
\documentclass[11pt,a4paper]{article}
\usepackage[spanish,es-noquoting]{babel}
\usepackage[utf8]{inputenc}
\usepackage[T1]{fontenc}
\usepackage{lmodern}
\usepackage{amsmath,amssymb,mathtools}
\usepackage{geometry}
\usepackage{enumitem}
\usepackage{multicol}
\usepackage{fancyhdr}
\usepackage{xcolor}
\usepackage{hyperref}
\usepackage{titlesec}
\usepackage{microtype}
\geometry{margin=1.55cm}
\hypersetup{colorlinks=true,linkcolor=black,urlcolor=blue}
\setlength{\parindent}{0pt}
\setlength{\parskip}{4pt}
\setlist[enumerate]{leftmargin=*, itemsep=3pt, topsep=3pt}
\allowdisplaybreaks
\pagestyle{fancy}
\fancyhf{}
\lhead{Práctico de cálculo}
\rhead{Funciones, derivadas y límites}
\cfoot{\thepage}
\newcommand{\R}{\mathbb{R}}
\titleformat{\section}{\Large\bfseries}{\thesection.}{0.5em}{}
\titleformat{\subsection}{\large\bfseries}{\thesubsection.}{0.5em}{}
\begin{document}
\begin{titlepage}
\centering
\vspace*{2cm}
{\Huge\bfseries Práctico Extenso de Funciones, Derivadas y Límites\\[0.4cm]}
{\Large Dominio, imagen, asíntotas, crecimiento, concavidad y límites\\[1cm]}
\vfill
{\large Documento de práctica progresiva\\}
{\large De nivel básico a nivel imposible\\[0.5cm]}
{\large Con solucionario desarrollado paso a paso\\}
\vfill
{\large Elaborado en \LaTeX\\}
\end{titlepage}
\tableofcontents
\newpage
\section*{Indicaciones}
\addcontentsline{toc}{section}{Indicaciones}
\begin{itemize}
\item Cada bloque tiene cinco niveles de dificultad.
\item Cada nivel contiene veinte ejercicios.
\item Primero resuelve sin mirar el solucionario.
\item El solucionario usa desarrollo matemático paso a paso, con texto solo cuando ayuda a ubicar una idea.
\end{itemize}
\newpage
\section{Dominio, dominio de imagen y asíntotas horizontales}

\subsection{Nivel 1 - Básico}

\begin{enumerate}[label=\textbf{E\arabic*.},start=1]

\item Determina dominio, imagen y asíntotas horizontales de $f(x)=1x-9$.

\item Determina dominio, imagen y asíntotas horizontales de $f(x)=(x+2)^2-1$.

\item Determina dominio, imagen y asíntotas horizontales de $f(x)=\sqrt{x+2}+0$.

\item Determina dominio, imagen y asíntota horizontal de $f(x)=\frac{1}{x+1}+1$.

\item Determina dominio, imagen y asíntota horizontal de $f(x)=2^{x-2}+3$.

\item Determina dominio, imagen y asíntotas horizontales de $f(x)=2x-4$.

\item Determina dominio, imagen y asíntotas horizontales de $f(x)=(x-3)^2-2$.

\item Determina dominio, imagen y asíntotas horizontales de $f(x)=\sqrt{x-3}+0$.

\item Determina dominio, imagen y asíntota horizontal de $f(x)=\frac{1}{x-4}+1$.

\item Determina dominio, imagen y asíntota horizontal de $f(x)=2^{x+0}+3$.

\item Determina dominio, imagen y asíntotas horizontales de $f(x)=3x+1$.

\item Determina dominio, imagen y asíntotas horizontales de $f(x)=(x-1)^2+3$.

\item Determina dominio, imagen y asíntotas horizontales de $f(x)=\sqrt{x+0}+0$.

\item Determina dominio, imagen y asíntota horizontal de $f(x)=\frac{1}{x+0}+1$.

\item Determina dominio, imagen y asíntota horizontal de $f(x)=2^{x+2}+3$.

\item Determina dominio, imagen y asíntotas horizontales de $f(x)=4x+6$.

\item Determina dominio, imagen y asíntotas horizontales de $f(x)=(x+1)^2+2$.

\item Determina dominio, imagen y asíntotas horizontales de $f(x)=\sqrt{x+3}+0$.

\item Determina dominio, imagen y asíntota horizontal de $f(x)=\frac{1}{x+4}+1$.

\item Determina dominio, imagen y asíntota horizontal de $f(x)=2^{x-3}+3$.

\end{enumerate}

\subsection{Nivel 2 - Intermedio bajo}

\begin{enumerate}[label=\textbf{E\arabic*.},start=21]

\item Determina dominio, imagen y asíntota horizontal de $f(x)=\frac{1x-3}{x+1}$.

\item Determina dominio, imagen y asíntotas horizontales de $f(x)=\ln(x+1)$.

\item Determina dominio, imagen y asíntota horizontal de $f(x)=-3e^{x+0}-1$.

\item Determina dominio, imagen y asíntota horizontal de $f(x)=\frac{4}{(x-2)^2}+1$.

\item Determina dominio, imagen y asíntota horizontal de $f(x)=\ln\left(\frac{x-5}{x-8}\right)$.

\item Determina dominio, imagen y asíntota horizontal de $f(x)=\frac{1x+2}{x+6}$.

\item Determina dominio, imagen y asíntotas horizontales de $f(x)=\ln(x+2)$.

\item Determina dominio, imagen y asíntota horizontal de $f(x)=-4e^{x+0}-3$.

\item Determina dominio, imagen y asíntota horizontal de $f(x)=\frac{4}{(x+1)^2}+0$.

\item Determina dominio, imagen y asíntota horizontal de $f(x)=\ln\left(\frac{x-5}{x-8}\right)$.

\item Determina dominio, imagen y asíntota horizontal de $f(x)=\frac{1x+0}{x+5}$.

\item Determina dominio, imagen y asíntotas horizontales de $f(x)=\ln(x-3)$.

\item Determina dominio, imagen y asíntota horizontal de $f(x)=-1e^{x+0}+2$.

\item Determina dominio, imagen y asíntota horizontal de $f(x)=\frac{4}{(x+0)^2}-1$.

\item Determina dominio, imagen y asíntota horizontal de $f(x)=\ln\left(\frac{x-5}{x-8}\right)$.

\item Determina dominio, imagen y asíntota horizontal de $f(x)=\frac{1x-2}{x+4}$.

\item Determina dominio, imagen y asíntotas horizontales de $f(x)=\ln(x-2)$.

\item Determina dominio, imagen y asíntota horizontal de $f(x)=-2e^{x+0}+0$.

\item Determina dominio, imagen y asíntota horizontal de $f(x)=\frac{4}{(x-1)^2}-2$.

\item Determina dominio, imagen y asíntota horizontal de $f(x)=\ln\left(\frac{x-5}{x-8}\right)$.

\end{enumerate}

\newpage

\subsection{Nivel 3 - Intermedio}

\begin{enumerate}[label=\textbf{E\arabic*.},start=41]

\item Determina dominio, imagen y asíntota horizontal de $f(x)=\sqrt{\frac{x-1}{x-3}}-1$.

\item Determina dominio, imagen y asíntota horizontal de $f(x)=\frac{2^x+2}{2^x+4}$. 

\item Determina dominio, imagen y asíntota horizontal de $f(x)=\ln\left(\frac{x^2-9}{x^2-25}\right)$.

\item Determina dominio, imagen y asíntotas horizontales de $f(x)=\arctan(x-2)$.

\item Determina dominio, imagen y asíntotas horizontales de $f(x)=\frac{e^x-e^{-x}}{e^x+e^{-x}}+-1$.

\item Determina dominio, imagen y asíntota horizontal de $f(x)=\sqrt{\frac{x-1}{x-3}}+0$.

\item Determina dominio, imagen y asíntota horizontal de $f(x)=\frac{2^x+2}{2^x+4}$. 

\item Determina dominio, imagen y asíntota horizontal de $f(x)=\ln\left(\frac{x^2-9}{x^2-25}\right)$.

\item Determina dominio, imagen y asíntotas horizontales de $f(x)=\arctan(x+1)$.

\item Determina dominio, imagen y asíntotas horizontales de $f(x)=\frac{e^x-e^{-x}}{e^x+e^{-x}}+0$.

\item Determina dominio, imagen y asíntota horizontal de $f(x)=\sqrt{\frac{x-1}{x-3}}+1$.

\item Determina dominio, imagen y asíntota horizontal de $f(x)=\frac{2^x+2}{2^x+4}$. 

\item Determina dominio, imagen y asíntota horizontal de $f(x)=\ln\left(\frac{x^2-9}{x^2-25}\right)$.

\item Determina dominio, imagen y asíntotas horizontales de $f(x)=\arctan(x+0)$.

\item Determina dominio, imagen y asíntotas horizontales de $f(x)=\frac{e^x-e^{-x}}{e^x+e^{-x}}+1$.

\item Determina dominio, imagen y asíntota horizontal de $f(x)=\sqrt{\frac{x-1}{x-3}}+2$.

\item Determina dominio, imagen y asíntota horizontal de $f(x)=\frac{2^x+2}{2^x+4}$. 

\item Determina dominio, imagen y asíntota horizontal de $f(x)=\ln\left(\frac{x^2-9}{x^2-25}\right)$.

\item Determina dominio, imagen y asíntotas horizontales de $f(x)=\arctan(x-1)$.

\item Determina dominio, imagen y asíntotas horizontales de $f(x)=\frac{e^x-e^{-x}}{e^x+e^{-x}}+2$.

\end{enumerate}

\subsection{Nivel 4 - Avanzado}

\begin{enumerate}[label=\textbf{E\arabic*.},start=61]

\item Determina dominio, imagen y asíntota horizontal de $f(x)=\sqrt{\frac{x-1}{x-3}}-1$.

\item Determina dominio, imagen y asíntota horizontal de $f(x)=\frac{2^x+2}{2^x+4}$. 

\item Determina dominio, imagen y asíntota horizontal de $f(x)=\ln\left(\frac{x^2-9}{x^2-25}\right)$.

\item Determina dominio, imagen y asíntotas horizontales de $f(x)=\arctan(x-2)$.

\item Determina dominio, imagen y asíntotas horizontales de $f(x)=\frac{e^x-e^{-x}}{e^x+e^{-x}}+-1$.

\item Determina dominio, imagen y asíntota horizontal de $f(x)=\sqrt{\frac{x-1}{x-3}}+0$.

\item Determina dominio, imagen y asíntota horizontal de $f(x)=\frac{2^x+2}{2^x+4}$. 

\item Determina dominio, imagen y asíntota horizontal de $f(x)=\ln\left(\frac{x^2-9}{x^2-25}\right)$.

\item Determina dominio, imagen y asíntotas horizontales de $f(x)=\arctan(x+1)$.

\item Determina dominio, imagen y asíntotas horizontales de $f(x)=\frac{e^x-e^{-x}}{e^x+e^{-x}}+0$.

\item Determina dominio, imagen y asíntota horizontal de $f(x)=\sqrt{\frac{x-1}{x-3}}+1$.

\item Determina dominio, imagen y asíntota horizontal de $f(x)=\frac{2^x+2}{2^x+4}$. 

\item Determina dominio, imagen y asíntota horizontal de $f(x)=\ln\left(\frac{x^2-9}{x^2-25}\right)$.

\item Determina dominio, imagen y asíntotas horizontales de $f(x)=\arctan(x+0)$.

\item Determina dominio, imagen y asíntotas horizontales de $f(x)=\frac{e^x-e^{-x}}{e^x+e^{-x}}+1$.

\item Determina dominio, imagen y asíntota horizontal de $f(x)=\sqrt{\frac{x-1}{x-3}}+2$.

\item Determina dominio, imagen y asíntota horizontal de $f(x)=\frac{2^x+2}{2^x+4}$. 

\item Determina dominio, imagen y asíntota horizontal de $f(x)=\ln\left(\frac{x^2-9}{x^2-25}\right)$.

\item Determina dominio, imagen y asíntotas horizontales de $f(x)=\arctan(x-1)$.

\item Determina dominio, imagen y asíntotas horizontales de $f(x)=\frac{e^x-e^{-x}}{e^x+e^{-x}}+2$.

\end{enumerate}

\newpage

\subsection{Nivel 5 - Imposible}

\begin{enumerate}[label=\textbf{E\arabic*.},start=81]

\item Determina dominio, imagen y asíntota horizontal de $f(x)=\sqrt{\frac{x-1}{x-3}}-1$.

\item Determina dominio, imagen y asíntota horizontal de $f(x)=\frac{2^x+2}{2^x+4}$. 

\item Determina dominio, imagen y asíntota horizontal de $f(x)=\ln\left(\frac{x^2-9}{x^2-25}\right)$.

\item Determina dominio, imagen y asíntotas horizontales de $f(x)=\arctan(x-2)$.

\item Determina dominio, imagen y asíntotas horizontales de $f(x)=\frac{e^x-e^{-x}}{e^x+e^{-x}}+-1$.

\item Determina dominio, imagen y asíntota horizontal de $f(x)=\sqrt{\frac{x-1}{x-3}}+0$.

\item Determina dominio, imagen y asíntota horizontal de $f(x)=\frac{2^x+2}{2^x+4}$. 

\item Determina dominio, imagen y asíntota horizontal de $f(x)=\ln\left(\frac{x^2-9}{x^2-25}\right)$.

\item Determina dominio, imagen y asíntotas horizontales de $f(x)=\arctan(x+1)$.

\item Determina dominio, imagen y asíntotas horizontales de $f(x)=\frac{e^x-e^{-x}}{e^x+e^{-x}}+0$.

\item Determina dominio, imagen y asíntota horizontal de $f(x)=\sqrt{\frac{x-1}{x-3}}+1$.

\item Determina dominio, imagen y asíntota horizontal de $f(x)=\frac{2^x+2}{2^x+4}$. 

\item Determina dominio, imagen y asíntota horizontal de $f(x)=\ln\left(\frac{x^2-9}{x^2-25}\right)$.

\item Determina dominio, imagen y asíntotas horizontales de $f(x)=\arctan(x+0)$.

\item Determina dominio, imagen y asíntotas horizontales de $f(x)=\frac{e^x-e^{-x}}{e^x+e^{-x}}+1$.

\item Determina dominio, imagen y asíntota horizontal de $f(x)=\sqrt{\frac{x-1}{x-3}}+2$.

\item Determina dominio, imagen y asíntota horizontal de $f(x)=\frac{2^x+2}{2^x+4}$. 

\item Determina dominio, imagen y asíntota horizontal de $f(x)=\ln\left(\frac{x^2-9}{x^2-25}\right)$.

\item Determina dominio, imagen y asíntotas horizontales de $f(x)=\arctan(x-1)$.

\item Determina dominio, imagen y asíntotas horizontales de $f(x)=\frac{e^x-e^{-x}}{e^x+e^{-x}}+2$.

\end{enumerate}

\newpage

\section{Crecimiento, puntos críticos y concavidad}

\subsection{Nivel 1 - Básico}

\begin{enumerate}[label=\textbf{E\arabic*.},start=101]

\item Para $f(x)=(x+4)^2-3$, determina crecimiento, puntos críticos y concavidad.

\item Para $f(x)=-(x+3)^2-2$, determina crecimiento, puntos críticos y concavidad.

\item Para $f(x)=(x+2)^2-1$, determina crecimiento, puntos críticos y concavidad.

\item Para $f(x)=-(x+1)^2+0$, determina crecimiento, puntos críticos y concavidad.

\item Para $f(x)=(x+0)^2+1$, determina crecimiento, puntos críticos y concavidad.

\item Para $f(x)=-(x-1)^2+2$, determina crecimiento, puntos críticos y concavidad.

\item Para $f(x)=(x-2)^2+3$, determina crecimiento, puntos críticos y concavidad.

\item Para $f(x)=-(x-3)^2-3$, determina crecimiento, puntos críticos y concavidad.

\item Para $f(x)=(x-4)^2-2$, determina crecimiento, puntos críticos y concavidad.

\item Para $f(x)=-(x+4)^2-1$, determina crecimiento, puntos críticos y concavidad.

\item Para $f(x)=(x+3)^2+0$, determina crecimiento, puntos críticos y concavidad.

\item Para $f(x)=-(x+2)^2+1$, determina crecimiento, puntos críticos y concavidad.

\item Para $f(x)=(x+1)^2+2$, determina crecimiento, puntos críticos y concavidad.

\item Para $f(x)=-(x+0)^2+3$, determina crecimiento, puntos críticos y concavidad.

\item Para $f(x)=(x-1)^2-3$, determina crecimiento, puntos críticos y concavidad.

\item Para $f(x)=-(x-2)^2-2$, determina crecimiento, puntos críticos y concavidad.

\item Para $f(x)=(x-3)^2-1$, determina crecimiento, puntos críticos y concavidad.

\item Para $f(x)=-(x-4)^2+0$, determina crecimiento, puntos críticos y concavidad.

\item Para $f(x)=(x+4)^2+1$, determina crecimiento, puntos críticos y concavidad.

\item Para $f(x)=-(x+3)^2+2$, determina crecimiento, puntos críticos y concavidad.

\end{enumerate}

\subsection{Nivel 2 - Intermedio bajo}

\begin{enumerate}[label=\textbf{E\arabic*.},start=121]

\item Para $f(x)=x^3-3(1)^2x$, determina crecimiento, puntos críticos y concavidad.

\item Para $f(x)=(x+2)^3$, determina crecimiento, puntos críticos y concavidad.

\item Para $f(x)=x^3-3(3)^2x$, determina crecimiento, puntos críticos y concavidad.

\item Para $f(x)=(x+0)^3$, determina crecimiento, puntos críticos y concavidad.

\item Para $f(x)=x^3-3(5)^2x$, determina crecimiento, puntos críticos y concavidad.

\item Para $f(x)=(x-2)^3$, determina crecimiento, puntos críticos y concavidad.

\item Para $f(x)=x^3-3(2)^2x$, determina crecimiento, puntos críticos y concavidad.

\item Para $f(x)=(x+3)^3$, determina crecimiento, puntos críticos y concavidad.

\item Para $f(x)=x^3-3(4)^2x$, determina crecimiento, puntos críticos y concavidad.

\item Para $f(x)=(x+1)^3$, determina crecimiento, puntos críticos y concavidad.

\item Para $f(x)=x^3-3(1)^2x$, determina crecimiento, puntos críticos y concavidad.

\item Para $f(x)=(x-1)^3$, determina crecimiento, puntos críticos y concavidad.

\item Para $f(x)=x^3-3(3)^2x$, determina crecimiento, puntos críticos y concavidad.

\item Para $f(x)=(x-3)^3$, determina crecimiento, puntos críticos y concavidad.

\item Para $f(x)=x^3-3(5)^2x$, determina crecimiento, puntos críticos y concavidad.

\item Para $f(x)=(x+2)^3$, determina crecimiento, puntos críticos y concavidad.

\item Para $f(x)=x^3-3(2)^2x$, determina crecimiento, puntos críticos y concavidad.

\item Para $f(x)=(x+0)^3$, determina crecimiento, puntos críticos y concavidad.

\item Para $f(x)=x^3-3(4)^2x$, determina crecimiento, puntos críticos y concavidad.

\item Para $f(x)=(x-2)^3$, determina crecimiento, puntos críticos y concavidad.

\end{enumerate}

\newpage

\subsection{Nivel 3 - Intermedio}

\begin{enumerate}[label=\textbf{E\arabic*.},start=141]

\item Para $f(x)=x+\frac{1}{x}$, determina crecimiento, puntos críticos y concavidad.

\item Para $f(x)=(x+1)^4$, determina crecimiento, puntos críticos y concavidad.

\item Para $f(x)=x+\frac{3}{x}$, determina crecimiento, puntos críticos y concavidad.

\item Para $f(x)=(x-1)^4$, determina crecimiento, puntos críticos y concavidad.

\item Para $f(x)=x+\frac{5}{x}$, determina crecimiento, puntos críticos y concavidad.

\item Para $f(x)=(x+2)^4$, determina crecimiento, puntos críticos y concavidad.

\item Para $f(x)=x+\frac{2}{x}$, determina crecimiento, puntos críticos y concavidad.

\item Para $f(x)=(x+0)^4$, determina crecimiento, puntos críticos y concavidad.

\item Para $f(x)=x+\frac{4}{x}$, determina crecimiento, puntos críticos y concavidad.

\item Para $f(x)=(x-2)^4$, determina crecimiento, puntos críticos y concavidad.

\item Para $f(x)=x+\frac{1}{x}$, determina crecimiento, puntos críticos y concavidad.

\item Para $f(x)=(x+1)^4$, determina crecimiento, puntos críticos y concavidad.

\item Para $f(x)=x+\frac{3}{x}$, determina crecimiento, puntos críticos y concavidad.

\item Para $f(x)=(x-1)^4$, determina crecimiento, puntos críticos y concavidad.

\item Para $f(x)=x+\frac{5}{x}$, determina crecimiento, puntos críticos y concavidad.

\item Para $f(x)=(x+2)^4$, determina crecimiento, puntos críticos y concavidad.

\item Para $f(x)=x+\frac{2}{x}$, determina crecimiento, puntos críticos y concavidad.

\item Para $f(x)=(x+0)^4$, determina crecimiento, puntos críticos y concavidad.

\item Para $f(x)=x+\frac{4}{x}$, determina crecimiento, puntos críticos y concavidad.

\item Para $f(x)=(x-2)^4$, determina crecimiento, puntos críticos y concavidad.

\end{enumerate}

\subsection{Nivel 4 - Avanzado}

\begin{enumerate}[label=\textbf{E\arabic*.},start=161]

\item Para $f(x)=e^x-2x$, determina crecimiento, puntos críticos y concavidad.

\item Para $f(x)=x^2\ln x$, determina crecimiento, puntos críticos y concavidad.

\item Para $f(x)=e^x-4x$, determina crecimiento, puntos críticos y concavidad.

\item Para $f(x)=x^2\ln x$, determina crecimiento, puntos críticos y concavidad.

\item Para $f(x)=e^x-6x$, determina crecimiento, puntos críticos y concavidad.

\item Para $f(x)=x^2\ln x$, determina crecimiento, puntos críticos y concavidad.

\item Para $f(x)=e^x-3x$, determina crecimiento, puntos críticos y concavidad.

\item Para $f(x)=x^2\ln x$, determina crecimiento, puntos críticos y concavidad.

\item Para $f(x)=e^x-5x$, determina crecimiento, puntos críticos y concavidad.

\item Para $f(x)=x^2\ln x$, determina crecimiento, puntos críticos y concavidad.

\item Para $f(x)=e^x-2x$, determina crecimiento, puntos críticos y concavidad.

\item Para $f(x)=x^2\ln x$, determina crecimiento, puntos críticos y concavidad.

\item Para $f(x)=e^x-4x$, determina crecimiento, puntos críticos y concavidad.

\item Para $f(x)=x^2\ln x$, determina crecimiento, puntos críticos y concavidad.

\item Para $f(x)=e^x-6x$, determina crecimiento, puntos críticos y concavidad.

\item Para $f(x)=x^2\ln x$, determina crecimiento, puntos críticos y concavidad.

\item Para $f(x)=e^x-3x$, determina crecimiento, puntos críticos y concavidad.

\item Para $f(x)=x^2\ln x$, determina crecimiento, puntos críticos y concavidad.

\item Para $f(x)=e^x-5x$, determina crecimiento, puntos críticos y concavidad.

\item Para $f(x)=x^2\ln x$, determina crecimiento, puntos críticos y concavidad.

\end{enumerate}

\newpage

\subsection{Nivel 5 - Imposible}

\begin{enumerate}[label=\textbf{E\arabic*.},start=181]

\item Para $f(x)=\frac{\ln x}{x}$, determina crecimiento, puntos críticos y concavidad.

\item Para $f(x)=x^2e^{-x}$, determina crecimiento, puntos críticos y concavidad.

\item Para $f(x)=\frac{x^2+1}{x-1}$, determina crecimiento, puntos críticos y concavidad.

\item Para $f(x)=x+\frac{1}{x}+\frac{1}{x^2}$, determina crecimiento, puntos críticos y concavidad.

\item Para $f(x)=\ln(1+x^2)$, determina crecimiento, puntos críticos y concavidad.

\item Para $f(x)=\frac{\ln x}{x}$, determina crecimiento, puntos críticos y concavidad.

\item Para $f(x)=x^2e^{-x}$, determina crecimiento, puntos críticos y concavidad.

\item Para $f(x)=\frac{x^2+1}{x-1}$, determina crecimiento, puntos críticos y concavidad.

\item Para $f(x)=x+\frac{1}{x}+\frac{1}{x^2}$, determina crecimiento, puntos críticos y concavidad.

\item Para $f(x)=\ln(1+x^2)$, determina crecimiento, puntos críticos y concavidad.

\item Para $f(x)=\frac{\ln x}{x}$, determina crecimiento, puntos críticos y concavidad.

\item Para $f(x)=x^2e^{-x}$, determina crecimiento, puntos críticos y concavidad.

\item Para $f(x)=\frac{x^2+1}{x-1}$, determina crecimiento, puntos críticos y concavidad.

\item Para $f(x)=x+\frac{1}{x}+\frac{1}{x^2}$, determina crecimiento, puntos críticos y concavidad.

\item Para $f(x)=\ln(1+x^2)$, determina crecimiento, puntos críticos y concavidad.

\item Para $f(x)=\frac{\ln x}{x}$, determina crecimiento, puntos críticos y concavidad.

\item Para $f(x)=x^2e^{-x}$, determina crecimiento, puntos críticos y concavidad.

\item Para $f(x)=\frac{x^2+1}{x-1}$, determina crecimiento, puntos críticos y concavidad.

\item Para $f(x)=x+\frac{1}{x}+\frac{1}{x^2}$, determina crecimiento, puntos críticos y concavidad.

\item Para $f(x)=\ln(1+x^2)$, determina crecimiento, puntos críticos y concavidad.

\end{enumerate}

\newpage

\section{Límites de fracciones algebraicas y fracción sintética}

\subsection{Nivel 1 - Básico}

\begin{enumerate}[label=\textbf{E\arabic*.},start=201]

\item Calcula $\displaystyle\lim_{x\to -3}\frac{1x-2}{x+1}$. 

\item Calcula $\displaystyle\lim_{x\to -2}\frac{2x-1}{x+2}$. 

\item Calcula $\displaystyle\lim_{x\to -1}\frac{3x+0}{x+3}$. 

\item Calcula $\displaystyle\lim_{x\to 2}\frac{4x+1}{x+1}$. 

\item Calcula $\displaystyle\lim_{x\to 1}\frac{1x+2}{x+2}$. 

\item Calcula $\displaystyle\lim_{x\to 2}\frac{2x-2}{x+3}$. 

\item Calcula $\displaystyle\lim_{x\to 3}\frac{3x-1}{x+1}$. 

\item Calcula $\displaystyle\lim_{x\to -3}\frac{4x+0}{x+2}$. 

\item Calcula $\displaystyle\lim_{x\to -2}\frac{1x+1}{x+3}$. 

\item Calcula $\displaystyle\lim_{x\to -1}\frac{2x+2}{x+1}$. 

\item Calcula $\displaystyle\lim_{x\to 2}\frac{3x-2}{x+2}$. 

\item Calcula $\displaystyle\lim_{x\to 1}\frac{4x-1}{x+3}$. 

\item Calcula $\displaystyle\lim_{x\to 2}\frac{1x+0}{x+1}$. 

\item Calcula $\displaystyle\lim_{x\to 3}\frac{2x+1}{x+2}$. 

\item Calcula $\displaystyle\lim_{x\to -3}\frac{3x+2}{x+3}$. 

\item Calcula $\displaystyle\lim_{x\to -2}\frac{4x-2}{x+1}$. 

\item Calcula $\displaystyle\lim_{x\to -1}\frac{1x-1}{x+2}$. 

\item Calcula $\displaystyle\lim_{x\to 2}\frac{2x+0}{x+3}$. 

\item Calcula $\displaystyle\lim_{x\to 1}\frac{3x+1}{x+1}$. 

\item Calcula $\displaystyle\lim_{x\to 2}\frac{4x+2}{x+2}$. 

\end{enumerate}

\subsection{Nivel 2 - Intermedio bajo}

\begin{enumerate}[label=\textbf{E\arabic*.},start=221]

\item Calcula $\displaystyle\lim_{x\to -3}\frac{x^2--5x+6}{x--3}$. 

\item Calcula $\displaystyle\lim_{x\to -2}\frac{x^2--2x+0}{x--2}$. 

\item Calcula $\displaystyle\lim_{x\to -1}\frac{x^2-1x+-2}{x--1}$. 

\item Calcula $\displaystyle\lim_{x\to 2}\frac{x^2-8x+12}{x-2}$. 

\item Calcula $\displaystyle\lim_{x\to 1}\frac{x^2-7x+6}{x-1}$. 

\item Calcula $\displaystyle\lim_{x\to 2}\frac{x^2-5x+6}{x-2}$. 

\item Calcula $\displaystyle\lim_{x\to 3}\frac{x^2-8x+15}{x-3}$. 

\item Calcula $\displaystyle\lim_{x\to -3}\frac{x^2--3x+0}{x--3}$. 

\item Calcula $\displaystyle\lim_{x\to -2}\frac{x^2-0x+-4}{x--2}$. 

\item Calcula $\displaystyle\lim_{x\to -1}\frac{x^2-3x+-4}{x--1}$. 

\item Calcula $\displaystyle\lim_{x\to 2}\frac{x^2-5x+6}{x-2}$. 

\item Calcula $\displaystyle\lim_{x\to 1}\frac{x^2-4x+3}{x-1}$. 

\item Calcula $\displaystyle\lim_{x\to 2}\frac{x^2-7x+10}{x-2}$. 

\item Calcula $\displaystyle\lim_{x\to 3}\frac{x^2-10x+21}{x-3}$. 

\item Calcula $\displaystyle\lim_{x\to -3}\frac{x^2--1x+-6}{x--3}$. 

\item Calcula $\displaystyle\lim_{x\to -2}\frac{x^2--3x+2}{x--2}$. 

\item Calcula $\displaystyle\lim_{x\to -1}\frac{x^2-0x+-1}{x--1}$. 

\item Calcula $\displaystyle\lim_{x\to 2}\frac{x^2-7x+10}{x-2}$. 

\item Calcula $\displaystyle\lim_{x\to 1}\frac{x^2-6x+5}{x-1}$. 

\item Calcula $\displaystyle\lim_{x\to 2}\frac{x^2-9x+14}{x-2}$. 

\end{enumerate}

\newpage

\subsection{Nivel 3 - Intermedio}

\begin{enumerate}[label=\textbf{E\arabic*.},start=241]

\item Calcula usando fracción sintética $\displaystyle\lim_{x\to -4}\frac{x^3 + x^2 - 11x + 4}{x--4}$. 

\item Calcula usando fracción sintética $\displaystyle\lim_{x\to -3}\frac{x^3 + x^2 - 4x + 6}{x--3}$. 

\item Calcula usando fracción sintética $\displaystyle\lim_{x\to -2}\frac{x^3 + x^2 + x + 6}{x--2}$. 

\item Calcula usando fracción sintética $\displaystyle\lim_{x\to -1}\frac{x^3 + x^2 + 4x + 4}{x--1}$. 

\item Calcula usando fracción sintética $\displaystyle\lim_{x\to 3}\frac{x^3 - 2x^2 + 2x - 15}{x-3}$. 

\item Calcula usando fracción sintética $\displaystyle\lim_{x\to 1}\frac{x^3 + x^2 - x - 1}{x-1}$. 

\item Calcula usando fracción sintética $\displaystyle\lim_{x\to 2}\frac{x^3 + x^2 - 4x - 4}{x-2}$. 

\item Calcula usando fracción sintética $\displaystyle\lim_{x\to 3}\frac{x^3 - 6x^2 + 12x - 9}{x-3}$. 

\item Calcula usando fracción sintética $\displaystyle\lim_{x\to 4}\frac{x^3 - 6x^2 + 12x - 16}{x-4}$. 

\item Calcula usando fracción sintética $\displaystyle\lim_{x\to -4}\frac{x^3 + 3x^2 + x + 20}{x--4}$. 

\item Calcula usando fracción sintética $\displaystyle\lim_{x\to -3}\frac{x^3 + 3x^2 + x + 3}{x--3}$. 

\item Calcula usando fracción sintética $\displaystyle\lim_{x\to -2}\frac{x^3 + 3x^2 + 4x + 4}{x--2}$. 

\item Calcula usando fracción sintética $\displaystyle\lim_{x\to -1}\frac{x^3 + 3x^2 + 5x + 3}{x--1}$. 

\item Calcula usando fracción sintética $\displaystyle\lim_{x\to 3}\frac{x^3 - 5x - 12}{x-3}$. 

\item Calcula usando fracción sintética $\displaystyle\lim_{x\to 1}\frac{x^3 - 4x^2 + 8x - 5}{x-1}$. 

\item Calcula usando fracción sintética $\displaystyle\lim_{x\to 2}\frac{x^3 - 4x^2 + 5x - 2}{x-2}$. 

\item Calcula usando fracción sintética $\displaystyle\lim_{x\to 3}\frac{x^3 - 4x^2 + 5x - 6}{x-3}$. 

\item Calcula usando fracción sintética $\displaystyle\lim_{x\to 4}\frac{x^3 - 4x^2 + 3x - 12}{x-4}$. 

\item Calcula usando fracción sintética $\displaystyle\lim_{x\to -4}\frac{x^3 + 5x^2 + 8x + 16}{x--4}$. 

\item Calcula usando fracción sintética $\displaystyle\lim_{x\to -3}\frac{x^3 + 5x^2 + 11x + 15}{x--3}$. 

\end{enumerate}

\subsection{Nivel 4 - Avanzado}

\begin{enumerate}[label=\textbf{E\arabic*.},start=261]

\item Calcula $\displaystyle\lim_{x\to\infty}\frac{1x^3+2x^2-1}{2x^3-1x+1}$. 

\item Calcula $\displaystyle\lim_{x\to -2}\frac{(x--2)(x-0)(x-2)}{(x--2)(x+3)}$. 

\item Calcula $\displaystyle\lim_{x\to\infty}\frac{3x^3+4x^2-3}{4x^3-3x+3}$. 

\item Calcula $\displaystyle\lim_{x\to 2}\frac{(x-2)(x-4)(x-6)}{(x-2)(x+3)}$. 

\item Calcula $\displaystyle\lim_{x\to\infty}\frac{5x^3+2x^2-2}{2x^3-5x+2}$. 

\item Calcula $\displaystyle\lim_{x\to 2}\frac{(x-2)(x-4)(x-6)}{(x-2)(x+3)}$. 

\item Calcula $\displaystyle\lim_{x\to\infty}\frac{2x^3+4x^2-1}{4x^3-2x+1}$. 

\item Calcula $\displaystyle\lim_{x\to -3}\frac{(x--3)(x--1)(x-1)}{(x--3)(x+4)}$. 

\item Calcula $\displaystyle\lim_{x\to\infty}\frac{4x^3+2x^2-3}{2x^3-4x+3}$. 

\item Calcula $\displaystyle\lim_{x\to -1}\frac{(x--1)(x-1)(x-3)}{(x--1)(x+2)}$. 

\item Calcula $\displaystyle\lim_{x\to\infty}\frac{1x^3+4x^2-2}{4x^3-1x+2}$. 

\item Calcula $\displaystyle\lim_{x\to 1}\frac{(x-1)(x-3)(x-5)}{(x-1)(x+2)}$. 

\item Calcula $\displaystyle\lim_{x\to\infty}\frac{3x^3+2x^2-1}{2x^3-3x+1}$. 

\item Calcula $\displaystyle\lim_{x\to 3}\frac{(x-3)(x-5)(x-7)}{(x-3)(x+4)}$. 

\item Calcula $\displaystyle\lim_{x\to\infty}\frac{5x^3+4x^2-3}{4x^3-5x+3}$. 

\item Calcula $\displaystyle\lim_{x\to -2}\frac{(x--2)(x-0)(x-2)}{(x--2)(x+3)}$. 

\item Calcula $\displaystyle\lim_{x\to\infty}\frac{2x^3+2x^2-2}{2x^3-2x+2}$. 

\item Calcula $\displaystyle\lim_{x\to 2}\frac{(x-2)(x-4)(x-6)}{(x-2)(x+3)}$. 

\item Calcula $\displaystyle\lim_{x\to\infty}\frac{4x^3+4x^2-1}{4x^3-4x+1}$. 

\item Calcula $\displaystyle\lim_{x\to 2}\frac{(x-2)(x-4)(x-6)}{(x-2)(x+3)}$. 

\end{enumerate}

\newpage

\subsection{Nivel 5 - Imposible}

\begin{enumerate}[label=\textbf{E\arabic*.},start=281]

\item Calcula usando fracción sintética $\displaystyle\lim_{x\to -4}\frac{x^4 + 2x^3 - 7x^2 + 6x + 8}{x--4}$. 

\item Calcula usando fracción sintética $\displaystyle\lim_{x\to -3}\frac{x^4 + 2x^3 - x^2 + 9x + 9}{x--3}$. 

\item Calcula usando fracción sintética $\displaystyle\lim_{x\to -2}\frac{x^4 + 2x^3 + 3x^2 + 10x + 8}{x--2}$. 

\item Calcula usando fracción sintética $\displaystyle\lim_{x\to -1}\frac{x^4 + 2x^3 + 5x^2 + 9x + 5}{x--1}$. 

\item Calcula usando fracción sintética $\displaystyle\lim_{x\to 5}\frac{x^4 - 3x^3 - 5x^2 - 23x - 10}{x-5}$. 

\item Calcula usando fracción sintética $\displaystyle\lim_{x\to 1}\frac{x^4 - 3x^3 + 8x^2 - 3x - 3}{x-1}$. 

\item Calcula usando fracción sintética $\displaystyle\lim_{x\to 2}\frac{x^4 - 3x^3 + 9x^2 - 10x - 8}{x-2}$. 

\item Calcula usando fracción sintética $\displaystyle\lim_{x\to 3}\frac{x^4 - 3x^3 + x^2 + 2x - 15}{x-3}$. 

\item Calcula usando fracción sintética $\displaystyle\lim_{x\to 4}\frac{x^4 - 3x^3 - 2x^2 - 6x - 8}{x-4}$. 

\item Calcula usando fracción sintética $\displaystyle\lim_{x\to -4}\frac{x^4 + 6x^3 + 11x^2 + 15x + 12}{x--4}$. 

\item Calcula usando fracción sintética $\displaystyle\lim_{x\to -3}\frac{x^4 + x^3 - 2x^2 + 16x + 12}{x--3}$. 

\item Calcula usando fracción sintética $\displaystyle\lim_{x\to -2}\frac{x^4 + x^3 + 3x^2 + 15x + 10}{x--2}$. 

\item Calcula usando fracción sintética $\displaystyle\lim_{x\to -1}\frac{x^4 + x^3 + 6x^2 + 8x + 2}{x--1}$. 

\item Calcula usando fracción sintética $\displaystyle\lim_{x\to 5}\frac{x^4 - 4x^3 + 2x^2 - 32x - 15}{x-5}$. 

\item Calcula usando fracción sintética $\displaystyle\lim_{x\to 1}\frac{x^4 + x^3 - x^2 + 3x - 4}{x-1}$. 

\item Calcula usando fracción sintética $\displaystyle\lim_{x\to 2}\frac{x^4 - 4x^3 + 6x^2 + x - 10}{x-2}$. 

\item Calcula usando fracción sintética $\displaystyle\lim_{x\to 3}\frac{x^4 - 4x^3 + 6x^2 - 7x - 6}{x-3}$. 

\item Calcula usando fracción sintética $\displaystyle\lim_{x\to 4}\frac{x^4 - 4x^3 + 4x^2 - 13x - 12}{x-4}$. 

\item Calcula usando fracción sintética $\displaystyle\lim_{x\to -4}\frac{x^4 + 5x^3 + 9x^2 + 24x + 16}{x--4}$. 

\item Calcula $\displaystyle\lim_{x\to1}\frac{x^5-5x+4}{x^4-4x+3}$ usando fracción sintética las veces necesarias.

\end{enumerate}

\newpage

\section{Límites exponenciales}

\subsection{Nivel 1 - Básico}

\begin{enumerate}[label=\textbf{E\arabic*.},start=301]

\item Calcula $\displaystyle\lim_{x\to -2}2^x$.

\item Calcula $\displaystyle\lim_{x\to -1}3^x$.

\item Calcula $\displaystyle\lim_{x\to 0}4^x$.

\item Calcula $\displaystyle\lim_{x\to 1}5^x$.

\item Calcula $\displaystyle\lim_{x\to 2}2^x$.

\item Calcula $\displaystyle\lim_{x\to -2}3^x$.

\item Calcula $\displaystyle\lim_{x\to -1}4^x$.

\item Calcula $\displaystyle\lim_{x\to 0}5^x$.

\item Calcula $\displaystyle\lim_{x\to 1}2^x$.

\item Calcula $\displaystyle\lim_{x\to 2}3^x$.

\item Calcula $\displaystyle\lim_{x\to -2}4^x$.

\item Calcula $\displaystyle\lim_{x\to -1}5^x$.

\item Calcula $\displaystyle\lim_{x\to 0}2^x$.

\item Calcula $\displaystyle\lim_{x\to 1}3^x$.

\item Calcula $\displaystyle\lim_{x\to 2}4^x$.

\item Calcula $\displaystyle\lim_{x\to -2}5^x$.

\item Calcula $\displaystyle\lim_{x\to -1}2^x$.

\item Calcula $\displaystyle\lim_{x\to 0}3^x$.

\item Calcula $\displaystyle\lim_{x\to 1}4^x$.

\item Calcula $\displaystyle\lim_{x\to 2}5^x$.

\end{enumerate}

\subsection{Nivel 2 - Intermedio bajo}

\begin{enumerate}[label=\textbf{E\arabic*.},start=321]

\item Calcula $\displaystyle\lim_{x\to\infty}\frac{1}{2^x}$.

\item Calcula $\displaystyle\lim_{x\to-\infty}3^x$.

\item Calcula $\displaystyle\lim_{x\to\infty}\frac{1}{4^x}$.

\item Calcula $\displaystyle\lim_{x\to-\infty}5^x$.

\item Calcula $\displaystyle\lim_{x\to\infty}\frac{1}{2^x}$.

\item Calcula $\displaystyle\lim_{x\to-\infty}3^x$.

\item Calcula $\displaystyle\lim_{x\to\infty}\frac{1}{4^x}$.

\item Calcula $\displaystyle\lim_{x\to-\infty}5^x$.

\item Calcula $\displaystyle\lim_{x\to\infty}\frac{1}{2^x}$.

\item Calcula $\displaystyle\lim_{x\to-\infty}3^x$.

\item Calcula $\displaystyle\lim_{x\to\infty}\frac{1}{4^x}$.

\item Calcula $\displaystyle\lim_{x\to-\infty}5^x$.

\item Calcula $\displaystyle\lim_{x\to\infty}\frac{1}{2^x}$.

\item Calcula $\displaystyle\lim_{x\to-\infty}3^x$.

\item Calcula $\displaystyle\lim_{x\to\infty}\frac{1}{4^x}$.

\item Calcula $\displaystyle\lim_{x\to-\infty}5^x$.

\item Calcula $\displaystyle\lim_{x\to\infty}\frac{1}{2^x}$.

\item Calcula $\displaystyle\lim_{x\to-\infty}3^x$.

\item Calcula $\displaystyle\lim_{x\to\infty}\frac{1}{4^x}$.

\item Calcula $\displaystyle\lim_{x\to-\infty}5^x$.

\end{enumerate}

\newpage

\subsection{Nivel 3 - Intermedio}

\begin{enumerate}[label=\textbf{E\arabic*.},start=341]

\item Calcula $\displaystyle\lim_{x\to0}\frac{2^{1x}-1}{x}$.

\item Calcula $\displaystyle\lim_{x\to0}\frac{3^{2x}-1}{x}$.

\item Calcula $\displaystyle\lim_{x\to0}\frac{4^{3x}-1}{x}$.

\item Calcula $\displaystyle\lim_{x\to0}\frac{5^{4x}-1}{x}$.

\item Calcula $\displaystyle\lim_{x\to0}\frac{6^{5x}-1}{x}$.

\item Calcula $\displaystyle\lim_{x\to0}\frac{7^{1x}-1}{x}$.

\item Calcula $\displaystyle\lim_{x\to0}\frac{8^{2x}-1}{x}$.

\item Calcula $\displaystyle\lim_{x\to0}\frac{2^{3x}-1}{x}$.

\item Calcula $\displaystyle\lim_{x\to0}\frac{3^{4x}-1}{x}$.

\item Calcula $\displaystyle\lim_{x\to0}\frac{4^{5x}-1}{x}$.

\item Calcula $\displaystyle\lim_{x\to0}\frac{5^{1x}-1}{x}$.

\item Calcula $\displaystyle\lim_{x\to0}\frac{6^{2x}-1}{x}$.

\item Calcula $\displaystyle\lim_{x\to0}\frac{7^{3x}-1}{x}$.

\item Calcula $\displaystyle\lim_{x\to0}\frac{8^{4x}-1}{x}$.

\item Calcula $\displaystyle\lim_{x\to0}\frac{2^{5x}-1}{x}$.

\item Calcula $\displaystyle\lim_{x\to0}\frac{3^{1x}-1}{x}$.

\item Calcula $\displaystyle\lim_{x\to0}\frac{4^{2x}-1}{x}$.

\item Calcula $\displaystyle\lim_{x\to0}\frac{5^{3x}-1}{x}$.

\item Calcula $\displaystyle\lim_{x\to0}\frac{6^{4x}-1}{x}$.

\item Calcula $\displaystyle\lim_{x\to0}\frac{7^{5x}-1}{x}$.

\end{enumerate}

\subsection{Nivel 4 - Avanzado}

\begin{enumerate}[label=\textbf{E\arabic*.},start=361]

\item Calcula $\displaystyle\lim_{x\to0}\frac{e^{1x}-e^{1x}}{x}$.

\item Calcula $\displaystyle\lim_{x\to\infty}\frac{e^{2x}+2}{e^{2x}-2}$.

\item Calcula $\displaystyle\lim_{x\to0}\frac{e^{3x}-e^{3x}}{x}$.

\item Calcula $\displaystyle\lim_{x\to\infty}\frac{e^{4x}+4}{e^{4x}-4}$.

\item Calcula $\displaystyle\lim_{x\to0}\frac{e^{5x}-e^{5x}}{x}$.

\item Calcula $\displaystyle\lim_{x\to\infty}\frac{e^{6x}+1}{e^{6x}-1}$.

\item Calcula $\displaystyle\lim_{x\to0}\frac{e^{7x}-e^{2x}}{x}$.

\item Calcula $\displaystyle\lim_{x\to\infty}\frac{e^{1x}+3}{e^{1x}-3}$.

\item Calcula $\displaystyle\lim_{x\to0}\frac{e^{2x}-e^{4x}}{x}$.

\item Calcula $\displaystyle\lim_{x\to\infty}\frac{e^{3x}+5}{e^{3x}-5}$.

\item Calcula $\displaystyle\lim_{x\to0}\frac{e^{4x}-e^{1x}}{x}$.

\item Calcula $\displaystyle\lim_{x\to\infty}\frac{e^{5x}+2}{e^{5x}-2}$.

\item Calcula $\displaystyle\lim_{x\to0}\frac{e^{6x}-e^{3x}}{x}$.

\item Calcula $\displaystyle\lim_{x\to\infty}\frac{e^{7x}+4}{e^{7x}-4}$.

\item Calcula $\displaystyle\lim_{x\to0}\frac{e^{1x}-e^{5x}}{x}$.

\item Calcula $\displaystyle\lim_{x\to\infty}\frac{e^{2x}+1}{e^{2x}-1}$.

\item Calcula $\displaystyle\lim_{x\to0}\frac{e^{3x}-e^{2x}}{x}$.

\item Calcula $\displaystyle\lim_{x\to\infty}\frac{e^{4x}+3}{e^{4x}-3}$.

\item Calcula $\displaystyle\lim_{x\to0}\frac{e^{5x}-e^{4x}}{x}$.

\item Calcula $\displaystyle\lim_{x\to\infty}\frac{e^{6x}+5}{e^{6x}-5}$.

\end{enumerate}

\newpage

\subsection{Nivel 5 - Imposible}

\begin{enumerate}[label=\textbf{E\arabic*.},start=381]

\item Calcula $\displaystyle \lim_{x\to0}\frac{2^x-3^x}{x}$.

\item Calcula $\displaystyle \lim_{x\to0}\frac{a^x-b^x}{x}\quad(a,b>0)$.

\item Calcula $\displaystyle \lim_{x\to0}\left(\frac{e^x-1}{x}\right)^{1/x}$.

\item Calcula $\displaystyle \lim_{x\to\infty}\left(1+\frac{3}{x}\right)^{2x}$.

\item Calcula $\displaystyle \lim_{x\to0}\frac{e^{\sin x}-e^x}{x^3}$.

\item Calcula $\displaystyle \lim_{x\to0}\frac{2^x-3^x}{x}$.

\item Calcula $\displaystyle \lim_{x\to0}\frac{a^x-b^x}{x}\quad(a,b>0)$.

\item Calcula $\displaystyle \lim_{x\to0}\left(\frac{e^x-1}{x}\right)^{1/x}$.

\item Calcula $\displaystyle \lim_{x\to\infty}\left(1+\frac{3}{x}\right)^{2x}$.

\item Calcula $\displaystyle \lim_{x\to0}\frac{e^{\sin x}-e^x}{x^3}$.

\item Calcula $\displaystyle \lim_{x\to0}\frac{2^x-3^x}{x}$.

\item Calcula $\displaystyle \lim_{x\to0}\frac{a^x-b^x}{x}\quad(a,b>0)$.

\item Calcula $\displaystyle \lim_{x\to0}\left(\frac{e^x-1}{x}\right)^{1/x}$.

\item Calcula $\displaystyle \lim_{x\to\infty}\left(1+\frac{3}{x}\right)^{2x}$.

\item Calcula $\displaystyle \lim_{x\to0}\frac{e^{\sin x}-e^x}{x^3}$.

\item Calcula $\displaystyle \lim_{x\to0}\frac{2^x-3^x}{x}$.

\item Calcula $\displaystyle \lim_{x\to0}\frac{a^x-b^x}{x}\quad(a,b>0)$.

\item Calcula $\displaystyle \lim_{x\to0}\left(\frac{e^x-1}{x}\right)^{1/x}$.

\item Calcula $\displaystyle \lim_{x\to\infty}\left(1+\frac{3}{x}\right)^{2x}$.

\item Calcula $\displaystyle \lim_{x\to0}\frac{e^{\sin x}-e^x}{x^3}$.

\end{enumerate}

\newpage

\section{Límites trigonométricos}

\subsection{Nivel 1 - Básico}

\begin{enumerate}[label=\textbf{E\arabic*.},start=401]

\item Calcula $\displaystyle\lim_{x\to0}\frac{\sin(1x)}{1x}$.

\item Calcula $\displaystyle\lim_{x\to0}\frac{\sin(2x)}{2x}$.

\item Calcula $\displaystyle\lim_{x\to0}\frac{\sin(3x)}{3x}$.

\item Calcula $\displaystyle\lim_{x\to0}\frac{\sin(4x)}{4x}$.

\item Calcula $\displaystyle\lim_{x\to0}\frac{\sin(5x)}{5x}$.

\item Calcula $\displaystyle\lim_{x\to0}\frac{\sin(6x)}{1x}$.

\item Calcula $\displaystyle\lim_{x\to0}\frac{\sin(7x)}{2x}$.

\item Calcula $\displaystyle\lim_{x\to0}\frac{\sin(8x)}{3x}$.

\item Calcula $\displaystyle\lim_{x\to0}\frac{\sin(9x)}{4x}$.

\item Calcula $\displaystyle\lim_{x\to0}\frac{\sin(1x)}{5x}$.

\item Calcula $\displaystyle\lim_{x\to0}\frac{\sin(2x)}{1x}$.

\item Calcula $\displaystyle\lim_{x\to0}\frac{\sin(3x)}{2x}$.

\item Calcula $\displaystyle\lim_{x\to0}\frac{\sin(4x)}{3x}$.

\item Calcula $\displaystyle\lim_{x\to0}\frac{\sin(5x)}{4x}$.

\item Calcula $\displaystyle\lim_{x\to0}\frac{\sin(6x)}{5x}$.

\item Calcula $\displaystyle\lim_{x\to0}\frac{\sin(7x)}{1x}$.

\item Calcula $\displaystyle\lim_{x\to0}\frac{\sin(8x)}{2x}$.

\item Calcula $\displaystyle\lim_{x\to0}\frac{\sin(9x)}{3x}$.

\item Calcula $\displaystyle\lim_{x\to0}\frac{\sin(1x)}{4x}$.

\item Calcula $\displaystyle\lim_{x\to0}\frac{\sin(2x)}{5x}$.

\end{enumerate}

\subsection{Nivel 2 - Intermedio bajo}

\begin{enumerate}[label=\textbf{E\arabic*.},start=421]

\item Calcula $\displaystyle\lim_{x\to0}\frac{\tan(1x)}{1x}$.

\item Calcula $\displaystyle\lim_{x\to0}\frac{\sin(2x)}{\sin(2x)}$.

\item Calcula $\displaystyle\lim_{x\to0}\frac{\tan(3x)}{3x}$.

\item Calcula $\displaystyle\lim_{x\to0}\frac{\sin(4x)}{\sin(4x)}$.

\item Calcula $\displaystyle\lim_{x\to0}\frac{\tan(5x)}{5x}$.

\item Calcula $\displaystyle\lim_{x\to0}\frac{\sin(6x)}{\sin(6x)}$.

\item Calcula $\displaystyle\lim_{x\to0}\frac{\tan(7x)}{1x}$.

\item Calcula $\displaystyle\lim_{x\to0}\frac{\sin(1x)}{\sin(2x)}$.

\item Calcula $\displaystyle\lim_{x\to0}\frac{\tan(2x)}{3x}$.

\item Calcula $\displaystyle\lim_{x\to0}\frac{\sin(3x)}{\sin(4x)}$.

\item Calcula $\displaystyle\lim_{x\to0}\frac{\tan(4x)}{5x}$.

\item Calcula $\displaystyle\lim_{x\to0}\frac{\sin(5x)}{\sin(6x)}$.

\item Calcula $\displaystyle\lim_{x\to0}\frac{\tan(6x)}{1x}$.

\item Calcula $\displaystyle\lim_{x\to0}\frac{\sin(7x)}{\sin(2x)}$.

\item Calcula $\displaystyle\lim_{x\to0}\frac{\tan(1x)}{3x}$.

\item Calcula $\displaystyle\lim_{x\to0}\frac{\sin(2x)}{\sin(4x)}$.

\item Calcula $\displaystyle\lim_{x\to0}\frac{\tan(3x)}{5x}$.

\item Calcula $\displaystyle\lim_{x\to0}\frac{\sin(4x)}{\sin(6x)}$.

\item Calcula $\displaystyle\lim_{x\to0}\frac{\tan(5x)}{1x}$.

\item Calcula $\displaystyle\lim_{x\to0}\frac{\sin(6x)}{\sin(2x)}$.

\end{enumerate}

\newpage

\subsection{Nivel 3 - Intermedio}

\begin{enumerate}[label=\textbf{E\arabic*.},start=441]

\item Calcula $\displaystyle\lim_{x\to0}\frac{1-\cos(1x)}{1x^2}$.

\item Calcula $\displaystyle\lim_{x\to0}\frac{1-\cos(2x)}{2x^2}$.

\item Calcula $\displaystyle\lim_{x\to0}\frac{1-\cos(3x)}{3x^2}$.

\item Calcula $\displaystyle\lim_{x\to0}\frac{1-\cos(4x)}{4x^2}$.

\item Calcula $\displaystyle\lim_{x\to0}\frac{1-\cos(5x)}{5x^2}$.

\item Calcula $\displaystyle\lim_{x\to0}\frac{1-\cos(6x)}{1x^2}$.

\item Calcula $\displaystyle\lim_{x\to0}\frac{1-\cos(7x)}{2x^2}$.

\item Calcula $\displaystyle\lim_{x\to0}\frac{1-\cos(8x)}{3x^2}$.

\item Calcula $\displaystyle\lim_{x\to0}\frac{1-\cos(9x)}{4x^2}$.

\item Calcula $\displaystyle\lim_{x\to0}\frac{1-\cos(1x)}{5x^2}$.

\item Calcula $\displaystyle\lim_{x\to0}\frac{1-\cos(2x)}{1x^2}$.

\item Calcula $\displaystyle\lim_{x\to0}\frac{1-\cos(3x)}{2x^2}$.

\item Calcula $\displaystyle\lim_{x\to0}\frac{1-\cos(4x)}{3x^2}$.

\item Calcula $\displaystyle\lim_{x\to0}\frac{1-\cos(5x)}{4x^2}$.

\item Calcula $\displaystyle\lim_{x\to0}\frac{1-\cos(6x)}{5x^2}$.

\item Calcula $\displaystyle\lim_{x\to0}\frac{1-\cos(7x)}{1x^2}$.

\item Calcula $\displaystyle\lim_{x\to0}\frac{1-\cos(8x)}{2x^2}$.

\item Calcula $\displaystyle\lim_{x\to0}\frac{1-\cos(9x)}{3x^2}$.

\item Calcula $\displaystyle\lim_{x\to0}\frac{1-\cos(1x)}{4x^2}$.

\item Calcula $\displaystyle\lim_{x\to0}\frac{1-\cos(2x)}{5x^2}$.

\end{enumerate}

\subsection{Nivel 4 - Avanzado}

\begin{enumerate}[label=\textbf{E\arabic*.},start=461]

\item Calcula $\displaystyle\lim_{x\to0}\frac{\sin(1x)\tan(2x)}{x^2}$.

\item Calcula $\displaystyle\lim_{x\to0}\frac{\sin(2x)-\sin(3x)}{x}$.

\item Calcula $\displaystyle\lim_{x\to0}\frac{\sin(3x)\tan(4x)}{x^2}$.

\item Calcula $\displaystyle\lim_{x\to0}\frac{\sin(4x)-\sin(5x)}{x}$.

\item Calcula $\displaystyle\lim_{x\to0}\frac{\sin(5x)\tan(6x)}{x^2}$.

\item Calcula $\displaystyle\lim_{x\to0}\frac{\sin(1x)-\sin(7x)}{x}$.

\item Calcula $\displaystyle\lim_{x\to0}\frac{\sin(2x)\tan(2x)}{x^2}$.

\item Calcula $\displaystyle\lim_{x\to0}\frac{\sin(3x)-\sin(3x)}{x}$.

\item Calcula $\displaystyle\lim_{x\to0}\frac{\sin(4x)\tan(4x)}{x^2}$.

\item Calcula $\displaystyle\lim_{x\to0}\frac{\sin(5x)-\sin(5x)}{x}$.

\item Calcula $\displaystyle\lim_{x\to0}\frac{\sin(1x)\tan(6x)}{x^2}$.

\item Calcula $\displaystyle\lim_{x\to0}\frac{\sin(2x)-\sin(7x)}{x}$.

\item Calcula $\displaystyle\lim_{x\to0}\frac{\sin(3x)\tan(2x)}{x^2}$.

\item Calcula $\displaystyle\lim_{x\to0}\frac{\sin(4x)-\sin(3x)}{x}$.

\item Calcula $\displaystyle\lim_{x\to0}\frac{\sin(5x)\tan(4x)}{x^2}$.

\item Calcula $\displaystyle\lim_{x\to0}\frac{\sin(1x)-\sin(5x)}{x}$.

\item Calcula $\displaystyle\lim_{x\to0}\frac{\sin(2x)\tan(6x)}{x^2}$.

\item Calcula $\displaystyle\lim_{x\to0}\frac{\sin(3x)-\sin(7x)}{x}$.

\item Calcula $\displaystyle\lim_{x\to0}\frac{\sin(4x)\tan(2x)}{x^2}$.

\item Calcula $\displaystyle\lim_{x\to0}\frac{\sin(5x)-\sin(3x)}{x}$.

\end{enumerate}

\newpage

\subsection{Nivel 5 - Imposible}

\begin{enumerate}[label=\textbf{E\arabic*.},start=481]

\item Calcula $\displaystyle \lim_{x\to0}\frac{\sin x-x\cos x}{x^3}$.

\item Calcula $\displaystyle \lim_{x\to0}\frac{\tan x-x}{x^3}$.

\item Calcula $\displaystyle \lim_{x\to0}\frac{\sin(3x)-3\sin x}{x^3}$.

\item Calcula $\displaystyle \lim_{x\to0}\frac{1-\cos x}{\sin^2x}$.

\item Calcula $\displaystyle \lim_{x\to0}\frac{\sin x\sin2x\sin3x}{x^3}$.

\item Calcula $\displaystyle \lim_{x\to0}\frac{\sin x-x\cos x}{x^3}$.

\item Calcula $\displaystyle \lim_{x\to0}\frac{\tan x-x}{x^3}$.

\item Calcula $\displaystyle \lim_{x\to0}\frac{\sin(3x)-3\sin x}{x^3}$.

\item Calcula $\displaystyle \lim_{x\to0}\frac{1-\cos x}{\sin^2x}$.

\item Calcula $\displaystyle \lim_{x\to0}\frac{\sin x\sin2x\sin3x}{x^3}$.

\item Calcula $\displaystyle \lim_{x\to0}\frac{\sin x-x\cos x}{x^3}$.

\item Calcula $\displaystyle \lim_{x\to0}\frac{\tan x-x}{x^3}$.

\item Calcula $\displaystyle \lim_{x\to0}\frac{\sin(3x)-3\sin x}{x^3}$.

\item Calcula $\displaystyle \lim_{x\to0}\frac{1-\cos x}{\sin^2x}$.

\item Calcula $\displaystyle \lim_{x\to0}\frac{\sin x\sin2x\sin3x}{x^3}$.

\item Calcula $\displaystyle \lim_{x\to0}\frac{\sin x-x\cos x}{x^3}$.

\item Calcula $\displaystyle \lim_{x\to0}\frac{\tan x-x}{x^3}$.

\item Calcula $\displaystyle \lim_{x\to0}\frac{\sin(3x)-3\sin x}{x^3}$.

\item Calcula $\displaystyle \lim_{x\to0}\frac{1-\cos x}{\sin^2x}$.

\item Calcula $\displaystyle \lim_{x\to0}\frac{\sin x\sin2x\sin3x}{x^3}$.

\end{enumerate}

\newpage

\section{Límites con conjugados normales y trigonométricos}

\subsection{Nivel 1 - Básico}

\begin{enumerate}[label=\textbf{E\arabic*.},start=501]

\item Calcula $\displaystyle\lim_{x\to0}\frac{\sqrt{x+1}-\sqrt{1}}{x}$.

\item Calcula $\displaystyle\lim_{x\to0}\frac{\sqrt{x+2}-\sqrt{2}}{x}$.

\item Calcula $\displaystyle\lim_{x\to0}\frac{\sqrt{x+3}-\sqrt{3}}{x}$.

\item Calcula $\displaystyle\lim_{x\to0}\frac{\sqrt{x+4}-\sqrt{4}}{x}$.

\item Calcula $\displaystyle\lim_{x\to0}\frac{\sqrt{x+5}-\sqrt{5}}{x}$.

\item Calcula $\displaystyle\lim_{x\to0}\frac{\sqrt{x+6}-\sqrt{6}}{x}$.

\item Calcula $\displaystyle\lim_{x\to0}\frac{\sqrt{x+7}-\sqrt{7}}{x}$.

\item Calcula $\displaystyle\lim_{x\to0}\frac{\sqrt{x+1}-\sqrt{1}}{x}$.

\item Calcula $\displaystyle\lim_{x\to0}\frac{\sqrt{x+2}-\sqrt{2}}{x}$.

\item Calcula $\displaystyle\lim_{x\to0}\frac{\sqrt{x+3}-\sqrt{3}}{x}$.

\item Calcula $\displaystyle\lim_{x\to0}\frac{\sqrt{x+4}-\sqrt{4}}{x}$.

\item Calcula $\displaystyle\lim_{x\to0}\frac{\sqrt{x+5}-\sqrt{5}}{x}$.

\item Calcula $\displaystyle\lim_{x\to0}\frac{\sqrt{x+6}-\sqrt{6}}{x}$.

\item Calcula $\displaystyle\lim_{x\to0}\frac{\sqrt{x+7}-\sqrt{7}}{x}$.

\item Calcula $\displaystyle\lim_{x\to0}\frac{\sqrt{x+1}-\sqrt{1}}{x}$.

\item Calcula $\displaystyle\lim_{x\to0}\frac{\sqrt{x+2}-\sqrt{2}}{x}$.

\item Calcula $\displaystyle\lim_{x\to0}\frac{\sqrt{x+3}-\sqrt{3}}{x}$.

\item Calcula $\displaystyle\lim_{x\to0}\frac{\sqrt{x+4}-\sqrt{4}}{x}$.

\item Calcula $\displaystyle\lim_{x\to0}\frac{\sqrt{x+5}-\sqrt{5}}{x}$.

\item Calcula $\displaystyle\lim_{x\to0}\frac{\sqrt{x+6}-\sqrt{6}}{x}$.

\end{enumerate}

\subsection{Nivel 2 - Intermedio bajo}

\begin{enumerate}[label=\textbf{E\arabic*.},start=521]

\item Calcula $\displaystyle\lim_{x\to -2}\frac{\sqrt{x+1}-\sqrt{-1}}{x--2}$.

\item Calcula $\displaystyle\lim_{x\to -1}\frac{\sqrt{x+2}-\sqrt{1}}{x--1}$.

\item Calcula $\displaystyle\lim_{x\to 0}\frac{\sqrt{x+3}-\sqrt{3}}{x-0}$.

\item Calcula $\displaystyle\lim_{x\to 1}\frac{\sqrt{x+4}-\sqrt{5}}{x-1}$.

\item Calcula $\displaystyle\lim_{x\to 2}\frac{\sqrt{x+5}-\sqrt{7}}{x-2}$.

\item Calcula $\displaystyle\lim_{x\to -2}\frac{\sqrt{x+6}-\sqrt{4}}{x--2}$.

\item Calcula $\displaystyle\lim_{x\to -1}\frac{\sqrt{x+7}-\sqrt{6}}{x--1}$.

\item Calcula $\displaystyle\lim_{x\to 0}\frac{\sqrt{x+1}-\sqrt{1}}{x-0}$.

\item Calcula $\displaystyle\lim_{x\to 1}\frac{\sqrt{x+2}-\sqrt{3}}{x-1}$.

\item Calcula $\displaystyle\lim_{x\to 2}\frac{\sqrt{x+3}-\sqrt{5}}{x-2}$.

\item Calcula $\displaystyle\lim_{x\to -2}\frac{\sqrt{x+4}-\sqrt{2}}{x--2}$.

\item Calcula $\displaystyle\lim_{x\to -1}\frac{\sqrt{x+5}-\sqrt{4}}{x--1}$.

\item Calcula $\displaystyle\lim_{x\to 0}\frac{\sqrt{x+6}-\sqrt{6}}{x-0}$.

\item Calcula $\displaystyle\lim_{x\to 1}\frac{\sqrt{x+7}-\sqrt{8}}{x-1}$.

\item Calcula $\displaystyle\lim_{x\to 2}\frac{\sqrt{x+1}-\sqrt{3}}{x-2}$.

\item Calcula $\displaystyle\lim_{x\to -2}\frac{\sqrt{x+2}-\sqrt{0}}{x--2}$.

\item Calcula $\displaystyle\lim_{x\to -1}\frac{\sqrt{x+3}-\sqrt{2}}{x--1}$.

\item Calcula $\displaystyle\lim_{x\to 0}\frac{\sqrt{x+4}-\sqrt{4}}{x-0}$.

\item Calcula $\displaystyle\lim_{x\to 1}\frac{\sqrt{x+5}-\sqrt{6}}{x-1}$.

\item Calcula $\displaystyle\lim_{x\to 2}\frac{\sqrt{x+6}-\sqrt{8}}{x-2}$.

\end{enumerate}

\newpage

\subsection{Nivel 3 - Intermedio}

\begin{enumerate}[label=\textbf{E\arabic*.},start=541]

\item Calcula $\displaystyle\lim_{x\to0}\frac{1-\cos(1x)}{1x^2}$ usando conjugado trigonométrico.

\item Calcula $\displaystyle\lim_{x\to0}\frac{1-\cos(2x)}{2x^2}$ usando conjugado trigonométrico.

\item Calcula $\displaystyle\lim_{x\to0}\frac{1-\cos(3x)}{3x^2}$ usando conjugado trigonométrico.

\item Calcula $\displaystyle\lim_{x\to0}\frac{1-\cos(4x)}{4x^2}$ usando conjugado trigonométrico.

\item Calcula $\displaystyle\lim_{x\to0}\frac{1-\cos(5x)}{5x^2}$ usando conjugado trigonométrico.

\item Calcula $\displaystyle\lim_{x\to0}\frac{1-\cos(6x)}{1x^2}$ usando conjugado trigonométrico.

\item Calcula $\displaystyle\lim_{x\to0}\frac{1-\cos(7x)}{2x^2}$ usando conjugado trigonométrico.

\item Calcula $\displaystyle\lim_{x\to0}\frac{1-\cos(8x)}{3x^2}$ usando conjugado trigonométrico.

\item Calcula $\displaystyle\lim_{x\to0}\frac{1-\cos(1x)}{4x^2}$ usando conjugado trigonométrico.

\item Calcula $\displaystyle\lim_{x\to0}\frac{1-\cos(2x)}{5x^2}$ usando conjugado trigonométrico.

\item Calcula $\displaystyle\lim_{x\to0}\frac{1-\cos(3x)}{1x^2}$ usando conjugado trigonométrico.

\item Calcula $\displaystyle\lim_{x\to0}\frac{1-\cos(4x)}{2x^2}$ usando conjugado trigonométrico.

\item Calcula $\displaystyle\lim_{x\to0}\frac{1-\cos(5x)}{3x^2}$ usando conjugado trigonométrico.

\item Calcula $\displaystyle\lim_{x\to0}\frac{1-\cos(6x)}{4x^2}$ usando conjugado trigonométrico.

\item Calcula $\displaystyle\lim_{x\to0}\frac{1-\cos(7x)}{5x^2}$ usando conjugado trigonométrico.

\item Calcula $\displaystyle\lim_{x\to0}\frac{1-\cos(8x)}{1x^2}$ usando conjugado trigonométrico.

\item Calcula $\displaystyle\lim_{x\to0}\frac{1-\cos(1x)}{2x^2}$ usando conjugado trigonométrico.

\item Calcula $\displaystyle\lim_{x\to0}\frac{1-\cos(2x)}{3x^2}$ usando conjugado trigonométrico.

\item Calcula $\displaystyle\lim_{x\to0}\frac{1-\cos(3x)}{4x^2}$ usando conjugado trigonométrico.

\item Calcula $\displaystyle\lim_{x\to0}\frac{1-\cos(4x)}{5x^2}$ usando conjugado trigonométrico.

\end{enumerate}

\subsection{Nivel 4 - Avanzado}

\begin{enumerate}[label=\textbf{E\arabic*.},start=561]

\item Calcula $\displaystyle\lim_{x\to0}\frac{\sqrt{1+\sin(1x)}-1}{x}$.

\item Calcula $\displaystyle\lim_{x\to0}\frac{\sqrt{1+\sin(2x)}-1}{x}$.

\item Calcula $\displaystyle\lim_{x\to0}\frac{\sqrt{1+\sin(3x)}-1}{x}$.

\item Calcula $\displaystyle\lim_{x\to0}\frac{\sqrt{1+\sin(4x)}-1}{x}$.

\item Calcula $\displaystyle\lim_{x\to0}\frac{\sqrt{1+\sin(5x)}-1}{x}$.

\item Calcula $\displaystyle\lim_{x\to0}\frac{\sqrt{1+\sin(6x)}-1}{x}$.

\item Calcula $\displaystyle\lim_{x\to0}\frac{\sqrt{1+\sin(1x)}-1}{x}$.

\item Calcula $\displaystyle\lim_{x\to0}\frac{\sqrt{1+\sin(2x)}-1}{x}$.

\item Calcula $\displaystyle\lim_{x\to0}\frac{\sqrt{1+\sin(3x)}-1}{x}$.

\item Calcula $\displaystyle\lim_{x\to0}\frac{\sqrt{1+\sin(4x)}-1}{x}$.

\item Calcula $\displaystyle\lim_{x\to0}\frac{\sqrt{1+\sin(5x)}-1}{x}$.

\item Calcula $\displaystyle\lim_{x\to0}\frac{\sqrt{1+\sin(6x)}-1}{x}$.

\item Calcula $\displaystyle\lim_{x\to0}\frac{\sqrt{1+\sin(1x)}-1}{x}$.

\item Calcula $\displaystyle\lim_{x\to0}\frac{\sqrt{1+\sin(2x)}-1}{x}$.

\item Calcula $\displaystyle\lim_{x\to0}\frac{\sqrt{1+\sin(3x)}-1}{x}$.

\item Calcula $\displaystyle\lim_{x\to0}\frac{\sqrt{1+\sin(4x)}-1}{x}$.

\item Calcula $\displaystyle\lim_{x\to0}\frac{\sqrt{1+\sin(5x)}-1}{x}$.

\item Calcula $\displaystyle\lim_{x\to0}\frac{\sqrt{1+\sin(6x)}-1}{x}$.

\item Calcula $\displaystyle\lim_{x\to0}\frac{\sqrt{1+\sin(1x)}-1}{x}$.

\item Calcula $\displaystyle\lim_{x\to0}\frac{\sqrt{1+\sin(2x)}-1}{x}$.

\end{enumerate}

\newpage

\subsection{Nivel 5 - Imposible}

\begin{enumerate}[label=\textbf{E\arabic*.},start=581]

\item Calcula $\displaystyle \lim_{x\to0}\frac{\sqrt{1+\sin x}-\sqrt{1-\sin x}}{x}$.

\item Calcula $\displaystyle \lim_{x\to0}\frac{\sqrt{\cos x}-1}{x^2}$.

\item Calcula $\displaystyle \lim_{x\to0}\frac{\sin x}{\sqrt{1+x}-\sqrt{1-x}}$.

\item Calcula $\displaystyle \lim_{x\to0}\frac{1-\cos x}{x(\sqrt{1+x}-1)}$.

\item Calcula $\displaystyle \lim_{x\to0}\frac{\sqrt{1+\tan x}-\sqrt{1-\tan x}}{\tan(2x)}$.

\item Calcula $\displaystyle \lim_{x\to0}\frac{\sqrt{1+\sin x}-\sqrt{1-\sin x}}{x}$.

\item Calcula $\displaystyle \lim_{x\to0}\frac{\sqrt{\cos x}-1}{x^2}$.

\item Calcula $\displaystyle \lim_{x\to0}\frac{\sin x}{\sqrt{1+x}-\sqrt{1-x}}$.

\item Calcula $\displaystyle \lim_{x\to0}\frac{1-\cos x}{x(\sqrt{1+x}-1)}$.

\item Calcula $\displaystyle \lim_{x\to0}\frac{\sqrt{1+\tan x}-\sqrt{1-\tan x}}{\tan(2x)}$.

\item Calcula $\displaystyle \lim_{x\to0}\frac{\sqrt{1+\sin x}-\sqrt{1-\sin x}}{x}$.

\item Calcula $\displaystyle \lim_{x\to0}\frac{\sqrt{\cos x}-1}{x^2}$.

\item Calcula $\displaystyle \lim_{x\to0}\frac{\sin x}{\sqrt{1+x}-\sqrt{1-x}}$.

\item Calcula $\displaystyle \lim_{x\to0}\frac{1-\cos x}{x(\sqrt{1+x}-1)}$.

\item Calcula $\displaystyle \lim_{x\to0}\frac{\sqrt{1+\tan x}-\sqrt{1-\tan x}}{\tan(2x)}$.

\item Calcula $\displaystyle \lim_{x\to0}\frac{\sqrt{1+\sin x}-\sqrt{1-\sin x}}{x}$.

\item Calcula $\displaystyle \lim_{x\to0}\frac{\sqrt{\cos x}-1}{x^2}$.

\item Calcula $\displaystyle \lim_{x\to0}\frac{\sin x}{\sqrt{1+x}-\sqrt{1-x}}$.

\item Calcula $\displaystyle \lim_{x\to0}\frac{1-\cos x}{x(\sqrt{1+x}-1)}$.

\item Calcula $\displaystyle \lim_{x\to0}\frac{\sqrt{1+\tan x}-\sqrt{1-\tan x}}{\tan(2x)}$.

\end{enumerate}

\newpage

\section{Otros límites}

\subsection{Nivel 1 - Básico}

\begin{enumerate}[label=\textbf{E\arabic*.},start=601]

\item Calcula $\displaystyle\lim_{x\to -3}(x^2+0x+0)$.

\item Calcula $\displaystyle\lim_{x\to -2}(x^2+1x+1)$.

\item Calcula $\displaystyle\lim_{x\to -1}(x^2+2x+2)$.

\item Calcula $\displaystyle\lim_{x\to 0}(x^2+3x+3)$.

\item Calcula $\displaystyle\lim_{x\to 1}(x^2+4x+0)$.

\item Calcula $\displaystyle\lim_{x\to 2}(x^2+0x+1)$.

\item Calcula $\displaystyle\lim_{x\to 3}(x^2+1x+2)$.

\item Calcula $\displaystyle\lim_{x\to -3}(x^2+2x+3)$.

\item Calcula $\displaystyle\lim_{x\to -2}(x^2+3x+0)$.

\item Calcula $\displaystyle\lim_{x\to -1}(x^2+4x+1)$.

\item Calcula $\displaystyle\lim_{x\to 0}(x^2+0x+2)$.

\item Calcula $\displaystyle\lim_{x\to 1}(x^2+1x+3)$.

\item Calcula $\displaystyle\lim_{x\to 2}(x^2+2x+0)$.

\item Calcula $\displaystyle\lim_{x\to 3}(x^2+3x+1)$.

\item Calcula $\displaystyle\lim_{x\to -3}(x^2+4x+2)$.

\item Calcula $\displaystyle\lim_{x\to -2}(x^2+0x+3)$.

\item Calcula $\displaystyle\lim_{x\to -1}(x^2+1x+0)$.

\item Calcula $\displaystyle\lim_{x\to 0}(x^2+2x+1)$.

\item Calcula $\displaystyle\lim_{x\to 1}(x^2+3x+2)$.

\item Calcula $\displaystyle\lim_{x\to 2}(x^2+4x+3)$.

\end{enumerate}

\subsection{Nivel 2 - Intermedio bajo}

\begin{enumerate}[label=\textbf{E\arabic*.},start=621]

\item Calcula $\displaystyle\lim_{x\to -3}\frac{|x--3|}{x--3}$. 

\item Calcula $\displaystyle\lim_{x\to -2}\frac{|x--2|}{x--2}$. 

\item Calcula $\displaystyle\lim_{x\to -1}\frac{|x--1|}{x--1}$. 

\item Calcula $\displaystyle\lim_{x\to 2}\frac{|x-2|}{x-2}$. 

\item Calcula $\displaystyle\lim_{x\to 1}\frac{|x-1|}{x-1}$. 

\item Calcula $\displaystyle\lim_{x\to 2}\frac{|x-2|}{x-2}$. 

\item Calcula $\displaystyle\lim_{x\to 3}\frac{|x-3|}{x-3}$. 

\item Calcula $\displaystyle\lim_{x\to -3}\frac{|x--3|}{x--3}$. 

\item Calcula $\displaystyle\lim_{x\to -2}\frac{|x--2|}{x--2}$. 

\item Calcula $\displaystyle\lim_{x\to -1}\frac{|x--1|}{x--1}$. 

\item Calcula $\displaystyle\lim_{x\to 2}\frac{|x-2|}{x-2}$. 

\item Calcula $\displaystyle\lim_{x\to 1}\frac{|x-1|}{x-1}$. 

\item Calcula $\displaystyle\lim_{x\to 2}\frac{|x-2|}{x-2}$. 

\item Calcula $\displaystyle\lim_{x\to 3}\frac{|x-3|}{x-3}$. 

\item Calcula $\displaystyle\lim_{x\to -3}\frac{|x--3|}{x--3}$. 

\item Calcula $\displaystyle\lim_{x\to -2}\frac{|x--2|}{x--2}$. 

\item Calcula $\displaystyle\lim_{x\to -1}\frac{|x--1|}{x--1}$. 

\item Calcula $\displaystyle\lim_{x\to 2}\frac{|x-2|}{x-2}$. 

\item Calcula $\displaystyle\lim_{x\to 1}\frac{|x-1|}{x-1}$. 

\item Calcula $\displaystyle\lim_{x\to 2}\frac{|x-2|}{x-2}$. 

\end{enumerate}

\newpage

\subsection{Nivel 3 - Intermedio}

\begin{enumerate}[label=\textbf{E\arabic*.},start=641]

\item Calcula $\displaystyle\lim_{x\to0}x^2\sin\left(\frac{1}{x}\right)$.

\item Calcula $\displaystyle\lim_{x\to0}x^2\sin\left(\frac{2}{x}\right)$.

\item Calcula $\displaystyle\lim_{x\to0}x^2\sin\left(\frac{3}{x}\right)$.

\item Calcula $\displaystyle\lim_{x\to0}x^2\sin\left(\frac{4}{x}\right)$.

\item Calcula $\displaystyle\lim_{x\to0}x^2\sin\left(\frac{5}{x}\right)$.

\item Calcula $\displaystyle\lim_{x\to0}x^2\sin\left(\frac{6}{x}\right)$.

\item Calcula $\displaystyle\lim_{x\to0}x^2\sin\left(\frac{7}{x}\right)$.

\item Calcula $\displaystyle\lim_{x\to0}x^2\sin\left(\frac{1}{x}\right)$.

\item Calcula $\displaystyle\lim_{x\to0}x^2\sin\left(\frac{2}{x}\right)$.

\item Calcula $\displaystyle\lim_{x\to0}x^2\sin\left(\frac{3}{x}\right)$.

\item Calcula $\displaystyle\lim_{x\to0}x^2\sin\left(\frac{4}{x}\right)$.

\item Calcula $\displaystyle\lim_{x\to0}x^2\sin\left(\frac{5}{x}\right)$.

\item Calcula $\displaystyle\lim_{x\to0}x^2\sin\left(\frac{6}{x}\right)$.

\item Calcula $\displaystyle\lim_{x\to0}x^2\sin\left(\frac{7}{x}\right)$.

\item Calcula $\displaystyle\lim_{x\to0}x^2\sin\left(\frac{1}{x}\right)$.

\item Calcula $\displaystyle\lim_{x\to0}x^2\sin\left(\frac{2}{x}\right)$.

\item Calcula $\displaystyle\lim_{x\to0}x^2\sin\left(\frac{3}{x}\right)$.

\item Calcula $\displaystyle\lim_{x\to0}x^2\sin\left(\frac{4}{x}\right)$.

\item Calcula $\displaystyle\lim_{x\to0}x^2\sin\left(\frac{5}{x}\right)$.

\item Calcula $\displaystyle\lim_{x\to0}x^2\sin\left(\frac{6}{x}\right)$.

\end{enumerate}

\subsection{Nivel 4 - Avanzado}

\begin{enumerate}[label=\textbf{E\arabic*.},start=661]

\item Calcula $\displaystyle\lim_{x\to0}\frac{\ln(1+x)}{x}$.

\item Calcula $\displaystyle\lim_{x\to0}\frac{e^x-1-\sin x}{x^2}$.

\item Calcula $\displaystyle\lim_{x\to0}\frac{\ln(1+x)}{x}$.

\item Calcula $\displaystyle\lim_{x\to0}\frac{e^x-1-\sin x}{x^2}$.

\item Calcula $\displaystyle\lim_{x\to0}\frac{\ln(1+x)}{x}$.

\item Calcula $\displaystyle\lim_{x\to0}\frac{e^x-1-\sin x}{x^2}$.

\item Calcula $\displaystyle\lim_{x\to0}\frac{\ln(1+x)}{x}$.

\item Calcula $\displaystyle\lim_{x\to0}\frac{e^x-1-\sin x}{x^2}$.

\item Calcula $\displaystyle\lim_{x\to0}\frac{\ln(1+x)}{x}$.

\item Calcula $\displaystyle\lim_{x\to0}\frac{e^x-1-\sin x}{x^2}$.

\item Calcula $\displaystyle\lim_{x\to0}\frac{\ln(1+x)}{x}$.

\item Calcula $\displaystyle\lim_{x\to0}\frac{e^x-1-\sin x}{x^2}$.

\item Calcula $\displaystyle\lim_{x\to0}\frac{\ln(1+x)}{x}$.

\item Calcula $\displaystyle\lim_{x\to0}\frac{e^x-1-\sin x}{x^2}$.

\item Calcula $\displaystyle\lim_{x\to0}\frac{\ln(1+x)}{x}$.

\item Calcula $\displaystyle\lim_{x\to0}\frac{e^x-1-\sin x}{x^2}$.

\item Calcula $\displaystyle\lim_{x\to0}\frac{\ln(1+x)}{x}$.

\item Calcula $\displaystyle\lim_{x\to0}\frac{e^x-1-\sin x}{x^2}$.

\item Calcula $\displaystyle\lim_{x\to0}\frac{\ln(1+x)}{x}$.

\item Calcula $\displaystyle\lim_{x\to0}\frac{e^x-1-\sin x}{x^2}$.

\end{enumerate}

\newpage

\subsection{Nivel 5 - Imposible}

\begin{enumerate}[label=\textbf{E\arabic*.},start=681]

\item Calcula $\displaystyle \lim_{x\to0}\frac{\ln(1+x)-x+\frac{x^2}{2}}{x^3}$.

\item Calcula $\displaystyle \lim_{x\to0}\frac{e^x-\cos x-\sin x}{x^2}$.

\item Calcula $\displaystyle \lim_{x\to0^+}x\ln x$.

\item Calcula $\displaystyle \lim_{x\to0}\frac{|x|}{x}$.

\item Calcula $\displaystyle \lim_{x\to0}\frac{\sqrt[3]{1+x}-1-\frac{x}{3}}{x^2}$.

\item Calcula $\displaystyle \lim_{x\to0}\frac{\ln(1+x)-x+\frac{x^2}{2}}{x^3}$.

\item Calcula $\displaystyle \lim_{x\to0}\frac{e^x-\cos x-\sin x}{x^2}$.

\item Calcula $\displaystyle \lim_{x\to0^+}x\ln x$.

\item Calcula $\displaystyle \lim_{x\to0}\frac{|x|}{x}$.

\item Calcula $\displaystyle \lim_{x\to0}\frac{\sqrt[3]{1+x}-1-\frac{x}{3}}{x^2}$.

\item Calcula $\displaystyle \lim_{x\to0}\frac{\ln(1+x)-x+\frac{x^2}{2}}{x^3}$.

\item Calcula $\displaystyle \lim_{x\to0}\frac{e^x-\cos x-\sin x}{x^2}$.

\item Calcula $\displaystyle \lim_{x\to0^+}x\ln x$.

\item Calcula $\displaystyle \lim_{x\to0}\frac{|x|}{x}$.

\item Calcula $\displaystyle \lim_{x\to0}\frac{\sqrt[3]{1+x}-1-\frac{x}{3}}{x^2}$.

\item Calcula $\displaystyle \lim_{x\to0}\frac{\ln(1+x)-x+\frac{x^2}{2}}{x^3}$.

\item Calcula $\displaystyle \lim_{x\to0}\frac{e^x-\cos x-\sin x}{x^2}$.

\item Calcula $\displaystyle \lim_{x\to0^+}x\ln x$.

\item Calcula $\displaystyle \lim_{x\to0}\frac{|x|}{x}$.

\item Calcula $\displaystyle \lim_{x\to0}\frac{\sqrt[3]{1+x}-1-\frac{x}{3}}{x^2}$.

\end{enumerate}

\newpage

\section{Solucionario desarrollado}

\subsection*{Dominio, dominio de imagen y asíntotas horizontales}

\addcontentsline{toc}{subsection}{Soluciones - Dominio, dominio de imagen y asíntotas horizontales}

\subsubsection*{Nivel 1 - Básico}

\paragraph{Ejercicio 1.} Determina dominio, imagen y asíntotas horizontales de $f(x)=1x-9$.

\[
\begin{aligned}
D_f=\mathbb{R}\\
\operatorname{Im}(f)=\mathbb{R}\\
\text{{AH: no existe}}
\end{aligned}
\]


\paragraph{Ejercicio 2.} Determina dominio, imagen y asíntotas horizontales de $f(x)=(x+2)^2-1$.

\[
\begin{aligned}
D_f=\mathbb{R}\\
(x+2)^2\ge 0\\
f(x)\ge -1\\
\operatorname{Im}(f)=[-1,\infty)\\
\text{{AH: no existe}}
\end{aligned}
\]


\paragraph{Ejercicio 3.} Determina dominio, imagen y asíntotas horizontales de $f(x)=\sqrt{x+2}+0$.

\[
\begin{aligned}
x+2\ge0\\
x\ge -2\\
D_f=[-2,\infty)\\
\sqrt{x+2}\ge0\\
\operatorname{Im}(f)=[0,\infty)\\
\text{{AH: no existe}}
\end{aligned}
\]


\paragraph{Ejercicio 4.} Determina dominio, imagen y asíntota horizontal de $f(x)=\frac{1}{x+1}+1$.

\[
\begin{aligned}
x+1\ne0\\
x\ne -1\\
D_f=\mathbb{R}\setminus\{-1\}\\
\frac{1}{x+1}\ne0\\
\operatorname{Im}(f)=\mathbb{R}\setminus\{1\}\\
\text{AH: }y=1
\end{aligned}
\]


\paragraph{Ejercicio 5.} Determina dominio, imagen y asíntota horizontal de $f(x)=2^{x-2}+3$.

\[
\begin{aligned}
2^{x-2}>0\\
D_f=\mathbb{R}\\
\operatorname{Im}(f)=(3,\infty)\\
\text{AH: }y=3
\end{aligned}
\]


\paragraph{Ejercicio 6.} Determina dominio, imagen y asíntotas horizontales de $f(x)=2x-4$.

\[
\begin{aligned}
D_f=\mathbb{R}\\
\operatorname{Im}(f)=\mathbb{R}\\
\text{{AH: no existe}}
\end{aligned}
\]


\paragraph{Ejercicio 7.} Determina dominio, imagen y asíntotas horizontales de $f(x)=(x-3)^2-2$.

\[
\begin{aligned}
D_f=\mathbb{R}\\
(x-3)^2\ge 0\\
f(x)\ge -2\\
\operatorname{Im}(f)=[-2,\infty)\\
\text{{AH: no existe}}
\end{aligned}
\]


\paragraph{Ejercicio 8.} Determina dominio, imagen y asíntotas horizontales de $f(x)=\sqrt{x-3}+0$.

\[
\begin{aligned}
x-3\ge0\\
x\ge 3\\
D_f=[3,\infty)\\
\sqrt{x-3}\ge0\\
\operatorname{Im}(f)=[0,\infty)\\
\text{{AH: no existe}}
\end{aligned}
\]


\paragraph{Ejercicio 9.} Determina dominio, imagen y asíntota horizontal de $f(x)=\frac{1}{x-4}+1$.

\[
\begin{aligned}
x-4\ne0\\
x\ne 4\\
D_f=\mathbb{R}\setminus\{4\}\\
\frac{1}{x-4}\ne0\\
\operatorname{Im}(f)=\mathbb{R}\setminus\{1\}\\
\text{AH: }y=1
\end{aligned}
\]


\paragraph{Ejercicio 10.} Determina dominio, imagen y asíntota horizontal de $f(x)=2^{x+0}+3$.

\[
\begin{aligned}
2^{x+0}>0\\
D_f=\mathbb{R}\\
\operatorname{Im}(f)=(3,\infty)\\
\text{AH: }y=3
\end{aligned}
\]


\paragraph{Ejercicio 11.} Determina dominio, imagen y asíntotas horizontales de $f(x)=3x+1$.

\[
\begin{aligned}
D_f=\mathbb{R}\\
\operatorname{Im}(f)=\mathbb{R}\\
\text{{AH: no existe}}
\end{aligned}
\]


\paragraph{Ejercicio 12.} Determina dominio, imagen y asíntotas horizontales de $f(x)=(x-1)^2+3$.

\[
\begin{aligned}
D_f=\mathbb{R}\\
(x-1)^2\ge 0\\
f(x)\ge 3\\
\operatorname{Im}(f)=[3,\infty)\\
\text{{AH: no existe}}
\end{aligned}
\]


\paragraph{Ejercicio 13.} Determina dominio, imagen y asíntotas horizontales de $f(x)=\sqrt{x+0}+0$.

\[
\begin{aligned}
x+0\ge0\\
x\ge 0\\
D_f=[0,\infty)\\
\sqrt{x+0}\ge0\\
\operatorname{Im}(f)=[0,\infty)\\
\text{{AH: no existe}}
\end{aligned}
\]


\paragraph{Ejercicio 14.} Determina dominio, imagen y asíntota horizontal de $f(x)=\frac{1}{x+0}+1$.

\[
\begin{aligned}
x+0\ne0\\
x\ne 0\\
D_f=\mathbb{R}\setminus\{0\}\\
\frac{1}{x+0}\ne0\\
\operatorname{Im}(f)=\mathbb{R}\setminus\{1\}\\
\text{AH: }y=1
\end{aligned}
\]


\paragraph{Ejercicio 15.} Determina dominio, imagen y asíntota horizontal de $f(x)=2^{x+2}+3$.

\[
\begin{aligned}
2^{x+2}>0\\
D_f=\mathbb{R}\\
\operatorname{Im}(f)=(3,\infty)\\
\text{AH: }y=3
\end{aligned}
\]


\paragraph{Ejercicio 16.} Determina dominio, imagen y asíntotas horizontales de $f(x)=4x+6$.

\[
\begin{aligned}
D_f=\mathbb{R}\\
\operatorname{Im}(f)=\mathbb{R}\\
\text{{AH: no existe}}
\end{aligned}
\]


\paragraph{Ejercicio 17.} Determina dominio, imagen y asíntotas horizontales de $f(x)=(x+1)^2+2$.

\[
\begin{aligned}
D_f=\mathbb{R}\\
(x+1)^2\ge 0\\
f(x)\ge 2\\
\operatorname{Im}(f)=[2,\infty)\\
\text{{AH: no existe}}
\end{aligned}
\]


\paragraph{Ejercicio 18.} Determina dominio, imagen y asíntotas horizontales de $f(x)=\sqrt{x+3}+0$.

\[
\begin{aligned}
x+3\ge0\\
x\ge -3\\
D_f=[-3,\infty)\\
\sqrt{x+3}\ge0\\
\operatorname{Im}(f)=[0,\infty)\\
\text{{AH: no existe}}
\end{aligned}
\]


\paragraph{Ejercicio 19.} Determina dominio, imagen y asíntota horizontal de $f(x)=\frac{1}{x+4}+1$.

\[
\begin{aligned}
x+4\ne0\\
x\ne -4\\
D_f=\mathbb{R}\setminus\{-4\}\\
\frac{1}{x+4}\ne0\\
\operatorname{Im}(f)=\mathbb{R}\setminus\{1\}\\
\text{AH: }y=1
\end{aligned}
\]


\paragraph{Ejercicio 20.} Determina dominio, imagen y asíntota horizontal de $f(x)=2^{x-3}+3$.

\[
\begin{aligned}
2^{x-3}>0\\
D_f=\mathbb{R}\\
\operatorname{Im}(f)=(3,\infty)\\
\text{AH: }y=3
\end{aligned}
\]


\newpage

\subsubsection*{Nivel 2 - Intermedio bajo}

\paragraph{Ejercicio 21.} Determina dominio, imagen y asíntota horizontal de $f(x)=\frac{1x-3}{x+1}$.

\[
\begin{aligned}
x+1\ne0\\
x\ne -1\\
D_f=\mathbb{R}\setminus\{-1\}\\
y=\frac{1x-3}{x+1}\\
yx+1y=1x-3\\
x(y-1)=-3-1y\\
y\ne 1\\
\operatorname{Im}(f)=\mathbb{R}\setminus\{1\}\\
\text{AH: }y=1
\end{aligned}
\]


\paragraph{Ejercicio 22.} Determina dominio, imagen y asíntotas horizontales de $f(x)=\ln(x+1)$.

\[
\begin{aligned}
x+1>0\\
x>-1\\
D_f=(-1,\infty)\\
\operatorname{Im}(f)=\mathbb{R}\\
\text{{AH: no existe}}
\end{aligned}
\]


\paragraph{Ejercicio 23.} Determina dominio, imagen y asíntota horizontal de $f(x)=-3e^{x+0}-1$.

\[
\begin{aligned}
e^{x+0}>0\\
D_f=\mathbb{R}\\
-3e^{x+0}<0\\
\operatorname{Im}(f)=(-\infty,-1)\\
\text{AH: }y=-1
\end{aligned}
\]


\paragraph{Ejercicio 24.} Determina dominio, imagen y asíntota horizontal de $f(x)=\frac{4}{(x-2)^2}+1$.

\[
\begin{aligned}
(x-2)^2\ne0\\
x\ne 2\\
D_f=\mathbb{R}\setminus\{2\}\\
\frac{4}{(x-2)^2}>0\\
\operatorname{Im}(f)=(1,\infty)\\
\text{AH: }y=1
\end{aligned}
\]


\paragraph{Ejercicio 25.} Determina dominio, imagen y asíntota horizontal de $f(x)=\ln\left(\frac{x-5}{x-8}\right)$.

\[
\begin{aligned}
\frac{x-5}{x-8} >0\\
D_f=(-\infty,5)\cup(8,\infty)\\
\frac{x-5}{x-8}\ne1\\
\ln(1)=0\\
\operatorname{Im}(f)=\mathbb{R}\setminus\{0\}\\
\text{{AH: }}y=0
\end{aligned}
\]


\paragraph{Ejercicio 26.} Determina dominio, imagen y asíntota horizontal de $f(x)=\frac{1x+2}{x+6}$.

\[
\begin{aligned}
x+6\ne0\\
x\ne -6\\
D_f=\mathbb{R}\setminus\{-6\}\\
y=\frac{1x+2}{x+6}\\
yx+6y=1x+2\\
x(y-1)=2-6y\\
y\ne 1\\
\operatorname{Im}(f)=\mathbb{R}\setminus\{1\}\\
\text{AH: }y=1
\end{aligned}
\]


\paragraph{Ejercicio 27.} Determina dominio, imagen y asíntotas horizontales de $f(x)=\ln(x+2)$.

\[
\begin{aligned}
x+2>0\\
x>-2\\
D_f=(-2,\infty)\\
\operatorname{Im}(f)=\mathbb{R}\\
\text{{AH: no existe}}
\end{aligned}
\]


\paragraph{Ejercicio 28.} Determina dominio, imagen y asíntota horizontal de $f(x)=-4e^{x+0}-3$.

\[
\begin{aligned}
e^{x+0}>0\\
D_f=\mathbb{R}\\
-4e^{x+0}<0\\
\operatorname{Im}(f)=(-\infty,-3)\\
\text{AH: }y=-3
\end{aligned}
\]


\paragraph{Ejercicio 29.} Determina dominio, imagen y asíntota horizontal de $f(x)=\frac{4}{(x+1)^2}+0$.

\[
\begin{aligned}
(x+1)^2\ne0\\
x\ne -1\\
D_f=\mathbb{R}\setminus\{-1\}\\
\frac{4}{(x+1)^2}>0\\
\operatorname{Im}(f)=(0,\infty)\\
\text{AH: }y=0
\end{aligned}
\]


\paragraph{Ejercicio 30.} Determina dominio, imagen y asíntota horizontal de $f(x)=\ln\left(\frac{x-5}{x-8}\right)$.

\[
\begin{aligned}
\frac{x-5}{x-8} >0\\
D_f=(-\infty,5)\cup(8,\infty)\\
\frac{x-5}{x-8}\ne1\\
\ln(1)=0\\
\operatorname{Im}(f)=\mathbb{R}\setminus\{0\}\\
\text{{AH: }}y=0
\end{aligned}
\]


\paragraph{Ejercicio 31.} Determina dominio, imagen y asíntota horizontal de $f(x)=\frac{1x+0}{x+5}$.

\[
\begin{aligned}
x+5\ne0\\
x\ne -5\\
D_f=\mathbb{R}\setminus\{-5\}\\
y=\frac{1x+0}{x+5}\\
yx+5y=1x+0\\
x(y-1)=0-5y\\
y\ne 1\\
\operatorname{Im}(f)=\mathbb{R}\setminus\{1\}\\
\text{AH: }y=1
\end{aligned}
\]


\paragraph{Ejercicio 32.} Determina dominio, imagen y asíntotas horizontales de $f(x)=\ln(x-3)$.

\[
\begin{aligned}
x-3>0\\
x>3\\
D_f=(3,\infty)\\
\operatorname{Im}(f)=\mathbb{R}\\
\text{{AH: no existe}}
\end{aligned}
\]


\paragraph{Ejercicio 33.} Determina dominio, imagen y asíntota horizontal de $f(x)=-1e^{x+0}+2$.

\[
\begin{aligned}
e^{x+0}>0\\
D_f=\mathbb{R}\\
-1e^{x+0}<0\\
\operatorname{Im}(f)=(-\infty,2)\\
\text{AH: }y=2
\end{aligned}
\]


\paragraph{Ejercicio 34.} Determina dominio, imagen y asíntota horizontal de $f(x)=\frac{4}{(x+0)^2}-1$.

\[
\begin{aligned}
(x+0)^2\ne0\\
x\ne 0\\
D_f=\mathbb{R}\setminus\{0\}\\
\frac{4}{(x+0)^2}>0\\
\operatorname{Im}(f)=(-1,\infty)\\
\text{AH: }y=-1
\end{aligned}
\]


\paragraph{Ejercicio 35.} Determina dominio, imagen y asíntota horizontal de $f(x)=\ln\left(\frac{x-5}{x-8}\right)$.

\[
\begin{aligned}
\frac{x-5}{x-8} >0\\
D_f=(-\infty,5)\cup(8,\infty)\\
\frac{x-5}{x-8}\ne1\\
\ln(1)=0\\
\operatorname{Im}(f)=\mathbb{R}\setminus\{0\}\\
\text{{AH: }}y=0
\end{aligned}
\]


\paragraph{Ejercicio 36.} Determina dominio, imagen y asíntota horizontal de $f(x)=\frac{1x-2}{x+4}$.

\[
\begin{aligned}
x+4\ne0\\
x\ne -4\\
D_f=\mathbb{R}\setminus\{-4\}\\
y=\frac{1x-2}{x+4}\\
yx+4y=1x-2\\
x(y-1)=-2-4y\\
y\ne 1\\
\operatorname{Im}(f)=\mathbb{R}\setminus\{1\}\\
\text{AH: }y=1
\end{aligned}
\]


\paragraph{Ejercicio 37.} Determina dominio, imagen y asíntotas horizontales de $f(x)=\ln(x-2)$.

\[
\begin{aligned}
x-2>0\\
x>2\\
D_f=(2,\infty)\\
\operatorname{Im}(f)=\mathbb{R}\\
\text{{AH: no existe}}
\end{aligned}
\]


\paragraph{Ejercicio 38.} Determina dominio, imagen y asíntota horizontal de $f(x)=-2e^{x+0}+0$.

\[
\begin{aligned}
e^{x+0}>0\\
D_f=\mathbb{R}\\
-2e^{x+0}<0\\
\operatorname{Im}(f)=(-\infty,0)\\
\text{AH: }y=0
\end{aligned}
\]


\paragraph{Ejercicio 39.} Determina dominio, imagen y asíntota horizontal de $f(x)=\frac{4}{(x-1)^2}-2$.

\[
\begin{aligned}
(x-1)^2\ne0\\
x\ne 1\\
D_f=\mathbb{R}\setminus\{1\}\\
\frac{4}{(x-1)^2}>0\\
\operatorname{Im}(f)=(-2,\infty)\\
\text{AH: }y=-2
\end{aligned}
\]


\paragraph{Ejercicio 40.} Determina dominio, imagen y asíntota horizontal de $f(x)=\ln\left(\frac{x-5}{x-8}\right)$.

\[
\begin{aligned}
\frac{x-5}{x-8} >0\\
D_f=(-\infty,5)\cup(8,\infty)\\
\frac{x-5}{x-8}\ne1\\
\ln(1)=0\\
\operatorname{Im}(f)=\mathbb{R}\setminus\{0\}\\
\text{{AH: }}y=0
\end{aligned}
\]


\newpage

\subsubsection*{Nivel 3 - Intermedio}

\paragraph{Ejercicio 41.} Determina dominio, imagen y asíntota horizontal de $f(x)=\sqrt{\frac{x-1}{x-3}}-1$.

\[
\begin{aligned}
\frac{x-1}{x-3}\ge0\\
D_f=(-\infty,1]\cup(3,\infty)\\
\frac{x-1}{x-3}\ne1\\
\sqrt{\frac{x-1}{x-3}}\ne1\\
\operatorname{Im}(f)=[-1,\infty)\setminus\{0\}\\
\text{AH: }y=0
\end{aligned}
\]


\paragraph{Ejercicio 42.} Determina dominio, imagen y asíntota horizontal de $f(x)=\frac{2^x+2}{2^x+4}$. 

\[
\begin{aligned}
t=2^x>0\\
y=\frac{t+2}{t+4}\\
y(t+4)=t+2\\
t(y-1)=2-4y\\
t=\frac{2-4y}{y-1}>0\\
\operatorname{Im}(f)=(\frac{2}{4},1)\\
D_f=\mathbb{R}\\
\text{{AH: }}y=1
\end{aligned}
\]


\paragraph{Ejercicio 43.} Determina dominio, imagen y asíntota horizontal de $f(x)=\ln\left(\frac{x^2-9}{x^2-25}\right)$.

\[
\begin{aligned}
\frac{x^2-9}{x^2-25} >0\\
D_f=(-\infty,-5)\cup(-3,3)\cup(5,\infty)\\
\frac{x^2-9}{x^2-25}\ne1\\
\ln(1)=0\\
\operatorname{Im}(f)=\mathbb{R}\setminus\{0\}\\
\text{{AH: }}y=0
\end{aligned}
\]


\paragraph{Ejercicio 44.} Determina dominio, imagen y asíntotas horizontales de $f(x)=\arctan(x-2)$.

\[
\begin{aligned}
D_f=\mathbb{R}\\
-\frac{\pi}{2}<\arctan(x-2)<\frac{\pi}{2}\\
\operatorname{Im}(f)=\left(-\frac{\pi}{2},\frac{\pi}{2}\right)\\
\text{AH: }y=\frac{\pi}{2}\text{ y }y=-\frac{\pi}{2}
\end{aligned}
\]


\paragraph{Ejercicio 45.} Determina dominio, imagen y asíntotas horizontales de $f(x)=\frac{e^x-e^{-x}}{e^x+e^{-x}}+-1$.

\[
\begin{aligned}
D_f=\mathbb{R}\\
t=e^x>0\\
f(x)=\frac{t^2-1}{t^2+1}+-1\\
-1<\frac{t^2-1}{t^2+1}<1\\
\operatorname{Im}(f)=(-2,0)\\
\text{AH: }y=-2\text{ y }y=0
\end{aligned}
\]


\paragraph{Ejercicio 46.} Determina dominio, imagen y asíntota horizontal de $f(x)=\sqrt{\frac{x-1}{x-3}}+0$.

\[
\begin{aligned}
\frac{x-1}{x-3}\ge0\\
D_f=(-\infty,1]\cup(3,\infty)\\
\frac{x-1}{x-3}\ne1\\
\sqrt{\frac{x-1}{x-3}}\ne1\\
\operatorname{Im}(f)=[0,\infty)\setminus\{1\}\\
\text{AH: }y=1
\end{aligned}
\]


\paragraph{Ejercicio 47.} Determina dominio, imagen y asíntota horizontal de $f(x)=\frac{2^x+2}{2^x+4}$. 

\[
\begin{aligned}
t=2^x>0\\
y=\frac{t+2}{t+4}\\
y(t+4)=t+2\\
t(y-1)=2-4y\\
t=\frac{2-4y}{y-1}>0\\
\operatorname{Im}(f)=(\frac{2}{4},1)\\
D_f=\mathbb{R}\\
\text{{AH: }}y=1
\end{aligned}
\]


\paragraph{Ejercicio 48.} Determina dominio, imagen y asíntota horizontal de $f(x)=\ln\left(\frac{x^2-9}{x^2-25}\right)$.

\[
\begin{aligned}
\frac{x^2-9}{x^2-25} >0\\
D_f=(-\infty,-5)\cup(-3,3)\cup(5,\infty)\\
\frac{x^2-9}{x^2-25}\ne1\\
\ln(1)=0\\
\operatorname{Im}(f)=\mathbb{R}\setminus\{0\}\\
\text{{AH: }}y=0
\end{aligned}
\]


\paragraph{Ejercicio 49.} Determina dominio, imagen y asíntotas horizontales de $f(x)=\arctan(x+1)$.

\[
\begin{aligned}
D_f=\mathbb{R}\\
-\frac{\pi}{2}<\arctan(x+1)<\frac{\pi}{2}\\
\operatorname{Im}(f)=\left(-\frac{\pi}{2},\frac{\pi}{2}\right)\\
\text{AH: }y=\frac{\pi}{2}\text{ y }y=-\frac{\pi}{2}
\end{aligned}
\]


\paragraph{Ejercicio 50.} Determina dominio, imagen y asíntotas horizontales de $f(x)=\frac{e^x-e^{-x}}{e^x+e^{-x}}+0$.

\[
\begin{aligned}
D_f=\mathbb{R}\\
t=e^x>0\\
f(x)=\frac{t^2-1}{t^2+1}+0\\
-1<\frac{t^2-1}{t^2+1}<1\\
\operatorname{Im}(f)=(-1,1)\\
\text{AH: }y=-1\text{ y }y=1
\end{aligned}
\]


\paragraph{Ejercicio 51.} Determina dominio, imagen y asíntota horizontal de $f(x)=\sqrt{\frac{x-1}{x-3}}+1$.

\[
\begin{aligned}
\frac{x-1}{x-3}\ge0\\
D_f=(-\infty,1]\cup(3,\infty)\\
\frac{x-1}{x-3}\ne1\\
\sqrt{\frac{x-1}{x-3}}\ne1\\
\operatorname{Im}(f)=[1,\infty)\setminus\{2\}\\
\text{AH: }y=2
\end{aligned}
\]


\paragraph{Ejercicio 52.} Determina dominio, imagen y asíntota horizontal de $f(x)=\frac{2^x+2}{2^x+4}$. 

\[
\begin{aligned}
t=2^x>0\\
y=\frac{t+2}{t+4}\\
y(t+4)=t+2\\
t(y-1)=2-4y\\
t=\frac{2-4y}{y-1}>0\\
\operatorname{Im}(f)=(\frac{2}{4},1)\\
D_f=\mathbb{R}\\
\text{{AH: }}y=1
\end{aligned}
\]


\paragraph{Ejercicio 53.} Determina dominio, imagen y asíntota horizontal de $f(x)=\ln\left(\frac{x^2-9}{x^2-25}\right)$.

\[
\begin{aligned}
\frac{x^2-9}{x^2-25} >0\\
D_f=(-\infty,-5)\cup(-3,3)\cup(5,\infty)\\
\frac{x^2-9}{x^2-25}\ne1\\
\ln(1)=0\\
\operatorname{Im}(f)=\mathbb{R}\setminus\{0\}\\
\text{{AH: }}y=0
\end{aligned}
\]


\paragraph{Ejercicio 54.} Determina dominio, imagen y asíntotas horizontales de $f(x)=\arctan(x+0)$.

\[
\begin{aligned}
D_f=\mathbb{R}\\
-\frac{\pi}{2}<\arctan(x+0)<\frac{\pi}{2}\\
\operatorname{Im}(f)=\left(-\frac{\pi}{2},\frac{\pi}{2}\right)\\
\text{AH: }y=\frac{\pi}{2}\text{ y }y=-\frac{\pi}{2}
\end{aligned}
\]


\paragraph{Ejercicio 55.} Determina dominio, imagen y asíntotas horizontales de $f(x)=\frac{e^x-e^{-x}}{e^x+e^{-x}}+1$.

\[
\begin{aligned}
D_f=\mathbb{R}\\
t=e^x>0\\
f(x)=\frac{t^2-1}{t^2+1}+1\\
-1<\frac{t^2-1}{t^2+1}<1\\
\operatorname{Im}(f)=(0,2)\\
\text{AH: }y=0\text{ y }y=2
\end{aligned}
\]


\paragraph{Ejercicio 56.} Determina dominio, imagen y asíntota horizontal de $f(x)=\sqrt{\frac{x-1}{x-3}}+2$.

\[
\begin{aligned}
\frac{x-1}{x-3}\ge0\\
D_f=(-\infty,1]\cup(3,\infty)\\
\frac{x-1}{x-3}\ne1\\
\sqrt{\frac{x-1}{x-3}}\ne1\\
\operatorname{Im}(f)=[2,\infty)\setminus\{3\}\\
\text{AH: }y=3
\end{aligned}
\]


\paragraph{Ejercicio 57.} Determina dominio, imagen y asíntota horizontal de $f(x)=\frac{2^x+2}{2^x+4}$. 

\[
\begin{aligned}
t=2^x>0\\
y=\frac{t+2}{t+4}\\
y(t+4)=t+2\\
t(y-1)=2-4y\\
t=\frac{2-4y}{y-1}>0\\
\operatorname{Im}(f)=(\frac{2}{4},1)\\
D_f=\mathbb{R}\\
\text{{AH: }}y=1
\end{aligned}
\]


\paragraph{Ejercicio 58.} Determina dominio, imagen y asíntota horizontal de $f(x)=\ln\left(\frac{x^2-9}{x^2-25}\right)$.

\[
\begin{aligned}
\frac{x^2-9}{x^2-25} >0\\
D_f=(-\infty,-5)\cup(-3,3)\cup(5,\infty)\\
\frac{x^2-9}{x^2-25}\ne1\\
\ln(1)=0\\
\operatorname{Im}(f)=\mathbb{R}\setminus\{0\}\\
\text{{AH: }}y=0
\end{aligned}
\]


\paragraph{Ejercicio 59.} Determina dominio, imagen y asíntotas horizontales de $f(x)=\arctan(x-1)$.

\[
\begin{aligned}
D_f=\mathbb{R}\\
-\frac{\pi}{2}<\arctan(x-1)<\frac{\pi}{2}\\
\operatorname{Im}(f)=\left(-\frac{\pi}{2},\frac{\pi}{2}\right)\\
\text{AH: }y=\frac{\pi}{2}\text{ y }y=-\frac{\pi}{2}
\end{aligned}
\]


\paragraph{Ejercicio 60.} Determina dominio, imagen y asíntotas horizontales de $f(x)=\frac{e^x-e^{-x}}{e^x+e^{-x}}+2$.

\[
\begin{aligned}
D_f=\mathbb{R}\\
t=e^x>0\\
f(x)=\frac{t^2-1}{t^2+1}+2\\
-1<\frac{t^2-1}{t^2+1}<1\\
\operatorname{Im}(f)=(1,3)\\
\text{AH: }y=1\text{ y }y=3
\end{aligned}
\]


\newpage

\subsubsection*{Nivel 4 - Avanzado}

\paragraph{Ejercicio 61.} Determina dominio, imagen y asíntota horizontal de $f(x)=\sqrt{\frac{x-1}{x-3}}-1$.

\[
\begin{aligned}
\frac{x-1}{x-3}\ge0\\
D_f=(-\infty,1]\cup(3,\infty)\\
\frac{x-1}{x-3}\ne1\\
\sqrt{\frac{x-1}{x-3}}\ne1\\
\operatorname{Im}(f)=[-1,\infty)\setminus\{0\}\\
\text{AH: }y=0
\end{aligned}
\]


\paragraph{Ejercicio 62.} Determina dominio, imagen y asíntota horizontal de $f(x)=\frac{2^x+2}{2^x+4}$. 

\[
\begin{aligned}
t=2^x>0\\
y=\frac{t+2}{t+4}\\
y(t+4)=t+2\\
t(y-1)=2-4y\\
t=\frac{2-4y}{y-1}>0\\
\operatorname{Im}(f)=(\frac{2}{4},1)\\
D_f=\mathbb{R}\\
\text{{AH: }}y=1
\end{aligned}
\]


\paragraph{Ejercicio 63.} Determina dominio, imagen y asíntota horizontal de $f(x)=\ln\left(\frac{x^2-9}{x^2-25}\right)$.

\[
\begin{aligned}
\frac{x^2-9}{x^2-25} >0\\
D_f=(-\infty,-5)\cup(-3,3)\cup(5,\infty)\\
\frac{x^2-9}{x^2-25}\ne1\\
\ln(1)=0\\
\operatorname{Im}(f)=\mathbb{R}\setminus\{0\}\\
\text{{AH: }}y=0
\end{aligned}
\]


\paragraph{Ejercicio 64.} Determina dominio, imagen y asíntotas horizontales de $f(x)=\arctan(x-2)$.

\[
\begin{aligned}
D_f=\mathbb{R}\\
-\frac{\pi}{2}<\arctan(x-2)<\frac{\pi}{2}\\
\operatorname{Im}(f)=\left(-\frac{\pi}{2},\frac{\pi}{2}\right)\\
\text{AH: }y=\frac{\pi}{2}\text{ y }y=-\frac{\pi}{2}
\end{aligned}
\]


\paragraph{Ejercicio 65.} Determina dominio, imagen y asíntotas horizontales de $f(x)=\frac{e^x-e^{-x}}{e^x+e^{-x}}+-1$.

\[
\begin{aligned}
D_f=\mathbb{R}\\
t=e^x>0\\
f(x)=\frac{t^2-1}{t^2+1}+-1\\
-1<\frac{t^2-1}{t^2+1}<1\\
\operatorname{Im}(f)=(-2,0)\\
\text{AH: }y=-2\text{ y }y=0
\end{aligned}
\]


\paragraph{Ejercicio 66.} Determina dominio, imagen y asíntota horizontal de $f(x)=\sqrt{\frac{x-1}{x-3}}+0$.

\[
\begin{aligned}
\frac{x-1}{x-3}\ge0\\
D_f=(-\infty,1]\cup(3,\infty)\\
\frac{x-1}{x-3}\ne1\\
\sqrt{\frac{x-1}{x-3}}\ne1\\
\operatorname{Im}(f)=[0,\infty)\setminus\{1\}\\
\text{AH: }y=1
\end{aligned}
\]


\paragraph{Ejercicio 67.} Determina dominio, imagen y asíntota horizontal de $f(x)=\frac{2^x+2}{2^x+4}$. 

\[
\begin{aligned}
t=2^x>0\\
y=\frac{t+2}{t+4}\\
y(t+4)=t+2\\
t(y-1)=2-4y\\
t=\frac{2-4y}{y-1}>0\\
\operatorname{Im}(f)=(\frac{2}{4},1)\\
D_f=\mathbb{R}\\
\text{{AH: }}y=1
\end{aligned}
\]


\paragraph{Ejercicio 68.} Determina dominio, imagen y asíntota horizontal de $f(x)=\ln\left(\frac{x^2-9}{x^2-25}\right)$.

\[
\begin{aligned}
\frac{x^2-9}{x^2-25} >0\\
D_f=(-\infty,-5)\cup(-3,3)\cup(5,\infty)\\
\frac{x^2-9}{x^2-25}\ne1\\
\ln(1)=0\\
\operatorname{Im}(f)=\mathbb{R}\setminus\{0\}\\
\text{{AH: }}y=0
\end{aligned}
\]


\paragraph{Ejercicio 69.} Determina dominio, imagen y asíntotas horizontales de $f(x)=\arctan(x+1)$.

\[
\begin{aligned}
D_f=\mathbb{R}\\
-\frac{\pi}{2}<\arctan(x+1)<\frac{\pi}{2}\\
\operatorname{Im}(f)=\left(-\frac{\pi}{2},\frac{\pi}{2}\right)\\
\text{AH: }y=\frac{\pi}{2}\text{ y }y=-\frac{\pi}{2}
\end{aligned}
\]


\paragraph{Ejercicio 70.} Determina dominio, imagen y asíntotas horizontales de $f(x)=\frac{e^x-e^{-x}}{e^x+e^{-x}}+0$.

\[
\begin{aligned}
D_f=\mathbb{R}\\
t=e^x>0\\
f(x)=\frac{t^2-1}{t^2+1}+0\\
-1<\frac{t^2-1}{t^2+1}<1\\
\operatorname{Im}(f)=(-1,1)\\
\text{AH: }y=-1\text{ y }y=1
\end{aligned}
\]


\paragraph{Ejercicio 71.} Determina dominio, imagen y asíntota horizontal de $f(x)=\sqrt{\frac{x-1}{x-3}}+1$.

\[
\begin{aligned}
\frac{x-1}{x-3}\ge0\\
D_f=(-\infty,1]\cup(3,\infty)\\
\frac{x-1}{x-3}\ne1\\
\sqrt{\frac{x-1}{x-3}}\ne1\\
\operatorname{Im}(f)=[1,\infty)\setminus\{2\}\\
\text{AH: }y=2
\end{aligned}
\]


\paragraph{Ejercicio 72.} Determina dominio, imagen y asíntota horizontal de $f(x)=\frac{2^x+2}{2^x+4}$. 

\[
\begin{aligned}
t=2^x>0\\
y=\frac{t+2}{t+4}\\
y(t+4)=t+2\\
t(y-1)=2-4y\\
t=\frac{2-4y}{y-1}>0\\
\operatorname{Im}(f)=(\frac{2}{4},1)\\
D_f=\mathbb{R}\\
\text{{AH: }}y=1
\end{aligned}
\]


\paragraph{Ejercicio 73.} Determina dominio, imagen y asíntota horizontal de $f(x)=\ln\left(\frac{x^2-9}{x^2-25}\right)$.

\[
\begin{aligned}
\frac{x^2-9}{x^2-25} >0\\
D_f=(-\infty,-5)\cup(-3,3)\cup(5,\infty)\\
\frac{x^2-9}{x^2-25}\ne1\\
\ln(1)=0\\
\operatorname{Im}(f)=\mathbb{R}\setminus\{0\}\\
\text{{AH: }}y=0
\end{aligned}
\]


\paragraph{Ejercicio 74.} Determina dominio, imagen y asíntotas horizontales de $f(x)=\arctan(x+0)$.

\[
\begin{aligned}
D_f=\mathbb{R}\\
-\frac{\pi}{2}<\arctan(x+0)<\frac{\pi}{2}\\
\operatorname{Im}(f)=\left(-\frac{\pi}{2},\frac{\pi}{2}\right)\\
\text{AH: }y=\frac{\pi}{2}\text{ y }y=-\frac{\pi}{2}
\end{aligned}
\]


\paragraph{Ejercicio 75.} Determina dominio, imagen y asíntotas horizontales de $f(x)=\frac{e^x-e^{-x}}{e^x+e^{-x}}+1$.

\[
\begin{aligned}
D_f=\mathbb{R}\\
t=e^x>0\\
f(x)=\frac{t^2-1}{t^2+1}+1\\
-1<\frac{t^2-1}{t^2+1}<1\\
\operatorname{Im}(f)=(0,2)\\
\text{AH: }y=0\text{ y }y=2
\end{aligned}
\]


\paragraph{Ejercicio 76.} Determina dominio, imagen y asíntota horizontal de $f(x)=\sqrt{\frac{x-1}{x-3}}+2$.

\[
\begin{aligned}
\frac{x-1}{x-3}\ge0\\
D_f=(-\infty,1]\cup(3,\infty)\\
\frac{x-1}{x-3}\ne1\\
\sqrt{\frac{x-1}{x-3}}\ne1\\
\operatorname{Im}(f)=[2,\infty)\setminus\{3\}\\
\text{AH: }y=3
\end{aligned}
\]


\paragraph{Ejercicio 77.} Determina dominio, imagen y asíntota horizontal de $f(x)=\frac{2^x+2}{2^x+4}$. 

\[
\begin{aligned}
t=2^x>0\\
y=\frac{t+2}{t+4}\\
y(t+4)=t+2\\
t(y-1)=2-4y\\
t=\frac{2-4y}{y-1}>0\\
\operatorname{Im}(f)=(\frac{2}{4},1)\\
D_f=\mathbb{R}\\
\text{{AH: }}y=1
\end{aligned}
\]


\paragraph{Ejercicio 78.} Determina dominio, imagen y asíntota horizontal de $f(x)=\ln\left(\frac{x^2-9}{x^2-25}\right)$.

\[
\begin{aligned}
\frac{x^2-9}{x^2-25} >0\\
D_f=(-\infty,-5)\cup(-3,3)\cup(5,\infty)\\
\frac{x^2-9}{x^2-25}\ne1\\
\ln(1)=0\\
\operatorname{Im}(f)=\mathbb{R}\setminus\{0\}\\
\text{{AH: }}y=0
\end{aligned}
\]


\paragraph{Ejercicio 79.} Determina dominio, imagen y asíntotas horizontales de $f(x)=\arctan(x-1)$.

\[
\begin{aligned}
D_f=\mathbb{R}\\
-\frac{\pi}{2}<\arctan(x-1)<\frac{\pi}{2}\\
\operatorname{Im}(f)=\left(-\frac{\pi}{2},\frac{\pi}{2}\right)\\
\text{AH: }y=\frac{\pi}{2}\text{ y }y=-\frac{\pi}{2}
\end{aligned}
\]


\paragraph{Ejercicio 80.} Determina dominio, imagen y asíntotas horizontales de $f(x)=\frac{e^x-e^{-x}}{e^x+e^{-x}}+2$.

\[
\begin{aligned}
D_f=\mathbb{R}\\
t=e^x>0\\
f(x)=\frac{t^2-1}{t^2+1}+2\\
-1<\frac{t^2-1}{t^2+1}<1\\
\operatorname{Im}(f)=(1,3)\\
\text{AH: }y=1\text{ y }y=3
\end{aligned}
\]


\newpage

\subsubsection*{Nivel 5 - Imposible}

\paragraph{Ejercicio 81.} Determina dominio, imagen y asíntota horizontal de $f(x)=\sqrt{\frac{x-1}{x-3}}-1$.

\[
\begin{aligned}
\frac{x-1}{x-3}\ge0\\
D_f=(-\infty,1]\cup(3,\infty)\\
\frac{x-1}{x-3}\ne1\\
\sqrt{\frac{x-1}{x-3}}\ne1\\
\operatorname{Im}(f)=[-1,\infty)\setminus\{0\}\\
\text{AH: }y=0
\end{aligned}
\]


\paragraph{Ejercicio 82.} Determina dominio, imagen y asíntota horizontal de $f(x)=\frac{2^x+2}{2^x+4}$. 

\[
\begin{aligned}
t=2^x>0\\
y=\frac{t+2}{t+4}\\
y(t+4)=t+2\\
t(y-1)=2-4y\\
t=\frac{2-4y}{y-1}>0\\
\operatorname{Im}(f)=(\frac{2}{4},1)\\
D_f=\mathbb{R}\\
\text{{AH: }}y=1
\end{aligned}
\]


\paragraph{Ejercicio 83.} Determina dominio, imagen y asíntota horizontal de $f(x)=\ln\left(\frac{x^2-9}{x^2-25}\right)$.

\[
\begin{aligned}
\frac{x^2-9}{x^2-25} >0\\
D_f=(-\infty,-5)\cup(-3,3)\cup(5,\infty)\\
\frac{x^2-9}{x^2-25}\ne1\\
\ln(1)=0\\
\operatorname{Im}(f)=\mathbb{R}\setminus\{0\}\\
\text{{AH: }}y=0
\end{aligned}
\]


\paragraph{Ejercicio 84.} Determina dominio, imagen y asíntotas horizontales de $f(x)=\arctan(x-2)$.

\[
\begin{aligned}
D_f=\mathbb{R}\\
-\frac{\pi}{2}<\arctan(x-2)<\frac{\pi}{2}\\
\operatorname{Im}(f)=\left(-\frac{\pi}{2},\frac{\pi}{2}\right)\\
\text{AH: }y=\frac{\pi}{2}\text{ y }y=-\frac{\pi}{2}
\end{aligned}
\]


\paragraph{Ejercicio 85.} Determina dominio, imagen y asíntotas horizontales de $f(x)=\frac{e^x-e^{-x}}{e^x+e^{-x}}+-1$.

\[
\begin{aligned}
D_f=\mathbb{R}\\
t=e^x>0\\
f(x)=\frac{t^2-1}{t^2+1}+-1\\
-1<\frac{t^2-1}{t^2+1}<1\\
\operatorname{Im}(f)=(-2,0)\\
\text{AH: }y=-2\text{ y }y=0
\end{aligned}
\]


\paragraph{Ejercicio 86.} Determina dominio, imagen y asíntota horizontal de $f(x)=\sqrt{\frac{x-1}{x-3}}+0$.

\[
\begin{aligned}
\frac{x-1}{x-3}\ge0\\
D_f=(-\infty,1]\cup(3,\infty)\\
\frac{x-1}{x-3}\ne1\\
\sqrt{\frac{x-1}{x-3}}\ne1\\
\operatorname{Im}(f)=[0,\infty)\setminus\{1\}\\
\text{AH: }y=1
\end{aligned}
\]


\paragraph{Ejercicio 87.} Determina dominio, imagen y asíntota horizontal de $f(x)=\frac{2^x+2}{2^x+4}$. 

\[
\begin{aligned}
t=2^x>0\\
y=\frac{t+2}{t+4}\\
y(t+4)=t+2\\
t(y-1)=2-4y\\
t=\frac{2-4y}{y-1}>0\\
\operatorname{Im}(f)=(\frac{2}{4},1)\\
D_f=\mathbb{R}\\
\text{{AH: }}y=1
\end{aligned}
\]


\paragraph{Ejercicio 88.} Determina dominio, imagen y asíntota horizontal de $f(x)=\ln\left(\frac{x^2-9}{x^2-25}\right)$.

\[
\begin{aligned}
\frac{x^2-9}{x^2-25} >0\\
D_f=(-\infty,-5)\cup(-3,3)\cup(5,\infty)\\
\frac{x^2-9}{x^2-25}\ne1\\
\ln(1)=0\\
\operatorname{Im}(f)=\mathbb{R}\setminus\{0\}\\
\text{{AH: }}y=0
\end{aligned}
\]


\paragraph{Ejercicio 89.} Determina dominio, imagen y asíntotas horizontales de $f(x)=\arctan(x+1)$.

\[
\begin{aligned}
D_f=\mathbb{R}\\
-\frac{\pi}{2}<\arctan(x+1)<\frac{\pi}{2}\\
\operatorname{Im}(f)=\left(-\frac{\pi}{2},\frac{\pi}{2}\right)\\
\text{AH: }y=\frac{\pi}{2}\text{ y }y=-\frac{\pi}{2}
\end{aligned}
\]


\paragraph{Ejercicio 90.} Determina dominio, imagen y asíntotas horizontales de $f(x)=\frac{e^x-e^{-x}}{e^x+e^{-x}}+0$.

\[
\begin{aligned}
D_f=\mathbb{R}\\
t=e^x>0\\
f(x)=\frac{t^2-1}{t^2+1}+0\\
-1<\frac{t^2-1}{t^2+1}<1\\
\operatorname{Im}(f)=(-1,1)\\
\text{AH: }y=-1\text{ y }y=1
\end{aligned}
\]


\paragraph{Ejercicio 91.} Determina dominio, imagen y asíntota horizontal de $f(x)=\sqrt{\frac{x-1}{x-3}}+1$.

\[
\begin{aligned}
\frac{x-1}{x-3}\ge0\\
D_f=(-\infty,1]\cup(3,\infty)\\
\frac{x-1}{x-3}\ne1\\
\sqrt{\frac{x-1}{x-3}}\ne1\\
\operatorname{Im}(f)=[1,\infty)\setminus\{2\}\\
\text{AH: }y=2
\end{aligned}
\]


\paragraph{Ejercicio 92.} Determina dominio, imagen y asíntota horizontal de $f(x)=\frac{2^x+2}{2^x+4}$. 

\[
\begin{aligned}
t=2^x>0\\
y=\frac{t+2}{t+4}\\
y(t+4)=t+2\\
t(y-1)=2-4y\\
t=\frac{2-4y}{y-1}>0\\
\operatorname{Im}(f)=(\frac{2}{4},1)\\
D_f=\mathbb{R}\\
\text{{AH: }}y=1
\end{aligned}
\]


\paragraph{Ejercicio 93.} Determina dominio, imagen y asíntota horizontal de $f(x)=\ln\left(\frac{x^2-9}{x^2-25}\right)$.

\[
\begin{aligned}
\frac{x^2-9}{x^2-25} >0\\
D_f=(-\infty,-5)\cup(-3,3)\cup(5,\infty)\\
\frac{x^2-9}{x^2-25}\ne1\\
\ln(1)=0\\
\operatorname{Im}(f)=\mathbb{R}\setminus\{0\}\\
\text{{AH: }}y=0
\end{aligned}
\]


\paragraph{Ejercicio 94.} Determina dominio, imagen y asíntotas horizontales de $f(x)=\arctan(x+0)$.

\[
\begin{aligned}
D_f=\mathbb{R}\\
-\frac{\pi}{2}<\arctan(x+0)<\frac{\pi}{2}\\
\operatorname{Im}(f)=\left(-\frac{\pi}{2},\frac{\pi}{2}\right)\\
\text{AH: }y=\frac{\pi}{2}\text{ y }y=-\frac{\pi}{2}
\end{aligned}
\]


\paragraph{Ejercicio 95.} Determina dominio, imagen y asíntotas horizontales de $f(x)=\frac{e^x-e^{-x}}{e^x+e^{-x}}+1$.

\[
\begin{aligned}
D_f=\mathbb{R}\\
t=e^x>0\\
f(x)=\frac{t^2-1}{t^2+1}+1\\
-1<\frac{t^2-1}{t^2+1}<1\\
\operatorname{Im}(f)=(0,2)\\
\text{AH: }y=0\text{ y }y=2
\end{aligned}
\]


\paragraph{Ejercicio 96.} Determina dominio, imagen y asíntota horizontal de $f(x)=\sqrt{\frac{x-1}{x-3}}+2$.

\[
\begin{aligned}
\frac{x-1}{x-3}\ge0\\
D_f=(-\infty,1]\cup(3,\infty)\\
\frac{x-1}{x-3}\ne1\\
\sqrt{\frac{x-1}{x-3}}\ne1\\
\operatorname{Im}(f)=[2,\infty)\setminus\{3\}\\
\text{AH: }y=3
\end{aligned}
\]


\paragraph{Ejercicio 97.} Determina dominio, imagen y asíntota horizontal de $f(x)=\frac{2^x+2}{2^x+4}$. 

\[
\begin{aligned}
t=2^x>0\\
y=\frac{t+2}{t+4}\\
y(t+4)=t+2\\
t(y-1)=2-4y\\
t=\frac{2-4y}{y-1}>0\\
\operatorname{Im}(f)=(\frac{2}{4},1)\\
D_f=\mathbb{R}\\
\text{{AH: }}y=1
\end{aligned}
\]


\paragraph{Ejercicio 98.} Determina dominio, imagen y asíntota horizontal de $f(x)=\ln\left(\frac{x^2-9}{x^2-25}\right)$.

\[
\begin{aligned}
\frac{x^2-9}{x^2-25} >0\\
D_f=(-\infty,-5)\cup(-3,3)\cup(5,\infty)\\
\frac{x^2-9}{x^2-25}\ne1\\
\ln(1)=0\\
\operatorname{Im}(f)=\mathbb{R}\setminus\{0\}\\
\text{{AH: }}y=0
\end{aligned}
\]


\paragraph{Ejercicio 99.} Determina dominio, imagen y asíntotas horizontales de $f(x)=\arctan(x-1)$.

\[
\begin{aligned}
D_f=\mathbb{R}\\
-\frac{\pi}{2}<\arctan(x-1)<\frac{\pi}{2}\\
\operatorname{Im}(f)=\left(-\frac{\pi}{2},\frac{\pi}{2}\right)\\
\text{AH: }y=\frac{\pi}{2}\text{ y }y=-\frac{\pi}{2}
\end{aligned}
\]


\paragraph{Ejercicio 100.} Determina dominio, imagen y asíntotas horizontales de $f(x)=\frac{e^x-e^{-x}}{e^x+e^{-x}}+2$.

\[
\begin{aligned}
D_f=\mathbb{R}\\
t=e^x>0\\
f(x)=\frac{t^2-1}{t^2+1}+2\\
-1<\frac{t^2-1}{t^2+1}<1\\
\operatorname{Im}(f)=(1,3)\\
\text{AH: }y=1\text{ y }y=3
\end{aligned}
\]


\newpage

\subsection*{Crecimiento, puntos críticos y concavidad}

\addcontentsline{toc}{subsection}{Soluciones - Crecimiento, puntos críticos y concavidad}

\subsubsection*{Nivel 1 - Básico}

\paragraph{Ejercicio 101.} Para $f(x)=(x+4)^2-3$, determina crecimiento, puntos críticos y concavidad.

\[
\begin{aligned}
f'(x)=2(x+4)\\
2(x+4)=0\Rightarrow x=-4\\
f(-4)=-3\\
f''(x)=2>0\\
(-\infty,-4)\downarrow,\quad(-4,\infty)\uparrow\\
(-4,-3)\text{ mínimo}\\
\text{{cóncava hacia arriba en }}\mathbb{R}
\end{aligned}
\]


\paragraph{Ejercicio 102.} Para $f(x)=-(x+3)^2-2$, determina crecimiento, puntos críticos y concavidad.

\[
\begin{aligned}
f'(x)=-2(x+3)\\
-2(x+3)=0\Rightarrow x=-3\\
f(-3)=-2\\
f''(x)=-2<0\\
(-\infty,-3)\uparrow,\quad(-3,\infty)\downarrow\\
(-3,-2)\text{ máximo}\\
\text{{cóncava hacia abajo en }}\mathbb{R}
\end{aligned}
\]


\paragraph{Ejercicio 103.} Para $f(x)=(x+2)^2-1$, determina crecimiento, puntos críticos y concavidad.

\[
\begin{aligned}
f'(x)=2(x+2)\\
2(x+2)=0\Rightarrow x=-2\\
f(-2)=-1\\
f''(x)=2>0\\
(-\infty,-2)\downarrow,\quad(-2,\infty)\uparrow\\
(-2,-1)\text{ mínimo}\\
\text{{cóncava hacia arriba en }}\mathbb{R}
\end{aligned}
\]


\paragraph{Ejercicio 104.} Para $f(x)=-(x+1)^2+0$, determina crecimiento, puntos críticos y concavidad.

\[
\begin{aligned}
f'(x)=-2(x+1)\\
-2(x+1)=0\Rightarrow x=-1\\
f(-1)=0\\
f''(x)=-2<0\\
(-\infty,-1)\uparrow,\quad(-1,\infty)\downarrow\\
(-1,0)\text{ máximo}\\
\text{{cóncava hacia abajo en }}\mathbb{R}
\end{aligned}
\]


\paragraph{Ejercicio 105.} Para $f(x)=(x+0)^2+1$, determina crecimiento, puntos críticos y concavidad.

\[
\begin{aligned}
f'(x)=2(x+0)\\
2(x+0)=0\Rightarrow x=0\\
f(0)=1\\
f''(x)=2>0\\
(-\infty,0)\downarrow,\quad(0,\infty)\uparrow\\
(0,1)\text{ mínimo}\\
\text{{cóncava hacia arriba en }}\mathbb{R}
\end{aligned}
\]


\paragraph{Ejercicio 106.} Para $f(x)=-(x-1)^2+2$, determina crecimiento, puntos críticos y concavidad.

\[
\begin{aligned}
f'(x)=-2(x-1)\\
-2(x-1)=0\Rightarrow x=1\\
f(1)=2\\
f''(x)=-2<0\\
(-\infty,1)\uparrow,\quad(1,\infty)\downarrow\\
(1,2)\text{ máximo}\\
\text{{cóncava hacia abajo en }}\mathbb{R}
\end{aligned}
\]


\paragraph{Ejercicio 107.} Para $f(x)=(x-2)^2+3$, determina crecimiento, puntos críticos y concavidad.

\[
\begin{aligned}
f'(x)=2(x-2)\\
2(x-2)=0\Rightarrow x=2\\
f(2)=3\\
f''(x)=2>0\\
(-\infty,2)\downarrow,\quad(2,\infty)\uparrow\\
(2,3)\text{ mínimo}\\
\text{{cóncava hacia arriba en }}\mathbb{R}
\end{aligned}
\]


\paragraph{Ejercicio 108.} Para $f(x)=-(x-3)^2-3$, determina crecimiento, puntos críticos y concavidad.

\[
\begin{aligned}
f'(x)=-2(x-3)\\
-2(x-3)=0\Rightarrow x=3\\
f(3)=-3\\
f''(x)=-2<0\\
(-\infty,3)\uparrow,\quad(3,\infty)\downarrow\\
(3,-3)\text{ máximo}\\
\text{{cóncava hacia abajo en }}\mathbb{R}
\end{aligned}
\]


\paragraph{Ejercicio 109.} Para $f(x)=(x-4)^2-2$, determina crecimiento, puntos críticos y concavidad.

\[
\begin{aligned}
f'(x)=2(x-4)\\
2(x-4)=0\Rightarrow x=4\\
f(4)=-2\\
f''(x)=2>0\\
(-\infty,4)\downarrow,\quad(4,\infty)\uparrow\\
(4,-2)\text{ mínimo}\\
\text{{cóncava hacia arriba en }}\mathbb{R}
\end{aligned}
\]


\paragraph{Ejercicio 110.} Para $f(x)=-(x+4)^2-1$, determina crecimiento, puntos críticos y concavidad.

\[
\begin{aligned}
f'(x)=-2(x+4)\\
-2(x+4)=0\Rightarrow x=-4\\
f(-4)=-1\\
f''(x)=-2<0\\
(-\infty,-4)\uparrow,\quad(-4,\infty)\downarrow\\
(-4,-1)\text{ máximo}\\
\text{{cóncava hacia abajo en }}\mathbb{R}
\end{aligned}
\]


\paragraph{Ejercicio 111.} Para $f(x)=(x+3)^2+0$, determina crecimiento, puntos críticos y concavidad.

\[
\begin{aligned}
f'(x)=2(x+3)\\
2(x+3)=0\Rightarrow x=-3\\
f(-3)=0\\
f''(x)=2>0\\
(-\infty,-3)\downarrow,\quad(-3,\infty)\uparrow\\
(-3,0)\text{ mínimo}\\
\text{{cóncava hacia arriba en }}\mathbb{R}
\end{aligned}
\]


\paragraph{Ejercicio 112.} Para $f(x)=-(x+2)^2+1$, determina crecimiento, puntos críticos y concavidad.

\[
\begin{aligned}
f'(x)=-2(x+2)\\
-2(x+2)=0\Rightarrow x=-2\\
f(-2)=1\\
f''(x)=-2<0\\
(-\infty,-2)\uparrow,\quad(-2,\infty)\downarrow\\
(-2,1)\text{ máximo}\\
\text{{cóncava hacia abajo en }}\mathbb{R}
\end{aligned}
\]


\paragraph{Ejercicio 113.} Para $f(x)=(x+1)^2+2$, determina crecimiento, puntos críticos y concavidad.

\[
\begin{aligned}
f'(x)=2(x+1)\\
2(x+1)=0\Rightarrow x=-1\\
f(-1)=2\\
f''(x)=2>0\\
(-\infty,-1)\downarrow,\quad(-1,\infty)\uparrow\\
(-1,2)\text{ mínimo}\\
\text{{cóncava hacia arriba en }}\mathbb{R}
\end{aligned}
\]


\paragraph{Ejercicio 114.} Para $f(x)=-(x+0)^2+3$, determina crecimiento, puntos críticos y concavidad.

\[
\begin{aligned}
f'(x)=-2(x+0)\\
-2(x+0)=0\Rightarrow x=0\\
f(0)=3\\
f''(x)=-2<0\\
(-\infty,0)\uparrow,\quad(0,\infty)\downarrow\\
(0,3)\text{ máximo}\\
\text{{cóncava hacia abajo en }}\mathbb{R}
\end{aligned}
\]


\paragraph{Ejercicio 115.} Para $f(x)=(x-1)^2-3$, determina crecimiento, puntos críticos y concavidad.

\[
\begin{aligned}
f'(x)=2(x-1)\\
2(x-1)=0\Rightarrow x=1\\
f(1)=-3\\
f''(x)=2>0\\
(-\infty,1)\downarrow,\quad(1,\infty)\uparrow\\
(1,-3)\text{ mínimo}\\
\text{{cóncava hacia arriba en }}\mathbb{R}
\end{aligned}
\]


\paragraph{Ejercicio 116.} Para $f(x)=-(x-2)^2-2$, determina crecimiento, puntos críticos y concavidad.

\[
\begin{aligned}
f'(x)=-2(x-2)\\
-2(x-2)=0\Rightarrow x=2\\
f(2)=-2\\
f''(x)=-2<0\\
(-\infty,2)\uparrow,\quad(2,\infty)\downarrow\\
(2,-2)\text{ máximo}\\
\text{{cóncava hacia abajo en }}\mathbb{R}
\end{aligned}
\]


\paragraph{Ejercicio 117.} Para $f(x)=(x-3)^2-1$, determina crecimiento, puntos críticos y concavidad.

\[
\begin{aligned}
f'(x)=2(x-3)\\
2(x-3)=0\Rightarrow x=3\\
f(3)=-1\\
f''(x)=2>0\\
(-\infty,3)\downarrow,\quad(3,\infty)\uparrow\\
(3,-1)\text{ mínimo}\\
\text{{cóncava hacia arriba en }}\mathbb{R}
\end{aligned}
\]


\paragraph{Ejercicio 118.} Para $f(x)=-(x-4)^2+0$, determina crecimiento, puntos críticos y concavidad.

\[
\begin{aligned}
f'(x)=-2(x-4)\\
-2(x-4)=0\Rightarrow x=4\\
f(4)=0\\
f''(x)=-2<0\\
(-\infty,4)\uparrow,\quad(4,\infty)\downarrow\\
(4,0)\text{ máximo}\\
\text{{cóncava hacia abajo en }}\mathbb{R}
\end{aligned}
\]


\paragraph{Ejercicio 119.} Para $f(x)=(x+4)^2+1$, determina crecimiento, puntos críticos y concavidad.

\[
\begin{aligned}
f'(x)=2(x+4)\\
2(x+4)=0\Rightarrow x=-4\\
f(-4)=1\\
f''(x)=2>0\\
(-\infty,-4)\downarrow,\quad(-4,\infty)\uparrow\\
(-4,1)\text{ mínimo}\\
\text{{cóncava hacia arriba en }}\mathbb{R}
\end{aligned}
\]


\paragraph{Ejercicio 120.} Para $f(x)=-(x+3)^2+2$, determina crecimiento, puntos críticos y concavidad.

\[
\begin{aligned}
f'(x)=-2(x+3)\\
-2(x+3)=0\Rightarrow x=-3\\
f(-3)=2\\
f''(x)=-2<0\\
(-\infty,-3)\uparrow,\quad(-3,\infty)\downarrow\\
(-3,2)\text{ máximo}\\
\text{{cóncava hacia abajo en }}\mathbb{R}
\end{aligned}
\]


\newpage

\subsubsection*{Nivel 2 - Intermedio bajo}

\paragraph{Ejercicio 121.} Para $f(x)=x^3-3(1)^2x$, determina crecimiento, puntos críticos y concavidad.

\[
\begin{aligned}
f'(x)=3x^2-3(1)^2=3(x-1)(x+1)\\
f'(x)=0\Rightarrow x=-1,\ 1\\
f''(x)=6x\\
(-\infty,-1)\uparrow,\quad(-1,1)\downarrow,\quad(1,\infty)\uparrow\\
x=-1:\text{ máximo local}\\
x=1:\text{ mínimo local}\\
(-\infty,0)\text{{ cóncava abajo}},\quad(0,\infty)\text{{ cóncava arriba}}\\
x=0\text{{ punto de inflexión}}
\end{aligned}
\]


\paragraph{Ejercicio 122.} Para $f(x)=(x+2)^3$, determina crecimiento, puntos críticos y concavidad.

\[
\begin{aligned}
f'(x)=3(x+2)^2\\
f'(x)=0\Rightarrow x=-2\\
f''(x)=6(x+2)\\
f'(x)\ge0\\
\mathbb{R}\uparrow\\
x=-2:\text{ crítico sin extremo}\\
(-\infty,-2)\text{ cóncava abajo},\quad(-2,\infty)\text{ cóncava arriba}\\
x=-2\text{ punto de inflexión}
\end{aligned}
\]


\paragraph{Ejercicio 123.} Para $f(x)=x^3-3(3)^2x$, determina crecimiento, puntos críticos y concavidad.

\[
\begin{aligned}
f'(x)=3x^2-3(3)^2=3(x-3)(x+3)\\
f'(x)=0\Rightarrow x=-3,\ 3\\
f''(x)=6x\\
(-\infty,-3)\uparrow,\quad(-3,3)\downarrow,\quad(3,\infty)\uparrow\\
x=-3:\text{ máximo local}\\
x=3:\text{ mínimo local}\\
(-\infty,0)\text{{ cóncava abajo}},\quad(0,\infty)\text{{ cóncava arriba}}\\
x=0\text{{ punto de inflexión}}
\end{aligned}
\]


\paragraph{Ejercicio 124.} Para $f(x)=(x+0)^3$, determina crecimiento, puntos críticos y concavidad.

\[
\begin{aligned}
f'(x)=3(x+0)^2\\
f'(x)=0\Rightarrow x=0\\
f''(x)=6(x+0)\\
f'(x)\ge0\\
\mathbb{R}\uparrow\\
x=0:\text{ crítico sin extremo}\\
(-\infty,0)\text{ cóncava abajo},\quad(0,\infty)\text{ cóncava arriba}\\
x=0\text{ punto de inflexión}
\end{aligned}
\]


\paragraph{Ejercicio 125.} Para $f(x)=x^3-3(5)^2x$, determina crecimiento, puntos críticos y concavidad.

\[
\begin{aligned}
f'(x)=3x^2-3(5)^2=3(x-5)(x+5)\\
f'(x)=0\Rightarrow x=-5,\ 5\\
f''(x)=6x\\
(-\infty,-5)\uparrow,\quad(-5,5)\downarrow,\quad(5,\infty)\uparrow\\
x=-5:\text{ máximo local}\\
x=5:\text{ mínimo local}\\
(-\infty,0)\text{{ cóncava abajo}},\quad(0,\infty)\text{{ cóncava arriba}}\\
x=0\text{{ punto de inflexión}}
\end{aligned}
\]


\paragraph{Ejercicio 126.} Para $f(x)=(x-2)^3$, determina crecimiento, puntos críticos y concavidad.

\[
\begin{aligned}
f'(x)=3(x-2)^2\\
f'(x)=0\Rightarrow x=2\\
f''(x)=6(x-2)\\
f'(x)\ge0\\
\mathbb{R}\uparrow\\
x=2:\text{ crítico sin extremo}\\
(-\infty,2)\text{ cóncava abajo},\quad(2,\infty)\text{ cóncava arriba}\\
x=2\text{ punto de inflexión}
\end{aligned}
\]


\paragraph{Ejercicio 127.} Para $f(x)=x^3-3(2)^2x$, determina crecimiento, puntos críticos y concavidad.

\[
\begin{aligned}
f'(x)=3x^2-3(2)^2=3(x-2)(x+2)\\
f'(x)=0\Rightarrow x=-2,\ 2\\
f''(x)=6x\\
(-\infty,-2)\uparrow,\quad(-2,2)\downarrow,\quad(2,\infty)\uparrow\\
x=-2:\text{ máximo local}\\
x=2:\text{ mínimo local}\\
(-\infty,0)\text{{ cóncava abajo}},\quad(0,\infty)\text{{ cóncava arriba}}\\
x=0\text{{ punto de inflexión}}
\end{aligned}
\]


\paragraph{Ejercicio 128.} Para $f(x)=(x+3)^3$, determina crecimiento, puntos críticos y concavidad.

\[
\begin{aligned}
f'(x)=3(x+3)^2\\
f'(x)=0\Rightarrow x=-3\\
f''(x)=6(x+3)\\
f'(x)\ge0\\
\mathbb{R}\uparrow\\
x=-3:\text{ crítico sin extremo}\\
(-\infty,-3)\text{ cóncava abajo},\quad(-3,\infty)\text{ cóncava arriba}\\
x=-3\text{ punto de inflexión}
\end{aligned}
\]


\paragraph{Ejercicio 129.} Para $f(x)=x^3-3(4)^2x$, determina crecimiento, puntos críticos y concavidad.

\[
\begin{aligned}
f'(x)=3x^2-3(4)^2=3(x-4)(x+4)\\
f'(x)=0\Rightarrow x=-4,\ 4\\
f''(x)=6x\\
(-\infty,-4)\uparrow,\quad(-4,4)\downarrow,\quad(4,\infty)\uparrow\\
x=-4:\text{ máximo local}\\
x=4:\text{ mínimo local}\\
(-\infty,0)\text{{ cóncava abajo}},\quad(0,\infty)\text{{ cóncava arriba}}\\
x=0\text{{ punto de inflexión}}
\end{aligned}
\]


\paragraph{Ejercicio 130.} Para $f(x)=(x+1)^3$, determina crecimiento, puntos críticos y concavidad.

\[
\begin{aligned}
f'(x)=3(x+1)^2\\
f'(x)=0\Rightarrow x=-1\\
f''(x)=6(x+1)\\
f'(x)\ge0\\
\mathbb{R}\uparrow\\
x=-1:\text{ crítico sin extremo}\\
(-\infty,-1)\text{ cóncava abajo},\quad(-1,\infty)\text{ cóncava arriba}\\
x=-1\text{ punto de inflexión}
\end{aligned}
\]


\paragraph{Ejercicio 131.} Para $f(x)=x^3-3(1)^2x$, determina crecimiento, puntos críticos y concavidad.

\[
\begin{aligned}
f'(x)=3x^2-3(1)^2=3(x-1)(x+1)\\
f'(x)=0\Rightarrow x=-1,\ 1\\
f''(x)=6x\\
(-\infty,-1)\uparrow,\quad(-1,1)\downarrow,\quad(1,\infty)\uparrow\\
x=-1:\text{ máximo local}\\
x=1:\text{ mínimo local}\\
(-\infty,0)\text{{ cóncava abajo}},\quad(0,\infty)\text{{ cóncava arriba}}\\
x=0\text{{ punto de inflexión}}
\end{aligned}
\]


\paragraph{Ejercicio 132.} Para $f(x)=(x-1)^3$, determina crecimiento, puntos críticos y concavidad.

\[
\begin{aligned}
f'(x)=3(x-1)^2\\
f'(x)=0\Rightarrow x=1\\
f''(x)=6(x-1)\\
f'(x)\ge0\\
\mathbb{R}\uparrow\\
x=1:\text{ crítico sin extremo}\\
(-\infty,1)\text{ cóncava abajo},\quad(1,\infty)\text{ cóncava arriba}\\
x=1\text{ punto de inflexión}
\end{aligned}
\]


\paragraph{Ejercicio 133.} Para $f(x)=x^3-3(3)^2x$, determina crecimiento, puntos críticos y concavidad.

\[
\begin{aligned}
f'(x)=3x^2-3(3)^2=3(x-3)(x+3)\\
f'(x)=0\Rightarrow x=-3,\ 3\\
f''(x)=6x\\
(-\infty,-3)\uparrow,\quad(-3,3)\downarrow,\quad(3,\infty)\uparrow\\
x=-3:\text{ máximo local}\\
x=3:\text{ mínimo local}\\
(-\infty,0)\text{{ cóncava abajo}},\quad(0,\infty)\text{{ cóncava arriba}}\\
x=0\text{{ punto de inflexión}}
\end{aligned}
\]


\paragraph{Ejercicio 134.} Para $f(x)=(x-3)^3$, determina crecimiento, puntos críticos y concavidad.

\[
\begin{aligned}
f'(x)=3(x-3)^2\\
f'(x)=0\Rightarrow x=3\\
f''(x)=6(x-3)\\
f'(x)\ge0\\
\mathbb{R}\uparrow\\
x=3:\text{ crítico sin extremo}\\
(-\infty,3)\text{ cóncava abajo},\quad(3,\infty)\text{ cóncava arriba}\\
x=3\text{ punto de inflexión}
\end{aligned}
\]


\paragraph{Ejercicio 135.} Para $f(x)=x^3-3(5)^2x$, determina crecimiento, puntos críticos y concavidad.

\[
\begin{aligned}
f'(x)=3x^2-3(5)^2=3(x-5)(x+5)\\
f'(x)=0\Rightarrow x=-5,\ 5\\
f''(x)=6x\\
(-\infty,-5)\uparrow,\quad(-5,5)\downarrow,\quad(5,\infty)\uparrow\\
x=-5:\text{ máximo local}\\
x=5:\text{ mínimo local}\\
(-\infty,0)\text{{ cóncava abajo}},\quad(0,\infty)\text{{ cóncava arriba}}\\
x=0\text{{ punto de inflexión}}
\end{aligned}
\]


\paragraph{Ejercicio 136.} Para $f(x)=(x+2)^3$, determina crecimiento, puntos críticos y concavidad.

\[
\begin{aligned}
f'(x)=3(x+2)^2\\
f'(x)=0\Rightarrow x=-2\\
f''(x)=6(x+2)\\
f'(x)\ge0\\
\mathbb{R}\uparrow\\
x=-2:\text{ crítico sin extremo}\\
(-\infty,-2)\text{ cóncava abajo},\quad(-2,\infty)\text{ cóncava arriba}\\
x=-2\text{ punto de inflexión}
\end{aligned}
\]


\paragraph{Ejercicio 137.} Para $f(x)=x^3-3(2)^2x$, determina crecimiento, puntos críticos y concavidad.

\[
\begin{aligned}
f'(x)=3x^2-3(2)^2=3(x-2)(x+2)\\
f'(x)=0\Rightarrow x=-2,\ 2\\
f''(x)=6x\\
(-\infty,-2)\uparrow,\quad(-2,2)\downarrow,\quad(2,\infty)\uparrow\\
x=-2:\text{ máximo local}\\
x=2:\text{ mínimo local}\\
(-\infty,0)\text{{ cóncava abajo}},\quad(0,\infty)\text{{ cóncava arriba}}\\
x=0\text{{ punto de inflexión}}
\end{aligned}
\]


\paragraph{Ejercicio 138.} Para $f(x)=(x+0)^3$, determina crecimiento, puntos críticos y concavidad.

\[
\begin{aligned}
f'(x)=3(x+0)^2\\
f'(x)=0\Rightarrow x=0\\
f''(x)=6(x+0)\\
f'(x)\ge0\\
\mathbb{R}\uparrow\\
x=0:\text{ crítico sin extremo}\\
(-\infty,0)\text{ cóncava abajo},\quad(0,\infty)\text{ cóncava arriba}\\
x=0\text{ punto de inflexión}
\end{aligned}
\]


\paragraph{Ejercicio 139.} Para $f(x)=x^3-3(4)^2x$, determina crecimiento, puntos críticos y concavidad.

\[
\begin{aligned}
f'(x)=3x^2-3(4)^2=3(x-4)(x+4)\\
f'(x)=0\Rightarrow x=-4,\ 4\\
f''(x)=6x\\
(-\infty,-4)\uparrow,\quad(-4,4)\downarrow,\quad(4,\infty)\uparrow\\
x=-4:\text{ máximo local}\\
x=4:\text{ mínimo local}\\
(-\infty,0)\text{{ cóncava abajo}},\quad(0,\infty)\text{{ cóncava arriba}}\\
x=0\text{{ punto de inflexión}}
\end{aligned}
\]


\paragraph{Ejercicio 140.} Para $f(x)=(x-2)^3$, determina crecimiento, puntos críticos y concavidad.

\[
\begin{aligned}
f'(x)=3(x-2)^2\\
f'(x)=0\Rightarrow x=2\\
f''(x)=6(x-2)\\
f'(x)\ge0\\
\mathbb{R}\uparrow\\
x=2:\text{ crítico sin extremo}\\
(-\infty,2)\text{ cóncava abajo},\quad(2,\infty)\text{ cóncava arriba}\\
x=2\text{ punto de inflexión}
\end{aligned}
\]


\newpage

\subsubsection*{Nivel 3 - Intermedio}

\paragraph{Ejercicio 141.} Para $f(x)=x+\frac{1}{x}$, determina crecimiento, puntos críticos y concavidad.

\[
\begin{aligned}
D_f=\mathbb{R}\setminus\{0\}\\
f'(x)=1-\frac{1}{x^2}=\frac{x^2-1}{x^2}\\
f'(x)=0\Rightarrow x=\pm \sqrt{1}\\
f''(x)=\frac{2}{x^3}\\
(-\infty,-\sqrt{1})\uparrow,\quad(-\sqrt{1},0)\downarrow\\
(0,\sqrt{1})\downarrow,\quad(\sqrt{1},\infty)\uparrow\\
x=-\sqrt{1}:\text{ máximo},\quad x=\sqrt{1}:\text{ mínimo}\\
(-\infty,0)\text{{ cóncava abajo}},\quad(0,\infty)\text{{ cóncava arriba}}
\end{aligned}
\]


\paragraph{Ejercicio 142.} Para $f(x)=(x+1)^4$, determina crecimiento, puntos críticos y concavidad.

\[
\begin{aligned}
f'(x)=4(x+1)^3\\
f'(x)=0\Rightarrow x=-1\\
f''(x)=12(x+1)^2\ge0\\
(-\infty,-1)\downarrow,\quad(-1,\infty)\uparrow\\
x=-1:\text{ mínimo}\\
\text{{cóncava hacia arriba}}
\end{aligned}
\]


\paragraph{Ejercicio 143.} Para $f(x)=x+\frac{3}{x}$, determina crecimiento, puntos críticos y concavidad.

\[
\begin{aligned}
D_f=\mathbb{R}\setminus\{0\}\\
f'(x)=1-\frac{3}{x^2}=\frac{x^2-3}{x^2}\\
f'(x)=0\Rightarrow x=\pm \sqrt{3}\\
f''(x)=\frac{6}{x^3}\\
(-\infty,-\sqrt{3})\uparrow,\quad(-\sqrt{3},0)\downarrow\\
(0,\sqrt{3})\downarrow,\quad(\sqrt{3},\infty)\uparrow\\
x=-\sqrt{3}:\text{ máximo},\quad x=\sqrt{3}:\text{ mínimo}\\
(-\infty,0)\text{{ cóncava abajo}},\quad(0,\infty)\text{{ cóncava arriba}}
\end{aligned}
\]


\paragraph{Ejercicio 144.} Para $f(x)=(x-1)^4$, determina crecimiento, puntos críticos y concavidad.

\[
\begin{aligned}
f'(x)=4(x-1)^3\\
f'(x)=0\Rightarrow x=1\\
f''(x)=12(x-1)^2\ge0\\
(-\infty,1)\downarrow,\quad(1,\infty)\uparrow\\
x=1:\text{ mínimo}\\
\text{{cóncava hacia arriba}}
\end{aligned}
\]


\paragraph{Ejercicio 145.} Para $f(x)=x+\frac{5}{x}$, determina crecimiento, puntos críticos y concavidad.

\[
\begin{aligned}
D_f=\mathbb{R}\setminus\{0\}\\
f'(x)=1-\frac{5}{x^2}=\frac{x^2-5}{x^2}\\
f'(x)=0\Rightarrow x=\pm \sqrt{5}\\
f''(x)=\frac{10}{x^3}\\
(-\infty,-\sqrt{5})\uparrow,\quad(-\sqrt{5},0)\downarrow\\
(0,\sqrt{5})\downarrow,\quad(\sqrt{5},\infty)\uparrow\\
x=-\sqrt{5}:\text{ máximo},\quad x=\sqrt{5}:\text{ mínimo}\\
(-\infty,0)\text{{ cóncava abajo}},\quad(0,\infty)\text{{ cóncava arriba}}
\end{aligned}
\]


\paragraph{Ejercicio 146.} Para $f(x)=(x+2)^4$, determina crecimiento, puntos críticos y concavidad.

\[
\begin{aligned}
f'(x)=4(x+2)^3\\
f'(x)=0\Rightarrow x=-2\\
f''(x)=12(x+2)^2\ge0\\
(-\infty,-2)\downarrow,\quad(-2,\infty)\uparrow\\
x=-2:\text{ mínimo}\\
\text{{cóncava hacia arriba}}
\end{aligned}
\]


\paragraph{Ejercicio 147.} Para $f(x)=x+\frac{2}{x}$, determina crecimiento, puntos críticos y concavidad.

\[
\begin{aligned}
D_f=\mathbb{R}\setminus\{0\}\\
f'(x)=1-\frac{2}{x^2}=\frac{x^2-2}{x^2}\\
f'(x)=0\Rightarrow x=\pm \sqrt{2}\\
f''(x)=\frac{4}{x^3}\\
(-\infty,-\sqrt{2})\uparrow,\quad(-\sqrt{2},0)\downarrow\\
(0,\sqrt{2})\downarrow,\quad(\sqrt{2},\infty)\uparrow\\
x=-\sqrt{2}:\text{ máximo},\quad x=\sqrt{2}:\text{ mínimo}\\
(-\infty,0)\text{{ cóncava abajo}},\quad(0,\infty)\text{{ cóncava arriba}}
\end{aligned}
\]


\paragraph{Ejercicio 148.} Para $f(x)=(x+0)^4$, determina crecimiento, puntos críticos y concavidad.

\[
\begin{aligned}
f'(x)=4(x+0)^3\\
f'(x)=0\Rightarrow x=0\\
f''(x)=12(x+0)^2\ge0\\
(-\infty,0)\downarrow,\quad(0,\infty)\uparrow\\
x=0:\text{ mínimo}\\
\text{{cóncava hacia arriba}}
\end{aligned}
\]


\paragraph{Ejercicio 149.} Para $f(x)=x+\frac{4}{x}$, determina crecimiento, puntos críticos y concavidad.

\[
\begin{aligned}
D_f=\mathbb{R}\setminus\{0\}\\
f'(x)=1-\frac{4}{x^2}=\frac{x^2-4}{x^2}\\
f'(x)=0\Rightarrow x=\pm \sqrt{4}\\
f''(x)=\frac{8}{x^3}\\
(-\infty,-\sqrt{4})\uparrow,\quad(-\sqrt{4},0)\downarrow\\
(0,\sqrt{4})\downarrow,\quad(\sqrt{4},\infty)\uparrow\\
x=-\sqrt{4}:\text{ máximo},\quad x=\sqrt{4}:\text{ mínimo}\\
(-\infty,0)\text{{ cóncava abajo}},\quad(0,\infty)\text{{ cóncava arriba}}
\end{aligned}
\]


\paragraph{Ejercicio 150.} Para $f(x)=(x-2)^4$, determina crecimiento, puntos críticos y concavidad.

\[
\begin{aligned}
f'(x)=4(x-2)^3\\
f'(x)=0\Rightarrow x=2\\
f''(x)=12(x-2)^2\ge0\\
(-\infty,2)\downarrow,\quad(2,\infty)\uparrow\\
x=2:\text{ mínimo}\\
\text{{cóncava hacia arriba}}
\end{aligned}
\]


\paragraph{Ejercicio 151.} Para $f(x)=x+\frac{1}{x}$, determina crecimiento, puntos críticos y concavidad.

\[
\begin{aligned}
D_f=\mathbb{R}\setminus\{0\}\\
f'(x)=1-\frac{1}{x^2}=\frac{x^2-1}{x^2}\\
f'(x)=0\Rightarrow x=\pm \sqrt{1}\\
f''(x)=\frac{2}{x^3}\\
(-\infty,-\sqrt{1})\uparrow,\quad(-\sqrt{1},0)\downarrow\\
(0,\sqrt{1})\downarrow,\quad(\sqrt{1},\infty)\uparrow\\
x=-\sqrt{1}:\text{ máximo},\quad x=\sqrt{1}:\text{ mínimo}\\
(-\infty,0)\text{{ cóncava abajo}},\quad(0,\infty)\text{{ cóncava arriba}}
\end{aligned}
\]


\paragraph{Ejercicio 152.} Para $f(x)=(x+1)^4$, determina crecimiento, puntos críticos y concavidad.

\[
\begin{aligned}
f'(x)=4(x+1)^3\\
f'(x)=0\Rightarrow x=-1\\
f''(x)=12(x+1)^2\ge0\\
(-\infty,-1)\downarrow,\quad(-1,\infty)\uparrow\\
x=-1:\text{ mínimo}\\
\text{{cóncava hacia arriba}}
\end{aligned}
\]


\paragraph{Ejercicio 153.} Para $f(x)=x+\frac{3}{x}$, determina crecimiento, puntos críticos y concavidad.

\[
\begin{aligned}
D_f=\mathbb{R}\setminus\{0\}\\
f'(x)=1-\frac{3}{x^2}=\frac{x^2-3}{x^2}\\
f'(x)=0\Rightarrow x=\pm \sqrt{3}\\
f''(x)=\frac{6}{x^3}\\
(-\infty,-\sqrt{3})\uparrow,\quad(-\sqrt{3},0)\downarrow\\
(0,\sqrt{3})\downarrow,\quad(\sqrt{3},\infty)\uparrow\\
x=-\sqrt{3}:\text{ máximo},\quad x=\sqrt{3}:\text{ mínimo}\\
(-\infty,0)\text{{ cóncava abajo}},\quad(0,\infty)\text{{ cóncava arriba}}
\end{aligned}
\]


\paragraph{Ejercicio 154.} Para $f(x)=(x-1)^4$, determina crecimiento, puntos críticos y concavidad.

\[
\begin{aligned}
f'(x)=4(x-1)^3\\
f'(x)=0\Rightarrow x=1\\
f''(x)=12(x-1)^2\ge0\\
(-\infty,1)\downarrow,\quad(1,\infty)\uparrow\\
x=1:\text{ mínimo}\\
\text{{cóncava hacia arriba}}
\end{aligned}
\]


\paragraph{Ejercicio 155.} Para $f(x)=x+\frac{5}{x}$, determina crecimiento, puntos críticos y concavidad.

\[
\begin{aligned}
D_f=\mathbb{R}\setminus\{0\}\\
f'(x)=1-\frac{5}{x^2}=\frac{x^2-5}{x^2}\\
f'(x)=0\Rightarrow x=\pm \sqrt{5}\\
f''(x)=\frac{10}{x^3}\\
(-\infty,-\sqrt{5})\uparrow,\quad(-\sqrt{5},0)\downarrow\\
(0,\sqrt{5})\downarrow,\quad(\sqrt{5},\infty)\uparrow\\
x=-\sqrt{5}:\text{ máximo},\quad x=\sqrt{5}:\text{ mínimo}\\
(-\infty,0)\text{{ cóncava abajo}},\quad(0,\infty)\text{{ cóncava arriba}}
\end{aligned}
\]


\paragraph{Ejercicio 156.} Para $f(x)=(x+2)^4$, determina crecimiento, puntos críticos y concavidad.

\[
\begin{aligned}
f'(x)=4(x+2)^3\\
f'(x)=0\Rightarrow x=-2\\
f''(x)=12(x+2)^2\ge0\\
(-\infty,-2)\downarrow,\quad(-2,\infty)\uparrow\\
x=-2:\text{ mínimo}\\
\text{{cóncava hacia arriba}}
\end{aligned}
\]


\paragraph{Ejercicio 157.} Para $f(x)=x+\frac{2}{x}$, determina crecimiento, puntos críticos y concavidad.

\[
\begin{aligned}
D_f=\mathbb{R}\setminus\{0\}\\
f'(x)=1-\frac{2}{x^2}=\frac{x^2-2}{x^2}\\
f'(x)=0\Rightarrow x=\pm \sqrt{2}\\
f''(x)=\frac{4}{x^3}\\
(-\infty,-\sqrt{2})\uparrow,\quad(-\sqrt{2},0)\downarrow\\
(0,\sqrt{2})\downarrow,\quad(\sqrt{2},\infty)\uparrow\\
x=-\sqrt{2}:\text{ máximo},\quad x=\sqrt{2}:\text{ mínimo}\\
(-\infty,0)\text{{ cóncava abajo}},\quad(0,\infty)\text{{ cóncava arriba}}
\end{aligned}
\]


\paragraph{Ejercicio 158.} Para $f(x)=(x+0)^4$, determina crecimiento, puntos críticos y concavidad.

\[
\begin{aligned}
f'(x)=4(x+0)^3\\
f'(x)=0\Rightarrow x=0\\
f''(x)=12(x+0)^2\ge0\\
(-\infty,0)\downarrow,\quad(0,\infty)\uparrow\\
x=0:\text{ mínimo}\\
\text{{cóncava hacia arriba}}
\end{aligned}
\]


\paragraph{Ejercicio 159.} Para $f(x)=x+\frac{4}{x}$, determina crecimiento, puntos críticos y concavidad.

\[
\begin{aligned}
D_f=\mathbb{R}\setminus\{0\}\\
f'(x)=1-\frac{4}{x^2}=\frac{x^2-4}{x^2}\\
f'(x)=0\Rightarrow x=\pm \sqrt{4}\\
f''(x)=\frac{8}{x^3}\\
(-\infty,-\sqrt{4})\uparrow,\quad(-\sqrt{4},0)\downarrow\\
(0,\sqrt{4})\downarrow,\quad(\sqrt{4},\infty)\uparrow\\
x=-\sqrt{4}:\text{ máximo},\quad x=\sqrt{4}:\text{ mínimo}\\
(-\infty,0)\text{{ cóncava abajo}},\quad(0,\infty)\text{{ cóncava arriba}}
\end{aligned}
\]


\paragraph{Ejercicio 160.} Para $f(x)=(x-2)^4$, determina crecimiento, puntos críticos y concavidad.

\[
\begin{aligned}
f'(x)=4(x-2)^3\\
f'(x)=0\Rightarrow x=2\\
f''(x)=12(x-2)^2\ge0\\
(-\infty,2)\downarrow,\quad(2,\infty)\uparrow\\
x=2:\text{ mínimo}\\
\text{{cóncava hacia arriba}}
\end{aligned}
\]


\newpage

\subsubsection*{Nivel 4 - Avanzado}

\paragraph{Ejercicio 161.} Para $f(x)=e^x-2x$, determina crecimiento, puntos críticos y concavidad.

\[
\begin{aligned}
f'(x)=e^x-2\\
e^x=2\Rightarrow x=\ln 2\\
f''(x)=e^x>0\\
(-\infty,\ln 2)\downarrow,\quad(\ln 2,\infty)\uparrow\\
x=\ln 2:\text{ mínimo}\\
\text{{cóncava hacia arriba}}
\end{aligned}
\]


\paragraph{Ejercicio 162.} Para $f(x)=x^2\ln x$, determina crecimiento, puntos críticos y concavidad.

\[
\begin{aligned}
D_f=(0,\infty)\\
f'(x)=2x\ln x+x=x(2\ln x+1)\\
2\ln x+1=0\Rightarrow x=e^{-1/2}\\
f''(x)=2\ln x+3\\
f''(x)=0\Rightarrow x=e^{-3/2}\\
(0,e^{-1/2})\downarrow,\quad(e^{-1/2},\infty)\uparrow\\
x=e^{-1/2}:\text{{ mínimo}}\\
(0,e^{-3/2})\text{{ cóncava abajo}},\quad(e^{-3/2},\infty)\text{{ cóncava arriba}}
\end{aligned}
\]


\paragraph{Ejercicio 163.} Para $f(x)=e^x-4x$, determina crecimiento, puntos críticos y concavidad.

\[
\begin{aligned}
f'(x)=e^x-4\\
e^x=4\Rightarrow x=\ln 4\\
f''(x)=e^x>0\\
(-\infty,\ln 4)\downarrow,\quad(\ln 4,\infty)\uparrow\\
x=\ln 4:\text{ mínimo}\\
\text{{cóncava hacia arriba}}
\end{aligned}
\]


\paragraph{Ejercicio 164.} Para $f(x)=x^2\ln x$, determina crecimiento, puntos críticos y concavidad.

\[
\begin{aligned}
D_f=(0,\infty)\\
f'(x)=2x\ln x+x=x(2\ln x+1)\\
2\ln x+1=0\Rightarrow x=e^{-1/2}\\
f''(x)=2\ln x+3\\
f''(x)=0\Rightarrow x=e^{-3/2}\\
(0,e^{-1/2})\downarrow,\quad(e^{-1/2},\infty)\uparrow\\
x=e^{-1/2}:\text{{ mínimo}}\\
(0,e^{-3/2})\text{{ cóncava abajo}},\quad(e^{-3/2},\infty)\text{{ cóncava arriba}}
\end{aligned}
\]


\paragraph{Ejercicio 165.} Para $f(x)=e^x-6x$, determina crecimiento, puntos críticos y concavidad.

\[
\begin{aligned}
f'(x)=e^x-6\\
e^x=6\Rightarrow x=\ln 6\\
f''(x)=e^x>0\\
(-\infty,\ln 6)\downarrow,\quad(\ln 6,\infty)\uparrow\\
x=\ln 6:\text{ mínimo}\\
\text{{cóncava hacia arriba}}
\end{aligned}
\]


\paragraph{Ejercicio 166.} Para $f(x)=x^2\ln x$, determina crecimiento, puntos críticos y concavidad.

\[
\begin{aligned}
D_f=(0,\infty)\\
f'(x)=2x\ln x+x=x(2\ln x+1)\\
2\ln x+1=0\Rightarrow x=e^{-1/2}\\
f''(x)=2\ln x+3\\
f''(x)=0\Rightarrow x=e^{-3/2}\\
(0,e^{-1/2})\downarrow,\quad(e^{-1/2},\infty)\uparrow\\
x=e^{-1/2}:\text{{ mínimo}}\\
(0,e^{-3/2})\text{{ cóncava abajo}},\quad(e^{-3/2},\infty)\text{{ cóncava arriba}}
\end{aligned}
\]


\paragraph{Ejercicio 167.} Para $f(x)=e^x-3x$, determina crecimiento, puntos críticos y concavidad.

\[
\begin{aligned}
f'(x)=e^x-3\\
e^x=3\Rightarrow x=\ln 3\\
f''(x)=e^x>0\\
(-\infty,\ln 3)\downarrow,\quad(\ln 3,\infty)\uparrow\\
x=\ln 3:\text{ mínimo}\\
\text{{cóncava hacia arriba}}
\end{aligned}
\]


\paragraph{Ejercicio 168.} Para $f(x)=x^2\ln x$, determina crecimiento, puntos críticos y concavidad.

\[
\begin{aligned}
D_f=(0,\infty)\\
f'(x)=2x\ln x+x=x(2\ln x+1)\\
2\ln x+1=0\Rightarrow x=e^{-1/2}\\
f''(x)=2\ln x+3\\
f''(x)=0\Rightarrow x=e^{-3/2}\\
(0,e^{-1/2})\downarrow,\quad(e^{-1/2},\infty)\uparrow\\
x=e^{-1/2}:\text{{ mínimo}}\\
(0,e^{-3/2})\text{{ cóncava abajo}},\quad(e^{-3/2},\infty)\text{{ cóncava arriba}}
\end{aligned}
\]


\paragraph{Ejercicio 169.} Para $f(x)=e^x-5x$, determina crecimiento, puntos críticos y concavidad.

\[
\begin{aligned}
f'(x)=e^x-5\\
e^x=5\Rightarrow x=\ln 5\\
f''(x)=e^x>0\\
(-\infty,\ln 5)\downarrow,\quad(\ln 5,\infty)\uparrow\\
x=\ln 5:\text{ mínimo}\\
\text{{cóncava hacia arriba}}
\end{aligned}
\]


\paragraph{Ejercicio 170.} Para $f(x)=x^2\ln x$, determina crecimiento, puntos críticos y concavidad.

\[
\begin{aligned}
D_f=(0,\infty)\\
f'(x)=2x\ln x+x=x(2\ln x+1)\\
2\ln x+1=0\Rightarrow x=e^{-1/2}\\
f''(x)=2\ln x+3\\
f''(x)=0\Rightarrow x=e^{-3/2}\\
(0,e^{-1/2})\downarrow,\quad(e^{-1/2},\infty)\uparrow\\
x=e^{-1/2}:\text{{ mínimo}}\\
(0,e^{-3/2})\text{{ cóncava abajo}},\quad(e^{-3/2},\infty)\text{{ cóncava arriba}}
\end{aligned}
\]


\paragraph{Ejercicio 171.} Para $f(x)=e^x-2x$, determina crecimiento, puntos críticos y concavidad.

\[
\begin{aligned}
f'(x)=e^x-2\\
e^x=2\Rightarrow x=\ln 2\\
f''(x)=e^x>0\\
(-\infty,\ln 2)\downarrow,\quad(\ln 2,\infty)\uparrow\\
x=\ln 2:\text{ mínimo}\\
\text{{cóncava hacia arriba}}
\end{aligned}
\]


\paragraph{Ejercicio 172.} Para $f(x)=x^2\ln x$, determina crecimiento, puntos críticos y concavidad.

\[
\begin{aligned}
D_f=(0,\infty)\\
f'(x)=2x\ln x+x=x(2\ln x+1)\\
2\ln x+1=0\Rightarrow x=e^{-1/2}\\
f''(x)=2\ln x+3\\
f''(x)=0\Rightarrow x=e^{-3/2}\\
(0,e^{-1/2})\downarrow,\quad(e^{-1/2},\infty)\uparrow\\
x=e^{-1/2}:\text{{ mínimo}}\\
(0,e^{-3/2})\text{{ cóncava abajo}},\quad(e^{-3/2},\infty)\text{{ cóncava arriba}}
\end{aligned}
\]


\paragraph{Ejercicio 173.} Para $f(x)=e^x-4x$, determina crecimiento, puntos críticos y concavidad.

\[
\begin{aligned}
f'(x)=e^x-4\\
e^x=4\Rightarrow x=\ln 4\\
f''(x)=e^x>0\\
(-\infty,\ln 4)\downarrow,\quad(\ln 4,\infty)\uparrow\\
x=\ln 4:\text{ mínimo}\\
\text{{cóncava hacia arriba}}
\end{aligned}
\]


\paragraph{Ejercicio 174.} Para $f(x)=x^2\ln x$, determina crecimiento, puntos críticos y concavidad.

\[
\begin{aligned}
D_f=(0,\infty)\\
f'(x)=2x\ln x+x=x(2\ln x+1)\\
2\ln x+1=0\Rightarrow x=e^{-1/2}\\
f''(x)=2\ln x+3\\
f''(x)=0\Rightarrow x=e^{-3/2}\\
(0,e^{-1/2})\downarrow,\quad(e^{-1/2},\infty)\uparrow\\
x=e^{-1/2}:\text{{ mínimo}}\\
(0,e^{-3/2})\text{{ cóncava abajo}},\quad(e^{-3/2},\infty)\text{{ cóncava arriba}}
\end{aligned}
\]


\paragraph{Ejercicio 175.} Para $f(x)=e^x-6x$, determina crecimiento, puntos críticos y concavidad.

\[
\begin{aligned}
f'(x)=e^x-6\\
e^x=6\Rightarrow x=\ln 6\\
f''(x)=e^x>0\\
(-\infty,\ln 6)\downarrow,\quad(\ln 6,\infty)\uparrow\\
x=\ln 6:\text{ mínimo}\\
\text{{cóncava hacia arriba}}
\end{aligned}
\]


\paragraph{Ejercicio 176.} Para $f(x)=x^2\ln x$, determina crecimiento, puntos críticos y concavidad.

\[
\begin{aligned}
D_f=(0,\infty)\\
f'(x)=2x\ln x+x=x(2\ln x+1)\\
2\ln x+1=0\Rightarrow x=e^{-1/2}\\
f''(x)=2\ln x+3\\
f''(x)=0\Rightarrow x=e^{-3/2}\\
(0,e^{-1/2})\downarrow,\quad(e^{-1/2},\infty)\uparrow\\
x=e^{-1/2}:\text{{ mínimo}}\\
(0,e^{-3/2})\text{{ cóncava abajo}},\quad(e^{-3/2},\infty)\text{{ cóncava arriba}}
\end{aligned}
\]


\paragraph{Ejercicio 177.} Para $f(x)=e^x-3x$, determina crecimiento, puntos críticos y concavidad.

\[
\begin{aligned}
f'(x)=e^x-3\\
e^x=3\Rightarrow x=\ln 3\\
f''(x)=e^x>0\\
(-\infty,\ln 3)\downarrow,\quad(\ln 3,\infty)\uparrow\\
x=\ln 3:\text{ mínimo}\\
\text{{cóncava hacia arriba}}
\end{aligned}
\]


\paragraph{Ejercicio 178.} Para $f(x)=x^2\ln x$, determina crecimiento, puntos críticos y concavidad.

\[
\begin{aligned}
D_f=(0,\infty)\\
f'(x)=2x\ln x+x=x(2\ln x+1)\\
2\ln x+1=0\Rightarrow x=e^{-1/2}\\
f''(x)=2\ln x+3\\
f''(x)=0\Rightarrow x=e^{-3/2}\\
(0,e^{-1/2})\downarrow,\quad(e^{-1/2},\infty)\uparrow\\
x=e^{-1/2}:\text{{ mínimo}}\\
(0,e^{-3/2})\text{{ cóncava abajo}},\quad(e^{-3/2},\infty)\text{{ cóncava arriba}}
\end{aligned}
\]


\paragraph{Ejercicio 179.} Para $f(x)=e^x-5x$, determina crecimiento, puntos críticos y concavidad.

\[
\begin{aligned}
f'(x)=e^x-5\\
e^x=5\Rightarrow x=\ln 5\\
f''(x)=e^x>0\\
(-\infty,\ln 5)\downarrow,\quad(\ln 5,\infty)\uparrow\\
x=\ln 5:\text{ mínimo}\\
\text{{cóncava hacia arriba}}
\end{aligned}
\]


\paragraph{Ejercicio 180.} Para $f(x)=x^2\ln x$, determina crecimiento, puntos críticos y concavidad.

\[
\begin{aligned}
D_f=(0,\infty)\\
f'(x)=2x\ln x+x=x(2\ln x+1)\\
2\ln x+1=0\Rightarrow x=e^{-1/2}\\
f''(x)=2\ln x+3\\
f''(x)=0\Rightarrow x=e^{-3/2}\\
(0,e^{-1/2})\downarrow,\quad(e^{-1/2},\infty)\uparrow\\
x=e^{-1/2}:\text{{ mínimo}}\\
(0,e^{-3/2})\text{{ cóncava abajo}},\quad(e^{-3/2},\infty)\text{{ cóncava arriba}}
\end{aligned}
\]


\newpage

\subsubsection*{Nivel 5 - Imposible}

\paragraph{Ejercicio 181.} Para $f(x)=\frac{\ln x}{x}$, determina crecimiento, puntos críticos y concavidad.

\[
\begin{aligned}
D_f=(0,\infty)\\
f'(x)=\frac{1-\ln x}{x^2}\\
1-\ln x=0\Rightarrow x=e\\
f''(x)=\frac{2\ln x-3}{x^3}\\
2\ln x-3=0\Rightarrow x=e^{3/2}\\
(0,e)\uparrow,\quad(e,\infty)\downarrow\\
x=e:\text{{ máximo}}\\
(0,e^{3/2})\text{{ cóncava abajo}},\quad(e^{3/2},\infty)\text{{ cóncava arriba}}
\end{aligned}
\]


\paragraph{Ejercicio 182.} Para $f(x)=x^2e^{-x}$, determina crecimiento, puntos críticos y concavidad.

\[
\begin{aligned}
f'(x)=e^{-x}x(2-x)\\
f'(x)=0\Rightarrow x=0,\ 2\\
f''(x)=e^{-x}(x^2-4x+2)\\
x^2-4x+2=0\Rightarrow x=2\pm\sqrt2\\
(-\infty,0)\downarrow,\quad(0,2)\uparrow,\quad(2,\infty)\downarrow\\
x=0:\text{{ mínimo}},\quad x=2:\text{{ máximo}}\\
(-\infty,2-\sqrt2)\text{{ cóncava arriba}}\\
(2-\sqrt2,2+\sqrt2)\text{{ cóncava abajo}}\\
(2+\sqrt2,\infty)\text{{ cóncava arriba}}
\end{aligned}
\]


\paragraph{Ejercicio 183.} Para $f(x)=\frac{x^2+1}{x-1}$, determina crecimiento, puntos críticos y concavidad.

\[
\begin{aligned}
D_f=\mathbb{R}\setminus\{1\}\\
f'(x)=\frac{2x(x-1)-(x^2+1)}{(x-1)^2}\\
f'(x)=\frac{x^2-2x-1}{(x-1)^2}\\
x^2-2x-1=0\Rightarrow x=1\pm\sqrt2\\
f''(x)=\frac{4}{(x-1)^3}\\
(-\infty,1-\sqrt2)\uparrow,\quad(1-\sqrt2,1)\downarrow\\
(1,1+\sqrt2)\downarrow,\quad(1+\sqrt2,\infty)\uparrow\\
(-\infty,1)\text{{ cóncava abajo}},\quad(1,\infty)\text{{ cóncava arriba}}
\end{aligned}
\]


\paragraph{Ejercicio 184.} Para $f(x)=x+\frac{1}{x}+\frac{1}{x^2}$, determina crecimiento, puntos críticos y concavidad.

\[
\begin{aligned}
D_f=\mathbb{R}\setminus\{0\}\\
f'(x)=1-\frac{1}{x^2}-\frac{2}{x^3}=\frac{x^3-x-2}{x^3}\\
x^3-x-2=(x-2)(x+1)^2\\
f'(x)=\frac{(x-2)(x+1)^2}{x^3}\\
f'(x)=0\Rightarrow x=-1,\ 2\\
f''(x)=\frac{2(x+3)}{x^4}\\
f''(x)=0\Rightarrow x=-3\\
(-\infty,0)\uparrow,\quad(0,2)\downarrow,\quad(2,\infty)\uparrow
\end{aligned}
\]


\paragraph{Ejercicio 185.} Para $f(x)=\ln(1+x^2)$, determina crecimiento, puntos críticos y concavidad.

\[
\begin{aligned}
D_f=\mathbb{R}\\
f'(x)=\frac{2x}{1+x^2}\\
f'(x)=0\Rightarrow x=0\\
f''(x)=\frac{2(1-x^2)}{(1+x^2)^2}\\
f''(x)=0\Rightarrow x=\pm1\\
(-\infty,0)\downarrow,\quad(0,\infty)\uparrow\\
x=0:\text{{ mínimo}}\\
(-1,1)\text{{ cóncava arriba}}\\
(-\infty,-1)\cup(1,\infty)\text{{ cóncava abajo}}
\end{aligned}
\]


\paragraph{Ejercicio 186.} Para $f(x)=\frac{\ln x}{x}$, determina crecimiento, puntos críticos y concavidad.

\[
\begin{aligned}
D_f=(0,\infty)\\
f'(x)=\frac{1-\ln x}{x^2}\\
1-\ln x=0\Rightarrow x=e\\
f''(x)=\frac{2\ln x-3}{x^3}\\
2\ln x-3=0\Rightarrow x=e^{3/2}\\
(0,e)\uparrow,\quad(e,\infty)\downarrow\\
x=e:\text{{ máximo}}\\
(0,e^{3/2})\text{{ cóncava abajo}},\quad(e^{3/2},\infty)\text{{ cóncava arriba}}
\end{aligned}
\]


\paragraph{Ejercicio 187.} Para $f(x)=x^2e^{-x}$, determina crecimiento, puntos críticos y concavidad.

\[
\begin{aligned}
f'(x)=e^{-x}x(2-x)\\
f'(x)=0\Rightarrow x=0,\ 2\\
f''(x)=e^{-x}(x^2-4x+2)\\
x^2-4x+2=0\Rightarrow x=2\pm\sqrt2\\
(-\infty,0)\downarrow,\quad(0,2)\uparrow,\quad(2,\infty)\downarrow\\
x=0:\text{{ mínimo}},\quad x=2:\text{{ máximo}}\\
(-\infty,2-\sqrt2)\text{{ cóncava arriba}}\\
(2-\sqrt2,2+\sqrt2)\text{{ cóncava abajo}}\\
(2+\sqrt2,\infty)\text{{ cóncava arriba}}
\end{aligned}
\]


\paragraph{Ejercicio 188.} Para $f(x)=\frac{x^2+1}{x-1}$, determina crecimiento, puntos críticos y concavidad.

\[
\begin{aligned}
D_f=\mathbb{R}\setminus\{1\}\\
f'(x)=\frac{2x(x-1)-(x^2+1)}{(x-1)^2}\\
f'(x)=\frac{x^2-2x-1}{(x-1)^2}\\
x^2-2x-1=0\Rightarrow x=1\pm\sqrt2\\
f''(x)=\frac{4}{(x-1)^3}\\
(-\infty,1-\sqrt2)\uparrow,\quad(1-\sqrt2,1)\downarrow\\
(1,1+\sqrt2)\downarrow,\quad(1+\sqrt2,\infty)\uparrow\\
(-\infty,1)\text{{ cóncava abajo}},\quad(1,\infty)\text{{ cóncava arriba}}
\end{aligned}
\]


\paragraph{Ejercicio 189.} Para $f(x)=x+\frac{1}{x}+\frac{1}{x^2}$, determina crecimiento, puntos críticos y concavidad.

\[
\begin{aligned}
D_f=\mathbb{R}\setminus\{0\}\\
f'(x)=1-\frac{1}{x^2}-\frac{2}{x^3}=\frac{x^3-x-2}{x^3}\\
x^3-x-2=(x-2)(x+1)^2\\
f'(x)=\frac{(x-2)(x+1)^2}{x^3}\\
f'(x)=0\Rightarrow x=-1,\ 2\\
f''(x)=\frac{2(x+3)}{x^4}\\
f''(x)=0\Rightarrow x=-3\\
(-\infty,0)\uparrow,\quad(0,2)\downarrow,\quad(2,\infty)\uparrow
\end{aligned}
\]


\paragraph{Ejercicio 190.} Para $f(x)=\ln(1+x^2)$, determina crecimiento, puntos críticos y concavidad.

\[
\begin{aligned}
D_f=\mathbb{R}\\
f'(x)=\frac{2x}{1+x^2}\\
f'(x)=0\Rightarrow x=0\\
f''(x)=\frac{2(1-x^2)}{(1+x^2)^2}\\
f''(x)=0\Rightarrow x=\pm1\\
(-\infty,0)\downarrow,\quad(0,\infty)\uparrow\\
x=0:\text{{ mínimo}}\\
(-1,1)\text{{ cóncava arriba}}\\
(-\infty,-1)\cup(1,\infty)\text{{ cóncava abajo}}
\end{aligned}
\]


\paragraph{Ejercicio 191.} Para $f(x)=\frac{\ln x}{x}$, determina crecimiento, puntos críticos y concavidad.

\[
\begin{aligned}
D_f=(0,\infty)\\
f'(x)=\frac{1-\ln x}{x^2}\\
1-\ln x=0\Rightarrow x=e\\
f''(x)=\frac{2\ln x-3}{x^3}\\
2\ln x-3=0\Rightarrow x=e^{3/2}\\
(0,e)\uparrow,\quad(e,\infty)\downarrow\\
x=e:\text{{ máximo}}\\
(0,e^{3/2})\text{{ cóncava abajo}},\quad(e^{3/2},\infty)\text{{ cóncava arriba}}
\end{aligned}
\]


\paragraph{Ejercicio 192.} Para $f(x)=x^2e^{-x}$, determina crecimiento, puntos críticos y concavidad.

\[
\begin{aligned}
f'(x)=e^{-x}x(2-x)\\
f'(x)=0\Rightarrow x=0,\ 2\\
f''(x)=e^{-x}(x^2-4x+2)\\
x^2-4x+2=0\Rightarrow x=2\pm\sqrt2\\
(-\infty,0)\downarrow,\quad(0,2)\uparrow,\quad(2,\infty)\downarrow\\
x=0:\text{{ mínimo}},\quad x=2:\text{{ máximo}}\\
(-\infty,2-\sqrt2)\text{{ cóncava arriba}}\\
(2-\sqrt2,2+\sqrt2)\text{{ cóncava abajo}}\\
(2+\sqrt2,\infty)\text{{ cóncava arriba}}
\end{aligned}
\]


\paragraph{Ejercicio 193.} Para $f(x)=\frac{x^2+1}{x-1}$, determina crecimiento, puntos críticos y concavidad.

\[
\begin{aligned}
D_f=\mathbb{R}\setminus\{1\}\\
f'(x)=\frac{2x(x-1)-(x^2+1)}{(x-1)^2}\\
f'(x)=\frac{x^2-2x-1}{(x-1)^2}\\
x^2-2x-1=0\Rightarrow x=1\pm\sqrt2\\
f''(x)=\frac{4}{(x-1)^3}\\
(-\infty,1-\sqrt2)\uparrow,\quad(1-\sqrt2,1)\downarrow\\
(1,1+\sqrt2)\downarrow,\quad(1+\sqrt2,\infty)\uparrow\\
(-\infty,1)\text{{ cóncava abajo}},\quad(1,\infty)\text{{ cóncava arriba}}
\end{aligned}
\]


\paragraph{Ejercicio 194.} Para $f(x)=x+\frac{1}{x}+\frac{1}{x^2}$, determina crecimiento, puntos críticos y concavidad.

\[
\begin{aligned}
D_f=\mathbb{R}\setminus\{0\}\\
f'(x)=1-\frac{1}{x^2}-\frac{2}{x^3}=\frac{x^3-x-2}{x^3}\\
x^3-x-2=(x-2)(x+1)^2\\
f'(x)=\frac{(x-2)(x+1)^2}{x^3}\\
f'(x)=0\Rightarrow x=-1,\ 2\\
f''(x)=\frac{2(x+3)}{x^4}\\
f''(x)=0\Rightarrow x=-3\\
(-\infty,0)\uparrow,\quad(0,2)\downarrow,\quad(2,\infty)\uparrow
\end{aligned}
\]


\paragraph{Ejercicio 195.} Para $f(x)=\ln(1+x^2)$, determina crecimiento, puntos críticos y concavidad.

\[
\begin{aligned}
D_f=\mathbb{R}\\
f'(x)=\frac{2x}{1+x^2}\\
f'(x)=0\Rightarrow x=0\\
f''(x)=\frac{2(1-x^2)}{(1+x^2)^2}\\
f''(x)=0\Rightarrow x=\pm1\\
(-\infty,0)\downarrow,\quad(0,\infty)\uparrow\\
x=0:\text{{ mínimo}}\\
(-1,1)\text{{ cóncava arriba}}\\
(-\infty,-1)\cup(1,\infty)\text{{ cóncava abajo}}
\end{aligned}
\]


\paragraph{Ejercicio 196.} Para $f(x)=\frac{\ln x}{x}$, determina crecimiento, puntos críticos y concavidad.

\[
\begin{aligned}
D_f=(0,\infty)\\
f'(x)=\frac{1-\ln x}{x^2}\\
1-\ln x=0\Rightarrow x=e\\
f''(x)=\frac{2\ln x-3}{x^3}\\
2\ln x-3=0\Rightarrow x=e^{3/2}\\
(0,e)\uparrow,\quad(e,\infty)\downarrow\\
x=e:\text{{ máximo}}\\
(0,e^{3/2})\text{{ cóncava abajo}},\quad(e^{3/2},\infty)\text{{ cóncava arriba}}
\end{aligned}
\]


\paragraph{Ejercicio 197.} Para $f(x)=x^2e^{-x}$, determina crecimiento, puntos críticos y concavidad.

\[
\begin{aligned}
f'(x)=e^{-x}x(2-x)\\
f'(x)=0\Rightarrow x=0,\ 2\\
f''(x)=e^{-x}(x^2-4x+2)\\
x^2-4x+2=0\Rightarrow x=2\pm\sqrt2\\
(-\infty,0)\downarrow,\quad(0,2)\uparrow,\quad(2,\infty)\downarrow\\
x=0:\text{{ mínimo}},\quad x=2:\text{{ máximo}}\\
(-\infty,2-\sqrt2)\text{{ cóncava arriba}}\\
(2-\sqrt2,2+\sqrt2)\text{{ cóncava abajo}}\\
(2+\sqrt2,\infty)\text{{ cóncava arriba}}
\end{aligned}
\]


\paragraph{Ejercicio 198.} Para $f(x)=\frac{x^2+1}{x-1}$, determina crecimiento, puntos críticos y concavidad.

\[
\begin{aligned}
D_f=\mathbb{R}\setminus\{1\}\\
f'(x)=\frac{2x(x-1)-(x^2+1)}{(x-1)^2}\\
f'(x)=\frac{x^2-2x-1}{(x-1)^2}\\
x^2-2x-1=0\Rightarrow x=1\pm\sqrt2\\
f''(x)=\frac{4}{(x-1)^3}\\
(-\infty,1-\sqrt2)\uparrow,\quad(1-\sqrt2,1)\downarrow\\
(1,1+\sqrt2)\downarrow,\quad(1+\sqrt2,\infty)\uparrow\\
(-\infty,1)\text{{ cóncava abajo}},\quad(1,\infty)\text{{ cóncava arriba}}
\end{aligned}
\]


\paragraph{Ejercicio 199.} Para $f(x)=x+\frac{1}{x}+\frac{1}{x^2}$, determina crecimiento, puntos críticos y concavidad.

\[
\begin{aligned}
D_f=\mathbb{R}\setminus\{0\}\\
f'(x)=1-\frac{1}{x^2}-\frac{2}{x^3}=\frac{x^3-x-2}{x^3}\\
x^3-x-2=(x-2)(x+1)^2\\
f'(x)=\frac{(x-2)(x+1)^2}{x^3}\\
f'(x)=0\Rightarrow x=-1,\ 2\\
f''(x)=\frac{2(x+3)}{x^4}\\
f''(x)=0\Rightarrow x=-3\\
(-\infty,0)\uparrow,\quad(0,2)\downarrow,\quad(2,\infty)\uparrow
\end{aligned}
\]


\paragraph{Ejercicio 200.} Para $f(x)=\ln(1+x^2)$, determina crecimiento, puntos críticos y concavidad.

\[
\begin{aligned}
D_f=\mathbb{R}\\
f'(x)=\frac{2x}{1+x^2}\\
f'(x)=0\Rightarrow x=0\\
f''(x)=\frac{2(1-x^2)}{(1+x^2)^2}\\
f''(x)=0\Rightarrow x=\pm1\\
(-\infty,0)\downarrow,\quad(0,\infty)\uparrow\\
x=0:\text{{ mínimo}}\\
(-1,1)\text{{ cóncava arriba}}\\
(-\infty,-1)\cup(1,\infty)\text{{ cóncava abajo}}
\end{aligned}
\]


\newpage

\subsection*{Límites de fracciones algebraicas y fracción sintética}

\addcontentsline{toc}{subsection}{Soluciones - Límites de fracciones algebraicas y fracción sintética}

\subsubsection*{Nivel 1 - Básico}

\paragraph{Ejercicio 201.} Calcula $\displaystyle\lim_{x\to -3}\frac{1x-2}{x+1}$. 

\[
\begin{aligned}
\frac{1x-2}{x+1}\Big|_{x=-3}=\frac{1(-3)-2}{-2}\\
=\frac{-5}{-2}
\end{aligned}
\]
\[\boxed{\frac{-5}{-2}}\]


\paragraph{Ejercicio 202.} Calcula $\displaystyle\lim_{x\to -2}\frac{2x-1}{x+2}$. 

\[
\begin{aligned}
\frac{2x-1}{x+2}\Big|_{x=-2}=\frac{2(-2)-1}{0}\\
=\frac{-5}{0}
\end{aligned}
\]
\[\boxed{\frac{-5}{0}}\]


\paragraph{Ejercicio 203.} Calcula $\displaystyle\lim_{x\to -1}\frac{3x+0}{x+3}$. 

\[
\begin{aligned}
\frac{3x+0}{x+3}\Big|_{x=-1}=\frac{3(-1)+0}{2}\\
=\frac{-3}{2}
\end{aligned}
\]
\[\boxed{\frac{-3}{2}}\]


\paragraph{Ejercicio 204.} Calcula $\displaystyle\lim_{x\to 2}\frac{4x+1}{x+1}$. 

\[
\begin{aligned}
\frac{4x+1}{x+1}\Big|_{x=2}=\frac{4(2)+1}{3}\\
=\frac{9}{3}
\end{aligned}
\]
\[\boxed{\frac{9}{3}}\]


\paragraph{Ejercicio 205.} Calcula $\displaystyle\lim_{x\to 1}\frac{1x+2}{x+2}$. 

\[
\begin{aligned}
\frac{1x+2}{x+2}\Big|_{x=1}=\frac{1(1)+2}{3}\\
=\frac{3}{3}
\end{aligned}
\]
\[\boxed{\frac{3}{3}}\]


\paragraph{Ejercicio 206.} Calcula $\displaystyle\lim_{x\to 2}\frac{2x-2}{x+3}$. 

\[
\begin{aligned}
\frac{2x-2}{x+3}\Big|_{x=2}=\frac{2(2)-2}{5}\\
=\frac{2}{5}
\end{aligned}
\]
\[\boxed{\frac{2}{5}}\]


\paragraph{Ejercicio 207.} Calcula $\displaystyle\lim_{x\to 3}\frac{3x-1}{x+1}$. 

\[
\begin{aligned}
\frac{3x-1}{x+1}\Big|_{x=3}=\frac{3(3)-1}{4}\\
=\frac{8}{4}
\end{aligned}
\]
\[\boxed{\frac{8}{4}}\]


\paragraph{Ejercicio 208.} Calcula $\displaystyle\lim_{x\to -3}\frac{4x+0}{x+2}$. 

\[
\begin{aligned}
\frac{4x+0}{x+2}\Big|_{x=-3}=\frac{4(-3)+0}{-1}\\
=\frac{-12}{-1}
\end{aligned}
\]
\[\boxed{\frac{-12}{-1}}\]


\paragraph{Ejercicio 209.} Calcula $\displaystyle\lim_{x\to -2}\frac{1x+1}{x+3}$. 

\[
\begin{aligned}
\frac{1x+1}{x+3}\Big|_{x=-2}=\frac{1(-2)+1}{1}\\
=\frac{-1}{1}
\end{aligned}
\]
\[\boxed{\frac{-1}{1}}\]


\paragraph{Ejercicio 210.} Calcula $\displaystyle\lim_{x\to -1}\frac{2x+2}{x+1}$. 

\[
\begin{aligned}
\frac{2x+2}{x+1}\Big|_{x=-1}=\frac{2(-1)+2}{0}\\
=\frac{0}{0}
\end{aligned}
\]
\[\boxed{\frac{0}{0}}\]


\paragraph{Ejercicio 211.} Calcula $\displaystyle\lim_{x\to 2}\frac{3x-2}{x+2}$. 

\[
\begin{aligned}
\frac{3x-2}{x+2}\Big|_{x=2}=\frac{3(2)-2}{4}\\
=\frac{4}{4}
\end{aligned}
\]
\[\boxed{\frac{4}{4}}\]


\paragraph{Ejercicio 212.} Calcula $\displaystyle\lim_{x\to 1}\frac{4x-1}{x+3}$. 

\[
\begin{aligned}
\frac{4x-1}{x+3}\Big|_{x=1}=\frac{4(1)-1}{4}\\
=\frac{3}{4}
\end{aligned}
\]
\[\boxed{\frac{3}{4}}\]


\paragraph{Ejercicio 213.} Calcula $\displaystyle\lim_{x\to 2}\frac{1x+0}{x+1}$. 

\[
\begin{aligned}
\frac{1x+0}{x+1}\Big|_{x=2}=\frac{1(2)+0}{3}\\
=\frac{2}{3}
\end{aligned}
\]
\[\boxed{\frac{2}{3}}\]


\paragraph{Ejercicio 214.} Calcula $\displaystyle\lim_{x\to 3}\frac{2x+1}{x+2}$. 

\[
\begin{aligned}
\frac{2x+1}{x+2}\Big|_{x=3}=\frac{2(3)+1}{5}\\
=\frac{7}{5}
\end{aligned}
\]
\[\boxed{\frac{7}{5}}\]


\paragraph{Ejercicio 215.} Calcula $\displaystyle\lim_{x\to -3}\frac{3x+2}{x+3}$. 

\[
\begin{aligned}
\frac{3x+2}{x+3}\Big|_{x=-3}=\frac{3(-3)+2}{0}\\
=\frac{-7}{0}
\end{aligned}
\]
\[\boxed{\frac{-7}{0}}\]


\paragraph{Ejercicio 216.} Calcula $\displaystyle\lim_{x\to -2}\frac{4x-2}{x+1}$. 

\[
\begin{aligned}
\frac{4x-2}{x+1}\Big|_{x=-2}=\frac{4(-2)-2}{-1}\\
=\frac{-10}{-1}
\end{aligned}
\]
\[\boxed{\frac{-10}{-1}}\]


\paragraph{Ejercicio 217.} Calcula $\displaystyle\lim_{x\to -1}\frac{1x-1}{x+2}$. 

\[
\begin{aligned}
\frac{1x-1}{x+2}\Big|_{x=-1}=\frac{1(-1)-1}{1}\\
=\frac{-2}{1}
\end{aligned}
\]
\[\boxed{\frac{-2}{1}}\]


\paragraph{Ejercicio 218.} Calcula $\displaystyle\lim_{x\to 2}\frac{2x+0}{x+3}$. 

\[
\begin{aligned}
\frac{2x+0}{x+3}\Big|_{x=2}=\frac{2(2)+0}{5}\\
=\frac{4}{5}
\end{aligned}
\]
\[\boxed{\frac{4}{5}}\]


\paragraph{Ejercicio 219.} Calcula $\displaystyle\lim_{x\to 1}\frac{3x+1}{x+1}$. 

\[
\begin{aligned}
\frac{3x+1}{x+1}\Big|_{x=1}=\frac{3(1)+1}{2}\\
=\frac{4}{2}
\end{aligned}
\]
\[\boxed{\frac{4}{2}}\]


\paragraph{Ejercicio 220.} Calcula $\displaystyle\lim_{x\to 2}\frac{4x+2}{x+2}$. 

\[
\begin{aligned}
\frac{4x+2}{x+2}\Big|_{x=2}=\frac{4(2)+2}{4}\\
=\frac{10}{4}
\end{aligned}
\]
\[\boxed{\frac{10}{4}}\]


\newpage

\subsubsection*{Nivel 2 - Intermedio bajo}

\paragraph{Ejercicio 221.} Calcula $\displaystyle\lim_{x\to -3}\frac{x^2--5x+6}{x--3}$. 

\[
\begin{aligned}
x^2--5x+6=(x--3)(x--2)\\
\frac{(x--3)(x--2)}{x--3}=x--2\\
\lim_{x\to -3}(x--2)=-1
\end{aligned}
\]
\[\boxed{-1}\]


\paragraph{Ejercicio 222.} Calcula $\displaystyle\lim_{x\to -2}\frac{x^2--2x+0}{x--2}$. 

\[
\begin{aligned}
x^2--2x+0=(x--2)(x-0)\\
\frac{(x--2)(x-0)}{x--2}=x-0\\
\lim_{x\to -2}(x-0)=-2
\end{aligned}
\]
\[\boxed{-2}\]


\paragraph{Ejercicio 223.} Calcula $\displaystyle\lim_{x\to -1}\frac{x^2-1x+-2}{x--1}$. 

\[
\begin{aligned}
x^2-1x+-2=(x--1)(x-2)\\
\frac{(x--1)(x-2)}{x--1}=x-2\\
\lim_{x\to -1}(x-2)=-3
\end{aligned}
\]
\[\boxed{-3}\]


\paragraph{Ejercicio 224.} Calcula $\displaystyle\lim_{x\to 2}\frac{x^2-8x+12}{x-2}$. 

\[
\begin{aligned}
x^2-8x+12=(x-2)(x-6)\\
\frac{(x-2)(x-6)}{x-2}=x-6\\
\lim_{x\to 2}(x-6)=-4
\end{aligned}
\]
\[\boxed{-4}\]


\paragraph{Ejercicio 225.} Calcula $\displaystyle\lim_{x\to 1}\frac{x^2-7x+6}{x-1}$. 

\[
\begin{aligned}
x^2-7x+6=(x-1)(x-6)\\
\frac{(x-1)(x-6)}{x-1}=x-6\\
\lim_{x\to 1}(x-6)=-5
\end{aligned}
\]
\[\boxed{-5}\]


\paragraph{Ejercicio 226.} Calcula $\displaystyle\lim_{x\to 2}\frac{x^2-5x+6}{x-2}$. 

\[
\begin{aligned}
x^2-5x+6=(x-2)(x-3)\\
\frac{(x-2)(x-3)}{x-2}=x-3\\
\lim_{x\to 2}(x-3)=-1
\end{aligned}
\]
\[\boxed{-1}\]


\paragraph{Ejercicio 227.} Calcula $\displaystyle\lim_{x\to 3}\frac{x^2-8x+15}{x-3}$. 

\[
\begin{aligned}
x^2-8x+15=(x-3)(x-5)\\
\frac{(x-3)(x-5)}{x-3}=x-5\\
\lim_{x\to 3}(x-5)=-2
\end{aligned}
\]
\[\boxed{-2}\]


\paragraph{Ejercicio 228.} Calcula $\displaystyle\lim_{x\to -3}\frac{x^2--3x+0}{x--3}$. 

\[
\begin{aligned}
x^2--3x+0=(x--3)(x-0)\\
\frac{(x--3)(x-0)}{x--3}=x-0\\
\lim_{x\to -3}(x-0)=-3
\end{aligned}
\]
\[\boxed{-3}\]


\paragraph{Ejercicio 229.} Calcula $\displaystyle\lim_{x\to -2}\frac{x^2-0x+-4}{x--2}$. 

\[
\begin{aligned}
x^2-0x+-4=(x--2)(x-2)\\
\frac{(x--2)(x-2)}{x--2}=x-2\\
\lim_{x\to -2}(x-2)=-4
\end{aligned}
\]
\[\boxed{-4}\]


\paragraph{Ejercicio 230.} Calcula $\displaystyle\lim_{x\to -1}\frac{x^2-3x+-4}{x--1}$. 

\[
\begin{aligned}
x^2-3x+-4=(x--1)(x-4)\\
\frac{(x--1)(x-4)}{x--1}=x-4\\
\lim_{x\to -1}(x-4)=-5
\end{aligned}
\]
\[\boxed{-5}\]


\paragraph{Ejercicio 231.} Calcula $\displaystyle\lim_{x\to 2}\frac{x^2-5x+6}{x-2}$. 

\[
\begin{aligned}
x^2-5x+6=(x-2)(x-3)\\
\frac{(x-2)(x-3)}{x-2}=x-3\\
\lim_{x\to 2}(x-3)=-1
\end{aligned}
\]
\[\boxed{-1}\]


\paragraph{Ejercicio 232.} Calcula $\displaystyle\lim_{x\to 1}\frac{x^2-4x+3}{x-1}$. 

\[
\begin{aligned}
x^2-4x+3=(x-1)(x-3)\\
\frac{(x-1)(x-3)}{x-1}=x-3\\
\lim_{x\to 1}(x-3)=-2
\end{aligned}
\]
\[\boxed{-2}\]


\paragraph{Ejercicio 233.} Calcula $\displaystyle\lim_{x\to 2}\frac{x^2-7x+10}{x-2}$. 

\[
\begin{aligned}
x^2-7x+10=(x-2)(x-5)\\
\frac{(x-2)(x-5)}{x-2}=x-5\\
\lim_{x\to 2}(x-5)=-3
\end{aligned}
\]
\[\boxed{-3}\]


\paragraph{Ejercicio 234.} Calcula $\displaystyle\lim_{x\to 3}\frac{x^2-10x+21}{x-3}$. 

\[
\begin{aligned}
x^2-10x+21=(x-3)(x-7)\\
\frac{(x-3)(x-7)}{x-3}=x-7\\
\lim_{x\to 3}(x-7)=-4
\end{aligned}
\]
\[\boxed{-4}\]


\paragraph{Ejercicio 235.} Calcula $\displaystyle\lim_{x\to -3}\frac{x^2--1x+-6}{x--3}$. 

\[
\begin{aligned}
x^2--1x+-6=(x--3)(x-2)\\
\frac{(x--3)(x-2)}{x--3}=x-2\\
\lim_{x\to -3}(x-2)=-5
\end{aligned}
\]
\[\boxed{-5}\]


\paragraph{Ejercicio 236.} Calcula $\displaystyle\lim_{x\to -2}\frac{x^2--3x+2}{x--2}$. 

\[
\begin{aligned}
x^2--3x+2=(x--2)(x--1)\\
\frac{(x--2)(x--1)}{x--2}=x--1\\
\lim_{x\to -2}(x--1)=-1
\end{aligned}
\]
\[\boxed{-1}\]


\paragraph{Ejercicio 237.} Calcula $\displaystyle\lim_{x\to -1}\frac{x^2-0x+-1}{x--1}$. 

\[
\begin{aligned}
x^2-0x+-1=(x--1)(x-1)\\
\frac{(x--1)(x-1)}{x--1}=x-1\\
\lim_{x\to -1}(x-1)=-2
\end{aligned}
\]
\[\boxed{-2}\]


\paragraph{Ejercicio 238.} Calcula $\displaystyle\lim_{x\to 2}\frac{x^2-7x+10}{x-2}$. 

\[
\begin{aligned}
x^2-7x+10=(x-2)(x-5)\\
\frac{(x-2)(x-5)}{x-2}=x-5\\
\lim_{x\to 2}(x-5)=-3
\end{aligned}
\]
\[\boxed{-3}\]


\paragraph{Ejercicio 239.} Calcula $\displaystyle\lim_{x\to 1}\frac{x^2-6x+5}{x-1}$. 

\[
\begin{aligned}
x^2-6x+5=(x-1)(x-5)\\
\frac{(x-1)(x-5)}{x-1}=x-5\\
\lim_{x\to 1}(x-5)=-4
\end{aligned}
\]
\[\boxed{-4}\]


\paragraph{Ejercicio 240.} Calcula $\displaystyle\lim_{x\to 2}\frac{x^2-9x+14}{x-2}$. 

\[
\begin{aligned}
x^2-9x+14=(x-2)(x-7)\\
\frac{(x-2)(x-7)}{x-2}=x-7\\
\lim_{x\to 2}(x-7)=-5
\end{aligned}
\]
\[\boxed{-5}\]


\newpage

\subsubsection*{Nivel 3 - Intermedio}

\paragraph{Ejercicio 241.} Calcula usando fracción sintética $\displaystyle\lim_{x\to -4}\frac{x^3 + x^2 - 11x + 4}{x--4}$. 

\[
\begin{aligned}
P(-4)=0\\
\text{Sintética con }-4:\quad P(x)=(x--4)(x^2 - 3x + 1)\\
\frac{x^3 + x^2 - 11x + 4}{x--4}=x^2 - 3x + 1\\
\lim_{x\to -4}(x^2 - 3x + 1)=29
\end{aligned}
\]
\[\boxed{29}\]


\paragraph{Ejercicio 242.} Calcula usando fracción sintética $\displaystyle\lim_{x\to -3}\frac{x^3 + x^2 - 4x + 6}{x--3}$. 

\[
\begin{aligned}
P(-3)=0\\
\text{Sintética con }-3:\quad P(x)=(x--3)(x^2 - 2x + 2)\\
\frac{x^3 + x^2 - 4x + 6}{x--3}=x^2 - 2x + 2\\
\lim_{x\to -3}(x^2 - 2x + 2)=17
\end{aligned}
\]
\[\boxed{17}\]


\paragraph{Ejercicio 243.} Calcula usando fracción sintética $\displaystyle\lim_{x\to -2}\frac{x^3 + x^2 + x + 6}{x--2}$. 

\[
\begin{aligned}
P(-2)=0\\
\text{Sintética con }-2:\quad P(x)=(x--2)(x^2 - x + 3)\\
\frac{x^3 + x^2 + x + 6}{x--2}=x^2 - x + 3\\
\lim_{x\to -2}(x^2 - x + 3)=9
\end{aligned}
\]
\[\boxed{9}\]


\paragraph{Ejercicio 244.} Calcula usando fracción sintética $\displaystyle\lim_{x\to -1}\frac{x^3 + x^2 + 4x + 4}{x--1}$. 

\[
\begin{aligned}
P(-1)=0\\
\text{Sintética con }-1:\quad P(x)=(x--1)(x^2 + 4)\\
\frac{x^3 + x^2 + 4x + 4}{x--1}=x^2 + 4\\
\lim_{x\to -1}(x^2 + 4)=5
\end{aligned}
\]
\[\boxed{5}\]


\paragraph{Ejercicio 245.} Calcula usando fracción sintética $\displaystyle\lim_{x\to 3}\frac{x^3 - 2x^2 + 2x - 15}{x-3}$. 

\[
\begin{aligned}
P(3)=0\\
\text{Sintética con }3:\quad P(x)=(x-3)(x^2 + x + 5)\\
\frac{x^3 - 2x^2 + 2x - 15}{x-3}=x^2 + x + 5\\
\lim_{x\to 3}(x^2 + x + 5)=17
\end{aligned}
\]
\[\boxed{17}\]


\paragraph{Ejercicio 246.} Calcula usando fracción sintética $\displaystyle\lim_{x\to 1}\frac{x^3 + x^2 - x - 1}{x-1}$. 

\[
\begin{aligned}
P(1)=0\\
\text{Sintética con }1:\quad P(x)=(x-1)(x^2 + 2x + 1)\\
\frac{x^3 + x^2 - x - 1}{x-1}=x^2 + 2x + 1\\
\lim_{x\to 1}(x^2 + 2x + 1)=4
\end{aligned}
\]
\[\boxed{4}\]


\paragraph{Ejercicio 247.} Calcula usando fracción sintética $\displaystyle\lim_{x\to 2}\frac{x^3 + x^2 - 4x - 4}{x-2}$. 

\[
\begin{aligned}
P(2)=0\\
\text{Sintética con }2:\quad P(x)=(x-2)(x^2 + 3x + 2)\\
\frac{x^3 + x^2 - 4x - 4}{x-2}=x^2 + 3x + 2\\
\lim_{x\to 2}(x^2 + 3x + 2)=12
\end{aligned}
\]
\[\boxed{12}\]


\paragraph{Ejercicio 248.} Calcula usando fracción sintética $\displaystyle\lim_{x\to 3}\frac{x^3 - 6x^2 + 12x - 9}{x-3}$. 

\[
\begin{aligned}
P(3)=0\\
\text{Sintética con }3:\quad P(x)=(x-3)(x^2 - 3x + 3)\\
\frac{x^3 - 6x^2 + 12x - 9}{x-3}=x^2 - 3x + 3\\
\lim_{x\to 3}(x^2 - 3x + 3)=3
\end{aligned}
\]
\[\boxed{3}\]


\paragraph{Ejercicio 249.} Calcula usando fracción sintética $\displaystyle\lim_{x\to 4}\frac{x^3 - 6x^2 + 12x - 16}{x-4}$. 

\[
\begin{aligned}
P(4)=0\\
\text{Sintética con }4:\quad P(x)=(x-4)(x^2 - 2x + 4)\\
\frac{x^3 - 6x^2 + 12x - 16}{x-4}=x^2 - 2x + 4\\
\lim_{x\to 4}(x^2 - 2x + 4)=12
\end{aligned}
\]
\[\boxed{12}\]


\paragraph{Ejercicio 250.} Calcula usando fracción sintética $\displaystyle\lim_{x\to -4}\frac{x^3 + 3x^2 + x + 20}{x--4}$. 

\[
\begin{aligned}
P(-4)=0\\
\text{Sintética con }-4:\quad P(x)=(x--4)(x^2 - x + 5)\\
\frac{x^3 + 3x^2 + x + 20}{x--4}=x^2 - x + 5\\
\lim_{x\to -4}(x^2 - x + 5)=25
\end{aligned}
\]
\[\boxed{25}\]


\paragraph{Ejercicio 251.} Calcula usando fracción sintética $\displaystyle\lim_{x\to -3}\frac{x^3 + 3x^2 + x + 3}{x--3}$. 

\[
\begin{aligned}
P(-3)=0\\
\text{Sintética con }-3:\quad P(x)=(x--3)(x^2 + 1)\\
\frac{x^3 + 3x^2 + x + 3}{x--3}=x^2 + 1\\
\lim_{x\to -3}(x^2 + 1)=10
\end{aligned}
\]
\[\boxed{10}\]


\paragraph{Ejercicio 252.} Calcula usando fracción sintética $\displaystyle\lim_{x\to -2}\frac{x^3 + 3x^2 + 4x + 4}{x--2}$. 

\[
\begin{aligned}
P(-2)=0\\
\text{Sintética con }-2:\quad P(x)=(x--2)(x^2 + x + 2)\\
\frac{x^3 + 3x^2 + 4x + 4}{x--2}=x^2 + x + 2\\
\lim_{x\to -2}(x^2 + x + 2)=4
\end{aligned}
\]
\[\boxed{4}\]


\paragraph{Ejercicio 253.} Calcula usando fracción sintética $\displaystyle\lim_{x\to -1}\frac{x^3 + 3x^2 + 5x + 3}{x--1}$. 

\[
\begin{aligned}
P(-1)=0\\
\text{Sintética con }-1:\quad P(x)=(x--1)(x^2 + 2x + 3)\\
\frac{x^3 + 3x^2 + 5x + 3}{x--1}=x^2 + 2x + 3\\
\lim_{x\to -1}(x^2 + 2x + 3)=2
\end{aligned}
\]
\[\boxed{2}\]


\paragraph{Ejercicio 254.} Calcula usando fracción sintética $\displaystyle\lim_{x\to 3}\frac{x^3 - 5x - 12}{x-3}$. 

\[
\begin{aligned}
P(3)=0\\
\text{Sintética con }3:\quad P(x)=(x-3)(x^2 + 3x + 4)\\
\frac{x^3 - 5x - 12}{x-3}=x^2 + 3x + 4\\
\lim_{x\to 3}(x^2 + 3x + 4)=22
\end{aligned}
\]
\[\boxed{22}\]


\paragraph{Ejercicio 255.} Calcula usando fracción sintética $\displaystyle\lim_{x\to 1}\frac{x^3 - 4x^2 + 8x - 5}{x-1}$. 

\[
\begin{aligned}
P(1)=0\\
\text{Sintética con }1:\quad P(x)=(x-1)(x^2 - 3x + 5)\\
\frac{x^3 - 4x^2 + 8x - 5}{x-1}=x^2 - 3x + 5\\
\lim_{x\to 1}(x^2 - 3x + 5)=3
\end{aligned}
\]
\[\boxed{3}\]


\paragraph{Ejercicio 256.} Calcula usando fracción sintética $\displaystyle\lim_{x\to 2}\frac{x^3 - 4x^2 + 5x - 2}{x-2}$. 

\[
\begin{aligned}
P(2)=0\\
\text{Sintética con }2:\quad P(x)=(x-2)(x^2 - 2x + 1)\\
\frac{x^3 - 4x^2 + 5x - 2}{x-2}=x^2 - 2x + 1\\
\lim_{x\to 2}(x^2 - 2x + 1)=1
\end{aligned}
\]
\[\boxed{1}\]


\paragraph{Ejercicio 257.} Calcula usando fracción sintética $\displaystyle\lim_{x\to 3}\frac{x^3 - 4x^2 + 5x - 6}{x-3}$. 

\[
\begin{aligned}
P(3)=0\\
\text{Sintética con }3:\quad P(x)=(x-3)(x^2 - x + 2)\\
\frac{x^3 - 4x^2 + 5x - 6}{x-3}=x^2 - x + 2\\
\lim_{x\to 3}(x^2 - x + 2)=8
\end{aligned}
\]
\[\boxed{8}\]


\paragraph{Ejercicio 258.} Calcula usando fracción sintética $\displaystyle\lim_{x\to 4}\frac{x^3 - 4x^2 + 3x - 12}{x-4}$. 

\[
\begin{aligned}
P(4)=0\\
\text{Sintética con }4:\quad P(x)=(x-4)(x^2 + 3)\\
\frac{x^3 - 4x^2 + 3x - 12}{x-4}=x^2 + 3\\
\lim_{x\to 4}(x^2 + 3)=19
\end{aligned}
\]
\[\boxed{19}\]


\paragraph{Ejercicio 259.} Calcula usando fracción sintética $\displaystyle\lim_{x\to -4}\frac{x^3 + 5x^2 + 8x + 16}{x--4}$. 

\[
\begin{aligned}
P(-4)=0\\
\text{Sintética con }-4:\quad P(x)=(x--4)(x^2 + x + 4)\\
\frac{x^3 + 5x^2 + 8x + 16}{x--4}=x^2 + x + 4\\
\lim_{x\to -4}(x^2 + x + 4)=16
\end{aligned}
\]
\[\boxed{16}\]


\paragraph{Ejercicio 260.} Calcula usando fracción sintética $\displaystyle\lim_{x\to -3}\frac{x^3 + 5x^2 + 11x + 15}{x--3}$. 

\[
\begin{aligned}
P(-3)=0\\
\text{Sintética con }-3:\quad P(x)=(x--3)(x^2 + 2x + 5)\\
\frac{x^3 + 5x^2 + 11x + 15}{x--3}=x^2 + 2x + 5\\
\lim_{x\to -3}(x^2 + 2x + 5)=8
\end{aligned}
\]
\[\boxed{8}\]


\newpage

\subsubsection*{Nivel 4 - Avanzado}

\paragraph{Ejercicio 261.} Calcula $\displaystyle\lim_{x\to\infty}\frac{1x^3+2x^2-1}{2x^3-1x+1}$. 

\[
\begin{aligned}
\frac{1x^3+2x^2-1}{2x^3-1x+1}=\frac{x^3(1+\frac{2}{x}-\frac{1}{x^3})}{x^3(2-\frac{1}{x^2}+\frac{1}{x^3})}\\
\lim=\frac{1}{2}
\end{aligned}
\]
\[\boxed{\frac{1}{2}}\]


\paragraph{Ejercicio 262.} Calcula $\displaystyle\lim_{x\to -2}\frac{(x--2)(x-0)(x-2)}{(x--2)(x+3)}$. 

\[
\begin{aligned}
\frac{(x--2)(x-0)(x-2)}{(x--2)(x+3)}=\frac{(x-0)(x-2)}{x+3}\\
\lim=\frac{(-2)(-4)}{1}=\frac{8}{1}
\end{aligned}
\]
\[\boxed{\frac{8}{1}}\]


\paragraph{Ejercicio 263.} Calcula $\displaystyle\lim_{x\to\infty}\frac{3x^3+4x^2-3}{4x^3-3x+3}$. 

\[
\begin{aligned}
\frac{3x^3+4x^2-3}{4x^3-3x+3}=\frac{x^3(3+\frac{4}{x}-\frac{3}{x^3})}{x^3(4-\frac{3}{x^2}+\frac{3}{x^3})}\\
\lim=\frac{3}{4}
\end{aligned}
\]
\[\boxed{\frac{3}{4}}\]


\paragraph{Ejercicio 264.} Calcula $\displaystyle\lim_{x\to 2}\frac{(x-2)(x-4)(x-6)}{(x-2)(x+3)}$. 

\[
\begin{aligned}
\frac{(x-2)(x-4)(x-6)}{(x-2)(x+3)}=\frac{(x-4)(x-6)}{x+3}\\
\lim=\frac{(-2)(-4)}{5}=\frac{8}{5}
\end{aligned}
\]
\[\boxed{\frac{8}{5}}\]


\paragraph{Ejercicio 265.} Calcula $\displaystyle\lim_{x\to\infty}\frac{5x^3+2x^2-2}{2x^3-5x+2}$. 

\[
\begin{aligned}
\frac{5x^3+2x^2-2}{2x^3-5x+2}=\frac{x^3(5+\frac{2}{x}-\frac{2}{x^3})}{x^3(2-\frac{5}{x^2}+\frac{2}{x^3})}\\
\lim=\frac{5}{2}
\end{aligned}
\]
\[\boxed{\frac{5}{2}}\]


\paragraph{Ejercicio 266.} Calcula $\displaystyle\lim_{x\to 2}\frac{(x-2)(x-4)(x-6)}{(x-2)(x+3)}$. 

\[
\begin{aligned}
\frac{(x-2)(x-4)(x-6)}{(x-2)(x+3)}=\frac{(x-4)(x-6)}{x+3}\\
\lim=\frac{(-2)(-4)}{5}=\frac{8}{5}
\end{aligned}
\]
\[\boxed{\frac{8}{5}}\]


\paragraph{Ejercicio 267.} Calcula $\displaystyle\lim_{x\to\infty}\frac{2x^3+4x^2-1}{4x^3-2x+1}$. 

\[
\begin{aligned}
\frac{2x^3+4x^2-1}{4x^3-2x+1}=\frac{x^3(2+\frac{4}{x}-\frac{1}{x^3})}{x^3(4-\frac{2}{x^2}+\frac{1}{x^3})}\\
\lim=\frac{2}{4}
\end{aligned}
\]
\[\boxed{\frac{2}{4}}\]


\paragraph{Ejercicio 268.} Calcula $\displaystyle\lim_{x\to -3}\frac{(x--3)(x--1)(x-1)}{(x--3)(x+4)}$. 

\[
\begin{aligned}
\frac{(x--3)(x--1)(x-1)}{(x--3)(x+4)}=\frac{(x--1)(x-1)}{x+4}\\
\lim=\frac{(-2)(-4)}{1}=\frac{8}{1}
\end{aligned}
\]
\[\boxed{\frac{8}{1}}\]


\paragraph{Ejercicio 269.} Calcula $\displaystyle\lim_{x\to\infty}\frac{4x^3+2x^2-3}{2x^3-4x+3}$. 

\[
\begin{aligned}
\frac{4x^3+2x^2-3}{2x^3-4x+3}=\frac{x^3(4+\frac{2}{x}-\frac{3}{x^3})}{x^3(2-\frac{4}{x^2}+\frac{3}{x^3})}\\
\lim=\frac{4}{2}
\end{aligned}
\]
\[\boxed{\frac{4}{2}}\]


\paragraph{Ejercicio 270.} Calcula $\displaystyle\lim_{x\to -1}\frac{(x--1)(x-1)(x-3)}{(x--1)(x+2)}$. 

\[
\begin{aligned}
\frac{(x--1)(x-1)(x-3)}{(x--1)(x+2)}=\frac{(x-1)(x-3)}{x+2}\\
\lim=\frac{(-2)(-4)}{1}=\frac{8}{1}
\end{aligned}
\]
\[\boxed{\frac{8}{1}}\]


\paragraph{Ejercicio 271.} Calcula $\displaystyle\lim_{x\to\infty}\frac{1x^3+4x^2-2}{4x^3-1x+2}$. 

\[
\begin{aligned}
\frac{1x^3+4x^2-2}{4x^3-1x+2}=\frac{x^3(1+\frac{4}{x}-\frac{2}{x^3})}{x^3(4-\frac{1}{x^2}+\frac{2}{x^3})}\\
\lim=\frac{1}{4}
\end{aligned}
\]
\[\boxed{\frac{1}{4}}\]


\paragraph{Ejercicio 272.} Calcula $\displaystyle\lim_{x\to 1}\frac{(x-1)(x-3)(x-5)}{(x-1)(x+2)}$. 

\[
\begin{aligned}
\frac{(x-1)(x-3)(x-5)}{(x-1)(x+2)}=\frac{(x-3)(x-5)}{x+2}\\
\lim=\frac{(-2)(-4)}{3}=\frac{8}{3}
\end{aligned}
\]
\[\boxed{\frac{8}{3}}\]


\paragraph{Ejercicio 273.} Calcula $\displaystyle\lim_{x\to\infty}\frac{3x^3+2x^2-1}{2x^3-3x+1}$. 

\[
\begin{aligned}
\frac{3x^3+2x^2-1}{2x^3-3x+1}=\frac{x^3(3+\frac{2}{x}-\frac{1}{x^3})}{x^3(2-\frac{3}{x^2}+\frac{1}{x^3})}\\
\lim=\frac{3}{2}
\end{aligned}
\]
\[\boxed{\frac{3}{2}}\]


\paragraph{Ejercicio 274.} Calcula $\displaystyle\lim_{x\to 3}\frac{(x-3)(x-5)(x-7)}{(x-3)(x+4)}$. 

\[
\begin{aligned}
\frac{(x-3)(x-5)(x-7)}{(x-3)(x+4)}=\frac{(x-5)(x-7)}{x+4}\\
\lim=\frac{(-2)(-4)}{7}=\frac{8}{7}
\end{aligned}
\]
\[\boxed{\frac{8}{7}}\]


\paragraph{Ejercicio 275.} Calcula $\displaystyle\lim_{x\to\infty}\frac{5x^3+4x^2-3}{4x^3-5x+3}$. 

\[
\begin{aligned}
\frac{5x^3+4x^2-3}{4x^3-5x+3}=\frac{x^3(5+\frac{4}{x}-\frac{3}{x^3})}{x^3(4-\frac{5}{x^2}+\frac{3}{x^3})}\\
\lim=\frac{5}{4}
\end{aligned}
\]
\[\boxed{\frac{5}{4}}\]


\paragraph{Ejercicio 276.} Calcula $\displaystyle\lim_{x\to -2}\frac{(x--2)(x-0)(x-2)}{(x--2)(x+3)}$. 

\[
\begin{aligned}
\frac{(x--2)(x-0)(x-2)}{(x--2)(x+3)}=\frac{(x-0)(x-2)}{x+3}\\
\lim=\frac{(-2)(-4)}{1}=\frac{8}{1}
\end{aligned}
\]
\[\boxed{\frac{8}{1}}\]


\paragraph{Ejercicio 277.} Calcula $\displaystyle\lim_{x\to\infty}\frac{2x^3+2x^2-2}{2x^3-2x+2}$. 

\[
\begin{aligned}
\frac{2x^3+2x^2-2}{2x^3-2x+2}=\frac{x^3(2+\frac{2}{x}-\frac{2}{x^3})}{x^3(2-\frac{2}{x^2}+\frac{2}{x^3})}\\
\lim=\frac{2}{2}
\end{aligned}
\]
\[\boxed{\frac{2}{2}}\]


\paragraph{Ejercicio 278.} Calcula $\displaystyle\lim_{x\to 2}\frac{(x-2)(x-4)(x-6)}{(x-2)(x+3)}$. 

\[
\begin{aligned}
\frac{(x-2)(x-4)(x-6)}{(x-2)(x+3)}=\frac{(x-4)(x-6)}{x+3}\\
\lim=\frac{(-2)(-4)}{5}=\frac{8}{5}
\end{aligned}
\]
\[\boxed{\frac{8}{5}}\]


\paragraph{Ejercicio 279.} Calcula $\displaystyle\lim_{x\to\infty}\frac{4x^3+4x^2-1}{4x^3-4x+1}$. 

\[
\begin{aligned}
\frac{4x^3+4x^2-1}{4x^3-4x+1}=\frac{x^3(4+\frac{4}{x}-\frac{1}{x^3})}{x^3(4-\frac{4}{x^2}+\frac{1}{x^3})}\\
\lim=\frac{4}{4}
\end{aligned}
\]
\[\boxed{\frac{4}{4}}\]


\paragraph{Ejercicio 280.} Calcula $\displaystyle\lim_{x\to 2}\frac{(x-2)(x-4)(x-6)}{(x-2)(x+3)}$. 

\[
\begin{aligned}
\frac{(x-2)(x-4)(x-6)}{(x-2)(x+3)}=\frac{(x-4)(x-6)}{x+3}\\
\lim=\frac{(-2)(-4)}{5}=\frac{8}{5}
\end{aligned}
\]
\[\boxed{\frac{8}{5}}\]


\newpage

\subsubsection*{Nivel 5 - Imposible}

\paragraph{Ejercicio 281.} Calcula usando fracción sintética $\displaystyle\lim_{x\to -4}\frac{x^4 + 2x^3 - 7x^2 + 6x + 8}{x--4}$. 

\[
\begin{aligned}
P(-4)=0\\
P(x)=(x--4)(x^3 - 2x^2 + x + 2)\\
\frac{x^4 + 2x^3 - 7x^2 + 6x + 8}{x--4}=x^3 - 2x^2 + x + 2\\
\lim_{x\to -4}(x^3 - 2x^2 + x + 2)=-98
\end{aligned}
\]
\[\boxed{-98}\]


\paragraph{Ejercicio 282.} Calcula usando fracción sintética $\displaystyle\lim_{x\to -3}\frac{x^4 + 2x^3 - x^2 + 9x + 9}{x--3}$. 

\[
\begin{aligned}
P(-3)=0\\
P(x)=(x--3)(x^3 - x^2 + 2x + 3)\\
\frac{x^4 + 2x^3 - x^2 + 9x + 9}{x--3}=x^3 - x^2 + 2x + 3\\
\lim_{x\to -3}(x^3 - x^2 + 2x + 3)=-39
\end{aligned}
\]
\[\boxed{-39}\]


\paragraph{Ejercicio 283.} Calcula usando fracción sintética $\displaystyle\lim_{x\to -2}\frac{x^4 + 2x^3 + 3x^2 + 10x + 8}{x--2}$. 

\[
\begin{aligned}
P(-2)=0\\
P(x)=(x--2)(x^3 + 3x + 4)\\
\frac{x^4 + 2x^3 + 3x^2 + 10x + 8}{x--2}=x^3 + 3x + 4\\
\lim_{x\to -2}(x^3 + 3x + 4)=-10
\end{aligned}
\]
\[\boxed{-10}\]


\paragraph{Ejercicio 284.} Calcula usando fracción sintética $\displaystyle\lim_{x\to -1}\frac{x^4 + 2x^3 + 5x^2 + 9x + 5}{x--1}$. 

\[
\begin{aligned}
P(-1)=0\\
P(x)=(x--1)(x^3 + x^2 + 4x + 5)\\
\frac{x^4 + 2x^3 + 5x^2 + 9x + 5}{x--1}=x^3 + x^2 + 4x + 5\\
\lim_{x\to -1}(x^3 + x^2 + 4x + 5)=1
\end{aligned}
\]
\[\boxed{1}\]


\paragraph{Ejercicio 285.} Calcula usando fracción sintética $\displaystyle\lim_{x\to 5}\frac{x^4 - 3x^3 - 5x^2 - 23x - 10}{x-5}$. 

\[
\begin{aligned}
P(5)=0\\
P(x)=(x-5)(x^3 + 2x^2 + 5x + 2)\\
\frac{x^4 - 3x^3 - 5x^2 - 23x - 10}{x-5}=x^3 + 2x^2 + 5x + 2\\
\lim_{x\to 5}(x^3 + 2x^2 + 5x + 2)=202
\end{aligned}
\]
\[\boxed{202}\]


\paragraph{Ejercicio 286.} Calcula usando fracción sintética $\displaystyle\lim_{x\to 1}\frac{x^4 - 3x^3 + 8x^2 - 3x - 3}{x-1}$. 

\[
\begin{aligned}
P(1)=0\\
P(x)=(x-1)(x^3 - 2x^2 + 6x + 3)\\
\frac{x^4 - 3x^3 + 8x^2 - 3x - 3}{x-1}=x^3 - 2x^2 + 6x + 3\\
\lim_{x\to 1}(x^3 - 2x^2 + 6x + 3)=8
\end{aligned}
\]
\[\boxed{8}\]


\paragraph{Ejercicio 287.} Calcula usando fracción sintética $\displaystyle\lim_{x\to 2}\frac{x^4 - 3x^3 + 9x^2 - 10x - 8}{x-2}$. 

\[
\begin{aligned}
P(2)=0\\
P(x)=(x-2)(x^3 - x^2 + 7x + 4)\\
\frac{x^4 - 3x^3 + 9x^2 - 10x - 8}{x-2}=x^3 - x^2 + 7x + 4\\
\lim_{x\to 2}(x^3 - x^2 + 7x + 4)=22
\end{aligned}
\]
\[\boxed{22}\]


\paragraph{Ejercicio 288.} Calcula usando fracción sintética $\displaystyle\lim_{x\to 3}\frac{x^4 - 3x^3 + x^2 + 2x - 15}{x-3}$. 

\[
\begin{aligned}
P(3)=0\\
P(x)=(x-3)(x^3 + x + 5)\\
\frac{x^4 - 3x^3 + x^2 + 2x - 15}{x-3}=x^3 + x + 5\\
\lim_{x\to 3}(x^3 + x + 5)=35
\end{aligned}
\]
\[\boxed{35}\]


\paragraph{Ejercicio 289.} Calcula usando fracción sintética $\displaystyle\lim_{x\to 4}\frac{x^4 - 3x^3 - 2x^2 - 6x - 8}{x-4}$. 

\[
\begin{aligned}
P(4)=0\\
P(x)=(x-4)(x^3 + x^2 + 2x + 2)\\
\frac{x^4 - 3x^3 - 2x^2 - 6x - 8}{x-4}=x^3 + x^2 + 2x + 2\\
\lim_{x\to 4}(x^3 + x^2 + 2x + 2)=90
\end{aligned}
\]
\[\boxed{90}\]


\paragraph{Ejercicio 290.} Calcula usando fracción sintética $\displaystyle\lim_{x\to -4}\frac{x^4 + 6x^3 + 11x^2 + 15x + 12}{x--4}$. 

\[
\begin{aligned}
P(-4)=0\\
P(x)=(x--4)(x^3 + 2x^2 + 3x + 3)\\
\frac{x^4 + 6x^3 + 11x^2 + 15x + 12}{x--4}=x^3 + 2x^2 + 3x + 3\\
\lim_{x\to -4}(x^3 + 2x^2 + 3x + 3)=-41
\end{aligned}
\]
\[\boxed{-41}\]


\paragraph{Ejercicio 291.} Calcula usando fracción sintética $\displaystyle\lim_{x\to -3}\frac{x^4 + x^3 - 2x^2 + 16x + 12}{x--3}$. 

\[
\begin{aligned}
P(-3)=0\\
P(x)=(x--3)(x^3 - 2x^2 + 4x + 4)\\
\frac{x^4 + x^3 - 2x^2 + 16x + 12}{x--3}=x^3 - 2x^2 + 4x + 4\\
\lim_{x\to -3}(x^3 - 2x^2 + 4x + 4)=-53
\end{aligned}
\]
\[\boxed{-53}\]


\paragraph{Ejercicio 292.} Calcula usando fracción sintética $\displaystyle\lim_{x\to -2}\frac{x^4 + x^3 + 3x^2 + 15x + 10}{x--2}$. 

\[
\begin{aligned}
P(-2)=0\\
P(x)=(x--2)(x^3 - x^2 + 5x + 5)\\
\frac{x^4 + x^3 + 3x^2 + 15x + 10}{x--2}=x^3 - x^2 + 5x + 5\\
\lim_{x\to -2}(x^3 - x^2 + 5x + 5)=-17
\end{aligned}
\]
\[\boxed{-17}\]


\paragraph{Ejercicio 293.} Calcula usando fracción sintética $\displaystyle\lim_{x\to -1}\frac{x^4 + x^3 + 6x^2 + 8x + 2}{x--1}$. 

\[
\begin{aligned}
P(-1)=0\\
P(x)=(x--1)(x^3 + 6x + 2)\\
\frac{x^4 + x^3 + 6x^2 + 8x + 2}{x--1}=x^3 + 6x + 2\\
\lim_{x\to -1}(x^3 + 6x + 2)=-5
\end{aligned}
\]
\[\boxed{-5}\]


\paragraph{Ejercicio 294.} Calcula usando fracción sintética $\displaystyle\lim_{x\to 5}\frac{x^4 - 4x^3 + 2x^2 - 32x - 15}{x-5}$. 

\[
\begin{aligned}
P(5)=0\\
P(x)=(x-5)(x^3 + x^2 + 7x + 3)\\
\frac{x^4 - 4x^3 + 2x^2 - 32x - 15}{x-5}=x^3 + x^2 + 7x + 3\\
\lim_{x\to 5}(x^3 + x^2 + 7x + 3)=188
\end{aligned}
\]
\[\boxed{188}\]


\paragraph{Ejercicio 295.} Calcula usando fracción sintética $\displaystyle\lim_{x\to 1}\frac{x^4 + x^3 - x^2 + 3x - 4}{x-1}$. 

\[
\begin{aligned}
P(1)=0\\
P(x)=(x-1)(x^3 + 2x^2 + x + 4)\\
\frac{x^4 + x^3 - x^2 + 3x - 4}{x-1}=x^3 + 2x^2 + x + 4\\
\lim_{x\to 1}(x^3 + 2x^2 + x + 4)=8
\end{aligned}
\]
\[\boxed{8}\]


\paragraph{Ejercicio 296.} Calcula usando fracción sintética $\displaystyle\lim_{x\to 2}\frac{x^4 - 4x^3 + 6x^2 + x - 10}{x-2}$. 

\[
\begin{aligned}
P(2)=0\\
P(x)=(x-2)(x^3 - 2x^2 + 2x + 5)\\
\frac{x^4 - 4x^3 + 6x^2 + x - 10}{x-2}=x^3 - 2x^2 + 2x + 5\\
\lim_{x\to 2}(x^3 - 2x^2 + 2x + 5)=9
\end{aligned}
\]
\[\boxed{9}\]


\paragraph{Ejercicio 297.} Calcula usando fracción sintética $\displaystyle\lim_{x\to 3}\frac{x^4 - 4x^3 + 6x^2 - 7x - 6}{x-3}$. 

\[
\begin{aligned}
P(3)=0\\
P(x)=(x-3)(x^3 - x^2 + 3x + 2)\\
\frac{x^4 - 4x^3 + 6x^2 - 7x - 6}{x-3}=x^3 - x^2 + 3x + 2\\
\lim_{x\to 3}(x^3 - x^2 + 3x + 2)=29
\end{aligned}
\]
\[\boxed{29}\]


\paragraph{Ejercicio 298.} Calcula usando fracción sintética $\displaystyle\lim_{x\to 4}\frac{x^4 - 4x^3 + 4x^2 - 13x - 12}{x-4}$. 

\[
\begin{aligned}
P(4)=0\\
P(x)=(x-4)(x^3 + 4x + 3)\\
\frac{x^4 - 4x^3 + 4x^2 - 13x - 12}{x-4}=x^3 + 4x + 3\\
\lim_{x\to 4}(x^3 + 4x + 3)=83
\end{aligned}
\]
\[\boxed{83}\]


\paragraph{Ejercicio 299.} Calcula usando fracción sintética $\displaystyle\lim_{x\to -4}\frac{x^4 + 5x^3 + 9x^2 + 24x + 16}{x--4}$. 

\[
\begin{aligned}
P(-4)=0\\
P(x)=(x--4)(x^3 + x^2 + 5x + 4)\\
\frac{x^4 + 5x^3 + 9x^2 + 24x + 16}{x--4}=x^3 + x^2 + 5x + 4\\
\lim_{x\to -4}(x^3 + x^2 + 5x + 4)=-64
\end{aligned}
\]
\[\boxed{-64}\]


\paragraph{Ejercicio 300.} Calcula $\displaystyle\lim_{x\to1}\frac{x^5-5x+4}{x^4-4x+3}$ usando fracción sintética las veces necesarias.

\[
\begin{aligned}
x^5-5x+4=(x-1)^2(x^3+2x^2+3x+4)\\
x^4-4x+3=(x-1)^2(x^2+2x+3)\\
\frac{x^5-5x+4}{x^4-4x+3}=\frac{x^3+2x^2+3x+4}{x^2+2x+3}\\
\lim_{x\to1}=\frac{1+2+3+4}{1+2+3}=\frac{5}{3}
\end{aligned}
\]
\[\boxed{\frac{5}{3}}\]


\newpage

\subsection*{Límites exponenciales}

\addcontentsline{toc}{subsection}{Soluciones - Límites exponenciales}

\subsubsection*{Nivel 1 - Básico}

\paragraph{Ejercicio 301.} Calcula $\displaystyle\lim_{x\to -2}2^x$.

\[
\begin{aligned}
\lim_{x\to -2}2^x=2^{-2}
\end{aligned}
\]
\[\boxed{2^{-2}}\]


\paragraph{Ejercicio 302.} Calcula $\displaystyle\lim_{x\to -1}3^x$.

\[
\begin{aligned}
\lim_{x\to -1}3^x=3^{-1}
\end{aligned}
\]
\[\boxed{3^{-1}}\]


\paragraph{Ejercicio 303.} Calcula $\displaystyle\lim_{x\to 0}4^x$.

\[
\begin{aligned}
\lim_{x\to 0}4^x=4^{0}
\end{aligned}
\]
\[\boxed{4^{0}}\]


\paragraph{Ejercicio 304.} Calcula $\displaystyle\lim_{x\to 1}5^x$.

\[
\begin{aligned}
\lim_{x\to 1}5^x=5^{1}
\end{aligned}
\]
\[\boxed{5^{1}}\]


\paragraph{Ejercicio 305.} Calcula $\displaystyle\lim_{x\to 2}2^x$.

\[
\begin{aligned}
\lim_{x\to 2}2^x=2^{2}
\end{aligned}
\]
\[\boxed{2^{2}}\]


\paragraph{Ejercicio 306.} Calcula $\displaystyle\lim_{x\to -2}3^x$.

\[
\begin{aligned}
\lim_{x\to -2}3^x=3^{-2}
\end{aligned}
\]
\[\boxed{3^{-2}}\]


\paragraph{Ejercicio 307.} Calcula $\displaystyle\lim_{x\to -1}4^x$.

\[
\begin{aligned}
\lim_{x\to -1}4^x=4^{-1}
\end{aligned}
\]
\[\boxed{4^{-1}}\]


\paragraph{Ejercicio 308.} Calcula $\displaystyle\lim_{x\to 0}5^x$.

\[
\begin{aligned}
\lim_{x\to 0}5^x=5^{0}
\end{aligned}
\]
\[\boxed{5^{0}}\]


\paragraph{Ejercicio 309.} Calcula $\displaystyle\lim_{x\to 1}2^x$.

\[
\begin{aligned}
\lim_{x\to 1}2^x=2^{1}
\end{aligned}
\]
\[\boxed{2^{1}}\]


\paragraph{Ejercicio 310.} Calcula $\displaystyle\lim_{x\to 2}3^x$.

\[
\begin{aligned}
\lim_{x\to 2}3^x=3^{2}
\end{aligned}
\]
\[\boxed{3^{2}}\]


\paragraph{Ejercicio 311.} Calcula $\displaystyle\lim_{x\to -2}4^x$.

\[
\begin{aligned}
\lim_{x\to -2}4^x=4^{-2}
\end{aligned}
\]
\[\boxed{4^{-2}}\]


\paragraph{Ejercicio 312.} Calcula $\displaystyle\lim_{x\to -1}5^x$.

\[
\begin{aligned}
\lim_{x\to -1}5^x=5^{-1}
\end{aligned}
\]
\[\boxed{5^{-1}}\]


\paragraph{Ejercicio 313.} Calcula $\displaystyle\lim_{x\to 0}2^x$.

\[
\begin{aligned}
\lim_{x\to 0}2^x=2^{0}
\end{aligned}
\]
\[\boxed{2^{0}}\]


\paragraph{Ejercicio 314.} Calcula $\displaystyle\lim_{x\to 1}3^x$.

\[
\begin{aligned}
\lim_{x\to 1}3^x=3^{1}
\end{aligned}
\]
\[\boxed{3^{1}}\]


\paragraph{Ejercicio 315.} Calcula $\displaystyle\lim_{x\to 2}4^x$.

\[
\begin{aligned}
\lim_{x\to 2}4^x=4^{2}
\end{aligned}
\]
\[\boxed{4^{2}}\]


\paragraph{Ejercicio 316.} Calcula $\displaystyle\lim_{x\to -2}5^x$.

\[
\begin{aligned}
\lim_{x\to -2}5^x=5^{-2}
\end{aligned}
\]
\[\boxed{5^{-2}}\]


\paragraph{Ejercicio 317.} Calcula $\displaystyle\lim_{x\to -1}2^x$.

\[
\begin{aligned}
\lim_{x\to -1}2^x=2^{-1}
\end{aligned}
\]
\[\boxed{2^{-1}}\]


\paragraph{Ejercicio 318.} Calcula $\displaystyle\lim_{x\to 0}3^x$.

\[
\begin{aligned}
\lim_{x\to 0}3^x=3^{0}
\end{aligned}
\]
\[\boxed{3^{0}}\]


\paragraph{Ejercicio 319.} Calcula $\displaystyle\lim_{x\to 1}4^x$.

\[
\begin{aligned}
\lim_{x\to 1}4^x=4^{1}
\end{aligned}
\]
\[\boxed{4^{1}}\]


\paragraph{Ejercicio 320.} Calcula $\displaystyle\lim_{x\to 2}5^x$.

\[
\begin{aligned}
\lim_{x\to 2}5^x=5^{2}
\end{aligned}
\]
\[\boxed{5^{2}}\]


\newpage

\subsubsection*{Nivel 2 - Intermedio bajo}

\paragraph{Ejercicio 321.} Calcula $\displaystyle\lim_{x\to\infty}\frac{1}{2^x}$.

\[
\begin{aligned}
2^x\to\infty\\
\frac{1}{2^x}\to0
\end{aligned}
\]
\[\boxed{0}\]


\paragraph{Ejercicio 322.} Calcula $\displaystyle\lim_{x\to-\infty}3^x$.

\[
\begin{aligned}
3^x=\frac{1}{3^{-x}}\\
-x\to\infty\\
\lim=0
\end{aligned}
\]
\[\boxed{0}\]


\paragraph{Ejercicio 323.} Calcula $\displaystyle\lim_{x\to\infty}\frac{1}{4^x}$.

\[
\begin{aligned}
4^x\to\infty\\
\frac{1}{4^x}\to0
\end{aligned}
\]
\[\boxed{0}\]


\paragraph{Ejercicio 324.} Calcula $\displaystyle\lim_{x\to-\infty}5^x$.

\[
\begin{aligned}
5^x=\frac{1}{5^{-x}}\\
-x\to\infty\\
\lim=0
\end{aligned}
\]
\[\boxed{0}\]


\paragraph{Ejercicio 325.} Calcula $\displaystyle\lim_{x\to\infty}\frac{1}{2^x}$.

\[
\begin{aligned}
2^x\to\infty\\
\frac{1}{2^x}\to0
\end{aligned}
\]
\[\boxed{0}\]


\paragraph{Ejercicio 326.} Calcula $\displaystyle\lim_{x\to-\infty}3^x$.

\[
\begin{aligned}
3^x=\frac{1}{3^{-x}}\\
-x\to\infty\\
\lim=0
\end{aligned}
\]
\[\boxed{0}\]


\paragraph{Ejercicio 327.} Calcula $\displaystyle\lim_{x\to\infty}\frac{1}{4^x}$.

\[
\begin{aligned}
4^x\to\infty\\
\frac{1}{4^x}\to0
\end{aligned}
\]
\[\boxed{0}\]


\paragraph{Ejercicio 328.} Calcula $\displaystyle\lim_{x\to-\infty}5^x$.

\[
\begin{aligned}
5^x=\frac{1}{5^{-x}}\\
-x\to\infty\\
\lim=0
\end{aligned}
\]
\[\boxed{0}\]


\paragraph{Ejercicio 329.} Calcula $\displaystyle\lim_{x\to\infty}\frac{1}{2^x}$.

\[
\begin{aligned}
2^x\to\infty\\
\frac{1}{2^x}\to0
\end{aligned}
\]
\[\boxed{0}\]


\paragraph{Ejercicio 330.} Calcula $\displaystyle\lim_{x\to-\infty}3^x$.

\[
\begin{aligned}
3^x=\frac{1}{3^{-x}}\\
-x\to\infty\\
\lim=0
\end{aligned}
\]
\[\boxed{0}\]


\paragraph{Ejercicio 331.} Calcula $\displaystyle\lim_{x\to\infty}\frac{1}{4^x}$.

\[
\begin{aligned}
4^x\to\infty\\
\frac{1}{4^x}\to0
\end{aligned}
\]
\[\boxed{0}\]


\paragraph{Ejercicio 332.} Calcula $\displaystyle\lim_{x\to-\infty}5^x$.

\[
\begin{aligned}
5^x=\frac{1}{5^{-x}}\\
-x\to\infty\\
\lim=0
\end{aligned}
\]
\[\boxed{0}\]


\paragraph{Ejercicio 333.} Calcula $\displaystyle\lim_{x\to\infty}\frac{1}{2^x}$.

\[
\begin{aligned}
2^x\to\infty\\
\frac{1}{2^x}\to0
\end{aligned}
\]
\[\boxed{0}\]


\paragraph{Ejercicio 334.} Calcula $\displaystyle\lim_{x\to-\infty}3^x$.

\[
\begin{aligned}
3^x=\frac{1}{3^{-x}}\\
-x\to\infty\\
\lim=0
\end{aligned}
\]
\[\boxed{0}\]


\paragraph{Ejercicio 335.} Calcula $\displaystyle\lim_{x\to\infty}\frac{1}{4^x}$.

\[
\begin{aligned}
4^x\to\infty\\
\frac{1}{4^x}\to0
\end{aligned}
\]
\[\boxed{0}\]


\paragraph{Ejercicio 336.} Calcula $\displaystyle\lim_{x\to-\infty}5^x$.

\[
\begin{aligned}
5^x=\frac{1}{5^{-x}}\\
-x\to\infty\\
\lim=0
\end{aligned}
\]
\[\boxed{0}\]


\paragraph{Ejercicio 337.} Calcula $\displaystyle\lim_{x\to\infty}\frac{1}{2^x}$.

\[
\begin{aligned}
2^x\to\infty\\
\frac{1}{2^x}\to0
\end{aligned}
\]
\[\boxed{0}\]


\paragraph{Ejercicio 338.} Calcula $\displaystyle\lim_{x\to-\infty}3^x$.

\[
\begin{aligned}
3^x=\frac{1}{3^{-x}}\\
-x\to\infty\\
\lim=0
\end{aligned}
\]
\[\boxed{0}\]


\paragraph{Ejercicio 339.} Calcula $\displaystyle\lim_{x\to\infty}\frac{1}{4^x}$.

\[
\begin{aligned}
4^x\to\infty\\
\frac{1}{4^x}\to0
\end{aligned}
\]
\[\boxed{0}\]


\paragraph{Ejercicio 340.} Calcula $\displaystyle\lim_{x\to-\infty}5^x$.

\[
\begin{aligned}
5^x=\frac{1}{5^{-x}}\\
-x\to\infty\\
\lim=0
\end{aligned}
\]
\[\boxed{0}\]


\newpage

\subsubsection*{Nivel 3 - Intermedio}

\paragraph{Ejercicio 341.} Calcula $\displaystyle\lim_{x\to0}\frac{2^{1x}-1}{x}$.

\[
\begin{aligned}
u=1x\\
x=\frac{u}{1}\\
\frac{2^{1x}-1}{x}=1\frac{2^u-1}{u}\\
\lim=1\ln 2
\end{aligned}
\]
\[\boxed{1\ln 2}\]


\paragraph{Ejercicio 342.} Calcula $\displaystyle\lim_{x\to0}\frac{3^{2x}-1}{x}$.

\[
\begin{aligned}
u=2x\\
x=\frac{u}{2}\\
\frac{3^{2x}-1}{x}=2\frac{3^u-1}{u}\\
\lim=2\ln 3
\end{aligned}
\]
\[\boxed{2\ln 3}\]


\paragraph{Ejercicio 343.} Calcula $\displaystyle\lim_{x\to0}\frac{4^{3x}-1}{x}$.

\[
\begin{aligned}
u=3x\\
x=\frac{u}{3}\\
\frac{4^{3x}-1}{x}=3\frac{4^u-1}{u}\\
\lim=3\ln 4
\end{aligned}
\]
\[\boxed{3\ln 4}\]


\paragraph{Ejercicio 344.} Calcula $\displaystyle\lim_{x\to0}\frac{5^{4x}-1}{x}$.

\[
\begin{aligned}
u=4x\\
x=\frac{u}{4}\\
\frac{5^{4x}-1}{x}=4\frac{5^u-1}{u}\\
\lim=4\ln 5
\end{aligned}
\]
\[\boxed{4\ln 5}\]


\paragraph{Ejercicio 345.} Calcula $\displaystyle\lim_{x\to0}\frac{6^{5x}-1}{x}$.

\[
\begin{aligned}
u=5x\\
x=\frac{u}{5}\\
\frac{6^{5x}-1}{x}=5\frac{6^u-1}{u}\\
\lim=5\ln 6
\end{aligned}
\]
\[\boxed{5\ln 6}\]


\paragraph{Ejercicio 346.} Calcula $\displaystyle\lim_{x\to0}\frac{7^{1x}-1}{x}$.

\[
\begin{aligned}
u=1x\\
x=\frac{u}{1}\\
\frac{7^{1x}-1}{x}=1\frac{7^u-1}{u}\\
\lim=1\ln 7
\end{aligned}
\]
\[\boxed{1\ln 7}\]


\paragraph{Ejercicio 347.} Calcula $\displaystyle\lim_{x\to0}\frac{8^{2x}-1}{x}$.

\[
\begin{aligned}
u=2x\\
x=\frac{u}{2}\\
\frac{8^{2x}-1}{x}=2\frac{8^u-1}{u}\\
\lim=2\ln 8
\end{aligned}
\]
\[\boxed{2\ln 8}\]


\paragraph{Ejercicio 348.} Calcula $\displaystyle\lim_{x\to0}\frac{2^{3x}-1}{x}$.

\[
\begin{aligned}
u=3x\\
x=\frac{u}{3}\\
\frac{2^{3x}-1}{x}=3\frac{2^u-1}{u}\\
\lim=3\ln 2
\end{aligned}
\]
\[\boxed{3\ln 2}\]


\paragraph{Ejercicio 349.} Calcula $\displaystyle\lim_{x\to0}\frac{3^{4x}-1}{x}$.

\[
\begin{aligned}
u=4x\\
x=\frac{u}{4}\\
\frac{3^{4x}-1}{x}=4\frac{3^u-1}{u}\\
\lim=4\ln 3
\end{aligned}
\]
\[\boxed{4\ln 3}\]


\paragraph{Ejercicio 350.} Calcula $\displaystyle\lim_{x\to0}\frac{4^{5x}-1}{x}$.

\[
\begin{aligned}
u=5x\\
x=\frac{u}{5}\\
\frac{4^{5x}-1}{x}=5\frac{4^u-1}{u}\\
\lim=5\ln 4
\end{aligned}
\]
\[\boxed{5\ln 4}\]


\paragraph{Ejercicio 351.} Calcula $\displaystyle\lim_{x\to0}\frac{5^{1x}-1}{x}$.

\[
\begin{aligned}
u=1x\\
x=\frac{u}{1}\\
\frac{5^{1x}-1}{x}=1\frac{5^u-1}{u}\\
\lim=1\ln 5
\end{aligned}
\]
\[\boxed{1\ln 5}\]


\paragraph{Ejercicio 352.} Calcula $\displaystyle\lim_{x\to0}\frac{6^{2x}-1}{x}$.

\[
\begin{aligned}
u=2x\\
x=\frac{u}{2}\\
\frac{6^{2x}-1}{x}=2\frac{6^u-1}{u}\\
\lim=2\ln 6
\end{aligned}
\]
\[\boxed{2\ln 6}\]


\paragraph{Ejercicio 353.} Calcula $\displaystyle\lim_{x\to0}\frac{7^{3x}-1}{x}$.

\[
\begin{aligned}
u=3x\\
x=\frac{u}{3}\\
\frac{7^{3x}-1}{x}=3\frac{7^u-1}{u}\\
\lim=3\ln 7
\end{aligned}
\]
\[\boxed{3\ln 7}\]


\paragraph{Ejercicio 354.} Calcula $\displaystyle\lim_{x\to0}\frac{8^{4x}-1}{x}$.

\[
\begin{aligned}
u=4x\\
x=\frac{u}{4}\\
\frac{8^{4x}-1}{x}=4\frac{8^u-1}{u}\\
\lim=4\ln 8
\end{aligned}
\]
\[\boxed{4\ln 8}\]


\paragraph{Ejercicio 355.} Calcula $\displaystyle\lim_{x\to0}\frac{2^{5x}-1}{x}$.

\[
\begin{aligned}
u=5x\\
x=\frac{u}{5}\\
\frac{2^{5x}-1}{x}=5\frac{2^u-1}{u}\\
\lim=5\ln 2
\end{aligned}
\]
\[\boxed{5\ln 2}\]


\paragraph{Ejercicio 356.} Calcula $\displaystyle\lim_{x\to0}\frac{3^{1x}-1}{x}$.

\[
\begin{aligned}
u=1x\\
x=\frac{u}{1}\\
\frac{3^{1x}-1}{x}=1\frac{3^u-1}{u}\\
\lim=1\ln 3
\end{aligned}
\]
\[\boxed{1\ln 3}\]


\paragraph{Ejercicio 357.} Calcula $\displaystyle\lim_{x\to0}\frac{4^{2x}-1}{x}$.

\[
\begin{aligned}
u=2x\\
x=\frac{u}{2}\\
\frac{4^{2x}-1}{x}=2\frac{4^u-1}{u}\\
\lim=2\ln 4
\end{aligned}
\]
\[\boxed{2\ln 4}\]


\paragraph{Ejercicio 358.} Calcula $\displaystyle\lim_{x\to0}\frac{5^{3x}-1}{x}$.

\[
\begin{aligned}
u=3x\\
x=\frac{u}{3}\\
\frac{5^{3x}-1}{x}=3\frac{5^u-1}{u}\\
\lim=3\ln 5
\end{aligned}
\]
\[\boxed{3\ln 5}\]


\paragraph{Ejercicio 359.} Calcula $\displaystyle\lim_{x\to0}\frac{6^{4x}-1}{x}$.

\[
\begin{aligned}
u=4x\\
x=\frac{u}{4}\\
\frac{6^{4x}-1}{x}=4\frac{6^u-1}{u}\\
\lim=4\ln 6
\end{aligned}
\]
\[\boxed{4\ln 6}\]


\paragraph{Ejercicio 360.} Calcula $\displaystyle\lim_{x\to0}\frac{7^{5x}-1}{x}$.

\[
\begin{aligned}
u=5x\\
x=\frac{u}{5}\\
\frac{7^{5x}-1}{x}=5\frac{7^u-1}{u}\\
\lim=5\ln 7
\end{aligned}
\]
\[\boxed{5\ln 7}\]


\newpage

\subsubsection*{Nivel 4 - Avanzado}

\paragraph{Ejercicio 361.} Calcula $\displaystyle\lim_{x\to0}\frac{e^{1x}-e^{1x}}{x}$.

\[
\begin{aligned}
e^{1x}=1+1x+o(x)\\
e^{1x}=1+1x+o(x)\\
\lim=0
\end{aligned}
\]
\[\boxed{0}\]


\paragraph{Ejercicio 362.} Calcula $\displaystyle\lim_{x\to\infty}\frac{e^{2x}+2}{e^{2x}-2}$.

\[
\begin{aligned}
\frac{e^{2x}+2}{e^{2x}-2}=\frac{1+2e^{-2x}}{1-2e^{-2x}}\\
e^{-2x}\to0\\
\lim=1
\end{aligned}
\]
\[\boxed{1}\]


\paragraph{Ejercicio 363.} Calcula $\displaystyle\lim_{x\to0}\frac{e^{3x}-e^{3x}}{x}$.

\[
\begin{aligned}
e^{3x}=1+3x+o(x)\\
e^{3x}=1+3x+o(x)\\
\lim=0
\end{aligned}
\]
\[\boxed{0}\]


\paragraph{Ejercicio 364.} Calcula $\displaystyle\lim_{x\to\infty}\frac{e^{4x}+4}{e^{4x}-4}$.

\[
\begin{aligned}
\frac{e^{4x}+4}{e^{4x}-4}=\frac{1+4e^{-4x}}{1-4e^{-4x}}\\
e^{-4x}\to0\\
\lim=1
\end{aligned}
\]
\[\boxed{1}\]


\paragraph{Ejercicio 365.} Calcula $\displaystyle\lim_{x\to0}\frac{e^{5x}-e^{5x}}{x}$.

\[
\begin{aligned}
e^{5x}=1+5x+o(x)\\
e^{5x}=1+5x+o(x)\\
\lim=0
\end{aligned}
\]
\[\boxed{0}\]


\paragraph{Ejercicio 366.} Calcula $\displaystyle\lim_{x\to\infty}\frac{e^{6x}+1}{e^{6x}-1}$.

\[
\begin{aligned}
\frac{e^{6x}+1}{e^{6x}-1}=\frac{1+1e^{-6x}}{1-1e^{-6x}}\\
e^{-6x}\to0\\
\lim=1
\end{aligned}
\]
\[\boxed{1}\]


\paragraph{Ejercicio 367.} Calcula $\displaystyle\lim_{x\to0}\frac{e^{7x}-e^{2x}}{x}$.

\[
\begin{aligned}
e^{7x}=1+7x+o(x)\\
e^{2x}=1+2x+o(x)\\
\lim=5
\end{aligned}
\]
\[\boxed{5}\]


\paragraph{Ejercicio 368.} Calcula $\displaystyle\lim_{x\to\infty}\frac{e^{1x}+3}{e^{1x}-3}$.

\[
\begin{aligned}
\frac{e^{1x}+3}{e^{1x}-3}=\frac{1+3e^{-1x}}{1-3e^{-1x}}\\
e^{-1x}\to0\\
\lim=1
\end{aligned}
\]
\[\boxed{1}\]


\paragraph{Ejercicio 369.} Calcula $\displaystyle\lim_{x\to0}\frac{e^{2x}-e^{4x}}{x}$.

\[
\begin{aligned}
e^{2x}=1+2x+o(x)\\
e^{4x}=1+4x+o(x)\\
\lim=-2
\end{aligned}
\]
\[\boxed{-2}\]


\paragraph{Ejercicio 370.} Calcula $\displaystyle\lim_{x\to\infty}\frac{e^{3x}+5}{e^{3x}-5}$.

\[
\begin{aligned}
\frac{e^{3x}+5}{e^{3x}-5}=\frac{1+5e^{-3x}}{1-5e^{-3x}}\\
e^{-3x}\to0\\
\lim=1
\end{aligned}
\]
\[\boxed{1}\]


\paragraph{Ejercicio 371.} Calcula $\displaystyle\lim_{x\to0}\frac{e^{4x}-e^{1x}}{x}$.

\[
\begin{aligned}
e^{4x}=1+4x+o(x)\\
e^{1x}=1+1x+o(x)\\
\lim=3
\end{aligned}
\]
\[\boxed{3}\]


\paragraph{Ejercicio 372.} Calcula $\displaystyle\lim_{x\to\infty}\frac{e^{5x}+2}{e^{5x}-2}$.

\[
\begin{aligned}
\frac{e^{5x}+2}{e^{5x}-2}=\frac{1+2e^{-5x}}{1-2e^{-5x}}\\
e^{-5x}\to0\\
\lim=1
\end{aligned}
\]
\[\boxed{1}\]


\paragraph{Ejercicio 373.} Calcula $\displaystyle\lim_{x\to0}\frac{e^{6x}-e^{3x}}{x}$.

\[
\begin{aligned}
e^{6x}=1+6x+o(x)\\
e^{3x}=1+3x+o(x)\\
\lim=3
\end{aligned}
\]
\[\boxed{3}\]


\paragraph{Ejercicio 374.} Calcula $\displaystyle\lim_{x\to\infty}\frac{e^{7x}+4}{e^{7x}-4}$.

\[
\begin{aligned}
\frac{e^{7x}+4}{e^{7x}-4}=\frac{1+4e^{-7x}}{1-4e^{-7x}}\\
e^{-7x}\to0\\
\lim=1
\end{aligned}
\]
\[\boxed{1}\]


\paragraph{Ejercicio 375.} Calcula $\displaystyle\lim_{x\to0}\frac{e^{1x}-e^{5x}}{x}$.

\[
\begin{aligned}
e^{1x}=1+1x+o(x)\\
e^{5x}=1+5x+o(x)\\
\lim=-4
\end{aligned}
\]
\[\boxed{-4}\]


\paragraph{Ejercicio 376.} Calcula $\displaystyle\lim_{x\to\infty}\frac{e^{2x}+1}{e^{2x}-1}$.

\[
\begin{aligned}
\frac{e^{2x}+1}{e^{2x}-1}=\frac{1+1e^{-2x}}{1-1e^{-2x}}\\
e^{-2x}\to0\\
\lim=1
\end{aligned}
\]
\[\boxed{1}\]


\paragraph{Ejercicio 377.} Calcula $\displaystyle\lim_{x\to0}\frac{e^{3x}-e^{2x}}{x}$.

\[
\begin{aligned}
e^{3x}=1+3x+o(x)\\
e^{2x}=1+2x+o(x)\\
\lim=1
\end{aligned}
\]
\[\boxed{1}\]


\paragraph{Ejercicio 378.} Calcula $\displaystyle\lim_{x\to\infty}\frac{e^{4x}+3}{e^{4x}-3}$.

\[
\begin{aligned}
\frac{e^{4x}+3}{e^{4x}-3}=\frac{1+3e^{-4x}}{1-3e^{-4x}}\\
e^{-4x}\to0\\
\lim=1
\end{aligned}
\]
\[\boxed{1}\]


\paragraph{Ejercicio 379.} Calcula $\displaystyle\lim_{x\to0}\frac{e^{5x}-e^{4x}}{x}$.

\[
\begin{aligned}
e^{5x}=1+5x+o(x)\\
e^{4x}=1+4x+o(x)\\
\lim=1
\end{aligned}
\]
\[\boxed{1}\]


\paragraph{Ejercicio 380.} Calcula $\displaystyle\lim_{x\to\infty}\frac{e^{6x}+5}{e^{6x}-5}$.

\[
\begin{aligned}
\frac{e^{6x}+5}{e^{6x}-5}=\frac{1+5e^{-6x}}{1-5e^{-6x}}\\
e^{-6x}\to0\\
\lim=1
\end{aligned}
\]
\[\boxed{1}\]


\newpage

\subsubsection*{Nivel 5 - Imposible}

\paragraph{Ejercicio 381.} Calcula $\displaystyle \lim_{x\to0}\frac{2^x-3^x}{x}$.

\[
\begin{aligned}
2^x=e^{x\ln2}\\
3^x=e^{x\ln3}\\
\frac{2^x-3^x}{x}=\frac{(1+x\ln2+o(x))-(1+x\ln3+o(x))}{x}\\
\lim=\ln2-\ln3=\ln\frac23
\end{aligned}
\]
\[\boxed{\ln\frac23}\]


\paragraph{Ejercicio 382.} Calcula $\displaystyle \lim_{x\to0}\frac{a^x-b^x}{x}\quad(a,b>0)$.

\[
\begin{aligned}
a^x=e^{x\ln a}\\
b^x=e^{x\ln b}\\
a^x-b^x=x(\ln a-\ln b)+o(x)\\
\lim=\ln\frac{a}{b}
\end{aligned}
\]
\[\boxed{\ln\frac{a}{b}}\]


\paragraph{Ejercicio 383.} Calcula $\displaystyle \lim_{x\to0}\left(\frac{e^x-1}{x}\right)^{1/x}$.

\[
\begin{aligned}
\frac{e^x-1}{x}=1+\frac{x}{2}+o(x)\\
L=\lim\left(1+\frac{x}{2}+o(x)\right)^{1/x}\\
\ln L=\lim\frac{\ln(1+\frac{x}{2}+o(x))}{x}\\
\ln L=\frac12\\
L=e^{1/2}
\end{aligned}
\]
\[\boxed{\sqrt e}\]


\paragraph{Ejercicio 384.} Calcula $\displaystyle \lim_{x\to\infty}\left(1+\frac{3}{x}\right)^{2x}$.

\[
\begin{aligned}
L=\lim\left(1+\frac{3}{x}\right)^{2x}\\
\ln L=\lim 2x\ln\left(1+\frac{3}{x}\right)\\
\ln L=6\\
L=e^6
\end{aligned}
\]
\[\boxed{e^6}\]


\paragraph{Ejercicio 385.} Calcula $\displaystyle \lim_{x\to0}\frac{e^{\sin x}-e^x}{x^3}$.

\[
\begin{aligned}
\sin x=x-\frac{x^3}{6}+o(x^3)\\
e^{\sin x}=e^{x-x^3/6+o(x^3)}=e^x\left(1-\frac{x^3}{6}+o(x^3)\right)\\
e^{\sin x}-e^x=-\frac{e^xx^3}{6}+o(x^3)\\
\lim=-\frac16
\end{aligned}
\]
\[\boxed{-\frac16}\]


\paragraph{Ejercicio 386.} Calcula $\displaystyle \lim_{x\to0}\frac{2^x-3^x}{x}$.

\[
\begin{aligned}
2^x=e^{x\ln2}\\
3^x=e^{x\ln3}\\
\frac{2^x-3^x}{x}=\frac{(1+x\ln2+o(x))-(1+x\ln3+o(x))}{x}\\
\lim=\ln2-\ln3=\ln\frac23
\end{aligned}
\]
\[\boxed{\ln\frac23}\]


\paragraph{Ejercicio 387.} Calcula $\displaystyle \lim_{x\to0}\frac{a^x-b^x}{x}\quad(a,b>0)$.

\[
\begin{aligned}
a^x=e^{x\ln a}\\
b^x=e^{x\ln b}\\
a^x-b^x=x(\ln a-\ln b)+o(x)\\
\lim=\ln\frac{a}{b}
\end{aligned}
\]
\[\boxed{\ln\frac{a}{b}}\]


\paragraph{Ejercicio 388.} Calcula $\displaystyle \lim_{x\to0}\left(\frac{e^x-1}{x}\right)^{1/x}$.

\[
\begin{aligned}
\frac{e^x-1}{x}=1+\frac{x}{2}+o(x)\\
L=\lim\left(1+\frac{x}{2}+o(x)\right)^{1/x}\\
\ln L=\lim\frac{\ln(1+\frac{x}{2}+o(x))}{x}\\
\ln L=\frac12\\
L=e^{1/2}
\end{aligned}
\]
\[\boxed{\sqrt e}\]


\paragraph{Ejercicio 389.} Calcula $\displaystyle \lim_{x\to\infty}\left(1+\frac{3}{x}\right)^{2x}$.

\[
\begin{aligned}
L=\lim\left(1+\frac{3}{x}\right)^{2x}\\
\ln L=\lim 2x\ln\left(1+\frac{3}{x}\right)\\
\ln L=6\\
L=e^6
\end{aligned}
\]
\[\boxed{e^6}\]


\paragraph{Ejercicio 390.} Calcula $\displaystyle \lim_{x\to0}\frac{e^{\sin x}-e^x}{x^3}$.

\[
\begin{aligned}
\sin x=x-\frac{x^3}{6}+o(x^3)\\
e^{\sin x}=e^{x-x^3/6+o(x^3)}=e^x\left(1-\frac{x^3}{6}+o(x^3)\right)\\
e^{\sin x}-e^x=-\frac{e^xx^3}{6}+o(x^3)\\
\lim=-\frac16
\end{aligned}
\]
\[\boxed{-\frac16}\]


\paragraph{Ejercicio 391.} Calcula $\displaystyle \lim_{x\to0}\frac{2^x-3^x}{x}$.

\[
\begin{aligned}
2^x=e^{x\ln2}\\
3^x=e^{x\ln3}\\
\frac{2^x-3^x}{x}=\frac{(1+x\ln2+o(x))-(1+x\ln3+o(x))}{x}\\
\lim=\ln2-\ln3=\ln\frac23
\end{aligned}
\]
\[\boxed{\ln\frac23}\]


\paragraph{Ejercicio 392.} Calcula $\displaystyle \lim_{x\to0}\frac{a^x-b^x}{x}\quad(a,b>0)$.

\[
\begin{aligned}
a^x=e^{x\ln a}\\
b^x=e^{x\ln b}\\
a^x-b^x=x(\ln a-\ln b)+o(x)\\
\lim=\ln\frac{a}{b}
\end{aligned}
\]
\[\boxed{\ln\frac{a}{b}}\]


\paragraph{Ejercicio 393.} Calcula $\displaystyle \lim_{x\to0}\left(\frac{e^x-1}{x}\right)^{1/x}$.

\[
\begin{aligned}
\frac{e^x-1}{x}=1+\frac{x}{2}+o(x)\\
L=\lim\left(1+\frac{x}{2}+o(x)\right)^{1/x}\\
\ln L=\lim\frac{\ln(1+\frac{x}{2}+o(x))}{x}\\
\ln L=\frac12\\
L=e^{1/2}
\end{aligned}
\]
\[\boxed{\sqrt e}\]


\paragraph{Ejercicio 394.} Calcula $\displaystyle \lim_{x\to\infty}\left(1+\frac{3}{x}\right)^{2x}$.

\[
\begin{aligned}
L=\lim\left(1+\frac{3}{x}\right)^{2x}\\
\ln L=\lim 2x\ln\left(1+\frac{3}{x}\right)\\
\ln L=6\\
L=e^6
\end{aligned}
\]
\[\boxed{e^6}\]


\paragraph{Ejercicio 395.} Calcula $\displaystyle \lim_{x\to0}\frac{e^{\sin x}-e^x}{x^3}$.

\[
\begin{aligned}
\sin x=x-\frac{x^3}{6}+o(x^3)\\
e^{\sin x}=e^{x-x^3/6+o(x^3)}=e^x\left(1-\frac{x^3}{6}+o(x^3)\right)\\
e^{\sin x}-e^x=-\frac{e^xx^3}{6}+o(x^3)\\
\lim=-\frac16
\end{aligned}
\]
\[\boxed{-\frac16}\]


\paragraph{Ejercicio 396.} Calcula $\displaystyle \lim_{x\to0}\frac{2^x-3^x}{x}$.

\[
\begin{aligned}
2^x=e^{x\ln2}\\
3^x=e^{x\ln3}\\
\frac{2^x-3^x}{x}=\frac{(1+x\ln2+o(x))-(1+x\ln3+o(x))}{x}\\
\lim=\ln2-\ln3=\ln\frac23
\end{aligned}
\]
\[\boxed{\ln\frac23}\]


\paragraph{Ejercicio 397.} Calcula $\displaystyle \lim_{x\to0}\frac{a^x-b^x}{x}\quad(a,b>0)$.

\[
\begin{aligned}
a^x=e^{x\ln a}\\
b^x=e^{x\ln b}\\
a^x-b^x=x(\ln a-\ln b)+o(x)\\
\lim=\ln\frac{a}{b}
\end{aligned}
\]
\[\boxed{\ln\frac{a}{b}}\]


\paragraph{Ejercicio 398.} Calcula $\displaystyle \lim_{x\to0}\left(\frac{e^x-1}{x}\right)^{1/x}$.

\[
\begin{aligned}
\frac{e^x-1}{x}=1+\frac{x}{2}+o(x)\\
L=\lim\left(1+\frac{x}{2}+o(x)\right)^{1/x}\\
\ln L=\lim\frac{\ln(1+\frac{x}{2}+o(x))}{x}\\
\ln L=\frac12\\
L=e^{1/2}
\end{aligned}
\]
\[\boxed{\sqrt e}\]


\paragraph{Ejercicio 399.} Calcula $\displaystyle \lim_{x\to\infty}\left(1+\frac{3}{x}\right)^{2x}$.

\[
\begin{aligned}
L=\lim\left(1+\frac{3}{x}\right)^{2x}\\
\ln L=\lim 2x\ln\left(1+\frac{3}{x}\right)\\
\ln L=6\\
L=e^6
\end{aligned}
\]
\[\boxed{e^6}\]


\paragraph{Ejercicio 400.} Calcula $\displaystyle \lim_{x\to0}\frac{e^{\sin x}-e^x}{x^3}$.

\[
\begin{aligned}
\sin x=x-\frac{x^3}{6}+o(x^3)\\
e^{\sin x}=e^{x-x^3/6+o(x^3)}=e^x\left(1-\frac{x^3}{6}+o(x^3)\right)\\
e^{\sin x}-e^x=-\frac{e^xx^3}{6}+o(x^3)\\
\lim=-\frac16
\end{aligned}
\]
\[\boxed{-\frac16}\]


\newpage

\subsection*{Límites trigonométricos}

\addcontentsline{toc}{subsection}{Soluciones - Límites trigonométricos}

\subsubsection*{Nivel 1 - Básico}

\paragraph{Ejercicio 401.} Calcula $\displaystyle\lim_{x\to0}\frac{\sin(1x)}{1x}$.

\[
\begin{aligned}
\frac{\sin(1x)}{1x}=\frac{1}{1}\frac{\sin(1x)}{1x}\\
\lim=\frac{1}{1}
\end{aligned}
\]
\[\boxed{\frac{1}{1}}\]


\paragraph{Ejercicio 402.} Calcula $\displaystyle\lim_{x\to0}\frac{\sin(2x)}{2x}$.

\[
\begin{aligned}
\frac{\sin(2x)}{2x}=\frac{2}{2}\frac{\sin(2x)}{2x}\\
\lim=\frac{2}{2}
\end{aligned}
\]
\[\boxed{\frac{2}{2}}\]


\paragraph{Ejercicio 403.} Calcula $\displaystyle\lim_{x\to0}\frac{\sin(3x)}{3x}$.

\[
\begin{aligned}
\frac{\sin(3x)}{3x}=\frac{3}{3}\frac{\sin(3x)}{3x}\\
\lim=\frac{3}{3}
\end{aligned}
\]
\[\boxed{\frac{3}{3}}\]


\paragraph{Ejercicio 404.} Calcula $\displaystyle\lim_{x\to0}\frac{\sin(4x)}{4x}$.

\[
\begin{aligned}
\frac{\sin(4x)}{4x}=\frac{4}{4}\frac{\sin(4x)}{4x}\\
\lim=\frac{4}{4}
\end{aligned}
\]
\[\boxed{\frac{4}{4}}\]


\paragraph{Ejercicio 405.} Calcula $\displaystyle\lim_{x\to0}\frac{\sin(5x)}{5x}$.

\[
\begin{aligned}
\frac{\sin(5x)}{5x}=\frac{5}{5}\frac{\sin(5x)}{5x}\\
\lim=\frac{5}{5}
\end{aligned}
\]
\[\boxed{\frac{5}{5}}\]


\paragraph{Ejercicio 406.} Calcula $\displaystyle\lim_{x\to0}\frac{\sin(6x)}{1x}$.

\[
\begin{aligned}
\frac{\sin(6x)}{1x}=\frac{6}{1}\frac{\sin(6x)}{6x}\\
\lim=\frac{6}{1}
\end{aligned}
\]
\[\boxed{\frac{6}{1}}\]


\paragraph{Ejercicio 407.} Calcula $\displaystyle\lim_{x\to0}\frac{\sin(7x)}{2x}$.

\[
\begin{aligned}
\frac{\sin(7x)}{2x}=\frac{7}{2}\frac{\sin(7x)}{7x}\\
\lim=\frac{7}{2}
\end{aligned}
\]
\[\boxed{\frac{7}{2}}\]


\paragraph{Ejercicio 408.} Calcula $\displaystyle\lim_{x\to0}\frac{\sin(8x)}{3x}$.

\[
\begin{aligned}
\frac{\sin(8x)}{3x}=\frac{8}{3}\frac{\sin(8x)}{8x}\\
\lim=\frac{8}{3}
\end{aligned}
\]
\[\boxed{\frac{8}{3}}\]


\paragraph{Ejercicio 409.} Calcula $\displaystyle\lim_{x\to0}\frac{\sin(9x)}{4x}$.

\[
\begin{aligned}
\frac{\sin(9x)}{4x}=\frac{9}{4}\frac{\sin(9x)}{9x}\\
\lim=\frac{9}{4}
\end{aligned}
\]
\[\boxed{\frac{9}{4}}\]


\paragraph{Ejercicio 410.} Calcula $\displaystyle\lim_{x\to0}\frac{\sin(1x)}{5x}$.

\[
\begin{aligned}
\frac{\sin(1x)}{5x}=\frac{1}{5}\frac{\sin(1x)}{1x}\\
\lim=\frac{1}{5}
\end{aligned}
\]
\[\boxed{\frac{1}{5}}\]


\paragraph{Ejercicio 411.} Calcula $\displaystyle\lim_{x\to0}\frac{\sin(2x)}{1x}$.

\[
\begin{aligned}
\frac{\sin(2x)}{1x}=\frac{2}{1}\frac{\sin(2x)}{2x}\\
\lim=\frac{2}{1}
\end{aligned}
\]
\[\boxed{\frac{2}{1}}\]


\paragraph{Ejercicio 412.} Calcula $\displaystyle\lim_{x\to0}\frac{\sin(3x)}{2x}$.

\[
\begin{aligned}
\frac{\sin(3x)}{2x}=\frac{3}{2}\frac{\sin(3x)}{3x}\\
\lim=\frac{3}{2}
\end{aligned}
\]
\[\boxed{\frac{3}{2}}\]


\paragraph{Ejercicio 413.} Calcula $\displaystyle\lim_{x\to0}\frac{\sin(4x)}{3x}$.

\[
\begin{aligned}
\frac{\sin(4x)}{3x}=\frac{4}{3}\frac{\sin(4x)}{4x}\\
\lim=\frac{4}{3}
\end{aligned}
\]
\[\boxed{\frac{4}{3}}\]


\paragraph{Ejercicio 414.} Calcula $\displaystyle\lim_{x\to0}\frac{\sin(5x)}{4x}$.

\[
\begin{aligned}
\frac{\sin(5x)}{4x}=\frac{5}{4}\frac{\sin(5x)}{5x}\\
\lim=\frac{5}{4}
\end{aligned}
\]
\[\boxed{\frac{5}{4}}\]


\paragraph{Ejercicio 415.} Calcula $\displaystyle\lim_{x\to0}\frac{\sin(6x)}{5x}$.

\[
\begin{aligned}
\frac{\sin(6x)}{5x}=\frac{6}{5}\frac{\sin(6x)}{6x}\\
\lim=\frac{6}{5}
\end{aligned}
\]
\[\boxed{\frac{6}{5}}\]


\paragraph{Ejercicio 416.} Calcula $\displaystyle\lim_{x\to0}\frac{\sin(7x)}{1x}$.

\[
\begin{aligned}
\frac{\sin(7x)}{1x}=\frac{7}{1}\frac{\sin(7x)}{7x}\\
\lim=\frac{7}{1}
\end{aligned}
\]
\[\boxed{\frac{7}{1}}\]


\paragraph{Ejercicio 417.} Calcula $\displaystyle\lim_{x\to0}\frac{\sin(8x)}{2x}$.

\[
\begin{aligned}
\frac{\sin(8x)}{2x}=\frac{8}{2}\frac{\sin(8x)}{8x}\\
\lim=\frac{8}{2}
\end{aligned}
\]
\[\boxed{\frac{8}{2}}\]


\paragraph{Ejercicio 418.} Calcula $\displaystyle\lim_{x\to0}\frac{\sin(9x)}{3x}$.

\[
\begin{aligned}
\frac{\sin(9x)}{3x}=\frac{9}{3}\frac{\sin(9x)}{9x}\\
\lim=\frac{9}{3}
\end{aligned}
\]
\[\boxed{\frac{9}{3}}\]


\paragraph{Ejercicio 419.} Calcula $\displaystyle\lim_{x\to0}\frac{\sin(1x)}{4x}$.

\[
\begin{aligned}
\frac{\sin(1x)}{4x}=\frac{1}{4}\frac{\sin(1x)}{1x}\\
\lim=\frac{1}{4}
\end{aligned}
\]
\[\boxed{\frac{1}{4}}\]


\paragraph{Ejercicio 420.} Calcula $\displaystyle\lim_{x\to0}\frac{\sin(2x)}{5x}$.

\[
\begin{aligned}
\frac{\sin(2x)}{5x}=\frac{2}{5}\frac{\sin(2x)}{2x}\\
\lim=\frac{2}{5}
\end{aligned}
\]
\[\boxed{\frac{2}{5}}\]


\newpage

\subsubsection*{Nivel 2 - Intermedio bajo}

\paragraph{Ejercicio 421.} Calcula $\displaystyle\lim_{x\to0}\frac{\tan(1x)}{1x}$.

\[
\begin{aligned}
\lim=\frac{1}{1}
\end{aligned}
\]
\[\boxed{\frac{1}{1}}\]


\paragraph{Ejercicio 422.} Calcula $\displaystyle\lim_{x\to0}\frac{\sin(2x)}{\sin(2x)}$.

\[
\begin{aligned}
\lim=\frac{2}{2}
\end{aligned}
\]
\[\boxed{\frac{2}{2}}\]


\paragraph{Ejercicio 423.} Calcula $\displaystyle\lim_{x\to0}\frac{\tan(3x)}{3x}$.

\[
\begin{aligned}
\lim=\frac{3}{3}
\end{aligned}
\]
\[\boxed{\frac{3}{3}}\]


\paragraph{Ejercicio 424.} Calcula $\displaystyle\lim_{x\to0}\frac{\sin(4x)}{\sin(4x)}$.

\[
\begin{aligned}
\lim=\frac{4}{4}
\end{aligned}
\]
\[\boxed{\frac{4}{4}}\]


\paragraph{Ejercicio 425.} Calcula $\displaystyle\lim_{x\to0}\frac{\tan(5x)}{5x}$.

\[
\begin{aligned}
\lim=\frac{5}{5}
\end{aligned}
\]
\[\boxed{\frac{5}{5}}\]


\paragraph{Ejercicio 426.} Calcula $\displaystyle\lim_{x\to0}\frac{\sin(6x)}{\sin(6x)}$.

\[
\begin{aligned}
\lim=\frac{6}{6}
\end{aligned}
\]
\[\boxed{\frac{6}{6}}\]


\paragraph{Ejercicio 427.} Calcula $\displaystyle\lim_{x\to0}\frac{\tan(7x)}{1x}$.

\[
\begin{aligned}
\lim=\frac{7}{1}
\end{aligned}
\]
\[\boxed{\frac{7}{1}}\]


\paragraph{Ejercicio 428.} Calcula $\displaystyle\lim_{x\to0}\frac{\sin(1x)}{\sin(2x)}$.

\[
\begin{aligned}
\lim=\frac{1}{2}
\end{aligned}
\]
\[\boxed{\frac{1}{2}}\]


\paragraph{Ejercicio 429.} Calcula $\displaystyle\lim_{x\to0}\frac{\tan(2x)}{3x}$.

\[
\begin{aligned}
\lim=\frac{2}{3}
\end{aligned}
\]
\[\boxed{\frac{2}{3}}\]


\paragraph{Ejercicio 430.} Calcula $\displaystyle\lim_{x\to0}\frac{\sin(3x)}{\sin(4x)}$.

\[
\begin{aligned}
\lim=\frac{3}{4}
\end{aligned}
\]
\[\boxed{\frac{3}{4}}\]


\paragraph{Ejercicio 431.} Calcula $\displaystyle\lim_{x\to0}\frac{\tan(4x)}{5x}$.

\[
\begin{aligned}
\lim=\frac{4}{5}
\end{aligned}
\]
\[\boxed{\frac{4}{5}}\]


\paragraph{Ejercicio 432.} Calcula $\displaystyle\lim_{x\to0}\frac{\sin(5x)}{\sin(6x)}$.

\[
\begin{aligned}
\lim=\frac{5}{6}
\end{aligned}
\]
\[\boxed{\frac{5}{6}}\]


\paragraph{Ejercicio 433.} Calcula $\displaystyle\lim_{x\to0}\frac{\tan(6x)}{1x}$.

\[
\begin{aligned}
\lim=\frac{6}{1}
\end{aligned}
\]
\[\boxed{\frac{6}{1}}\]


\paragraph{Ejercicio 434.} Calcula $\displaystyle\lim_{x\to0}\frac{\sin(7x)}{\sin(2x)}$.

\[
\begin{aligned}
\lim=\frac{7}{2}
\end{aligned}
\]
\[\boxed{\frac{7}{2}}\]


\paragraph{Ejercicio 435.} Calcula $\displaystyle\lim_{x\to0}\frac{\tan(1x)}{3x}$.

\[
\begin{aligned}
\lim=\frac{1}{3}
\end{aligned}
\]
\[\boxed{\frac{1}{3}}\]


\paragraph{Ejercicio 436.} Calcula $\displaystyle\lim_{x\to0}\frac{\sin(2x)}{\sin(4x)}$.

\[
\begin{aligned}
\lim=\frac{2}{4}
\end{aligned}
\]
\[\boxed{\frac{2}{4}}\]


\paragraph{Ejercicio 437.} Calcula $\displaystyle\lim_{x\to0}\frac{\tan(3x)}{5x}$.

\[
\begin{aligned}
\lim=\frac{3}{5}
\end{aligned}
\]
\[\boxed{\frac{3}{5}}\]


\paragraph{Ejercicio 438.} Calcula $\displaystyle\lim_{x\to0}\frac{\sin(4x)}{\sin(6x)}$.

\[
\begin{aligned}
\lim=\frac{4}{6}
\end{aligned}
\]
\[\boxed{\frac{4}{6}}\]


\paragraph{Ejercicio 439.} Calcula $\displaystyle\lim_{x\to0}\frac{\tan(5x)}{1x}$.

\[
\begin{aligned}
\lim=\frac{5}{1}
\end{aligned}
\]
\[\boxed{\frac{5}{1}}\]


\paragraph{Ejercicio 440.} Calcula $\displaystyle\lim_{x\to0}\frac{\sin(6x)}{\sin(2x)}$.

\[
\begin{aligned}
\lim=\frac{6}{2}
\end{aligned}
\]
\[\boxed{\frac{6}{2}}\]


\newpage

\subsubsection*{Nivel 3 - Intermedio}

\paragraph{Ejercicio 441.} Calcula $\displaystyle\lim_{x\to0}\frac{1-\cos(1x)}{1x^2}$.

\[
\begin{aligned}
1-\cos(1x)\sim \frac{(1x)^2}{2}\\
\lim=\frac{1}{2}
\end{aligned}
\]
\[\boxed{\frac{1}{2}}\]


\paragraph{Ejercicio 442.} Calcula $\displaystyle\lim_{x\to0}\frac{1-\cos(2x)}{2x^2}$.

\[
\begin{aligned}
1-\cos(2x)\sim \frac{(2x)^2}{2}\\
\lim=\frac{4}{4}
\end{aligned}
\]
\[\boxed{\frac{4}{4}}\]


\paragraph{Ejercicio 443.} Calcula $\displaystyle\lim_{x\to0}\frac{1-\cos(3x)}{3x^2}$.

\[
\begin{aligned}
1-\cos(3x)\sim \frac{(3x)^2}{2}\\
\lim=\frac{9}{6}
\end{aligned}
\]
\[\boxed{\frac{9}{6}}\]


\paragraph{Ejercicio 444.} Calcula $\displaystyle\lim_{x\to0}\frac{1-\cos(4x)}{4x^2}$.

\[
\begin{aligned}
1-\cos(4x)\sim \frac{(4x)^2}{2}\\
\lim=\frac{16}{8}
\end{aligned}
\]
\[\boxed{\frac{16}{8}}\]


\paragraph{Ejercicio 445.} Calcula $\displaystyle\lim_{x\to0}\frac{1-\cos(5x)}{5x^2}$.

\[
\begin{aligned}
1-\cos(5x)\sim \frac{(5x)^2}{2}\\
\lim=\frac{25}{10}
\end{aligned}
\]
\[\boxed{\frac{25}{10}}\]


\paragraph{Ejercicio 446.} Calcula $\displaystyle\lim_{x\to0}\frac{1-\cos(6x)}{1x^2}$.

\[
\begin{aligned}
1-\cos(6x)\sim \frac{(6x)^2}{2}\\
\lim=\frac{36}{2}
\end{aligned}
\]
\[\boxed{\frac{36}{2}}\]


\paragraph{Ejercicio 447.} Calcula $\displaystyle\lim_{x\to0}\frac{1-\cos(7x)}{2x^2}$.

\[
\begin{aligned}
1-\cos(7x)\sim \frac{(7x)^2}{2}\\
\lim=\frac{49}{4}
\end{aligned}
\]
\[\boxed{\frac{49}{4}}\]


\paragraph{Ejercicio 448.} Calcula $\displaystyle\lim_{x\to0}\frac{1-\cos(8x)}{3x^2}$.

\[
\begin{aligned}
1-\cos(8x)\sim \frac{(8x)^2}{2}\\
\lim=\frac{64}{6}
\end{aligned}
\]
\[\boxed{\frac{64}{6}}\]


\paragraph{Ejercicio 449.} Calcula $\displaystyle\lim_{x\to0}\frac{1-\cos(9x)}{4x^2}$.

\[
\begin{aligned}
1-\cos(9x)\sim \frac{(9x)^2}{2}\\
\lim=\frac{81}{8}
\end{aligned}
\]
\[\boxed{\frac{81}{8}}\]


\paragraph{Ejercicio 450.} Calcula $\displaystyle\lim_{x\to0}\frac{1-\cos(1x)}{5x^2}$.

\[
\begin{aligned}
1-\cos(1x)\sim \frac{(1x)^2}{2}\\
\lim=\frac{1}{10}
\end{aligned}
\]
\[\boxed{\frac{1}{10}}\]


\paragraph{Ejercicio 451.} Calcula $\displaystyle\lim_{x\to0}\frac{1-\cos(2x)}{1x^2}$.

\[
\begin{aligned}
1-\cos(2x)\sim \frac{(2x)^2}{2}\\
\lim=\frac{4}{2}
\end{aligned}
\]
\[\boxed{\frac{4}{2}}\]


\paragraph{Ejercicio 452.} Calcula $\displaystyle\lim_{x\to0}\frac{1-\cos(3x)}{2x^2}$.

\[
\begin{aligned}
1-\cos(3x)\sim \frac{(3x)^2}{2}\\
\lim=\frac{9}{4}
\end{aligned}
\]
\[\boxed{\frac{9}{4}}\]


\paragraph{Ejercicio 453.} Calcula $\displaystyle\lim_{x\to0}\frac{1-\cos(4x)}{3x^2}$.

\[
\begin{aligned}
1-\cos(4x)\sim \frac{(4x)^2}{2}\\
\lim=\frac{16}{6}
\end{aligned}
\]
\[\boxed{\frac{16}{6}}\]


\paragraph{Ejercicio 454.} Calcula $\displaystyle\lim_{x\to0}\frac{1-\cos(5x)}{4x^2}$.

\[
\begin{aligned}
1-\cos(5x)\sim \frac{(5x)^2}{2}\\
\lim=\frac{25}{8}
\end{aligned}
\]
\[\boxed{\frac{25}{8}}\]


\paragraph{Ejercicio 455.} Calcula $\displaystyle\lim_{x\to0}\frac{1-\cos(6x)}{5x^2}$.

\[
\begin{aligned}
1-\cos(6x)\sim \frac{(6x)^2}{2}\\
\lim=\frac{36}{10}
\end{aligned}
\]
\[\boxed{\frac{36}{10}}\]


\paragraph{Ejercicio 456.} Calcula $\displaystyle\lim_{x\to0}\frac{1-\cos(7x)}{1x^2}$.

\[
\begin{aligned}
1-\cos(7x)\sim \frac{(7x)^2}{2}\\
\lim=\frac{49}{2}
\end{aligned}
\]
\[\boxed{\frac{49}{2}}\]


\paragraph{Ejercicio 457.} Calcula $\displaystyle\lim_{x\to0}\frac{1-\cos(8x)}{2x^2}$.

\[
\begin{aligned}
1-\cos(8x)\sim \frac{(8x)^2}{2}\\
\lim=\frac{64}{4}
\end{aligned}
\]
\[\boxed{\frac{64}{4}}\]


\paragraph{Ejercicio 458.} Calcula $\displaystyle\lim_{x\to0}\frac{1-\cos(9x)}{3x^2}$.

\[
\begin{aligned}
1-\cos(9x)\sim \frac{(9x)^2}{2}\\
\lim=\frac{81}{6}
\end{aligned}
\]
\[\boxed{\frac{81}{6}}\]


\paragraph{Ejercicio 459.} Calcula $\displaystyle\lim_{x\to0}\frac{1-\cos(1x)}{4x^2}$.

\[
\begin{aligned}
1-\cos(1x)\sim \frac{(1x)^2}{2}\\
\lim=\frac{1}{8}
\end{aligned}
\]
\[\boxed{\frac{1}{8}}\]


\paragraph{Ejercicio 460.} Calcula $\displaystyle\lim_{x\to0}\frac{1-\cos(2x)}{5x^2}$.

\[
\begin{aligned}
1-\cos(2x)\sim \frac{(2x)^2}{2}\\
\lim=\frac{4}{10}
\end{aligned}
\]
\[\boxed{\frac{4}{10}}\]


\newpage

\subsubsection*{Nivel 4 - Avanzado}

\paragraph{Ejercicio 461.} Calcula $\displaystyle\lim_{x\to0}\frac{\sin(1x)\tan(2x)}{x^2}$.

\[
\begin{aligned}
\sin(1x)\sim 1x\\
\tan(2x)\sim 2x\\
\lim=2
\end{aligned}
\]
\[\boxed{2}\]


\paragraph{Ejercicio 462.} Calcula $\displaystyle\lim_{x\to0}\frac{\sin(2x)-\sin(3x)}{x}$.

\[
\begin{aligned}
\sin(2x)\sim 2x\\
\sin(3x)\sim 3x\\
\lim=-1
\end{aligned}
\]
\[\boxed{-1}\]


\paragraph{Ejercicio 463.} Calcula $\displaystyle\lim_{x\to0}\frac{\sin(3x)\tan(4x)}{x^2}$.

\[
\begin{aligned}
\sin(3x)\sim 3x\\
\tan(4x)\sim 4x\\
\lim=12
\end{aligned}
\]
\[\boxed{12}\]


\paragraph{Ejercicio 464.} Calcula $\displaystyle\lim_{x\to0}\frac{\sin(4x)-\sin(5x)}{x}$.

\[
\begin{aligned}
\sin(4x)\sim 4x\\
\sin(5x)\sim 5x\\
\lim=-1
\end{aligned}
\]
\[\boxed{-1}\]


\paragraph{Ejercicio 465.} Calcula $\displaystyle\lim_{x\to0}\frac{\sin(5x)\tan(6x)}{x^2}$.

\[
\begin{aligned}
\sin(5x)\sim 5x\\
\tan(6x)\sim 6x\\
\lim=30
\end{aligned}
\]
\[\boxed{30}\]


\paragraph{Ejercicio 466.} Calcula $\displaystyle\lim_{x\to0}\frac{\sin(1x)-\sin(7x)}{x}$.

\[
\begin{aligned}
\sin(1x)\sim 1x\\
\sin(7x)\sim 7x\\
\lim=-6
\end{aligned}
\]
\[\boxed{-6}\]


\paragraph{Ejercicio 467.} Calcula $\displaystyle\lim_{x\to0}\frac{\sin(2x)\tan(2x)}{x^2}$.

\[
\begin{aligned}
\sin(2x)\sim 2x\\
\tan(2x)\sim 2x\\
\lim=4
\end{aligned}
\]
\[\boxed{4}\]


\paragraph{Ejercicio 468.} Calcula $\displaystyle\lim_{x\to0}\frac{\sin(3x)-\sin(3x)}{x}$.

\[
\begin{aligned}
\sin(3x)\sim 3x\\
\sin(3x)\sim 3x\\
\lim=0
\end{aligned}
\]
\[\boxed{0}\]


\paragraph{Ejercicio 469.} Calcula $\displaystyle\lim_{x\to0}\frac{\sin(4x)\tan(4x)}{x^2}$.

\[
\begin{aligned}
\sin(4x)\sim 4x\\
\tan(4x)\sim 4x\\
\lim=16
\end{aligned}
\]
\[\boxed{16}\]


\paragraph{Ejercicio 470.} Calcula $\displaystyle\lim_{x\to0}\frac{\sin(5x)-\sin(5x)}{x}$.

\[
\begin{aligned}
\sin(5x)\sim 5x\\
\sin(5x)\sim 5x\\
\lim=0
\end{aligned}
\]
\[\boxed{0}\]


\paragraph{Ejercicio 471.} Calcula $\displaystyle\lim_{x\to0}\frac{\sin(1x)\tan(6x)}{x^2}$.

\[
\begin{aligned}
\sin(1x)\sim 1x\\
\tan(6x)\sim 6x\\
\lim=6
\end{aligned}
\]
\[\boxed{6}\]


\paragraph{Ejercicio 472.} Calcula $\displaystyle\lim_{x\to0}\frac{\sin(2x)-\sin(7x)}{x}$.

\[
\begin{aligned}
\sin(2x)\sim 2x\\
\sin(7x)\sim 7x\\
\lim=-5
\end{aligned}
\]
\[\boxed{-5}\]


\paragraph{Ejercicio 473.} Calcula $\displaystyle\lim_{x\to0}\frac{\sin(3x)\tan(2x)}{x^2}$.

\[
\begin{aligned}
\sin(3x)\sim 3x\\
\tan(2x)\sim 2x\\
\lim=6
\end{aligned}
\]
\[\boxed{6}\]


\paragraph{Ejercicio 474.} Calcula $\displaystyle\lim_{x\to0}\frac{\sin(4x)-\sin(3x)}{x}$.

\[
\begin{aligned}
\sin(4x)\sim 4x\\
\sin(3x)\sim 3x\\
\lim=1
\end{aligned}
\]
\[\boxed{1}\]


\paragraph{Ejercicio 475.} Calcula $\displaystyle\lim_{x\to0}\frac{\sin(5x)\tan(4x)}{x^2}$.

\[
\begin{aligned}
\sin(5x)\sim 5x\\
\tan(4x)\sim 4x\\
\lim=20
\end{aligned}
\]
\[\boxed{20}\]


\paragraph{Ejercicio 476.} Calcula $\displaystyle\lim_{x\to0}\frac{\sin(1x)-\sin(5x)}{x}$.

\[
\begin{aligned}
\sin(1x)\sim 1x\\
\sin(5x)\sim 5x\\
\lim=-4
\end{aligned}
\]
\[\boxed{-4}\]


\paragraph{Ejercicio 477.} Calcula $\displaystyle\lim_{x\to0}\frac{\sin(2x)\tan(6x)}{x^2}$.

\[
\begin{aligned}
\sin(2x)\sim 2x\\
\tan(6x)\sim 6x\\
\lim=12
\end{aligned}
\]
\[\boxed{12}\]


\paragraph{Ejercicio 478.} Calcula $\displaystyle\lim_{x\to0}\frac{\sin(3x)-\sin(7x)}{x}$.

\[
\begin{aligned}
\sin(3x)\sim 3x\\
\sin(7x)\sim 7x\\
\lim=-4
\end{aligned}
\]
\[\boxed{-4}\]


\paragraph{Ejercicio 479.} Calcula $\displaystyle\lim_{x\to0}\frac{\sin(4x)\tan(2x)}{x^2}$.

\[
\begin{aligned}
\sin(4x)\sim 4x\\
\tan(2x)\sim 2x\\
\lim=8
\end{aligned}
\]
\[\boxed{8}\]


\paragraph{Ejercicio 480.} Calcula $\displaystyle\lim_{x\to0}\frac{\sin(5x)-\sin(3x)}{x}$.

\[
\begin{aligned}
\sin(5x)\sim 5x\\
\sin(3x)\sim 3x\\
\lim=2
\end{aligned}
\]
\[\boxed{2}\]


\newpage

\subsubsection*{Nivel 5 - Imposible}

\paragraph{Ejercicio 481.} Calcula $\displaystyle \lim_{x\to0}\frac{\sin x-x\cos x}{x^3}$.

\[
\begin{aligned}
\sin x=x-\frac{x^3}{6}+o(x^3)\\
x\cos x=x\left(1-\frac{x^2}{2}+o(x^2)\right)=x-\frac{x^3}{2}+o(x^3)\\
\sin x-x\cos x=\frac{x^3}{3}+o(x^3)\\
\lim=\frac13
\end{aligned}
\]
\[\boxed{\frac13}\]


\paragraph{Ejercicio 482.} Calcula $\displaystyle \lim_{x\to0}\frac{\tan x-x}{x^3}$.

\[
\begin{aligned}
\tan x=x+\frac{x^3}{3}+o(x^3)\\
\tan x-x=\frac{x^3}{3}+o(x^3)\\
\lim=\frac13
\end{aligned}
\]
\[\boxed{\frac13}\]


\paragraph{Ejercicio 483.} Calcula $\displaystyle \lim_{x\to0}\frac{\sin(3x)-3\sin x}{x^3}$.

\[
\begin{aligned}
\sin(3x)=3x-\frac{27x^3}{6}+o(x^3)\\
3\sin x=3x-\frac{3x^3}{6}+o(x^3)\\
\sin(3x)-3\sin x=-4x^3+o(x^3)\\
\lim=-4
\end{aligned}
\]
\[\boxed{-4}\]


\paragraph{Ejercicio 484.} Calcula $\displaystyle \lim_{x\to0}\frac{1-\cos x}{\sin^2x}$.

\[
\begin{aligned}
1-\cos x=2\sin^2\frac{x}{2}\\
\sin x=2\sin\frac{x}{2}\cos\frac{x}{2}\\
\frac{1-\cos x}{\sin^2x}=\frac{1}{2\cos^2(x/2)}\\
\lim=\frac12
\end{aligned}
\]
\[\boxed{\frac12}\]


\paragraph{Ejercicio 485.} Calcula $\displaystyle \lim_{x\to0}\frac{\sin x\sin2x\sin3x}{x^3}$.

\[
\begin{aligned}
\sin x\sim x\\
\sin2x\sim2x\\
\sin3x\sim3x\\
\lim=6
\end{aligned}
\]
\[\boxed{6}\]


\paragraph{Ejercicio 486.} Calcula $\displaystyle \lim_{x\to0}\frac{\sin x-x\cos x}{x^3}$.

\[
\begin{aligned}
\sin x=x-\frac{x^3}{6}+o(x^3)\\
x\cos x=x\left(1-\frac{x^2}{2}+o(x^2)\right)=x-\frac{x^3}{2}+o(x^3)\\
\sin x-x\cos x=\frac{x^3}{3}+o(x^3)\\
\lim=\frac13
\end{aligned}
\]
\[\boxed{\frac13}\]


\paragraph{Ejercicio 487.} Calcula $\displaystyle \lim_{x\to0}\frac{\tan x-x}{x^3}$.

\[
\begin{aligned}
\tan x=x+\frac{x^3}{3}+o(x^3)\\
\tan x-x=\frac{x^3}{3}+o(x^3)\\
\lim=\frac13
\end{aligned}
\]
\[\boxed{\frac13}\]


\paragraph{Ejercicio 488.} Calcula $\displaystyle \lim_{x\to0}\frac{\sin(3x)-3\sin x}{x^3}$.

\[
\begin{aligned}
\sin(3x)=3x-\frac{27x^3}{6}+o(x^3)\\
3\sin x=3x-\frac{3x^3}{6}+o(x^3)\\
\sin(3x)-3\sin x=-4x^3+o(x^3)\\
\lim=-4
\end{aligned}
\]
\[\boxed{-4}\]


\paragraph{Ejercicio 489.} Calcula $\displaystyle \lim_{x\to0}\frac{1-\cos x}{\sin^2x}$.

\[
\begin{aligned}
1-\cos x=2\sin^2\frac{x}{2}\\
\sin x=2\sin\frac{x}{2}\cos\frac{x}{2}\\
\frac{1-\cos x}{\sin^2x}=\frac{1}{2\cos^2(x/2)}\\
\lim=\frac12
\end{aligned}
\]
\[\boxed{\frac12}\]


\paragraph{Ejercicio 490.} Calcula $\displaystyle \lim_{x\to0}\frac{\sin x\sin2x\sin3x}{x^3}$.

\[
\begin{aligned}
\sin x\sim x\\
\sin2x\sim2x\\
\sin3x\sim3x\\
\lim=6
\end{aligned}
\]
\[\boxed{6}\]


\paragraph{Ejercicio 491.} Calcula $\displaystyle \lim_{x\to0}\frac{\sin x-x\cos x}{x^3}$.

\[
\begin{aligned}
\sin x=x-\frac{x^3}{6}+o(x^3)\\
x\cos x=x\left(1-\frac{x^2}{2}+o(x^2)\right)=x-\frac{x^3}{2}+o(x^3)\\
\sin x-x\cos x=\frac{x^3}{3}+o(x^3)\\
\lim=\frac13
\end{aligned}
\]
\[\boxed{\frac13}\]


\paragraph{Ejercicio 492.} Calcula $\displaystyle \lim_{x\to0}\frac{\tan x-x}{x^3}$.

\[
\begin{aligned}
\tan x=x+\frac{x^3}{3}+o(x^3)\\
\tan x-x=\frac{x^3}{3}+o(x^3)\\
\lim=\frac13
\end{aligned}
\]
\[\boxed{\frac13}\]


\paragraph{Ejercicio 493.} Calcula $\displaystyle \lim_{x\to0}\frac{\sin(3x)-3\sin x}{x^3}$.

\[
\begin{aligned}
\sin(3x)=3x-\frac{27x^3}{6}+o(x^3)\\
3\sin x=3x-\frac{3x^3}{6}+o(x^3)\\
\sin(3x)-3\sin x=-4x^3+o(x^3)\\
\lim=-4
\end{aligned}
\]
\[\boxed{-4}\]


\paragraph{Ejercicio 494.} Calcula $\displaystyle \lim_{x\to0}\frac{1-\cos x}{\sin^2x}$.

\[
\begin{aligned}
1-\cos x=2\sin^2\frac{x}{2}\\
\sin x=2\sin\frac{x}{2}\cos\frac{x}{2}\\
\frac{1-\cos x}{\sin^2x}=\frac{1}{2\cos^2(x/2)}\\
\lim=\frac12
\end{aligned}
\]
\[\boxed{\frac12}\]


\paragraph{Ejercicio 495.} Calcula $\displaystyle \lim_{x\to0}\frac{\sin x\sin2x\sin3x}{x^3}$.

\[
\begin{aligned}
\sin x\sim x\\
\sin2x\sim2x\\
\sin3x\sim3x\\
\lim=6
\end{aligned}
\]
\[\boxed{6}\]


\paragraph{Ejercicio 496.} Calcula $\displaystyle \lim_{x\to0}\frac{\sin x-x\cos x}{x^3}$.

\[
\begin{aligned}
\sin x=x-\frac{x^3}{6}+o(x^3)\\
x\cos x=x\left(1-\frac{x^2}{2}+o(x^2)\right)=x-\frac{x^3}{2}+o(x^3)\\
\sin x-x\cos x=\frac{x^3}{3}+o(x^3)\\
\lim=\frac13
\end{aligned}
\]
\[\boxed{\frac13}\]


\paragraph{Ejercicio 497.} Calcula $\displaystyle \lim_{x\to0}\frac{\tan x-x}{x^3}$.

\[
\begin{aligned}
\tan x=x+\frac{x^3}{3}+o(x^3)\\
\tan x-x=\frac{x^3}{3}+o(x^3)\\
\lim=\frac13
\end{aligned}
\]
\[\boxed{\frac13}\]


\paragraph{Ejercicio 498.} Calcula $\displaystyle \lim_{x\to0}\frac{\sin(3x)-3\sin x}{x^3}$.

\[
\begin{aligned}
\sin(3x)=3x-\frac{27x^3}{6}+o(x^3)\\
3\sin x=3x-\frac{3x^3}{6}+o(x^3)\\
\sin(3x)-3\sin x=-4x^3+o(x^3)\\
\lim=-4
\end{aligned}
\]
\[\boxed{-4}\]


\paragraph{Ejercicio 499.} Calcula $\displaystyle \lim_{x\to0}\frac{1-\cos x}{\sin^2x}$.

\[
\begin{aligned}
1-\cos x=2\sin^2\frac{x}{2}\\
\sin x=2\sin\frac{x}{2}\cos\frac{x}{2}\\
\frac{1-\cos x}{\sin^2x}=\frac{1}{2\cos^2(x/2)}\\
\lim=\frac12
\end{aligned}
\]
\[\boxed{\frac12}\]


\paragraph{Ejercicio 500.} Calcula $\displaystyle \lim_{x\to0}\frac{\sin x\sin2x\sin3x}{x^3}$.

\[
\begin{aligned}
\sin x\sim x\\
\sin2x\sim2x\\
\sin3x\sim3x\\
\lim=6
\end{aligned}
\]
\[\boxed{6}\]


\newpage

\subsection*{Límites con conjugados normales y trigonométricos}

\addcontentsline{toc}{subsection}{Soluciones - Límites con conjugados normales y trigonométricos}

\subsubsection*{Nivel 1 - Básico}

\paragraph{Ejercicio 501.} Calcula $\displaystyle\lim_{x\to0}\frac{\sqrt{x+1}-\sqrt{1}}{x}$.

\[
\begin{aligned}
\frac{\sqrt{x+1}-\sqrt{1}}{x}\cdot \frac{\sqrt{x+1}+\sqrt{1}}{\sqrt{x+1}+\sqrt{1}}\\
=\frac{1}{\sqrt{x+1}+\sqrt{1}}\\
\lim=\frac{1}{2\sqrt{1}}
\end{aligned}
\]
\[\boxed{\frac{1}{2\sqrt{1}}}\]


\paragraph{Ejercicio 502.} Calcula $\displaystyle\lim_{x\to0}\frac{\sqrt{x+2}-\sqrt{2}}{x}$.

\[
\begin{aligned}
\frac{\sqrt{x+2}-\sqrt{2}}{x}\cdot \frac{\sqrt{x+2}+\sqrt{2}}{\sqrt{x+2}+\sqrt{2}}\\
=\frac{1}{\sqrt{x+2}+\sqrt{2}}\\
\lim=\frac{1}{2\sqrt{2}}
\end{aligned}
\]
\[\boxed{\frac{1}{2\sqrt{2}}}\]


\paragraph{Ejercicio 503.} Calcula $\displaystyle\lim_{x\to0}\frac{\sqrt{x+3}-\sqrt{3}}{x}$.

\[
\begin{aligned}
\frac{\sqrt{x+3}-\sqrt{3}}{x}\cdot \frac{\sqrt{x+3}+\sqrt{3}}{\sqrt{x+3}+\sqrt{3}}\\
=\frac{1}{\sqrt{x+3}+\sqrt{3}}\\
\lim=\frac{1}{2\sqrt{3}}
\end{aligned}
\]
\[\boxed{\frac{1}{2\sqrt{3}}}\]


\paragraph{Ejercicio 504.} Calcula $\displaystyle\lim_{x\to0}\frac{\sqrt{x+4}-\sqrt{4}}{x}$.

\[
\begin{aligned}
\frac{\sqrt{x+4}-\sqrt{4}}{x}\cdot \frac{\sqrt{x+4}+\sqrt{4}}{\sqrt{x+4}+\sqrt{4}}\\
=\frac{1}{\sqrt{x+4}+\sqrt{4}}\\
\lim=\frac{1}{2\sqrt{4}}
\end{aligned}
\]
\[\boxed{\frac{1}{2\sqrt{4}}}\]


\paragraph{Ejercicio 505.} Calcula $\displaystyle\lim_{x\to0}\frac{\sqrt{x+5}-\sqrt{5}}{x}$.

\[
\begin{aligned}
\frac{\sqrt{x+5}-\sqrt{5}}{x}\cdot \frac{\sqrt{x+5}+\sqrt{5}}{\sqrt{x+5}+\sqrt{5}}\\
=\frac{1}{\sqrt{x+5}+\sqrt{5}}\\
\lim=\frac{1}{2\sqrt{5}}
\end{aligned}
\]
\[\boxed{\frac{1}{2\sqrt{5}}}\]


\paragraph{Ejercicio 506.} Calcula $\displaystyle\lim_{x\to0}\frac{\sqrt{x+6}-\sqrt{6}}{x}$.

\[
\begin{aligned}
\frac{\sqrt{x+6}-\sqrt{6}}{x}\cdot \frac{\sqrt{x+6}+\sqrt{6}}{\sqrt{x+6}+\sqrt{6}}\\
=\frac{1}{\sqrt{x+6}+\sqrt{6}}\\
\lim=\frac{1}{2\sqrt{6}}
\end{aligned}
\]
\[\boxed{\frac{1}{2\sqrt{6}}}\]


\paragraph{Ejercicio 507.} Calcula $\displaystyle\lim_{x\to0}\frac{\sqrt{x+7}-\sqrt{7}}{x}$.

\[
\begin{aligned}
\frac{\sqrt{x+7}-\sqrt{7}}{x}\cdot \frac{\sqrt{x+7}+\sqrt{7}}{\sqrt{x+7}+\sqrt{7}}\\
=\frac{1}{\sqrt{x+7}+\sqrt{7}}\\
\lim=\frac{1}{2\sqrt{7}}
\end{aligned}
\]
\[\boxed{\frac{1}{2\sqrt{7}}}\]


\paragraph{Ejercicio 508.} Calcula $\displaystyle\lim_{x\to0}\frac{\sqrt{x+1}-\sqrt{1}}{x}$.

\[
\begin{aligned}
\frac{\sqrt{x+1}-\sqrt{1}}{x}\cdot \frac{\sqrt{x+1}+\sqrt{1}}{\sqrt{x+1}+\sqrt{1}}\\
=\frac{1}{\sqrt{x+1}+\sqrt{1}}\\
\lim=\frac{1}{2\sqrt{1}}
\end{aligned}
\]
\[\boxed{\frac{1}{2\sqrt{1}}}\]


\paragraph{Ejercicio 509.} Calcula $\displaystyle\lim_{x\to0}\frac{\sqrt{x+2}-\sqrt{2}}{x}$.

\[
\begin{aligned}
\frac{\sqrt{x+2}-\sqrt{2}}{x}\cdot \frac{\sqrt{x+2}+\sqrt{2}}{\sqrt{x+2}+\sqrt{2}}\\
=\frac{1}{\sqrt{x+2}+\sqrt{2}}\\
\lim=\frac{1}{2\sqrt{2}}
\end{aligned}
\]
\[\boxed{\frac{1}{2\sqrt{2}}}\]


\paragraph{Ejercicio 510.} Calcula $\displaystyle\lim_{x\to0}\frac{\sqrt{x+3}-\sqrt{3}}{x}$.

\[
\begin{aligned}
\frac{\sqrt{x+3}-\sqrt{3}}{x}\cdot \frac{\sqrt{x+3}+\sqrt{3}}{\sqrt{x+3}+\sqrt{3}}\\
=\frac{1}{\sqrt{x+3}+\sqrt{3}}\\
\lim=\frac{1}{2\sqrt{3}}
\end{aligned}
\]
\[\boxed{\frac{1}{2\sqrt{3}}}\]


\paragraph{Ejercicio 511.} Calcula $\displaystyle\lim_{x\to0}\frac{\sqrt{x+4}-\sqrt{4}}{x}$.

\[
\begin{aligned}
\frac{\sqrt{x+4}-\sqrt{4}}{x}\cdot \frac{\sqrt{x+4}+\sqrt{4}}{\sqrt{x+4}+\sqrt{4}}\\
=\frac{1}{\sqrt{x+4}+\sqrt{4}}\\
\lim=\frac{1}{2\sqrt{4}}
\end{aligned}
\]
\[\boxed{\frac{1}{2\sqrt{4}}}\]


\paragraph{Ejercicio 512.} Calcula $\displaystyle\lim_{x\to0}\frac{\sqrt{x+5}-\sqrt{5}}{x}$.

\[
\begin{aligned}
\frac{\sqrt{x+5}-\sqrt{5}}{x}\cdot \frac{\sqrt{x+5}+\sqrt{5}}{\sqrt{x+5}+\sqrt{5}}\\
=\frac{1}{\sqrt{x+5}+\sqrt{5}}\\
\lim=\frac{1}{2\sqrt{5}}
\end{aligned}
\]
\[\boxed{\frac{1}{2\sqrt{5}}}\]


\paragraph{Ejercicio 513.} Calcula $\displaystyle\lim_{x\to0}\frac{\sqrt{x+6}-\sqrt{6}}{x}$.

\[
\begin{aligned}
\frac{\sqrt{x+6}-\sqrt{6}}{x}\cdot \frac{\sqrt{x+6}+\sqrt{6}}{\sqrt{x+6}+\sqrt{6}}\\
=\frac{1}{\sqrt{x+6}+\sqrt{6}}\\
\lim=\frac{1}{2\sqrt{6}}
\end{aligned}
\]
\[\boxed{\frac{1}{2\sqrt{6}}}\]


\paragraph{Ejercicio 514.} Calcula $\displaystyle\lim_{x\to0}\frac{\sqrt{x+7}-\sqrt{7}}{x}$.

\[
\begin{aligned}
\frac{\sqrt{x+7}-\sqrt{7}}{x}\cdot \frac{\sqrt{x+7}+\sqrt{7}}{\sqrt{x+7}+\sqrt{7}}\\
=\frac{1}{\sqrt{x+7}+\sqrt{7}}\\
\lim=\frac{1}{2\sqrt{7}}
\end{aligned}
\]
\[\boxed{\frac{1}{2\sqrt{7}}}\]


\paragraph{Ejercicio 515.} Calcula $\displaystyle\lim_{x\to0}\frac{\sqrt{x+1}-\sqrt{1}}{x}$.

\[
\begin{aligned}
\frac{\sqrt{x+1}-\sqrt{1}}{x}\cdot \frac{\sqrt{x+1}+\sqrt{1}}{\sqrt{x+1}+\sqrt{1}}\\
=\frac{1}{\sqrt{x+1}+\sqrt{1}}\\
\lim=\frac{1}{2\sqrt{1}}
\end{aligned}
\]
\[\boxed{\frac{1}{2\sqrt{1}}}\]


\paragraph{Ejercicio 516.} Calcula $\displaystyle\lim_{x\to0}\frac{\sqrt{x+2}-\sqrt{2}}{x}$.

\[
\begin{aligned}
\frac{\sqrt{x+2}-\sqrt{2}}{x}\cdot \frac{\sqrt{x+2}+\sqrt{2}}{\sqrt{x+2}+\sqrt{2}}\\
=\frac{1}{\sqrt{x+2}+\sqrt{2}}\\
\lim=\frac{1}{2\sqrt{2}}
\end{aligned}
\]
\[\boxed{\frac{1}{2\sqrt{2}}}\]


\paragraph{Ejercicio 517.} Calcula $\displaystyle\lim_{x\to0}\frac{\sqrt{x+3}-\sqrt{3}}{x}$.

\[
\begin{aligned}
\frac{\sqrt{x+3}-\sqrt{3}}{x}\cdot \frac{\sqrt{x+3}+\sqrt{3}}{\sqrt{x+3}+\sqrt{3}}\\
=\frac{1}{\sqrt{x+3}+\sqrt{3}}\\
\lim=\frac{1}{2\sqrt{3}}
\end{aligned}
\]
\[\boxed{\frac{1}{2\sqrt{3}}}\]


\paragraph{Ejercicio 518.} Calcula $\displaystyle\lim_{x\to0}\frac{\sqrt{x+4}-\sqrt{4}}{x}$.

\[
\begin{aligned}
\frac{\sqrt{x+4}-\sqrt{4}}{x}\cdot \frac{\sqrt{x+4}+\sqrt{4}}{\sqrt{x+4}+\sqrt{4}}\\
=\frac{1}{\sqrt{x+4}+\sqrt{4}}\\
\lim=\frac{1}{2\sqrt{4}}
\end{aligned}
\]
\[\boxed{\frac{1}{2\sqrt{4}}}\]


\paragraph{Ejercicio 519.} Calcula $\displaystyle\lim_{x\to0}\frac{\sqrt{x+5}-\sqrt{5}}{x}$.

\[
\begin{aligned}
\frac{\sqrt{x+5}-\sqrt{5}}{x}\cdot \frac{\sqrt{x+5}+\sqrt{5}}{\sqrt{x+5}+\sqrt{5}}\\
=\frac{1}{\sqrt{x+5}+\sqrt{5}}\\
\lim=\frac{1}{2\sqrt{5}}
\end{aligned}
\]
\[\boxed{\frac{1}{2\sqrt{5}}}\]


\paragraph{Ejercicio 520.} Calcula $\displaystyle\lim_{x\to0}\frac{\sqrt{x+6}-\sqrt{6}}{x}$.

\[
\begin{aligned}
\frac{\sqrt{x+6}-\sqrt{6}}{x}\cdot \frac{\sqrt{x+6}+\sqrt{6}}{\sqrt{x+6}+\sqrt{6}}\\
=\frac{1}{\sqrt{x+6}+\sqrt{6}}\\
\lim=\frac{1}{2\sqrt{6}}
\end{aligned}
\]
\[\boxed{\frac{1}{2\sqrt{6}}}\]


\newpage

\subsubsection*{Nivel 2 - Intermedio bajo}

\paragraph{Ejercicio 521.} Calcula $\displaystyle\lim_{x\to -2}\frac{\sqrt{x+1}-\sqrt{-1}}{x--2}$.

\[
\begin{aligned}
\frac{\sqrt{x+1}-\sqrt{-1}}{x--2}\cdot \frac{\sqrt{x+1}+\sqrt{-1}}{\sqrt{x+1}+\sqrt{-1}}\\
=\frac{1}{\sqrt{x+1}+\sqrt{-1}}\\
\lim=\frac{1}{2\sqrt{-1}}
\end{aligned}
\]
\[\boxed{\frac{1}{2\sqrt{-1}}}\]


\paragraph{Ejercicio 522.} Calcula $\displaystyle\lim_{x\to -1}\frac{\sqrt{x+2}-\sqrt{1}}{x--1}$.

\[
\begin{aligned}
\frac{\sqrt{x+2}-\sqrt{1}}{x--1}\cdot \frac{\sqrt{x+2}+\sqrt{1}}{\sqrt{x+2}+\sqrt{1}}\\
=\frac{1}{\sqrt{x+2}+\sqrt{1}}\\
\lim=\frac{1}{2\sqrt{1}}
\end{aligned}
\]
\[\boxed{\frac{1}{2\sqrt{1}}}\]


\paragraph{Ejercicio 523.} Calcula $\displaystyle\lim_{x\to 0}\frac{\sqrt{x+3}-\sqrt{3}}{x-0}$.

\[
\begin{aligned}
\frac{\sqrt{x+3}-\sqrt{3}}{x-0}\cdot \frac{\sqrt{x+3}+\sqrt{3}}{\sqrt{x+3}+\sqrt{3}}\\
=\frac{1}{\sqrt{x+3}+\sqrt{3}}\\
\lim=\frac{1}{2\sqrt{3}}
\end{aligned}
\]
\[\boxed{\frac{1}{2\sqrt{3}}}\]


\paragraph{Ejercicio 524.} Calcula $\displaystyle\lim_{x\to 1}\frac{\sqrt{x+4}-\sqrt{5}}{x-1}$.

\[
\begin{aligned}
\frac{\sqrt{x+4}-\sqrt{5}}{x-1}\cdot \frac{\sqrt{x+4}+\sqrt{5}}{\sqrt{x+4}+\sqrt{5}}\\
=\frac{1}{\sqrt{x+4}+\sqrt{5}}\\
\lim=\frac{1}{2\sqrt{5}}
\end{aligned}
\]
\[\boxed{\frac{1}{2\sqrt{5}}}\]


\paragraph{Ejercicio 525.} Calcula $\displaystyle\lim_{x\to 2}\frac{\sqrt{x+5}-\sqrt{7}}{x-2}$.

\[
\begin{aligned}
\frac{\sqrt{x+5}-\sqrt{7}}{x-2}\cdot \frac{\sqrt{x+5}+\sqrt{7}}{\sqrt{x+5}+\sqrt{7}}\\
=\frac{1}{\sqrt{x+5}+\sqrt{7}}\\
\lim=\frac{1}{2\sqrt{7}}
\end{aligned}
\]
\[\boxed{\frac{1}{2\sqrt{7}}}\]


\paragraph{Ejercicio 526.} Calcula $\displaystyle\lim_{x\to -2}\frac{\sqrt{x+6}-\sqrt{4}}{x--2}$.

\[
\begin{aligned}
\frac{\sqrt{x+6}-\sqrt{4}}{x--2}\cdot \frac{\sqrt{x+6}+\sqrt{4}}{\sqrt{x+6}+\sqrt{4}}\\
=\frac{1}{\sqrt{x+6}+\sqrt{4}}\\
\lim=\frac{1}{2\sqrt{4}}
\end{aligned}
\]
\[\boxed{\frac{1}{2\sqrt{4}}}\]


\paragraph{Ejercicio 527.} Calcula $\displaystyle\lim_{x\to -1}\frac{\sqrt{x+7}-\sqrt{6}}{x--1}$.

\[
\begin{aligned}
\frac{\sqrt{x+7}-\sqrt{6}}{x--1}\cdot \frac{\sqrt{x+7}+\sqrt{6}}{\sqrt{x+7}+\sqrt{6}}\\
=\frac{1}{\sqrt{x+7}+\sqrt{6}}\\
\lim=\frac{1}{2\sqrt{6}}
\end{aligned}
\]
\[\boxed{\frac{1}{2\sqrt{6}}}\]


\paragraph{Ejercicio 528.} Calcula $\displaystyle\lim_{x\to 0}\frac{\sqrt{x+1}-\sqrt{1}}{x-0}$.

\[
\begin{aligned}
\frac{\sqrt{x+1}-\sqrt{1}}{x-0}\cdot \frac{\sqrt{x+1}+\sqrt{1}}{\sqrt{x+1}+\sqrt{1}}\\
=\frac{1}{\sqrt{x+1}+\sqrt{1}}\\
\lim=\frac{1}{2\sqrt{1}}
\end{aligned}
\]
\[\boxed{\frac{1}{2\sqrt{1}}}\]


\paragraph{Ejercicio 529.} Calcula $\displaystyle\lim_{x\to 1}\frac{\sqrt{x+2}-\sqrt{3}}{x-1}$.

\[
\begin{aligned}
\frac{\sqrt{x+2}-\sqrt{3}}{x-1}\cdot \frac{\sqrt{x+2}+\sqrt{3}}{\sqrt{x+2}+\sqrt{3}}\\
=\frac{1}{\sqrt{x+2}+\sqrt{3}}\\
\lim=\frac{1}{2\sqrt{3}}
\end{aligned}
\]
\[\boxed{\frac{1}{2\sqrt{3}}}\]


\paragraph{Ejercicio 530.} Calcula $\displaystyle\lim_{x\to 2}\frac{\sqrt{x+3}-\sqrt{5}}{x-2}$.

\[
\begin{aligned}
\frac{\sqrt{x+3}-\sqrt{5}}{x-2}\cdot \frac{\sqrt{x+3}+\sqrt{5}}{\sqrt{x+3}+\sqrt{5}}\\
=\frac{1}{\sqrt{x+3}+\sqrt{5}}\\
\lim=\frac{1}{2\sqrt{5}}
\end{aligned}
\]
\[\boxed{\frac{1}{2\sqrt{5}}}\]


\paragraph{Ejercicio 531.} Calcula $\displaystyle\lim_{x\to -2}\frac{\sqrt{x+4}-\sqrt{2}}{x--2}$.

\[
\begin{aligned}
\frac{\sqrt{x+4}-\sqrt{2}}{x--2}\cdot \frac{\sqrt{x+4}+\sqrt{2}}{\sqrt{x+4}+\sqrt{2}}\\
=\frac{1}{\sqrt{x+4}+\sqrt{2}}\\
\lim=\frac{1}{2\sqrt{2}}
\end{aligned}
\]
\[\boxed{\frac{1}{2\sqrt{2}}}\]


\paragraph{Ejercicio 532.} Calcula $\displaystyle\lim_{x\to -1}\frac{\sqrt{x+5}-\sqrt{4}}{x--1}$.

\[
\begin{aligned}
\frac{\sqrt{x+5}-\sqrt{4}}{x--1}\cdot \frac{\sqrt{x+5}+\sqrt{4}}{\sqrt{x+5}+\sqrt{4}}\\
=\frac{1}{\sqrt{x+5}+\sqrt{4}}\\
\lim=\frac{1}{2\sqrt{4}}
\end{aligned}
\]
\[\boxed{\frac{1}{2\sqrt{4}}}\]


\paragraph{Ejercicio 533.} Calcula $\displaystyle\lim_{x\to 0}\frac{\sqrt{x+6}-\sqrt{6}}{x-0}$.

\[
\begin{aligned}
\frac{\sqrt{x+6}-\sqrt{6}}{x-0}\cdot \frac{\sqrt{x+6}+\sqrt{6}}{\sqrt{x+6}+\sqrt{6}}\\
=\frac{1}{\sqrt{x+6}+\sqrt{6}}\\
\lim=\frac{1}{2\sqrt{6}}
\end{aligned}
\]
\[\boxed{\frac{1}{2\sqrt{6}}}\]


\paragraph{Ejercicio 534.} Calcula $\displaystyle\lim_{x\to 1}\frac{\sqrt{x+7}-\sqrt{8}}{x-1}$.

\[
\begin{aligned}
\frac{\sqrt{x+7}-\sqrt{8}}{x-1}\cdot \frac{\sqrt{x+7}+\sqrt{8}}{\sqrt{x+7}+\sqrt{8}}\\
=\frac{1}{\sqrt{x+7}+\sqrt{8}}\\
\lim=\frac{1}{2\sqrt{8}}
\end{aligned}
\]
\[\boxed{\frac{1}{2\sqrt{8}}}\]


\paragraph{Ejercicio 535.} Calcula $\displaystyle\lim_{x\to 2}\frac{\sqrt{x+1}-\sqrt{3}}{x-2}$.

\[
\begin{aligned}
\frac{\sqrt{x+1}-\sqrt{3}}{x-2}\cdot \frac{\sqrt{x+1}+\sqrt{3}}{\sqrt{x+1}+\sqrt{3}}\\
=\frac{1}{\sqrt{x+1}+\sqrt{3}}\\
\lim=\frac{1}{2\sqrt{3}}
\end{aligned}
\]
\[\boxed{\frac{1}{2\sqrt{3}}}\]


\paragraph{Ejercicio 536.} Calcula $\displaystyle\lim_{x\to -2}\frac{\sqrt{x+2}-\sqrt{0}}{x--2}$.

\[
\begin{aligned}
\frac{\sqrt{x+2}-\sqrt{0}}{x--2}\cdot \frac{\sqrt{x+2}+\sqrt{0}}{\sqrt{x+2}+\sqrt{0}}\\
=\frac{1}{\sqrt{x+2}+\sqrt{0}}\\
\lim=\frac{1}{2\sqrt{0}}
\end{aligned}
\]
\[\boxed{\frac{1}{2\sqrt{0}}}\]


\paragraph{Ejercicio 537.} Calcula $\displaystyle\lim_{x\to -1}\frac{\sqrt{x+3}-\sqrt{2}}{x--1}$.

\[
\begin{aligned}
\frac{\sqrt{x+3}-\sqrt{2}}{x--1}\cdot \frac{\sqrt{x+3}+\sqrt{2}}{\sqrt{x+3}+\sqrt{2}}\\
=\frac{1}{\sqrt{x+3}+\sqrt{2}}\\
\lim=\frac{1}{2\sqrt{2}}
\end{aligned}
\]
\[\boxed{\frac{1}{2\sqrt{2}}}\]


\paragraph{Ejercicio 538.} Calcula $\displaystyle\lim_{x\to 0}\frac{\sqrt{x+4}-\sqrt{4}}{x-0}$.

\[
\begin{aligned}
\frac{\sqrt{x+4}-\sqrt{4}}{x-0}\cdot \frac{\sqrt{x+4}+\sqrt{4}}{\sqrt{x+4}+\sqrt{4}}\\
=\frac{1}{\sqrt{x+4}+\sqrt{4}}\\
\lim=\frac{1}{2\sqrt{4}}
\end{aligned}
\]
\[\boxed{\frac{1}{2\sqrt{4}}}\]


\paragraph{Ejercicio 539.} Calcula $\displaystyle\lim_{x\to 1}\frac{\sqrt{x+5}-\sqrt{6}}{x-1}$.

\[
\begin{aligned}
\frac{\sqrt{x+5}-\sqrt{6}}{x-1}\cdot \frac{\sqrt{x+5}+\sqrt{6}}{\sqrt{x+5}+\sqrt{6}}\\
=\frac{1}{\sqrt{x+5}+\sqrt{6}}\\
\lim=\frac{1}{2\sqrt{6}}
\end{aligned}
\]
\[\boxed{\frac{1}{2\sqrt{6}}}\]


\paragraph{Ejercicio 540.} Calcula $\displaystyle\lim_{x\to 2}\frac{\sqrt{x+6}-\sqrt{8}}{x-2}$.

\[
\begin{aligned}
\frac{\sqrt{x+6}-\sqrt{8}}{x-2}\cdot \frac{\sqrt{x+6}+\sqrt{8}}{\sqrt{x+6}+\sqrt{8}}\\
=\frac{1}{\sqrt{x+6}+\sqrt{8}}\\
\lim=\frac{1}{2\sqrt{8}}
\end{aligned}
\]
\[\boxed{\frac{1}{2\sqrt{8}}}\]


\newpage

\subsubsection*{Nivel 3 - Intermedio}

\paragraph{Ejercicio 541.} Calcula $\displaystyle\lim_{x\to0}\frac{1-\cos(1x)}{1x^2}$ usando conjugado trigonométrico.

\[
\begin{aligned}
(1-\cos(1x))(1+\cos(1x))=\sin^2(1x)\\
\frac{1-\cos(1x)}{1x^2}=\frac{\sin^2(1x)}{1x^2(1+\cos(1x))}\\
\lim=\frac{1}{2}
\end{aligned}
\]
\[\boxed{\frac{1}{2}}\]


\paragraph{Ejercicio 542.} Calcula $\displaystyle\lim_{x\to0}\frac{1-\cos(2x)}{2x^2}$ usando conjugado trigonométrico.

\[
\begin{aligned}
(1-\cos(2x))(1+\cos(2x))=\sin^2(2x)\\
\frac{1-\cos(2x)}{2x^2}=\frac{\sin^2(2x)}{2x^2(1+\cos(2x))}\\
\lim=\frac{4}{4}
\end{aligned}
\]
\[\boxed{\frac{4}{4}}\]


\paragraph{Ejercicio 543.} Calcula $\displaystyle\lim_{x\to0}\frac{1-\cos(3x)}{3x^2}$ usando conjugado trigonométrico.

\[
\begin{aligned}
(1-\cos(3x))(1+\cos(3x))=\sin^2(3x)\\
\frac{1-\cos(3x)}{3x^2}=\frac{\sin^2(3x)}{3x^2(1+\cos(3x))}\\
\lim=\frac{9}{6}
\end{aligned}
\]
\[\boxed{\frac{9}{6}}\]


\paragraph{Ejercicio 544.} Calcula $\displaystyle\lim_{x\to0}\frac{1-\cos(4x)}{4x^2}$ usando conjugado trigonométrico.

\[
\begin{aligned}
(1-\cos(4x))(1+\cos(4x))=\sin^2(4x)\\
\frac{1-\cos(4x)}{4x^2}=\frac{\sin^2(4x)}{4x^2(1+\cos(4x))}\\
\lim=\frac{16}{8}
\end{aligned}
\]
\[\boxed{\frac{16}{8}}\]


\paragraph{Ejercicio 545.} Calcula $\displaystyle\lim_{x\to0}\frac{1-\cos(5x)}{5x^2}$ usando conjugado trigonométrico.

\[
\begin{aligned}
(1-\cos(5x))(1+\cos(5x))=\sin^2(5x)\\
\frac{1-\cos(5x)}{5x^2}=\frac{\sin^2(5x)}{5x^2(1+\cos(5x))}\\
\lim=\frac{25}{10}
\end{aligned}
\]
\[\boxed{\frac{25}{10}}\]


\paragraph{Ejercicio 546.} Calcula $\displaystyle\lim_{x\to0}\frac{1-\cos(6x)}{1x^2}$ usando conjugado trigonométrico.

\[
\begin{aligned}
(1-\cos(6x))(1+\cos(6x))=\sin^2(6x)\\
\frac{1-\cos(6x)}{1x^2}=\frac{\sin^2(6x)}{1x^2(1+\cos(6x))}\\
\lim=\frac{36}{2}
\end{aligned}
\]
\[\boxed{\frac{36}{2}}\]


\paragraph{Ejercicio 547.} Calcula $\displaystyle\lim_{x\to0}\frac{1-\cos(7x)}{2x^2}$ usando conjugado trigonométrico.

\[
\begin{aligned}
(1-\cos(7x))(1+\cos(7x))=\sin^2(7x)\\
\frac{1-\cos(7x)}{2x^2}=\frac{\sin^2(7x)}{2x^2(1+\cos(7x))}\\
\lim=\frac{49}{4}
\end{aligned}
\]
\[\boxed{\frac{49}{4}}\]


\paragraph{Ejercicio 548.} Calcula $\displaystyle\lim_{x\to0}\frac{1-\cos(8x)}{3x^2}$ usando conjugado trigonométrico.

\[
\begin{aligned}
(1-\cos(8x))(1+\cos(8x))=\sin^2(8x)\\
\frac{1-\cos(8x)}{3x^2}=\frac{\sin^2(8x)}{3x^2(1+\cos(8x))}\\
\lim=\frac{64}{6}
\end{aligned}
\]
\[\boxed{\frac{64}{6}}\]


\paragraph{Ejercicio 549.} Calcula $\displaystyle\lim_{x\to0}\frac{1-\cos(1x)}{4x^2}$ usando conjugado trigonométrico.

\[
\begin{aligned}
(1-\cos(1x))(1+\cos(1x))=\sin^2(1x)\\
\frac{1-\cos(1x)}{4x^2}=\frac{\sin^2(1x)}{4x^2(1+\cos(1x))}\\
\lim=\frac{1}{8}
\end{aligned}
\]
\[\boxed{\frac{1}{8}}\]


\paragraph{Ejercicio 550.} Calcula $\displaystyle\lim_{x\to0}\frac{1-\cos(2x)}{5x^2}$ usando conjugado trigonométrico.

\[
\begin{aligned}
(1-\cos(2x))(1+\cos(2x))=\sin^2(2x)\\
\frac{1-\cos(2x)}{5x^2}=\frac{\sin^2(2x)}{5x^2(1+\cos(2x))}\\
\lim=\frac{4}{10}
\end{aligned}
\]
\[\boxed{\frac{4}{10}}\]


\paragraph{Ejercicio 551.} Calcula $\displaystyle\lim_{x\to0}\frac{1-\cos(3x)}{1x^2}$ usando conjugado trigonométrico.

\[
\begin{aligned}
(1-\cos(3x))(1+\cos(3x))=\sin^2(3x)\\
\frac{1-\cos(3x)}{1x^2}=\frac{\sin^2(3x)}{1x^2(1+\cos(3x))}\\
\lim=\frac{9}{2}
\end{aligned}
\]
\[\boxed{\frac{9}{2}}\]


\paragraph{Ejercicio 552.} Calcula $\displaystyle\lim_{x\to0}\frac{1-\cos(4x)}{2x^2}$ usando conjugado trigonométrico.

\[
\begin{aligned}
(1-\cos(4x))(1+\cos(4x))=\sin^2(4x)\\
\frac{1-\cos(4x)}{2x^2}=\frac{\sin^2(4x)}{2x^2(1+\cos(4x))}\\
\lim=\frac{16}{4}
\end{aligned}
\]
\[\boxed{\frac{16}{4}}\]


\paragraph{Ejercicio 553.} Calcula $\displaystyle\lim_{x\to0}\frac{1-\cos(5x)}{3x^2}$ usando conjugado trigonométrico.

\[
\begin{aligned}
(1-\cos(5x))(1+\cos(5x))=\sin^2(5x)\\
\frac{1-\cos(5x)}{3x^2}=\frac{\sin^2(5x)}{3x^2(1+\cos(5x))}\\
\lim=\frac{25}{6}
\end{aligned}
\]
\[\boxed{\frac{25}{6}}\]


\paragraph{Ejercicio 554.} Calcula $\displaystyle\lim_{x\to0}\frac{1-\cos(6x)}{4x^2}$ usando conjugado trigonométrico.

\[
\begin{aligned}
(1-\cos(6x))(1+\cos(6x))=\sin^2(6x)\\
\frac{1-\cos(6x)}{4x^2}=\frac{\sin^2(6x)}{4x^2(1+\cos(6x))}\\
\lim=\frac{36}{8}
\end{aligned}
\]
\[\boxed{\frac{36}{8}}\]


\paragraph{Ejercicio 555.} Calcula $\displaystyle\lim_{x\to0}\frac{1-\cos(7x)}{5x^2}$ usando conjugado trigonométrico.

\[
\begin{aligned}
(1-\cos(7x))(1+\cos(7x))=\sin^2(7x)\\
\frac{1-\cos(7x)}{5x^2}=\frac{\sin^2(7x)}{5x^2(1+\cos(7x))}\\
\lim=\frac{49}{10}
\end{aligned}
\]
\[\boxed{\frac{49}{10}}\]


\paragraph{Ejercicio 556.} Calcula $\displaystyle\lim_{x\to0}\frac{1-\cos(8x)}{1x^2}$ usando conjugado trigonométrico.

\[
\begin{aligned}
(1-\cos(8x))(1+\cos(8x))=\sin^2(8x)\\
\frac{1-\cos(8x)}{1x^2}=\frac{\sin^2(8x)}{1x^2(1+\cos(8x))}\\
\lim=\frac{64}{2}
\end{aligned}
\]
\[\boxed{\frac{64}{2}}\]


\paragraph{Ejercicio 557.} Calcula $\displaystyle\lim_{x\to0}\frac{1-\cos(1x)}{2x^2}$ usando conjugado trigonométrico.

\[
\begin{aligned}
(1-\cos(1x))(1+\cos(1x))=\sin^2(1x)\\
\frac{1-\cos(1x)}{2x^2}=\frac{\sin^2(1x)}{2x^2(1+\cos(1x))}\\
\lim=\frac{1}{4}
\end{aligned}
\]
\[\boxed{\frac{1}{4}}\]


\paragraph{Ejercicio 558.} Calcula $\displaystyle\lim_{x\to0}\frac{1-\cos(2x)}{3x^2}$ usando conjugado trigonométrico.

\[
\begin{aligned}
(1-\cos(2x))(1+\cos(2x))=\sin^2(2x)\\
\frac{1-\cos(2x)}{3x^2}=\frac{\sin^2(2x)}{3x^2(1+\cos(2x))}\\
\lim=\frac{4}{6}
\end{aligned}
\]
\[\boxed{\frac{4}{6}}\]


\paragraph{Ejercicio 559.} Calcula $\displaystyle\lim_{x\to0}\frac{1-\cos(3x)}{4x^2}$ usando conjugado trigonométrico.

\[
\begin{aligned}
(1-\cos(3x))(1+\cos(3x))=\sin^2(3x)\\
\frac{1-\cos(3x)}{4x^2}=\frac{\sin^2(3x)}{4x^2(1+\cos(3x))}\\
\lim=\frac{9}{8}
\end{aligned}
\]
\[\boxed{\frac{9}{8}}\]


\paragraph{Ejercicio 560.} Calcula $\displaystyle\lim_{x\to0}\frac{1-\cos(4x)}{5x^2}$ usando conjugado trigonométrico.

\[
\begin{aligned}
(1-\cos(4x))(1+\cos(4x))=\sin^2(4x)\\
\frac{1-\cos(4x)}{5x^2}=\frac{\sin^2(4x)}{5x^2(1+\cos(4x))}\\
\lim=\frac{16}{10}
\end{aligned}
\]
\[\boxed{\frac{16}{10}}\]


\newpage

\subsubsection*{Nivel 4 - Avanzado}

\paragraph{Ejercicio 561.} Calcula $\displaystyle\lim_{x\to0}\frac{\sqrt{1+\sin(1x)}-1}{x}$.

\[
\begin{aligned}
\frac{\sqrt{1+\sin(1x)}-1}{x}\cdot \frac{\sqrt{1+\sin(1x)}+1}{\sqrt{1+\sin(1x)}+1}\\
=\frac{\sin(1x)}{x(\sqrt{1+\sin(1x)}+1)}\\
\lim=\frac{1}{2}
\end{aligned}
\]
\[\boxed{\frac{1}{2}}\]


\paragraph{Ejercicio 562.} Calcula $\displaystyle\lim_{x\to0}\frac{\sqrt{1+\sin(2x)}-1}{x}$.

\[
\begin{aligned}
\frac{\sqrt{1+\sin(2x)}-1}{x}\cdot \frac{\sqrt{1+\sin(2x)}+1}{\sqrt{1+\sin(2x)}+1}\\
=\frac{\sin(2x)}{x(\sqrt{1+\sin(2x)}+1)}\\
\lim=\frac{2}{2}
\end{aligned}
\]
\[\boxed{\frac{2}{2}}\]


\paragraph{Ejercicio 563.} Calcula $\displaystyle\lim_{x\to0}\frac{\sqrt{1+\sin(3x)}-1}{x}$.

\[
\begin{aligned}
\frac{\sqrt{1+\sin(3x)}-1}{x}\cdot \frac{\sqrt{1+\sin(3x)}+1}{\sqrt{1+\sin(3x)}+1}\\
=\frac{\sin(3x)}{x(\sqrt{1+\sin(3x)}+1)}\\
\lim=\frac{3}{2}
\end{aligned}
\]
\[\boxed{\frac{3}{2}}\]


\paragraph{Ejercicio 564.} Calcula $\displaystyle\lim_{x\to0}\frac{\sqrt{1+\sin(4x)}-1}{x}$.

\[
\begin{aligned}
\frac{\sqrt{1+\sin(4x)}-1}{x}\cdot \frac{\sqrt{1+\sin(4x)}+1}{\sqrt{1+\sin(4x)}+1}\\
=\frac{\sin(4x)}{x(\sqrt{1+\sin(4x)}+1)}\\
\lim=\frac{4}{2}
\end{aligned}
\]
\[\boxed{\frac{4}{2}}\]


\paragraph{Ejercicio 565.} Calcula $\displaystyle\lim_{x\to0}\frac{\sqrt{1+\sin(5x)}-1}{x}$.

\[
\begin{aligned}
\frac{\sqrt{1+\sin(5x)}-1}{x}\cdot \frac{\sqrt{1+\sin(5x)}+1}{\sqrt{1+\sin(5x)}+1}\\
=\frac{\sin(5x)}{x(\sqrt{1+\sin(5x)}+1)}\\
\lim=\frac{5}{2}
\end{aligned}
\]
\[\boxed{\frac{5}{2}}\]


\paragraph{Ejercicio 566.} Calcula $\displaystyle\lim_{x\to0}\frac{\sqrt{1+\sin(6x)}-1}{x}$.

\[
\begin{aligned}
\frac{\sqrt{1+\sin(6x)}-1}{x}\cdot \frac{\sqrt{1+\sin(6x)}+1}{\sqrt{1+\sin(6x)}+1}\\
=\frac{\sin(6x)}{x(\sqrt{1+\sin(6x)}+1)}\\
\lim=\frac{6}{2}
\end{aligned}
\]
\[\boxed{\frac{6}{2}}\]


\paragraph{Ejercicio 567.} Calcula $\displaystyle\lim_{x\to0}\frac{\sqrt{1+\sin(1x)}-1}{x}$.

\[
\begin{aligned}
\frac{\sqrt{1+\sin(1x)}-1}{x}\cdot \frac{\sqrt{1+\sin(1x)}+1}{\sqrt{1+\sin(1x)}+1}\\
=\frac{\sin(1x)}{x(\sqrt{1+\sin(1x)}+1)}\\
\lim=\frac{1}{2}
\end{aligned}
\]
\[\boxed{\frac{1}{2}}\]


\paragraph{Ejercicio 568.} Calcula $\displaystyle\lim_{x\to0}\frac{\sqrt{1+\sin(2x)}-1}{x}$.

\[
\begin{aligned}
\frac{\sqrt{1+\sin(2x)}-1}{x}\cdot \frac{\sqrt{1+\sin(2x)}+1}{\sqrt{1+\sin(2x)}+1}\\
=\frac{\sin(2x)}{x(\sqrt{1+\sin(2x)}+1)}\\
\lim=\frac{2}{2}
\end{aligned}
\]
\[\boxed{\frac{2}{2}}\]


\paragraph{Ejercicio 569.} Calcula $\displaystyle\lim_{x\to0}\frac{\sqrt{1+\sin(3x)}-1}{x}$.

\[
\begin{aligned}
\frac{\sqrt{1+\sin(3x)}-1}{x}\cdot \frac{\sqrt{1+\sin(3x)}+1}{\sqrt{1+\sin(3x)}+1}\\
=\frac{\sin(3x)}{x(\sqrt{1+\sin(3x)}+1)}\\
\lim=\frac{3}{2}
\end{aligned}
\]
\[\boxed{\frac{3}{2}}\]


\paragraph{Ejercicio 570.} Calcula $\displaystyle\lim_{x\to0}\frac{\sqrt{1+\sin(4x)}-1}{x}$.

\[
\begin{aligned}
\frac{\sqrt{1+\sin(4x)}-1}{x}\cdot \frac{\sqrt{1+\sin(4x)}+1}{\sqrt{1+\sin(4x)}+1}\\
=\frac{\sin(4x)}{x(\sqrt{1+\sin(4x)}+1)}\\
\lim=\frac{4}{2}
\end{aligned}
\]
\[\boxed{\frac{4}{2}}\]


\paragraph{Ejercicio 571.} Calcula $\displaystyle\lim_{x\to0}\frac{\sqrt{1+\sin(5x)}-1}{x}$.

\[
\begin{aligned}
\frac{\sqrt{1+\sin(5x)}-1}{x}\cdot \frac{\sqrt{1+\sin(5x)}+1}{\sqrt{1+\sin(5x)}+1}\\
=\frac{\sin(5x)}{x(\sqrt{1+\sin(5x)}+1)}\\
\lim=\frac{5}{2}
\end{aligned}
\]
\[\boxed{\frac{5}{2}}\]


\paragraph{Ejercicio 572.} Calcula $\displaystyle\lim_{x\to0}\frac{\sqrt{1+\sin(6x)}-1}{x}$.

\[
\begin{aligned}
\frac{\sqrt{1+\sin(6x)}-1}{x}\cdot \frac{\sqrt{1+\sin(6x)}+1}{\sqrt{1+\sin(6x)}+1}\\
=\frac{\sin(6x)}{x(\sqrt{1+\sin(6x)}+1)}\\
\lim=\frac{6}{2}
\end{aligned}
\]
\[\boxed{\frac{6}{2}}\]


\paragraph{Ejercicio 573.} Calcula $\displaystyle\lim_{x\to0}\frac{\sqrt{1+\sin(1x)}-1}{x}$.

\[
\begin{aligned}
\frac{\sqrt{1+\sin(1x)}-1}{x}\cdot \frac{\sqrt{1+\sin(1x)}+1}{\sqrt{1+\sin(1x)}+1}\\
=\frac{\sin(1x)}{x(\sqrt{1+\sin(1x)}+1)}\\
\lim=\frac{1}{2}
\end{aligned}
\]
\[\boxed{\frac{1}{2}}\]


\paragraph{Ejercicio 574.} Calcula $\displaystyle\lim_{x\to0}\frac{\sqrt{1+\sin(2x)}-1}{x}$.

\[
\begin{aligned}
\frac{\sqrt{1+\sin(2x)}-1}{x}\cdot \frac{\sqrt{1+\sin(2x)}+1}{\sqrt{1+\sin(2x)}+1}\\
=\frac{\sin(2x)}{x(\sqrt{1+\sin(2x)}+1)}\\
\lim=\frac{2}{2}
\end{aligned}
\]
\[\boxed{\frac{2}{2}}\]


\paragraph{Ejercicio 575.} Calcula $\displaystyle\lim_{x\to0}\frac{\sqrt{1+\sin(3x)}-1}{x}$.

\[
\begin{aligned}
\frac{\sqrt{1+\sin(3x)}-1}{x}\cdot \frac{\sqrt{1+\sin(3x)}+1}{\sqrt{1+\sin(3x)}+1}\\
=\frac{\sin(3x)}{x(\sqrt{1+\sin(3x)}+1)}\\
\lim=\frac{3}{2}
\end{aligned}
\]
\[\boxed{\frac{3}{2}}\]


\paragraph{Ejercicio 576.} Calcula $\displaystyle\lim_{x\to0}\frac{\sqrt{1+\sin(4x)}-1}{x}$.

\[
\begin{aligned}
\frac{\sqrt{1+\sin(4x)}-1}{x}\cdot \frac{\sqrt{1+\sin(4x)}+1}{\sqrt{1+\sin(4x)}+1}\\
=\frac{\sin(4x)}{x(\sqrt{1+\sin(4x)}+1)}\\
\lim=\frac{4}{2}
\end{aligned}
\]
\[\boxed{\frac{4}{2}}\]


\paragraph{Ejercicio 577.} Calcula $\displaystyle\lim_{x\to0}\frac{\sqrt{1+\sin(5x)}-1}{x}$.

\[
\begin{aligned}
\frac{\sqrt{1+\sin(5x)}-1}{x}\cdot \frac{\sqrt{1+\sin(5x)}+1}{\sqrt{1+\sin(5x)}+1}\\
=\frac{\sin(5x)}{x(\sqrt{1+\sin(5x)}+1)}\\
\lim=\frac{5}{2}
\end{aligned}
\]
\[\boxed{\frac{5}{2}}\]


\paragraph{Ejercicio 578.} Calcula $\displaystyle\lim_{x\to0}\frac{\sqrt{1+\sin(6x)}-1}{x}$.

\[
\begin{aligned}
\frac{\sqrt{1+\sin(6x)}-1}{x}\cdot \frac{\sqrt{1+\sin(6x)}+1}{\sqrt{1+\sin(6x)}+1}\\
=\frac{\sin(6x)}{x(\sqrt{1+\sin(6x)}+1)}\\
\lim=\frac{6}{2}
\end{aligned}
\]
\[\boxed{\frac{6}{2}}\]


\paragraph{Ejercicio 579.} Calcula $\displaystyle\lim_{x\to0}\frac{\sqrt{1+\sin(1x)}-1}{x}$.

\[
\begin{aligned}
\frac{\sqrt{1+\sin(1x)}-1}{x}\cdot \frac{\sqrt{1+\sin(1x)}+1}{\sqrt{1+\sin(1x)}+1}\\
=\frac{\sin(1x)}{x(\sqrt{1+\sin(1x)}+1)}\\
\lim=\frac{1}{2}
\end{aligned}
\]
\[\boxed{\frac{1}{2}}\]


\paragraph{Ejercicio 580.} Calcula $\displaystyle\lim_{x\to0}\frac{\sqrt{1+\sin(2x)}-1}{x}$.

\[
\begin{aligned}
\frac{\sqrt{1+\sin(2x)}-1}{x}\cdot \frac{\sqrt{1+\sin(2x)}+1}{\sqrt{1+\sin(2x)}+1}\\
=\frac{\sin(2x)}{x(\sqrt{1+\sin(2x)}+1)}\\
\lim=\frac{2}{2}
\end{aligned}
\]
\[\boxed{\frac{2}{2}}\]


\newpage

\subsubsection*{Nivel 5 - Imposible}

\paragraph{Ejercicio 581.} Calcula $\displaystyle \lim_{x\to0}\frac{\sqrt{1+\sin x}-\sqrt{1-\sin x}}{x}$.

\[
\begin{aligned}
\frac{\sqrt{1+\sin x}-\sqrt{1-\sin x}}{x}\cdot\frac{\sqrt{1+\sin x}+\sqrt{1-\sin x}}{\sqrt{1+\sin x}+\sqrt{1-\sin x}}\\
=\frac{2\sin x}{x(\sqrt{1+\sin x}+\sqrt{1-\sin x})}\\
\lim=\frac{2}{2}=1
\end{aligned}
\]
\[\boxed{1}\]


\paragraph{Ejercicio 582.} Calcula $\displaystyle \lim_{x\to0}\frac{\sqrt{\cos x}-1}{x^2}$.

\[
\begin{aligned}
\frac{\sqrt{\cos x}-1}{x^2}\cdot\frac{\sqrt{\cos x}+1}{\sqrt{\cos x}+1}\\
=\frac{\cos x-1}{x^2(\sqrt{\cos x}+1)}\\
=-\frac{1-\cos x}{x^2(\sqrt{\cos x}+1)}\\
\lim=-\frac14
\end{aligned}
\]
\[\boxed{-\frac14}\]


\paragraph{Ejercicio 583.} Calcula $\displaystyle \lim_{x\to0}\frac{\sin x}{\sqrt{1+x}-\sqrt{1-x}}$.

\[
\begin{aligned}
\sqrt{1+x}-\sqrt{1-x}=\frac{2x}{\sqrt{1+x}+\sqrt{1-x}}\\
\frac{\sin x}{\sqrt{1+x}-\sqrt{1-x}}=\frac{\sin x}{x}\frac{\sqrt{1+x}+\sqrt{1-x}}{2}\\
\lim=1
\end{aligned}
\]
\[\boxed{1}\]


\paragraph{Ejercicio 584.} Calcula $\displaystyle \lim_{x\to0}\frac{1-\cos x}{x(\sqrt{1+x}-1)}$.

\[
\begin{aligned}
1-\cos x\sim\frac{x^2}{2}\\
\sqrt{1+x}-1=\frac{x}{\sqrt{1+x}+1}\sim\frac{x}{2}\\
x(\sqrt{1+x}-1)\sim\frac{x^2}{2}\\
\lim=1
\end{aligned}
\]
\[\boxed{1}\]


\paragraph{Ejercicio 585.} Calcula $\displaystyle \lim_{x\to0}\frac{\sqrt{1+\tan x}-\sqrt{1-\tan x}}{\tan(2x)}$.

\[
\begin{aligned}
\frac{\sqrt{1+\tan x}-\sqrt{1-\tan x}}{\tan(2x)}\cdot\frac{\sqrt{1+\tan x}+\sqrt{1-\tan x}}{\sqrt{1+\tan x}+\sqrt{1-\tan x}}\\
=\frac{2\tan x}{\tan(2x)(\sqrt{1+\tan x}+\sqrt{1-\tan x})}\\
\frac{\tan x}{\tan(2x)}\to\frac12\\
\lim=\frac12
\end{aligned}
\]
\[\boxed{\frac12}\]


\paragraph{Ejercicio 586.} Calcula $\displaystyle \lim_{x\to0}\frac{\sqrt{1+\sin x}-\sqrt{1-\sin x}}{x}$.

\[
\begin{aligned}
\frac{\sqrt{1+\sin x}-\sqrt{1-\sin x}}{x}\cdot\frac{\sqrt{1+\sin x}+\sqrt{1-\sin x}}{\sqrt{1+\sin x}+\sqrt{1-\sin x}}\\
=\frac{2\sin x}{x(\sqrt{1+\sin x}+\sqrt{1-\sin x})}\\
\lim=\frac{2}{2}=1
\end{aligned}
\]
\[\boxed{1}\]


\paragraph{Ejercicio 587.} Calcula $\displaystyle \lim_{x\to0}\frac{\sqrt{\cos x}-1}{x^2}$.

\[
\begin{aligned}
\frac{\sqrt{\cos x}-1}{x^2}\cdot\frac{\sqrt{\cos x}+1}{\sqrt{\cos x}+1}\\
=\frac{\cos x-1}{x^2(\sqrt{\cos x}+1)}\\
=-\frac{1-\cos x}{x^2(\sqrt{\cos x}+1)}\\
\lim=-\frac14
\end{aligned}
\]
\[\boxed{-\frac14}\]


\paragraph{Ejercicio 588.} Calcula $\displaystyle \lim_{x\to0}\frac{\sin x}{\sqrt{1+x}-\sqrt{1-x}}$.

\[
\begin{aligned}
\sqrt{1+x}-\sqrt{1-x}=\frac{2x}{\sqrt{1+x}+\sqrt{1-x}}\\
\frac{\sin x}{\sqrt{1+x}-\sqrt{1-x}}=\frac{\sin x}{x}\frac{\sqrt{1+x}+\sqrt{1-x}}{2}\\
\lim=1
\end{aligned}
\]
\[\boxed{1}\]


\paragraph{Ejercicio 589.} Calcula $\displaystyle \lim_{x\to0}\frac{1-\cos x}{x(\sqrt{1+x}-1)}$.

\[
\begin{aligned}
1-\cos x\sim\frac{x^2}{2}\\
\sqrt{1+x}-1=\frac{x}{\sqrt{1+x}+1}\sim\frac{x}{2}\\
x(\sqrt{1+x}-1)\sim\frac{x^2}{2}\\
\lim=1
\end{aligned}
\]
\[\boxed{1}\]


\paragraph{Ejercicio 590.} Calcula $\displaystyle \lim_{x\to0}\frac{\sqrt{1+\tan x}-\sqrt{1-\tan x}}{\tan(2x)}$.

\[
\begin{aligned}
\frac{\sqrt{1+\tan x}-\sqrt{1-\tan x}}{\tan(2x)}\cdot\frac{\sqrt{1+\tan x}+\sqrt{1-\tan x}}{\sqrt{1+\tan x}+\sqrt{1-\tan x}}\\
=\frac{2\tan x}{\tan(2x)(\sqrt{1+\tan x}+\sqrt{1-\tan x})}\\
\frac{\tan x}{\tan(2x)}\to\frac12\\
\lim=\frac12
\end{aligned}
\]
\[\boxed{\frac12}\]


\paragraph{Ejercicio 591.} Calcula $\displaystyle \lim_{x\to0}\frac{\sqrt{1+\sin x}-\sqrt{1-\sin x}}{x}$.

\[
\begin{aligned}
\frac{\sqrt{1+\sin x}-\sqrt{1-\sin x}}{x}\cdot\frac{\sqrt{1+\sin x}+\sqrt{1-\sin x}}{\sqrt{1+\sin x}+\sqrt{1-\sin x}}\\
=\frac{2\sin x}{x(\sqrt{1+\sin x}+\sqrt{1-\sin x})}\\
\lim=\frac{2}{2}=1
\end{aligned}
\]
\[\boxed{1}\]


\paragraph{Ejercicio 592.} Calcula $\displaystyle \lim_{x\to0}\frac{\sqrt{\cos x}-1}{x^2}$.

\[
\begin{aligned}
\frac{\sqrt{\cos x}-1}{x^2}\cdot\frac{\sqrt{\cos x}+1}{\sqrt{\cos x}+1}\\
=\frac{\cos x-1}{x^2(\sqrt{\cos x}+1)}\\
=-\frac{1-\cos x}{x^2(\sqrt{\cos x}+1)}\\
\lim=-\frac14
\end{aligned}
\]
\[\boxed{-\frac14}\]


\paragraph{Ejercicio 593.} Calcula $\displaystyle \lim_{x\to0}\frac{\sin x}{\sqrt{1+x}-\sqrt{1-x}}$.

\[
\begin{aligned}
\sqrt{1+x}-\sqrt{1-x}=\frac{2x}{\sqrt{1+x}+\sqrt{1-x}}\\
\frac{\sin x}{\sqrt{1+x}-\sqrt{1-x}}=\frac{\sin x}{x}\frac{\sqrt{1+x}+\sqrt{1-x}}{2}\\
\lim=1
\end{aligned}
\]
\[\boxed{1}\]


\paragraph{Ejercicio 594.} Calcula $\displaystyle \lim_{x\to0}\frac{1-\cos x}{x(\sqrt{1+x}-1)}$.

\[
\begin{aligned}
1-\cos x\sim\frac{x^2}{2}\\
\sqrt{1+x}-1=\frac{x}{\sqrt{1+x}+1}\sim\frac{x}{2}\\
x(\sqrt{1+x}-1)\sim\frac{x^2}{2}\\
\lim=1
\end{aligned}
\]
\[\boxed{1}\]


\paragraph{Ejercicio 595.} Calcula $\displaystyle \lim_{x\to0}\frac{\sqrt{1+\tan x}-\sqrt{1-\tan x}}{\tan(2x)}$.

\[
\begin{aligned}
\frac{\sqrt{1+\tan x}-\sqrt{1-\tan x}}{\tan(2x)}\cdot\frac{\sqrt{1+\tan x}+\sqrt{1-\tan x}}{\sqrt{1+\tan x}+\sqrt{1-\tan x}}\\
=\frac{2\tan x}{\tan(2x)(\sqrt{1+\tan x}+\sqrt{1-\tan x})}\\
\frac{\tan x}{\tan(2x)}\to\frac12\\
\lim=\frac12
\end{aligned}
\]
\[\boxed{\frac12}\]


\paragraph{Ejercicio 596.} Calcula $\displaystyle \lim_{x\to0}\frac{\sqrt{1+\sin x}-\sqrt{1-\sin x}}{x}$.

\[
\begin{aligned}
\frac{\sqrt{1+\sin x}-\sqrt{1-\sin x}}{x}\cdot\frac{\sqrt{1+\sin x}+\sqrt{1-\sin x}}{\sqrt{1+\sin x}+\sqrt{1-\sin x}}\\
=\frac{2\sin x}{x(\sqrt{1+\sin x}+\sqrt{1-\sin x})}\\
\lim=\frac{2}{2}=1
\end{aligned}
\]
\[\boxed{1}\]


\paragraph{Ejercicio 597.} Calcula $\displaystyle \lim_{x\to0}\frac{\sqrt{\cos x}-1}{x^2}$.

\[
\begin{aligned}
\frac{\sqrt{\cos x}-1}{x^2}\cdot\frac{\sqrt{\cos x}+1}{\sqrt{\cos x}+1}\\
=\frac{\cos x-1}{x^2(\sqrt{\cos x}+1)}\\
=-\frac{1-\cos x}{x^2(\sqrt{\cos x}+1)}\\
\lim=-\frac14
\end{aligned}
\]
\[\boxed{-\frac14}\]


\paragraph{Ejercicio 598.} Calcula $\displaystyle \lim_{x\to0}\frac{\sin x}{\sqrt{1+x}-\sqrt{1-x}}$.

\[
\begin{aligned}
\sqrt{1+x}-\sqrt{1-x}=\frac{2x}{\sqrt{1+x}+\sqrt{1-x}}\\
\frac{\sin x}{\sqrt{1+x}-\sqrt{1-x}}=\frac{\sin x}{x}\frac{\sqrt{1+x}+\sqrt{1-x}}{2}\\
\lim=1
\end{aligned}
\]
\[\boxed{1}\]


\paragraph{Ejercicio 599.} Calcula $\displaystyle \lim_{x\to0}\frac{1-\cos x}{x(\sqrt{1+x}-1)}$.

\[
\begin{aligned}
1-\cos x\sim\frac{x^2}{2}\\
\sqrt{1+x}-1=\frac{x}{\sqrt{1+x}+1}\sim\frac{x}{2}\\
x(\sqrt{1+x}-1)\sim\frac{x^2}{2}\\
\lim=1
\end{aligned}
\]
\[\boxed{1}\]


\paragraph{Ejercicio 600.} Calcula $\displaystyle \lim_{x\to0}\frac{\sqrt{1+\tan x}-\sqrt{1-\tan x}}{\tan(2x)}$.

\[
\begin{aligned}
\frac{\sqrt{1+\tan x}-\sqrt{1-\tan x}}{\tan(2x)}\cdot\frac{\sqrt{1+\tan x}+\sqrt{1-\tan x}}{\sqrt{1+\tan x}+\sqrt{1-\tan x}}\\
=\frac{2\tan x}{\tan(2x)(\sqrt{1+\tan x}+\sqrt{1-\tan x})}\\
\frac{\tan x}{\tan(2x)}\to\frac12\\
\lim=\frac12
\end{aligned}
\]
\[\boxed{\frac12}\]


\newpage

\subsection*{Otros límites}

\addcontentsline{toc}{subsection}{Soluciones - Otros límites}

\subsubsection*{Nivel 1 - Básico}

\paragraph{Ejercicio 601.} Calcula $\displaystyle\lim_{x\to -3}(x^2+0x+0)$.

\[
\begin{aligned}
\lim=x^2+0x+0\big|_{x=-3}\\
=-3^2+0(-3)+0=9
\end{aligned}
\]
\[\boxed{9}\]


\paragraph{Ejercicio 602.} Calcula $\displaystyle\lim_{x\to -2}(x^2+1x+1)$.

\[
\begin{aligned}
\lim=x^2+1x+1\big|_{x=-2}\\
=-2^2+1(-2)+1=3
\end{aligned}
\]
\[\boxed{3}\]


\paragraph{Ejercicio 603.} Calcula $\displaystyle\lim_{x\to -1}(x^2+2x+2)$.

\[
\begin{aligned}
\lim=x^2+2x+2\big|_{x=-1}\\
=-1^2+2(-1)+2=1
\end{aligned}
\]
\[\boxed{1}\]


\paragraph{Ejercicio 604.} Calcula $\displaystyle\lim_{x\to 0}(x^2+3x+3)$.

\[
\begin{aligned}
\lim=x^2+3x+3\big|_{x=0}\\
=0^2+3(0)+3=3
\end{aligned}
\]
\[\boxed{3}\]


\paragraph{Ejercicio 605.} Calcula $\displaystyle\lim_{x\to 1}(x^2+4x+0)$.

\[
\begin{aligned}
\lim=x^2+4x+0\big|_{x=1}\\
=1^2+4(1)+0=5
\end{aligned}
\]
\[\boxed{5}\]


\paragraph{Ejercicio 606.} Calcula $\displaystyle\lim_{x\to 2}(x^2+0x+1)$.

\[
\begin{aligned}
\lim=x^2+0x+1\big|_{x=2}\\
=2^2+0(2)+1=5
\end{aligned}
\]
\[\boxed{5}\]


\paragraph{Ejercicio 607.} Calcula $\displaystyle\lim_{x\to 3}(x^2+1x+2)$.

\[
\begin{aligned}
\lim=x^2+1x+2\big|_{x=3}\\
=3^2+1(3)+2=14
\end{aligned}
\]
\[\boxed{14}\]


\paragraph{Ejercicio 608.} Calcula $\displaystyle\lim_{x\to -3}(x^2+2x+3)$.

\[
\begin{aligned}
\lim=x^2+2x+3\big|_{x=-3}\\
=-3^2+2(-3)+3=6
\end{aligned}
\]
\[\boxed{6}\]


\paragraph{Ejercicio 609.} Calcula $\displaystyle\lim_{x\to -2}(x^2+3x+0)$.

\[
\begin{aligned}
\lim=x^2+3x+0\big|_{x=-2}\\
=-2^2+3(-2)+0=-2
\end{aligned}
\]
\[\boxed{-2}\]


\paragraph{Ejercicio 610.} Calcula $\displaystyle\lim_{x\to -1}(x^2+4x+1)$.

\[
\begin{aligned}
\lim=x^2+4x+1\big|_{x=-1}\\
=-1^2+4(-1)+1=-2
\end{aligned}
\]
\[\boxed{-2}\]


\paragraph{Ejercicio 611.} Calcula $\displaystyle\lim_{x\to 0}(x^2+0x+2)$.

\[
\begin{aligned}
\lim=x^2+0x+2\big|_{x=0}\\
=0^2+0(0)+2=2
\end{aligned}
\]
\[\boxed{2}\]


\paragraph{Ejercicio 612.} Calcula $\displaystyle\lim_{x\to 1}(x^2+1x+3)$.

\[
\begin{aligned}
\lim=x^2+1x+3\big|_{x=1}\\
=1^2+1(1)+3=5
\end{aligned}
\]
\[\boxed{5}\]


\paragraph{Ejercicio 613.} Calcula $\displaystyle\lim_{x\to 2}(x^2+2x+0)$.

\[
\begin{aligned}
\lim=x^2+2x+0\big|_{x=2}\\
=2^2+2(2)+0=8
\end{aligned}
\]
\[\boxed{8}\]


\paragraph{Ejercicio 614.} Calcula $\displaystyle\lim_{x\to 3}(x^2+3x+1)$.

\[
\begin{aligned}
\lim=x^2+3x+1\big|_{x=3}\\
=3^2+3(3)+1=19
\end{aligned}
\]
\[\boxed{19}\]


\paragraph{Ejercicio 615.} Calcula $\displaystyle\lim_{x\to -3}(x^2+4x+2)$.

\[
\begin{aligned}
\lim=x^2+4x+2\big|_{x=-3}\\
=-3^2+4(-3)+2=-1
\end{aligned}
\]
\[\boxed{-1}\]


\paragraph{Ejercicio 616.} Calcula $\displaystyle\lim_{x\to -2}(x^2+0x+3)$.

\[
\begin{aligned}
\lim=x^2+0x+3\big|_{x=-2}\\
=-2^2+0(-2)+3=7
\end{aligned}
\]
\[\boxed{7}\]


\paragraph{Ejercicio 617.} Calcula $\displaystyle\lim_{x\to -1}(x^2+1x+0)$.

\[
\begin{aligned}
\lim=x^2+1x+0\big|_{x=-1}\\
=-1^2+1(-1)+0=0
\end{aligned}
\]
\[\boxed{0}\]


\paragraph{Ejercicio 618.} Calcula $\displaystyle\lim_{x\to 0}(x^2+2x+1)$.

\[
\begin{aligned}
\lim=x^2+2x+1\big|_{x=0}\\
=0^2+2(0)+1=1
\end{aligned}
\]
\[\boxed{1}\]


\paragraph{Ejercicio 619.} Calcula $\displaystyle\lim_{x\to 1}(x^2+3x+2)$.

\[
\begin{aligned}
\lim=x^2+3x+2\big|_{x=1}\\
=1^2+3(1)+2=6
\end{aligned}
\]
\[\boxed{6}\]


\paragraph{Ejercicio 620.} Calcula $\displaystyle\lim_{x\to 2}(x^2+4x+3)$.

\[
\begin{aligned}
\lim=x^2+4x+3\big|_{x=2}\\
=2^2+4(2)+3=15
\end{aligned}
\]
\[\boxed{15}\]


\newpage

\subsubsection*{Nivel 2 - Intermedio bajo}

\paragraph{Ejercicio 621.} Calcula $\displaystyle\lim_{x\to -3}\frac{|x--3|}{x--3}$. 

\[
\begin{aligned}
x\to -3^-:\quad \frac{|x--3|}{x--3}=-1\\
x\to -3^+:\quad \frac{|x--3|}{x--3}=1\\
-1\ne1\\
\lim\text{{ no existe}}
\end{aligned}
\]
\[\boxed{\text{{No existe}}}\]


\paragraph{Ejercicio 622.} Calcula $\displaystyle\lim_{x\to -2}\frac{|x--2|}{x--2}$. 

\[
\begin{aligned}
x\to -2^-:\quad \frac{|x--2|}{x--2}=-1\\
x\to -2^+:\quad \frac{|x--2|}{x--2}=1\\
-1\ne1\\
\lim\text{{ no existe}}
\end{aligned}
\]
\[\boxed{\text{{No existe}}}\]


\paragraph{Ejercicio 623.} Calcula $\displaystyle\lim_{x\to -1}\frac{|x--1|}{x--1}$. 

\[
\begin{aligned}
x\to -1^-:\quad \frac{|x--1|}{x--1}=-1\\
x\to -1^+:\quad \frac{|x--1|}{x--1}=1\\
-1\ne1\\
\lim\text{{ no existe}}
\end{aligned}
\]
\[\boxed{\text{{No existe}}}\]


\paragraph{Ejercicio 624.} Calcula $\displaystyle\lim_{x\to 2}\frac{|x-2|}{x-2}$. 

\[
\begin{aligned}
x\to 2^-:\quad \frac{|x-2|}{x-2}=-1\\
x\to 2^+:\quad \frac{|x-2|}{x-2}=1\\
-1\ne1\\
\lim\text{{ no existe}}
\end{aligned}
\]
\[\boxed{\text{{No existe}}}\]


\paragraph{Ejercicio 625.} Calcula $\displaystyle\lim_{x\to 1}\frac{|x-1|}{x-1}$. 

\[
\begin{aligned}
x\to 1^-:\quad \frac{|x-1|}{x-1}=-1\\
x\to 1^+:\quad \frac{|x-1|}{x-1}=1\\
-1\ne1\\
\lim\text{{ no existe}}
\end{aligned}
\]
\[\boxed{\text{{No existe}}}\]


\paragraph{Ejercicio 626.} Calcula $\displaystyle\lim_{x\to 2}\frac{|x-2|}{x-2}$. 

\[
\begin{aligned}
x\to 2^-:\quad \frac{|x-2|}{x-2}=-1\\
x\to 2^+:\quad \frac{|x-2|}{x-2}=1\\
-1\ne1\\
\lim\text{{ no existe}}
\end{aligned}
\]
\[\boxed{\text{{No existe}}}\]


\paragraph{Ejercicio 627.} Calcula $\displaystyle\lim_{x\to 3}\frac{|x-3|}{x-3}$. 

\[
\begin{aligned}
x\to 3^-:\quad \frac{|x-3|}{x-3}=-1\\
x\to 3^+:\quad \frac{|x-3|}{x-3}=1\\
-1\ne1\\
\lim\text{{ no existe}}
\end{aligned}
\]
\[\boxed{\text{{No existe}}}\]


\paragraph{Ejercicio 628.} Calcula $\displaystyle\lim_{x\to -3}\frac{|x--3|}{x--3}$. 

\[
\begin{aligned}
x\to -3^-:\quad \frac{|x--3|}{x--3}=-1\\
x\to -3^+:\quad \frac{|x--3|}{x--3}=1\\
-1\ne1\\
\lim\text{{ no existe}}
\end{aligned}
\]
\[\boxed{\text{{No existe}}}\]


\paragraph{Ejercicio 629.} Calcula $\displaystyle\lim_{x\to -2}\frac{|x--2|}{x--2}$. 

\[
\begin{aligned}
x\to -2^-:\quad \frac{|x--2|}{x--2}=-1\\
x\to -2^+:\quad \frac{|x--2|}{x--2}=1\\
-1\ne1\\
\lim\text{{ no existe}}
\end{aligned}
\]
\[\boxed{\text{{No existe}}}\]


\paragraph{Ejercicio 630.} Calcula $\displaystyle\lim_{x\to -1}\frac{|x--1|}{x--1}$. 

\[
\begin{aligned}
x\to -1^-:\quad \frac{|x--1|}{x--1}=-1\\
x\to -1^+:\quad \frac{|x--1|}{x--1}=1\\
-1\ne1\\
\lim\text{{ no existe}}
\end{aligned}
\]
\[\boxed{\text{{No existe}}}\]


\paragraph{Ejercicio 631.} Calcula $\displaystyle\lim_{x\to 2}\frac{|x-2|}{x-2}$. 

\[
\begin{aligned}
x\to 2^-:\quad \frac{|x-2|}{x-2}=-1\\
x\to 2^+:\quad \frac{|x-2|}{x-2}=1\\
-1\ne1\\
\lim\text{{ no existe}}
\end{aligned}
\]
\[\boxed{\text{{No existe}}}\]


\paragraph{Ejercicio 632.} Calcula $\displaystyle\lim_{x\to 1}\frac{|x-1|}{x-1}$. 

\[
\begin{aligned}
x\to 1^-:\quad \frac{|x-1|}{x-1}=-1\\
x\to 1^+:\quad \frac{|x-1|}{x-1}=1\\
-1\ne1\\
\lim\text{{ no existe}}
\end{aligned}
\]
\[\boxed{\text{{No existe}}}\]


\paragraph{Ejercicio 633.} Calcula $\displaystyle\lim_{x\to 2}\frac{|x-2|}{x-2}$. 

\[
\begin{aligned}
x\to 2^-:\quad \frac{|x-2|}{x-2}=-1\\
x\to 2^+:\quad \frac{|x-2|}{x-2}=1\\
-1\ne1\\
\lim\text{{ no existe}}
\end{aligned}
\]
\[\boxed{\text{{No existe}}}\]


\paragraph{Ejercicio 634.} Calcula $\displaystyle\lim_{x\to 3}\frac{|x-3|}{x-3}$. 

\[
\begin{aligned}
x\to 3^-:\quad \frac{|x-3|}{x-3}=-1\\
x\to 3^+:\quad \frac{|x-3|}{x-3}=1\\
-1\ne1\\
\lim\text{{ no existe}}
\end{aligned}
\]
\[\boxed{\text{{No existe}}}\]


\paragraph{Ejercicio 635.} Calcula $\displaystyle\lim_{x\to -3}\frac{|x--3|}{x--3}$. 

\[
\begin{aligned}
x\to -3^-:\quad \frac{|x--3|}{x--3}=-1\\
x\to -3^+:\quad \frac{|x--3|}{x--3}=1\\
-1\ne1\\
\lim\text{{ no existe}}
\end{aligned}
\]
\[\boxed{\text{{No existe}}}\]


\paragraph{Ejercicio 636.} Calcula $\displaystyle\lim_{x\to -2}\frac{|x--2|}{x--2}$. 

\[
\begin{aligned}
x\to -2^-:\quad \frac{|x--2|}{x--2}=-1\\
x\to -2^+:\quad \frac{|x--2|}{x--2}=1\\
-1\ne1\\
\lim\text{{ no existe}}
\end{aligned}
\]
\[\boxed{\text{{No existe}}}\]


\paragraph{Ejercicio 637.} Calcula $\displaystyle\lim_{x\to -1}\frac{|x--1|}{x--1}$. 

\[
\begin{aligned}
x\to -1^-:\quad \frac{|x--1|}{x--1}=-1\\
x\to -1^+:\quad \frac{|x--1|}{x--1}=1\\
-1\ne1\\
\lim\text{{ no existe}}
\end{aligned}
\]
\[\boxed{\text{{No existe}}}\]


\paragraph{Ejercicio 638.} Calcula $\displaystyle\lim_{x\to 2}\frac{|x-2|}{x-2}$. 

\[
\begin{aligned}
x\to 2^-:\quad \frac{|x-2|}{x-2}=-1\\
x\to 2^+:\quad \frac{|x-2|}{x-2}=1\\
-1\ne1\\
\lim\text{{ no existe}}
\end{aligned}
\]
\[\boxed{\text{{No existe}}}\]


\paragraph{Ejercicio 639.} Calcula $\displaystyle\lim_{x\to 1}\frac{|x-1|}{x-1}$. 

\[
\begin{aligned}
x\to 1^-:\quad \frac{|x-1|}{x-1}=-1\\
x\to 1^+:\quad \frac{|x-1|}{x-1}=1\\
-1\ne1\\
\lim\text{{ no existe}}
\end{aligned}
\]
\[\boxed{\text{{No existe}}}\]


\paragraph{Ejercicio 640.} Calcula $\displaystyle\lim_{x\to 2}\frac{|x-2|}{x-2}$. 

\[
\begin{aligned}
x\to 2^-:\quad \frac{|x-2|}{x-2}=-1\\
x\to 2^+:\quad \frac{|x-2|}{x-2}=1\\
-1\ne1\\
\lim\text{{ no existe}}
\end{aligned}
\]
\[\boxed{\text{{No existe}}}\]


\newpage

\subsubsection*{Nivel 3 - Intermedio}

\paragraph{Ejercicio 641.} Calcula $\displaystyle\lim_{x\to0}x^2\sin\left(\frac{1}{x}\right)$.

\[
\begin{aligned}
-1\le\sin\left(\frac{1}{x}\right)\le1\\
-x^2\le x^2\sin\left(\frac{1}{x}\right)\le x^2\\
\lim_{x\to0}(-x^2)=0\\
\lim_{x\to0}x^2=0
\end{aligned}
\]
\[\boxed{0}\]


\paragraph{Ejercicio 642.} Calcula $\displaystyle\lim_{x\to0}x^2\sin\left(\frac{2}{x}\right)$.

\[
\begin{aligned}
-1\le\sin\left(\frac{2}{x}\right)\le1\\
-x^2\le x^2\sin\left(\frac{2}{x}\right)\le x^2\\
\lim_{x\to0}(-x^2)=0\\
\lim_{x\to0}x^2=0
\end{aligned}
\]
\[\boxed{0}\]


\paragraph{Ejercicio 643.} Calcula $\displaystyle\lim_{x\to0}x^2\sin\left(\frac{3}{x}\right)$.

\[
\begin{aligned}
-1\le\sin\left(\frac{3}{x}\right)\le1\\
-x^2\le x^2\sin\left(\frac{3}{x}\right)\le x^2\\
\lim_{x\to0}(-x^2)=0\\
\lim_{x\to0}x^2=0
\end{aligned}
\]
\[\boxed{0}\]


\paragraph{Ejercicio 644.} Calcula $\displaystyle\lim_{x\to0}x^2\sin\left(\frac{4}{x}\right)$.

\[
\begin{aligned}
-1\le\sin\left(\frac{4}{x}\right)\le1\\
-x^2\le x^2\sin\left(\frac{4}{x}\right)\le x^2\\
\lim_{x\to0}(-x^2)=0\\
\lim_{x\to0}x^2=0
\end{aligned}
\]
\[\boxed{0}\]


\paragraph{Ejercicio 645.} Calcula $\displaystyle\lim_{x\to0}x^2\sin\left(\frac{5}{x}\right)$.

\[
\begin{aligned}
-1\le\sin\left(\frac{5}{x}\right)\le1\\
-x^2\le x^2\sin\left(\frac{5}{x}\right)\le x^2\\
\lim_{x\to0}(-x^2)=0\\
\lim_{x\to0}x^2=0
\end{aligned}
\]
\[\boxed{0}\]


\paragraph{Ejercicio 646.} Calcula $\displaystyle\lim_{x\to0}x^2\sin\left(\frac{6}{x}\right)$.

\[
\begin{aligned}
-1\le\sin\left(\frac{6}{x}\right)\le1\\
-x^2\le x^2\sin\left(\frac{6}{x}\right)\le x^2\\
\lim_{x\to0}(-x^2)=0\\
\lim_{x\to0}x^2=0
\end{aligned}
\]
\[\boxed{0}\]


\paragraph{Ejercicio 647.} Calcula $\displaystyle\lim_{x\to0}x^2\sin\left(\frac{7}{x}\right)$.

\[
\begin{aligned}
-1\le\sin\left(\frac{7}{x}\right)\le1\\
-x^2\le x^2\sin\left(\frac{7}{x}\right)\le x^2\\
\lim_{x\to0}(-x^2)=0\\
\lim_{x\to0}x^2=0
\end{aligned}
\]
\[\boxed{0}\]


\paragraph{Ejercicio 648.} Calcula $\displaystyle\lim_{x\to0}x^2\sin\left(\frac{1}{x}\right)$.

\[
\begin{aligned}
-1\le\sin\left(\frac{1}{x}\right)\le1\\
-x^2\le x^2\sin\left(\frac{1}{x}\right)\le x^2\\
\lim_{x\to0}(-x^2)=0\\
\lim_{x\to0}x^2=0
\end{aligned}
\]
\[\boxed{0}\]


\paragraph{Ejercicio 649.} Calcula $\displaystyle\lim_{x\to0}x^2\sin\left(\frac{2}{x}\right)$.

\[
\begin{aligned}
-1\le\sin\left(\frac{2}{x}\right)\le1\\
-x^2\le x^2\sin\left(\frac{2}{x}\right)\le x^2\\
\lim_{x\to0}(-x^2)=0\\
\lim_{x\to0}x^2=0
\end{aligned}
\]
\[\boxed{0}\]


\paragraph{Ejercicio 650.} Calcula $\displaystyle\lim_{x\to0}x^2\sin\left(\frac{3}{x}\right)$.

\[
\begin{aligned}
-1\le\sin\left(\frac{3}{x}\right)\le1\\
-x^2\le x^2\sin\left(\frac{3}{x}\right)\le x^2\\
\lim_{x\to0}(-x^2)=0\\
\lim_{x\to0}x^2=0
\end{aligned}
\]
\[\boxed{0}\]


\paragraph{Ejercicio 651.} Calcula $\displaystyle\lim_{x\to0}x^2\sin\left(\frac{4}{x}\right)$.

\[
\begin{aligned}
-1\le\sin\left(\frac{4}{x}\right)\le1\\
-x^2\le x^2\sin\left(\frac{4}{x}\right)\le x^2\\
\lim_{x\to0}(-x^2)=0\\
\lim_{x\to0}x^2=0
\end{aligned}
\]
\[\boxed{0}\]


\paragraph{Ejercicio 652.} Calcula $\displaystyle\lim_{x\to0}x^2\sin\left(\frac{5}{x}\right)$.

\[
\begin{aligned}
-1\le\sin\left(\frac{5}{x}\right)\le1\\
-x^2\le x^2\sin\left(\frac{5}{x}\right)\le x^2\\
\lim_{x\to0}(-x^2)=0\\
\lim_{x\to0}x^2=0
\end{aligned}
\]
\[\boxed{0}\]


\paragraph{Ejercicio 653.} Calcula $\displaystyle\lim_{x\to0}x^2\sin\left(\frac{6}{x}\right)$.

\[
\begin{aligned}
-1\le\sin\left(\frac{6}{x}\right)\le1\\
-x^2\le x^2\sin\left(\frac{6}{x}\right)\le x^2\\
\lim_{x\to0}(-x^2)=0\\
\lim_{x\to0}x^2=0
\end{aligned}
\]
\[\boxed{0}\]


\paragraph{Ejercicio 654.} Calcula $\displaystyle\lim_{x\to0}x^2\sin\left(\frac{7}{x}\right)$.

\[
\begin{aligned}
-1\le\sin\left(\frac{7}{x}\right)\le1\\
-x^2\le x^2\sin\left(\frac{7}{x}\right)\le x^2\\
\lim_{x\to0}(-x^2)=0\\
\lim_{x\to0}x^2=0
\end{aligned}
\]
\[\boxed{0}\]


\paragraph{Ejercicio 655.} Calcula $\displaystyle\lim_{x\to0}x^2\sin\left(\frac{1}{x}\right)$.

\[
\begin{aligned}
-1\le\sin\left(\frac{1}{x}\right)\le1\\
-x^2\le x^2\sin\left(\frac{1}{x}\right)\le x^2\\
\lim_{x\to0}(-x^2)=0\\
\lim_{x\to0}x^2=0
\end{aligned}
\]
\[\boxed{0}\]


\paragraph{Ejercicio 656.} Calcula $\displaystyle\lim_{x\to0}x^2\sin\left(\frac{2}{x}\right)$.

\[
\begin{aligned}
-1\le\sin\left(\frac{2}{x}\right)\le1\\
-x^2\le x^2\sin\left(\frac{2}{x}\right)\le x^2\\
\lim_{x\to0}(-x^2)=0\\
\lim_{x\to0}x^2=0
\end{aligned}
\]
\[\boxed{0}\]


\paragraph{Ejercicio 657.} Calcula $\displaystyle\lim_{x\to0}x^2\sin\left(\frac{3}{x}\right)$.

\[
\begin{aligned}
-1\le\sin\left(\frac{3}{x}\right)\le1\\
-x^2\le x^2\sin\left(\frac{3}{x}\right)\le x^2\\
\lim_{x\to0}(-x^2)=0\\
\lim_{x\to0}x^2=0
\end{aligned}
\]
\[\boxed{0}\]


\paragraph{Ejercicio 658.} Calcula $\displaystyle\lim_{x\to0}x^2\sin\left(\frac{4}{x}\right)$.

\[
\begin{aligned}
-1\le\sin\left(\frac{4}{x}\right)\le1\\
-x^2\le x^2\sin\left(\frac{4}{x}\right)\le x^2\\
\lim_{x\to0}(-x^2)=0\\
\lim_{x\to0}x^2=0
\end{aligned}
\]
\[\boxed{0}\]


\paragraph{Ejercicio 659.} Calcula $\displaystyle\lim_{x\to0}x^2\sin\left(\frac{5}{x}\right)$.

\[
\begin{aligned}
-1\le\sin\left(\frac{5}{x}\right)\le1\\
-x^2\le x^2\sin\left(\frac{5}{x}\right)\le x^2\\
\lim_{x\to0}(-x^2)=0\\
\lim_{x\to0}x^2=0
\end{aligned}
\]
\[\boxed{0}\]


\paragraph{Ejercicio 660.} Calcula $\displaystyle\lim_{x\to0}x^2\sin\left(\frac{6}{x}\right)$.

\[
\begin{aligned}
-1\le\sin\left(\frac{6}{x}\right)\le1\\
-x^2\le x^2\sin\left(\frac{6}{x}\right)\le x^2\\
\lim_{x\to0}(-x^2)=0\\
\lim_{x\to0}x^2=0
\end{aligned}
\]
\[\boxed{0}\]


\newpage

\subsubsection*{Nivel 4 - Avanzado}

\paragraph{Ejercicio 661.} Calcula $\displaystyle\lim_{x\to0}\frac{\ln(1+x)}{x}$.

\[
\begin{aligned}
\ln(1+x)=x+o(x)\\
\frac{\ln(1+x)}{x}=\frac{x+o(x)}{x}\\
\lim=1
\end{aligned}
\]
\[\boxed{1}\]


\paragraph{Ejercicio 662.} Calcula $\displaystyle\lim_{x\to0}\frac{e^x-1-\sin x}{x^2}$.

\[
\begin{aligned}
e^x=1+x+\frac{x^2}{2}+o(x^2)\\
\sin x=x+o(x^2)\\
e^x-1-\sin x=\frac{x^2}{2}+o(x^2)\\
\lim=\frac12
\end{aligned}
\]
\[\boxed{\frac12}\]


\paragraph{Ejercicio 663.} Calcula $\displaystyle\lim_{x\to0}\frac{\ln(1+x)}{x}$.

\[
\begin{aligned}
\ln(1+x)=x+o(x)\\
\frac{\ln(1+x)}{x}=\frac{x+o(x)}{x}\\
\lim=1
\end{aligned}
\]
\[\boxed{1}\]


\paragraph{Ejercicio 664.} Calcula $\displaystyle\lim_{x\to0}\frac{e^x-1-\sin x}{x^2}$.

\[
\begin{aligned}
e^x=1+x+\frac{x^2}{2}+o(x^2)\\
\sin x=x+o(x^2)\\
e^x-1-\sin x=\frac{x^2}{2}+o(x^2)\\
\lim=\frac12
\end{aligned}
\]
\[\boxed{\frac12}\]


\paragraph{Ejercicio 665.} Calcula $\displaystyle\lim_{x\to0}\frac{\ln(1+x)}{x}$.

\[
\begin{aligned}
\ln(1+x)=x+o(x)\\
\frac{\ln(1+x)}{x}=\frac{x+o(x)}{x}\\
\lim=1
\end{aligned}
\]
\[\boxed{1}\]


\paragraph{Ejercicio 666.} Calcula $\displaystyle\lim_{x\to0}\frac{e^x-1-\sin x}{x^2}$.

\[
\begin{aligned}
e^x=1+x+\frac{x^2}{2}+o(x^2)\\
\sin x=x+o(x^2)\\
e^x-1-\sin x=\frac{x^2}{2}+o(x^2)\\
\lim=\frac12
\end{aligned}
\]
\[\boxed{\frac12}\]


\paragraph{Ejercicio 667.} Calcula $\displaystyle\lim_{x\to0}\frac{\ln(1+x)}{x}$.

\[
\begin{aligned}
\ln(1+x)=x+o(x)\\
\frac{\ln(1+x)}{x}=\frac{x+o(x)}{x}\\
\lim=1
\end{aligned}
\]
\[\boxed{1}\]


\paragraph{Ejercicio 668.} Calcula $\displaystyle\lim_{x\to0}\frac{e^x-1-\sin x}{x^2}$.

\[
\begin{aligned}
e^x=1+x+\frac{x^2}{2}+o(x^2)\\
\sin x=x+o(x^2)\\
e^x-1-\sin x=\frac{x^2}{2}+o(x^2)\\
\lim=\frac12
\end{aligned}
\]
\[\boxed{\frac12}\]


\paragraph{Ejercicio 669.} Calcula $\displaystyle\lim_{x\to0}\frac{\ln(1+x)}{x}$.

\[
\begin{aligned}
\ln(1+x)=x+o(x)\\
\frac{\ln(1+x)}{x}=\frac{x+o(x)}{x}\\
\lim=1
\end{aligned}
\]
\[\boxed{1}\]


\paragraph{Ejercicio 670.} Calcula $\displaystyle\lim_{x\to0}\frac{e^x-1-\sin x}{x^2}$.

\[
\begin{aligned}
e^x=1+x+\frac{x^2}{2}+o(x^2)\\
\sin x=x+o(x^2)\\
e^x-1-\sin x=\frac{x^2}{2}+o(x^2)\\
\lim=\frac12
\end{aligned}
\]
\[\boxed{\frac12}\]


\paragraph{Ejercicio 671.} Calcula $\displaystyle\lim_{x\to0}\frac{\ln(1+x)}{x}$.

\[
\begin{aligned}
\ln(1+x)=x+o(x)\\
\frac{\ln(1+x)}{x}=\frac{x+o(x)}{x}\\
\lim=1
\end{aligned}
\]
\[\boxed{1}\]


\paragraph{Ejercicio 672.} Calcula $\displaystyle\lim_{x\to0}\frac{e^x-1-\sin x}{x^2}$.

\[
\begin{aligned}
e^x=1+x+\frac{x^2}{2}+o(x^2)\\
\sin x=x+o(x^2)\\
e^x-1-\sin x=\frac{x^2}{2}+o(x^2)\\
\lim=\frac12
\end{aligned}
\]
\[\boxed{\frac12}\]


\paragraph{Ejercicio 673.} Calcula $\displaystyle\lim_{x\to0}\frac{\ln(1+x)}{x}$.

\[
\begin{aligned}
\ln(1+x)=x+o(x)\\
\frac{\ln(1+x)}{x}=\frac{x+o(x)}{x}\\
\lim=1
\end{aligned}
\]
\[\boxed{1}\]


\paragraph{Ejercicio 674.} Calcula $\displaystyle\lim_{x\to0}\frac{e^x-1-\sin x}{x^2}$.

\[
\begin{aligned}
e^x=1+x+\frac{x^2}{2}+o(x^2)\\
\sin x=x+o(x^2)\\
e^x-1-\sin x=\frac{x^2}{2}+o(x^2)\\
\lim=\frac12
\end{aligned}
\]
\[\boxed{\frac12}\]


\paragraph{Ejercicio 675.} Calcula $\displaystyle\lim_{x\to0}\frac{\ln(1+x)}{x}$.

\[
\begin{aligned}
\ln(1+x)=x+o(x)\\
\frac{\ln(1+x)}{x}=\frac{x+o(x)}{x}\\
\lim=1
\end{aligned}
\]
\[\boxed{1}\]


\paragraph{Ejercicio 676.} Calcula $\displaystyle\lim_{x\to0}\frac{e^x-1-\sin x}{x^2}$.

\[
\begin{aligned}
e^x=1+x+\frac{x^2}{2}+o(x^2)\\
\sin x=x+o(x^2)\\
e^x-1-\sin x=\frac{x^2}{2}+o(x^2)\\
\lim=\frac12
\end{aligned}
\]
\[\boxed{\frac12}\]


\paragraph{Ejercicio 677.} Calcula $\displaystyle\lim_{x\to0}\frac{\ln(1+x)}{x}$.

\[
\begin{aligned}
\ln(1+x)=x+o(x)\\
\frac{\ln(1+x)}{x}=\frac{x+o(x)}{x}\\
\lim=1
\end{aligned}
\]
\[\boxed{1}\]


\paragraph{Ejercicio 678.} Calcula $\displaystyle\lim_{x\to0}\frac{e^x-1-\sin x}{x^2}$.

\[
\begin{aligned}
e^x=1+x+\frac{x^2}{2}+o(x^2)\\
\sin x=x+o(x^2)\\
e^x-1-\sin x=\frac{x^2}{2}+o(x^2)\\
\lim=\frac12
\end{aligned}
\]
\[\boxed{\frac12}\]


\paragraph{Ejercicio 679.} Calcula $\displaystyle\lim_{x\to0}\frac{\ln(1+x)}{x}$.

\[
\begin{aligned}
\ln(1+x)=x+o(x)\\
\frac{\ln(1+x)}{x}=\frac{x+o(x)}{x}\\
\lim=1
\end{aligned}
\]
\[\boxed{1}\]


\paragraph{Ejercicio 680.} Calcula $\displaystyle\lim_{x\to0}\frac{e^x-1-\sin x}{x^2}$.

\[
\begin{aligned}
e^x=1+x+\frac{x^2}{2}+o(x^2)\\
\sin x=x+o(x^2)\\
e^x-1-\sin x=\frac{x^2}{2}+o(x^2)\\
\lim=\frac12
\end{aligned}
\]
\[\boxed{\frac12}\]


\newpage

\subsubsection*{Nivel 5 - Imposible}

\paragraph{Ejercicio 681.} Calcula $\displaystyle \lim_{x\to0}\frac{\ln(1+x)-x+\frac{x^2}{2}}{x^3}$.

\[
\begin{aligned}
\ln(1+x)=x-\frac{x^2}{2}+\frac{x^3}{3}+o(x^3)\\
\ln(1+x)-x+\frac{x^2}{2}=\frac{x^3}{3}+o(x^3)\\
\lim=\frac13
\end{aligned}
\]
\[\boxed{\frac13}\]


\paragraph{Ejercicio 682.} Calcula $\displaystyle \lim_{x\to0}\frac{e^x-\cos x-\sin x}{x^2}$.

\[
\begin{aligned}
e^x=1+x+\frac{x^2}{2}+o(x^2)\\
\cos x=1-\frac{x^2}{2}+o(x^2)\\
\sin x=x+o(x^2)\\
e^x-\cos x-\sin x=x^2+o(x^2)\\
\lim=1
\end{aligned}
\]
\[\boxed{1}\]


\paragraph{Ejercicio 683.} Calcula $\displaystyle \lim_{x\to0^+}x\ln x$.

\[
\begin{aligned}
x=\frac1t\\
x\to0^+\Rightarrow t\to\infty\\
x\ln x=\frac1t\ln\frac1t=-\frac{\ln t}{t}\\
\lim=0
\end{aligned}
\]
\[\boxed{0}\]


\paragraph{Ejercicio 684.} Calcula $\displaystyle \lim_{x\to0}\frac{|x|}{x}$.

\[
\begin{aligned}
x\to0^-:\frac{|x|}{x}=-1\\
x\to0^+:\frac{|x|}{x}=1\\
-1\ne1\\
\lim\text{{ no existe}}
\end{aligned}
\]
\[\boxed{\text{{No existe}}}\]


\paragraph{Ejercicio 685.} Calcula $\displaystyle \lim_{x\to0}\frac{\sqrt[3]{1+x}-1-\frac{x}{3}}{x^2}$.

\[
\begin{aligned}
(1+x)^{1/3}=1+\frac{x}{3}-\frac{x^2}{9}+o(x^2)\\
\sqrt[3]{1+x}-1-\frac{x}{3}=-\frac{x^2}{9}+o(x^2)\\
\lim=-\frac19
\end{aligned}
\]
\[\boxed{-\frac19}\]


\paragraph{Ejercicio 686.} Calcula $\displaystyle \lim_{x\to0}\frac{\ln(1+x)-x+\frac{x^2}{2}}{x^3}$.

\[
\begin{aligned}
\ln(1+x)=x-\frac{x^2}{2}+\frac{x^3}{3}+o(x^3)\\
\ln(1+x)-x+\frac{x^2}{2}=\frac{x^3}{3}+o(x^3)\\
\lim=\frac13
\end{aligned}
\]
\[\boxed{\frac13}\]


\paragraph{Ejercicio 687.} Calcula $\displaystyle \lim_{x\to0}\frac{e^x-\cos x-\sin x}{x^2}$.

\[
\begin{aligned}
e^x=1+x+\frac{x^2}{2}+o(x^2)\\
\cos x=1-\frac{x^2}{2}+o(x^2)\\
\sin x=x+o(x^2)\\
e^x-\cos x-\sin x=x^2+o(x^2)\\
\lim=1
\end{aligned}
\]
\[\boxed{1}\]


\paragraph{Ejercicio 688.} Calcula $\displaystyle \lim_{x\to0^+}x\ln x$.

\[
\begin{aligned}
x=\frac1t\\
x\to0^+\Rightarrow t\to\infty\\
x\ln x=\frac1t\ln\frac1t=-\frac{\ln t}{t}\\
\lim=0
\end{aligned}
\]
\[\boxed{0}\]


\paragraph{Ejercicio 689.} Calcula $\displaystyle \lim_{x\to0}\frac{|x|}{x}$.

\[
\begin{aligned}
x\to0^-:\frac{|x|}{x}=-1\\
x\to0^+:\frac{|x|}{x}=1\\
-1\ne1\\
\lim\text{{ no existe}}
\end{aligned}
\]
\[\boxed{\text{{No existe}}}\]


\paragraph{Ejercicio 690.} Calcula $\displaystyle \lim_{x\to0}\frac{\sqrt[3]{1+x}-1-\frac{x}{3}}{x^2}$.

\[
\begin{aligned}
(1+x)^{1/3}=1+\frac{x}{3}-\frac{x^2}{9}+o(x^2)\\
\sqrt[3]{1+x}-1-\frac{x}{3}=-\frac{x^2}{9}+o(x^2)\\
\lim=-\frac19
\end{aligned}
\]
\[\boxed{-\frac19}\]


\paragraph{Ejercicio 691.} Calcula $\displaystyle \lim_{x\to0}\frac{\ln(1+x)-x+\frac{x^2}{2}}{x^3}$.

\[
\begin{aligned}
\ln(1+x)=x-\frac{x^2}{2}+\frac{x^3}{3}+o(x^3)\\
\ln(1+x)-x+\frac{x^2}{2}=\frac{x^3}{3}+o(x^3)\\
\lim=\frac13
\end{aligned}
\]
\[\boxed{\frac13}\]


\paragraph{Ejercicio 692.} Calcula $\displaystyle \lim_{x\to0}\frac{e^x-\cos x-\sin x}{x^2}$.

\[
\begin{aligned}
e^x=1+x+\frac{x^2}{2}+o(x^2)\\
\cos x=1-\frac{x^2}{2}+o(x^2)\\
\sin x=x+o(x^2)\\
e^x-\cos x-\sin x=x^2+o(x^2)\\
\lim=1
\end{aligned}
\]
\[\boxed{1}\]


\paragraph{Ejercicio 693.} Calcula $\displaystyle \lim_{x\to0^+}x\ln x$.

\[
\begin{aligned}
x=\frac1t\\
x\to0^+\Rightarrow t\to\infty\\
x\ln x=\frac1t\ln\frac1t=-\frac{\ln t}{t}\\
\lim=0
\end{aligned}
\]
\[\boxed{0}\]


\paragraph{Ejercicio 694.} Calcula $\displaystyle \lim_{x\to0}\frac{|x|}{x}$.

\[
\begin{aligned}
x\to0^-:\frac{|x|}{x}=-1\\
x\to0^+:\frac{|x|}{x}=1\\
-1\ne1\\
\lim\text{{ no existe}}
\end{aligned}
\]
\[\boxed{\text{{No existe}}}\]


\paragraph{Ejercicio 695.} Calcula $\displaystyle \lim_{x\to0}\frac{\sqrt[3]{1+x}-1-\frac{x}{3}}{x^2}$.

\[
\begin{aligned}
(1+x)^{1/3}=1+\frac{x}{3}-\frac{x^2}{9}+o(x^2)\\
\sqrt[3]{1+x}-1-\frac{x}{3}=-\frac{x^2}{9}+o(x^2)\\
\lim=-\frac19
\end{aligned}
\]
\[\boxed{-\frac19}\]


\paragraph{Ejercicio 696.} Calcula $\displaystyle \lim_{x\to0}\frac{\ln(1+x)-x+\frac{x^2}{2}}{x^3}$.

\[
\begin{aligned}
\ln(1+x)=x-\frac{x^2}{2}+\frac{x^3}{3}+o(x^3)\\
\ln(1+x)-x+\frac{x^2}{2}=\frac{x^3}{3}+o(x^3)\\
\lim=\frac13
\end{aligned}
\]
\[\boxed{\frac13}\]


\paragraph{Ejercicio 697.} Calcula $\displaystyle \lim_{x\to0}\frac{e^x-\cos x-\sin x}{x^2}$.

\[
\begin{aligned}
e^x=1+x+\frac{x^2}{2}+o(x^2)\\
\cos x=1-\frac{x^2}{2}+o(x^2)\\
\sin x=x+o(x^2)\\
e^x-\cos x-\sin x=x^2+o(x^2)\\
\lim=1
\end{aligned}
\]
\[\boxed{1}\]


\paragraph{Ejercicio 698.} Calcula $\displaystyle \lim_{x\to0^+}x\ln x$.

\[
\begin{aligned}
x=\frac1t\\
x\to0^+\Rightarrow t\to\infty\\
x\ln x=\frac1t\ln\frac1t=-\frac{\ln t}{t}\\
\lim=0
\end{aligned}
\]
\[\boxed{0}\]


\paragraph{Ejercicio 699.} Calcula $\displaystyle \lim_{x\to0}\frac{|x|}{x}$.

\[
\begin{aligned}
x\to0^-:\frac{|x|}{x}=-1\\
x\to0^+:\frac{|x|}{x}=1\\
-1\ne1\\
\lim\text{{ no existe}}
\end{aligned}
\]
\[\boxed{\text{{No existe}}}\]


\paragraph{Ejercicio 700.} Calcula $\displaystyle \lim_{x\to0}\frac{\sqrt[3]{1+x}-1-\frac{x}{3}}{x^2}$.

\[
\begin{aligned}
(1+x)^{1/3}=1+\frac{x}{3}-\frac{x^2}{9}+o(x^2)\\
\sqrt[3]{1+x}-1-\frac{x}{3}=-\frac{x^2}{9}+o(x^2)\\
\lim=-\frac19
\end{aligned}
\]
\[\boxed{-\frac19}\]


\newpage
\end{document}
